% !TeX program = xelatex
% !TeX encoding = UTF-8
\documentclass[a4paper,12pt,oneside,fontpreset=auto,theme=ocean,
  sectionstyle=classic,answers=true,solutions=appendix]{maki-notes}
\usepackage{functional-analysis}

\renewcommand{\LectureNotesTitle}{泛函分析}
\renewcommand{\AuthorInfo}{书稿框架与专题讲义}
\renewcommand{\CoverTitle}{泛函分析}
\renewcommand{\CoverAuthorBlock}{吉田耕作第六版主线\\广义函数专题扩展}
\renewcommand{\CoverFooterBlock}{以中译本为主，英文第六版为辅助\\逐章编写版 · 2026 年 10 月}
\hypersetup{pdftitle={泛函分析：吉田主线与广义函数专题},pdfauthor={书稿框架与专题讲义},
  colorlinks=true,linkcolor=LogicColor,citecolor=DefColor,urlcolor=ExColor}
\SetCourseInfo{course={泛函分析},audience={高年级本科生与研究生},
  term={第零至第七章正文与全书框架},date={2026-10-26}}

% 快速编译时取消下面一行的注释；章节编号沿用完整编译生成的 aux。
% \includeonly{books/03-distributions}
% 只编译吉田主线中的一章：
% \includeonly{books/yosida/01-seminorms}
\begin{document}
\frontmatter
\hypersetup{pageanchor=false}
\MakeTitlePage
\clearpage
\hypersetup{pageanchor=true}
\chapter*{前言与阅读约定}
\addcontentsline{toc}{chapter}{前言与阅读约定}
本书以吉田耕作《泛函分析》第六版中译本为主线，以英文第六版辅助核对术语、目录和页码。第一部分保留原书第零章至第十四章的章次与各节主题，包括第五章和第十章的原书附录。目前第零至第十一章共 110/132 个编号节已有首轮正文，附 186 道解答习题；第五章与第十章附录已写，第十二至十四章仍保留占位。已有正文的比例不代表阅读进度、工作量比例或最终审校完成率。

主线从预备知识、半范数与 Banach 空间基本定理进入 Hilbert 空间、Hahn--Banach、强弱收敛及对偶算子，再连接 Fourier 与微分方程、预解式和谱、半群及发展方程。第八章加入两类平均过程、殆周期均值与全纯演算；第九章展开生成元、Hille--Yosida、Lumer--Phillips、Stone、全纯半群、分数幂、Trotter--Kato 及对偶连续部分。第十章连接紧性判别、核算子、Sobolev 紧嵌入、Fredholm 谱论与 Dirichlet 弱解，并补写核空间附录。第十一章从赋范代数进入谱测度、无界自伴演算、紧群表示和 Wiener 逼近。每章核查记录区分正文覆盖、作者自查、排版验收与独立数学审校。

部分一般结论仍作为外部前提或待展开内容，例如 Riesz--Markov、一般 Eberlein--\v Smulian、Titchmarsh 与变系数正则性，以及部分局部凸推广、扇形预解式的反向构造和一般从属核。具体假设、直接查阅范围与证明边界在各节及 COVERAGE 中注明；不把首轮写成或编译通过等同于独立数学审校完成。

第二部分“广义函数专题扩展”接入此前已经写成的八章内容，章号顺接为第十五章至第二十二章。它以施瓦尔茨《广义函数论》和盖尔芳特、希烙夫《广义函数 III：微分方程理论的若干问题》为专题参考。吉田原书本身涉及广义函数的各节仍保留在第一部分，后续写作可引用专题中的详细证明，也可补充原书特有的论述。本次调整保持现有专题正文与习题，采用标签引用适应新的章号。

目录框架中的章节名与节名以中译本为准，英文版辅助核对人名与对应主题。每个主线占位节记录两版印刷起始页；附录另列 PDF 页序的换算。后续应逐节补充定义、定理、证明、例子和习题，并统一记号。下面的配对与 Fourier 约定目前适用于已写成的广义函数分部，主线正文写作时也应明确它与原书约定的差别。

\Needspace{10\baselineskip}
\section*{符号与配对}
\phantomsection
\addcontentsline{toc}{section}{符号与配对}
始终令 $\Omega\subset\R^n$ 为开集，$K\Subset\Omega$ 表示紧集 $K$ 包含于 $\Omega$。多重指标 $\alpha=(\alpha_1,\ldots,\alpha_n)\in\N_0^n$，并约定
\[
 |\alpha|=\sum_j\alpha_j,\quad \alpha!=\prod_j\alpha_j!,\quad
 x^\alpha=\prod_jx_j^{\alpha_j},\quad
 \partial^\alpha=\partial_{x_1}^{\alpha_1}\cdots\partial_{x_n}^{\alpha_n}.
\]
分布是复线性泛函，$\Pair{T}{\varphi}$ 对 $\varphi$ 线性，正则分布采用
\[
 \Pair{T_f}{\varphi}=\int_\Omega f(x)\varphi(x)\,\dd x,
\]
不在测试函数上加复共轭。Hilbert 空间的内积则另行采用 $ (u,v)=\int u\overline v$，对第一变量线性。这两种约定不可混同。

\begin{center}
\begin{tabular}{ll}
\toprule
符号 & 含义\\\midrule
$\Test(\Omega)$ & $C_c^\infty(\Omega)$，带标准 LF 拓扑\\
$\Dist(\Omega)$ & $\Test(\Omega)$ 的连续线性对偶\\
$\Smooth(\Omega)$ & $C^\infty(\Omega)$，带紧集上各阶一致收敛拓扑\\
$\CompactDist(\Omega)$ & 紧支集分布空间\\
$\Sch(\R^n),\Temp(\R^n)$ & Schwartz 空间与温和分布空间\\
$\supp T$ & 分布 $T$ 的支集\\
$T_j\to T$ & 默认指逐测试函数收敛，即弱对偶收敛\\
\bottomrule
\end{tabular}
\end{center}

傅里叶变换统一使用角频率约定：
\begin{equation}\label{eq:fourier-convention}
 \widehat f(\xi)=\int_{\R^n}e^{-ix\cdot\xi}f(x)\,\dd x,
 \qquad f(x)=(2\pi)^{-n}\int_{\R^n}e^{ix\cdot\xi}\widehat f(\xi)\,\dd\xi.
\end{equation}
施瓦尔茨原书部分公式使用 $2\pi$ 频率约定，盖尔芳特、希烙夫原书相关章节采用正号指数；查阅时须作符号与常数换算。

\section*{建议的阅读路径}
主线按吉田原书从预备知识读至发展方程。当前可直接阅读的专题正文从第 \ref{ch:motivation} 章起：第 \ref{ch:test}--\ref{ch:operations} 章建立测试空间与分布运算，第 \ref{ch:convolution}--\ref{ch:fourier} 章讨论卷积与 Fourier 变换，第 \ref{ch:pde} 章讨论微分方程模型，第 \ref{ch:weighted} 章介绍加权空间与广义固有函数。已写成的习题解答集中置于书后；新增占位节目前没有虚构的习题或证明。

\tableofcontents
\mainmatter
% 第一部分保留吉田原书的 0--14 章编号，不以分部重排原章次。
\part{泛函分析主线：吉田耕作第六版}
\setcounter{chapter}{-1}
% !TeX encoding = UTF-8
% 吉田第六版：中译本第 0 章，四节正文与参考文献整理。
\SetChapterInfo{subtitle={集合、拓扑、测度与线性结构},keywords={选择原理; 紧性; 函数逼近; 积分; 商空间},prereq={集合与映射、微积分、线性代数},goals={理解选择原理的等价性; 掌握紧性与类别方法; 证明积分收敛与交换定理; 建立线性算子和商空间}}
\chapter{预备知识}\label{ch:yosida-00}
\ChapterSource{本章保留吉田耕作第六版第零章的四节结构，原书位置为中译本与英文版印刷页 1--22，参见 \cite{yosida-zh,yosida-en}。正文按原书主题重组为讲义，并补充选择原理的等价性证明、积分构造及代数应用；使用现代测度术语，涉及原书与本稿的记号差别时另行说明。本章参考作者及其总书目条目已核对，正文不是逐句翻译。}

\section{集合论}\label{sec:yosida-00-01}
\noindent\textit{原书位置：中译本与英文版本节均从印刷第 1 页起。}
% !TeX encoding = UTF-8
泛函分析中常须把一个局部构造延拓到整个空间。例如，线性无关集能否扩充为基，或者已定义的线性泛函能否继续延拓？这类问题常通过“按包含关系排列所有候选对象，再取一个极大对象”解决。本节建立这一方法的集合论基础。

以下在 Zermelo--Fraenkel 集合论（ZF）中讨论；选择公理作为额外命题单独提出。所有偏序集、索引族及线性空间的底集都是集合。记 $\mathcal P(X)$ 为 $X$ 的幂集，$\varnothing$ 为空集。三个选择原理之间的证明不能预先借用其中另一个；为此，先交代良序与超限递归所需的基础工具。

\subsection{偏序、链与极大元}
\begin{definition}[偏序与全序]\label{def:y0-partial-order}
集合 $P$ 上的二元关系 $\le$ 称为\emph{偏序}，若它满足：
\begin{enumerate}
  \item 自反性：$x\le x$，对每个 $x\in P$ 成立；
  \item 反对称性：$x\le y$ 且 $y\le x$ 蕴含 $x=y$；
  \item 传递性：$x\le y$ 且 $y\le z$ 蕴含 $x\le z$。
\end{enumerate}
称 $(P,\le)$ 为偏序集，记 $x<y$ 表示 $x\le y$ 且 $x\ne y$。若任意两个元素都可比较，即 $x\le y$ 或 $y\le x$，则此偏序称为\emph{全序}。
\end{definition}

\begin{definition}[链、上界、最大元与极大元]\label{def:y0-chain}
设 $(P,\le)$ 为偏序集，$A\subseteq P$。
\begin{enumerate}
  \item 若 $C\subseteq P$ 中任意两元素可比较，则称 $C$ 为\emph{链}；允许 $C=\varnothing$。
  \item 若 $u\in P$ 且每个 $a\in A$ 都满足 $a\le u$，则称 $u$ 为 $A$ 在 $P$ 中的\emph{上界}。上界不必属于 $A$。上界中的最小元称为 $A$ 的\emph{上确界}，记为 $\sup A$。
  \item 若 $a_*\in A$ 且每个 $a\in A$ 都有 $a\le a_*$，则称 $a_*$ 为 $A$ 的\emph{最大元}。
  \item 若 $m\in P$ 且 $m\le x$ 蕴含 $x=m$，则称 $m$ 为 $P$ 的\emph{极大元}。
\end{enumerate}
反向不等式分别定义下界、下确界、最小元与极小元。讨论子集的极大元时，使用从 $P$ 限制得到的偏序。
\end{definition}

最大元必须支配所有元素，而极大元只要求不存在严格更大的元素。最大元若存在则唯一，并且是极大元；极大元可以不唯一。在全序集中，极大元也是最大元。空链的每个上界条件都是空条件，故它的上界恰为 $P$ 的全部元素。

\begin{example}[包含序与整除序]
幂集 $\mathcal P(X)$ 在包含关系 $\subseteq$ 下是偏序集。若 $X$ 至少有两个元素，它通常不是全序集。任意子族 $\mathcal A\subseteq\mathcal P(X)$ 的上确界为 $\bigcup\mathcal A$，下确界为 $\bigcap\mathcal A$；这里约定空族的交为 $X$。

在 $P=\{\{1\},\{2\}\}$ 的包含序下，两个元素均极大，却没有最大元。正整数集在整除关系下也是偏序集；例如 $2$ 与 $3$ 不可比较。通常的数值大小和整除关系是两种不同的偏序。
\end{example}

\subsection{良序与超限递归的基础}
\begin{definition}[良序与初始段]\label{def:y0-well-order}
全序集 $(X,\le)$ 称为\emph{良序集}，若每个非空子集 $A\subseteq X$ 都有最小元。对 $x\in X$，集合
\[
 X_{<x}=\{y\in X:y<x\}
\]
称为由 $x$ 决定的严格初始段。一般地，$I\subseteq X$ 若满足 $x\in I$ 且 $y<x$ 时 $y\in I$，则称 $I$ 为初始段。
\end{definition}

自然数集在通常次序下良序。整数集不是良序，因为整个整数集没有最小元；实数集不是良序，例如 $(0,1)$ 没有最小元。有限全序集总是良序。\emph{良序原则断言可给一个集合另赋某种良序，并非断言它原有的次序就是良序。}

\begin{proposition}[良序归纳]\label{prop:y0-well-order-induction}
设 $X$ 良序。若性质 $Q(x)$ 满足：对任意 $x\in X$，所有 $y<x$ 满足 $Q(y)$ 时，$Q(x)$ 也成立，则每个 $x\in X$ 都满足 $Q(x)$。
\end{proposition}
\begin{proof}
若反例集非空，取其中的最小元 $x_0$。所有 $y<x_0$ 都不是反例，由假设 $Q(x_0)$ 成立，矛盾。
\end{proof}

为精确说明“逐个选择，直到选完”何时停止，使用序数记号。\emph{序数}是由成员关系 $\in$ 严格良序的传递集合；传递指 $u\in\alpha$ 蕴含 $u\subseteq\alpha$。序数的严格次序记为 $\beta<\alpha$，即 $\beta\in\alpha$。最小序数为 $0=\varnothing$，后继序数为 $\alpha+1=\alpha\cup\{\alpha\}$；非零且不是后继的序数称为极限序数。自然数的序数集合记为 $\omega$。一族序数的并仍为序数，且是该族的上确界。

每个良序集都与唯一序数保序同构，此序数称为其\emph{序型}。这里“保序同构”指保持并反映次序的双射。其构造是沿良序定义
\[
 r(x)=\{r(y):y<x\}.
\]
良序归纳表明各 $r(x)$ 是序数，且 $y<x$ 当且仅当 $r(y)\in r(x)$；像集是传递的序数。唯一性也由良序归纳得到：两个保序的编号在每个位置都有相同的前驱编号，因而相等。这一构造使用下面的递归定理，它在 ZF 内成立，不需要选择公理。

\begin{lemma}[沿良序的递归]\label{lem:y0-transfinite-recursion}
设 $X$ 为良序集。若一条确定规则对每个 $x\in X$ 及已定义在 $X_{<x}$ 上的函数 $g$ 指定唯一集合值 $F(x,g)$，则存在唯一函数 $f$，定义域为 $X$，且
\[
 f(x)=F\bigl(x,f|_{X_{<x}}\bigr)\qquad(x\in X).
\]
在序数上使用此定理称为超限递归；在极限阶段，已有函数取并，得到此前全部取值。
\end{lemma}
\begin{proof}
称定义在某个初始段上并满足递归式的函数为近似解。任意两个近似解在公共定义域上一致：若有分歧，取最早分歧点；该点之前的函数相同，规则的唯一性迫使此点取值也相同，矛盾。

用良序归纳证明每个 $X_{<x}$ 上有唯一近似解。若所有更早的初始段都有解，对每个 $y<x$ 将 $X_{<y}$ 上的解延拓一步，定义 $y$ 处的取值为 $F(y,g)$。由替代公理，这些唯一确定的延拓组成一个集合；由相容性，其并给出 $X_{<x}$ 上的解。空初始段的解为空函数。最后，对所有 $x\in X$ 再将这些解延拓一步并取并，得到整个 $X$ 上的 $f$。相容性同时证明唯一性。这里收集的是唯一确定的对象，所用的是替代公理，而非选择公理。
\end{proof}

\begin{lemma}[Hartogs 引理]\label{lem:y0-hartogs}
对任意集合 $X$，存在序数 $\theta$，使得不存在从 $\theta$ 到 $X$ 的单射。此结论可在 ZF 中证明。
\end{lemma}
\begin{proof}
考虑 $X$ 的各个子集上的所有良序关系。每一项可表示为一对 $(A,R)$，其中 $A\subseteq X$，$R\subseteq X\times X$，故这些对象构成一个集合。对每个这样的良序取唯一序型，由替代公理得到一个序数集合 $T$。令
\[
 \theta=\bigcup_{\alpha\in T}(\alpha+1).
\]
这是序数，且对所有 $\alpha\in T$ 都有 $\alpha<\theta$。若有单射 $j:\theta\to X$，把 $\theta$ 的良序搬到像集 $j[\theta]$ 上，得到 $X$ 的一个子集上的良序，其序型正是 $\theta$。于是 $\theta\in T$，从而 $\theta<\theta$，矛盾。
\end{proof}

Hartogs 引理给出递归选择不能无限延长的集合论界限。只说“每次选一个，最后必能选完”并不构成证明；若只按自然数编号，更无法遍历任意不可数集合。

\subsection{选择公理、良序原则与 Zorn 引理}
\begin{definition}[选择函数与选择公理]\label{def:y0-choice}
设 $(A_i)_{i\in I}$ 为一族非空集合。函数 $f:I\to\bigcup_{i\in I}A_i$ 若满足 $f(i)\in A_i$，称为这族集合的\emph{选择函数}。\emph{选择公理}（AC）断言：每一族非空集合都有选择函数。等价地，每一这样的笛卡尔积都非空：
\[
 \prod_{i\in I}A_i
 =\left\{f:I\to\bigcup_{i\in I}A_i:
          f(i)\in A_i\text{ 对所有 }i\in I\right\}\ne\varnothing.
\]
若 $I=\varnothing$，空函数即为选择函数，此时乘积是单元素集 $\{\varnothing\}$。
\end{definition}

有限族非空集合的选择函数可用有限归纳在 ZF 中构造；一般无限族的结论是额外公理。若每个 $A_i$ 的元素由一条统一规则唯一确定，例如非空自然数子集的最小元，也不需要选择公理。选择公理保证函数存在，不提供可计算的选择算法。

\begin{theorem}[良序原则]\label{thm:y0-well-order-principle}
每个集合 $X$ 都可以赋予一个良序关系。空集的空关系也视为良序。
\end{theorem}

\begin{theorem}[Zorn 引理]\label{thm:y0-zorn}
设 $(P,\le)$ 是非空偏序集。如果每条链 $C\subseteq P$ 都在 $P$ 中有上界，则 $P$ 至少有一个极大元。
\end{theorem}

Zorn 引理只要求上界存在，不要求上界唯一，也不要求上确界存在。若把空链也包含在条件中，“每条链都有上界”本身已蕴含 $P\ne\varnothing$；若只考察非空链，则必须另列 $P\ne\varnothing$。

\begin{theorem}[三种选择原理的等价性]\label{thm:y0-choice-equivalence}
在 ZF 集合论中，以下三命题等价：选择公理、良序原则、Zorn 引理。
\end{theorem}
\begin{proof}
依次证明
\[
 \text{选择公理}\ \Longrightarrow\ \text{良序原则}
 \ \Longrightarrow\ \text{Zorn 引理}
 \ \Longrightarrow\ \text{选择公理}.
\]

\emph{第一步：选择公理推出良序原则。}
若 $X=\varnothing$，结论直接成立。设 $X\ne\varnothing$，对所有非空子集应用选择公理，固定函数
\[
 c:\mathcal P(X)\setminus\{\varnothing\}\longrightarrow X,
 \qquad c(A)\in A.
\]
由 Hartogs 引理取不能单射到 $X$ 的序数 $\theta$。取一个不在 $X$ 中的标记 $\ast$，沿 $\theta$ 递归定义
\[
 R_\alpha=X\setminus
       \{x_\beta:\beta<\alpha,\ x_\beta\in X\},\qquad
 x_\alpha=
 \begin{cases}
   c(R_\alpha),&R_\alpha\ne\varnothing,\\
   \ast,&R_\alpha=\varnothing.
 \end{cases}
\]
若每个 $\alpha<\theta$ 都有 $R_\alpha\ne\varnothing$，则 $\alpha\mapsto x_\alpha$ 是 $\theta$ 到 $X$ 的单射，因为每次所取元素均未在此前出现。这与 $\theta$ 的定义矛盾。

故存在最小的 $\gamma<\theta$ 使 $R_\gamma=\varnothing$。在 $\gamma$ 之前，取值两两不同，且已遍历 $X$，于是 $\alpha\mapsto x_\alpha$ 是 $\gamma$ 到 $X$ 的双射。定义 $x_\alpha\prec x_\beta$ 当且仅当 $\alpha<\beta$，便把序数 $\gamma$ 的良序搬到了 $X$ 上。

\emph{第二步：良序原则推出 Zorn 引理。}
设 $P$ 非空且每条链都有上界。由良序原则，将 $P$ 的全部元素无重复地写作 $(p_\alpha)_{\alpha<\kappa}$，其中 $\kappa$ 是某个非零序数。按这个编号逐个检查元素，只在新元素严格大于所有已保留元素时将它保留。精确地，令 $C_0=\varnothing$，并递归定义
\[
 C_{\alpha+1}=
 \begin{cases}
 C_\alpha\cup\{p_\alpha\},
       &x<p_\alpha\text{ 对每个 }x\in C_\alpha,\\
 C_\alpha,&\text{否则},
 \end{cases}
 \qquad
 C_\lambda=\bigcup_{\alpha<\lambda}C_\alpha
 \quad(\lambda\text{ 为极限序数}).
\]
良序归纳表明各 $C_\alpha$ 是链，且随 $\alpha$ 递增；极限阶段的并仍为链，因为任意两个元素已共同处于某个更早阶段。于是 $C_\kappa$ 是链，并且 $p_0$ 在第一步被保留，所以它非空。

取 $C_\kappa$ 的一个上界 $u\in P$。若 $u$ 不是极大元，则存在 $v\in P$ 使 $u<v$。设 $v=p_\beta$。在检查 $p_\beta$ 的阶段，对所有 $x\in C_\beta$ 都有 $x\le u<v$，所以 $v$ 被保留，属于 $C_\kappa$。上界条件又给出 $v\le u$，与 $u<v$ 矛盾。因此 $u$ 为极大元。

\emph{第三步：Zorn 引理推出选择公理。}
设 $(A_i)_{i\in I}$ 为非空集合族，令 $U=\bigcup_{i\in I}A_i$。考虑全部部分选择函数：
\[
 P=\{f:J\to U:J\subseteq I,\ f(i)\in A_i\ (i\in J)\},
\]
把函数看作图，即 $I\times U$ 的子集，以图的包含关系作为偏序。因空函数属于 $P$，此偏序集非空。

若 $\mathcal C\subseteq P$ 是链，令 $g=\bigcup_{f\in\mathcal C}f$。这是函数：同一指标若在两个成员的定义域中出现，这两个成员可比较，较大的函数延拓较小的函数，取值必然一致。且 $g(i)\in A_i$，所以 $g\in P$，并为 $\mathcal C$ 的上界。空链的并为空函数，也满足要求。

由 Zorn 引理存在极大部分选择函数 $f_*$。若其定义域 $J_*$ 不等于 $I$，取 $i_0\in I\setminus J_*$，再取一个 $a_0\in A_{i_0}$，则
\[
 f_*\cup\{(i_0,a_0)\}
\]
是严格延拓 $f_*$ 的部分选择函数，与极大性矛盾。因此 $J_*=I$，$f_*$ 即为所求选择函数。这里仅从一个非空集合取一个元素，不是在所有 $A_i$ 中同时作选择，故没有循环使用选择公理。
\end{proof}

上述闭环已经证明任意两命题互相蕴含。下面再给出两个常用的直接证明，帮助识别“良序化”和“取极大延拓”这两种构造。

\begin{proposition}[良序原则直接推出选择公理]\label{prop:y0-wo-ac}
若每个集合都可良序化，则每族非空集合都有选择函数。
\end{proposition}
\begin{proof}
将 $U=\bigcup_{i\in I}A_i$ 良序化，定义 $f(i)$ 为 $A_i$ 在这一良序下的最小元。非空性保证最小元存在，良序性保证唯一性；替代公理将这些唯一取值收集为函数。若 $I$ 为空，取空函数。
\end{proof}

\begin{proposition}[Zorn 引理直接推出良序原则]\label{prop:y0-zorn-wo}
若 Zorn 引理成立，则每个集合 $X$ 都可良序化。
\end{proposition}
\begin{proof}
令 $P$ 为 $X$ 的各个子集上的所有良序对 $(A,\le_A)$。规定
\[
 (A,\le_A)\preccurlyeq(B,\le_B)
\]
当且仅当 $A\subseteq B$，$\le_B$ 在 $A$ 上限制为 $\le_A$，且 $A$ 是 $B$ 的初始段。这是偏序，空集上的良序对使 $P$ 非空。

对非空链 $\mathcal C$，取底集的并 $B$ 和各个序关系的并 $\le_B$。任意有限个元素都共同处于某个链成员中，故 $\le_B$ 为全序。它还是良序：设 $S\subseteq B$ 非空，取 $s\in S$ 及含 $s$ 的链成员 $A$，令 $a$ 为 $S\cap A$ 在 $A$ 中的最小元。若 $t\in S\setminus A$，含 $t$ 的链成员必延拓 $A$，且 $A$ 是它的初始段，所以 $a<t$。于是 $a$ 是 $S$ 的最小元。每个链成员均为 $B$ 的初始段，故 $(B,\le_B)$ 为链的上界。空链由空良序对作为上界。

由 Zorn 引理取极大良序对 $(A_*,\le_*)$。若 $A_*\ne X$，取 $x\in X\setminus A_*$，把 $x$ 放在 $A_*$ 的全部元素之后。新次序仍是良序，且 $A_*$ 是其初始段，从而得到严格延拓，矛盾。所以 $A_*=X$。
\end{proof}

\begin{remark}[为什么需要初始段条件]
上一个证明若只要求两个良序的限制相容，其链的并未必良序。例如 $A_n=\{-n,\ldots,-1,0\}$ 各自在通常大小次序下良序，次序逐步延拓，但它们的并没有最小元。初始段条件禁止在已有元素之前不断插入新元素，正是并集良序性的关键。
\end{remark}

\subsection{极大链原理与泛函分析中的应用}
\begin{theorem}[Hausdorff 极大链原理]\label{thm:y0-maximal-chain}
每个偏序集的任意一条链都包含于一条按包含关系极大的链。此原理在 ZF 中与 Zorn 引理等价。
\end{theorem}
\begin{proof}
先假设 Zorn 引理。固定链 $C_0\subseteq P$，考虑全部包含 $C_0$ 的链组成的集合 $\mathcal Q$，以包含关系排序。它非空。$\mathcal Q$ 中一条链的并仍为 $P$ 中的链：任意两元素分别来自两个可比较的成员，故共同属于较大者。这个并包含 $C_0$，所以是 $\mathcal Q$ 中的上界；空链可用 $C_0$ 作上界。Zorn 引理给出所求极大链。

反过来，设 $P$ 非空且每条链有上界。用极大链原理将空链扩充为极大链 $C$，取其上界 $u$。由于 $C\cup\{u\}$ 仍为链，极大性迫使 $u\in C$。若 $u<v$，则 $C\cup\{v\}$ 也为链，故 $v\in C$；但 $u$ 是 $C$ 的上界，得到 $v\le u$，矛盾。所以 $u$ 为极大元。
\end{proof}

\begin{theorem}[线性无关集扩充为 Hamel 基]\label{thm:y0-hamel-basis}
假设选择公理，设 $V$ 是域 $\mathbb K$ 上的线性空间，$L\subseteq V$ 为线性无关集。则存在包含 $L$ 的 Hamel 基 $B$；也就是说，每个向量都唯一表示为 $B$ 中有限个向量的线性组合。特别地，每个线性空间都有 Hamel 基。
\end{theorem}
\begin{proof}
令 $P$ 为 $V$ 中全部包含 $L$ 的线性无关子集，以包含关系为偏序。$L\in P$，故 $P$ 非空。对非空链 $\mathcal C\subseteq P$，令 $M=\bigcup_{A\in\mathcal C}A$。任何有限个 $M$ 中的向量都共同处于某个链成员中：对有限个包含它们的成员逐次取较大者即可。因此它们的线性关系只能是平凡关系，$M$ 线性无关，并包含 $L$，故 $M\in P$。空链以 $L$ 为上界。

由 Zorn 引理取极大元素 $B$。若 $v\notin\operatorname{span}B$，则 $B\cup\{v\}$ 仍线性无关：任何涉及 $v$ 的有限线性关系，若 $v$ 的系数非零便把 $v$ 表示为 $B$ 的有限线性组合；若系数为零则由 $B$ 的线性无关性知其余系数全零。这样得到严格更大的成员，与极大性矛盾。因此 $\operatorname{span}B=V$。线性无关性保证表达唯一。取 $L=\varnothing$ 即得基的存在性；零空间的基为空集。
\end{proof}

\begin{remark}[代数基与拓扑基]
Hamel 基只涉及有限线性组合，不涉及收敛。Banach 空间中的 Schauder 基要求向量有收敛的级数展开；Hilbert 空间中的完全正交系也涉及内积与闭线性张成。Hamel 基存在定理不保证任意 Banach 空间具有 Schauder 基。
\end{remark}

对后续的 Hahn--Banach 延拓，候选对象则是定义在不同子空间上的线性泛函，以“定义域扩大且取值相容”排序。链的并给出一致的线性泛函，极大性再配合单步延拓将定义域推到整个空间。因此使用 Zorn 引理时，应依次明确：候选对象形成集合，偏序及非空性成立，每条链有合法上界，极大对象具有目标性质。

\subsection{补充习题与解答}
\begin{exercise}[极大元可以不唯一]
设 $X$ 为至少有两个元素的集合，$P=\mathcal P(X)\setminus\{X\}$，以包含关系排序。求 $P$ 的全部极大元，并判断 $P$ 是否有最大元。
\end{exercise}
\begin{solution}
极大元恰为 $X\setminus\{x\}$（$x\in X$）。这种集合再加入任何缺少的元素就变成 $X$，故极大。若 $A\in P$ 缺少两个不同元素 $x,y$，则 $A\cup\{x\}$ 仍是真子集，故 $A$ 不极大。若有最大元，它必须包含每个单元素子集，从而等于 $X$，与属于 $P$ 矛盾。
\end{solution}

\begin{exercise}[有界的数列并不足够]
在 Zorn 引理中，能否将“每条链都有上界”改为“每个可数链都有上界”？提示：考虑全体可数序数 $\omega_1$ 在通常序数次序下的偏序，并在选择公理下讨论。
\end{exercise}
\begin{solution}
不能。在选择公理下，可数个可数集合的并可数，所以任意可数子集 $C\subseteq\omega_1$ 的序数上确界 $\bigcup C$ 仍是可数序数，属于 $\omega_1$，并是其上界；空集的并为 $0$。但任意 $\alpha<\omega_1$ 的后继 $\alpha+1$ 仍在 $\omega_1$ 中，故没有极大元。这里可数并的结论使用了选择公理所蕴含的可数选择，不能在讨论 ZF 中的等价性时不加说明地借用。
\end{solution}

\begin{exercise}[已指定的局部选择可以延拓]
假设 Zorn 引理。对非空集合族 $(A_i)_{i\in I}$，若已给定 $J_0\subseteq I$ 上的部分选择函数 $f_0$，证明它能延拓为整个 $I$ 上的选择函数。
\end{exercise}
\begin{solution}
将定理 \ref{thm:y0-choice-equivalence} 第三步中的偏序集换为所有延拓 $f_0$ 的部分选择函数。它包含 $f_0$，因而非空；非空链的并仍延拓 $f_0$，空链以 $f_0$ 作上界。由 Zorn 引理取极大延拓，若有未定义指标，选其中一个指标及其集合中的一个元素即可继续延拓，矛盾。所以定义域为 $I$。
\end{solution}


\Needspace{8\baselineskip}
\section{拓扑空间}\label{sec:yosida-00-02}
\noindent\textit{原书位置：中译本印刷第 3--14 页；英文版第 3--15 页。}
% !TeX encoding = UTF-8
本节以开集组织邻近、连续与紧性，随后证明函数逼近和类别论中常用的基础定理。除特别说明外，拓扑空间均假设为 Hausdorff 空间；个别结论会明确指出不需要这一假设。紧性始终指开覆盖意义的紧性。

\subsection{开集、闭集、邻域与连续映射}
\begin{definition}[拓扑与子空间拓扑]\label{def:y0-topology}
集合 $X$ 上的拓扑是子集族 $\tau\subseteq\mathcal P(X)$，满足 $\varnothing,X\in\tau$，任意族成员的并属于 $\tau$，有限个成员的交属于 $\tau$。$\tau$ 的成员称为开集，补集为开集的集合称为闭集。子集 $A\subseteq X$ 的子空间拓扑为 $\{A\cap U:U\in\tau\}$。

若 $N$ 含有一个包含 $x$ 的开集，则称 $N$ 为 $x$ 的邻域。若任意两个不同点都有互不相交的开邻域，则空间称为 Hausdorff 空间。
\end{definition}

\begin{definition}[闭包、内部、边界与聚点]\label{def:y0-closure}
集合 $A$ 的闭包 $\overline A$ 是包含 $A$ 的全部闭集的交；内部 $A^\circ$ 是包含于 $A$ 的全部开集的并；边界为 $\partial A=\overline A\setminus A^\circ$。若 $\overline A=X$，称 $A$ 在 $X$ 中稠密。若每个 $x$ 的邻域都与 $A\setminus\{x\}$ 相交，称 $x$ 为 $A$ 的聚点。
\end{definition}

\begin{proposition}[闭包的邻域刻画]\label{prop:y0-closure}
$x\in\overline A$ 当且仅当 $x$ 的每个开邻域都与 $A$ 相交。因此 $A$ 闭当且仅当它包含自己的全部聚点，且
\[
 \overline A=X\setminus(X\setminus A)^\circ.
\]
\end{proposition}
\begin{proof}
若某开邻域 $U\ni x$ 与 $A$ 不交，则闭集 $X\setminus U$ 包含 $A$，而不含 $x$，所以 $x\notin\overline A$。反之，若 $x\notin\overline A$，则 $X\setminus\overline A$ 是所求开邻域。对 $x\notin A$，与 $A$ 相交等价于与 $A\setminus\{x\}$ 相交，故闭集的刻画随之成立。最后一个等式是邻域刻画的补集形式。
\end{proof}

\begin{definition}[拓扑基、邻域基与连续]\label{def:y0-continuity}
开集族 $\mathcal B$ 若使每个开集都是其成员的并，称为拓扑基。点 $x$ 的邻域族 $\mathcal B_x$ 若使每个 $x$ 的邻域都包含一个 $\mathcal B_x$ 的成员，称为 $x$ 的邻域基。每点有可数邻域基称为第一可数；有可数拓扑基称为第二可数。

映射 $f:X\to Y$ 在 $x$ 连续，指对每个 $f(x)$ 的邻域 $V$，存在 $x$ 的邻域 $U$ 使 $f(U)\subseteq V$；在每点连续称为连续。双射若其本身和逆映射都连续，称为同胚。
\end{definition}

\begin{proposition}[连续性的等价条件]\label{prop:y0-continuity}
映射 $f:X\to Y$ 连续当且仅当每个开集的原像为开集，也当且仅当每个闭集的原像为闭集。
\end{proposition}
\begin{proof}
若 $f$ 连续，对开集 $V\subseteq Y$ 及 $x\in f^{-1}(V)$，连续性给出含 $x$ 的开集 $U\subseteq f^{-1}(V)$，所以原像开。反之，给定邻域 $V\ni f(x)$，在其中取开邻域 $W$，则 $f^{-1}(W)$ 为 $x$ 的开邻域，且像包含于 $V$。闭集条件与开集条件由取补集互相转换。
\end{proof}

\begin{definition}[度量与收敛]\label{def:y0-metric}
函数 $d:X\times X\to[0,\infty)$ 若满足 $d(x,y)=0$ 当且仅当 $x=y$、$d(x,y)=d(y,x)$ 和三角不等式 $d(x,z)\le d(x,y)+d(y,z)$，称为度量。开球 $B(x,r)=\{y:d(x,y)<r\}$ 构成拓扑基，所定义的拓扑称为度量拓扑。

在拓扑空间中，$x_n\to x$ 指每个 $x$ 的邻域都包含数列的全部充分靠后项。在度量空间中，这等价于 $d(x_n,x)\to0$。
\end{definition}

度量拓扑是 Hausdorff 的：不同点之间距离为 $r>0$，半径 $r/3$ 的两球不交。一般 Hausdorff 空间中数列极限也唯一，否则两个不交的极限邻域都必须含有充分靠后的项。度量空间是第一可数的，因为 $B(x,1/n)$ 给出可数邻域基。

\begin{proposition}[第一可数空间中的数列判据]\label{prop:y0-sequential}
若 $X$ 第一可数，则 $x\in\overline A$ 当且仅当存在 $A$ 中的数列趋于 $x$；映射 $f:X\to Y$ 连续当且仅当它保持所有数列的收敛。
\end{proposition}
\begin{proof}
将 $x$ 的可数邻域基替换为有限交，使它递减，记为 $U_n$。若 $x\in\overline A$，取 $x_n\in A\cap U_n$，便有 $x_n\to x$。反向由闭包的邻域刻画成立。连续映射显然保持收敛；若 $f$ 在 $x$ 不连续，则有开邻域 $V\ni f(x)$ 使每个 $U_n$ 都含有点 $x_n$ 满足 $f(x_n)\notin V$，这给出反例数列。
\end{proof}

一般拓扑空间中数列未必刻画闭包或连续性。例如在不可数集的可数补拓扑中，每个收敛数列最终恒等于其极限，但所有不可数子集都稠密。此例不满足 Hausdorff 条件，仅用来说明不能省略数列判据的假设。

\subsection{紧性、局部紧性与单点紧化}
\begin{definition}[紧性与相对紧性]\label{def:y0-compact}
若子集 $K\subseteq X$ 的每个开覆盖都有有限子覆盖，称 $K$ 紧。若 $\overline A$ 紧，称 $A$ 相对紧。若每个点都具有紧邻域，称空间局部紧。

集合族具有有限交性质，指每个有限子族的交都非空，其中空族的交约定为整个底集。
\end{definition}

\begin{proposition}[紧性的基本性质]\label{prop:y0-compact}
连续像保持紧性；紧空间的闭子集紧；Hausdorff 空间的紧子集闭；紧性等价于每个具有有限交性质的闭集族的交非空。紧空间到 Hausdorff 空间的连续双射为同胚。
\end{proposition}
\begin{proof}
像空间的开覆盖取原像，再用有限子覆盖，即得第一条。覆盖闭子集时加上其补集，得到紧空间的覆盖，即得第二条。若 $x\notin K$ 且 $K$ 紧，对每个 $y\in K$ 取不交开邻域 $U_y\ni y$、$V_y\ni x$。有限个 $U_{y_j}$ 覆盖 $K$，则 $\bigcap_jV_{y_j}$ 与 $K$ 不交，故 $X\setminus K$ 开。有限交刻画由开覆盖取补集得到。最后，连续双射将闭集送到紧集，再送到闭集，故其逆映射连续。
\end{proof}

\begin{proposition}[紧 Hausdorff 空间的正规性]\label{prop:y0-normal}
紧 Hausdorff 空间是正规空间，即任意两个不交闭集都有互不相交的开邻域。它也是正则空间，即点与不含它的闭集可被不交开集分离。
\end{proposition}
\begin{proof}
设不交闭集 $A,B$，均紧。固定 $a\in A$，对每个 $b\in B$ 取不交开邻域 $U_{a,b}\ni a$、$V_{a,b}\ni b$。以有限个 $V_{a,b_j}$ 覆盖 $B$，令 $U_a=\bigcap_jU_{a,b_j}$、$V_a=\bigcup_jV_{a,b_j}$。则 $U_a\cap V_a=\varnothing$ 且 $B\subseteq V_a$。再以有限个 $U_{a_i}$ 覆盖 $A$，取 $U=\bigcup_iU_{a_i}$、$V=\bigcap_iV_{a_i}$，便有 $A\subseteq U$、$B\subseteq V$、$U\cap V=\varnothing$。单点在 Hausdorff 空间中闭，故正规性也给出正则性。
\end{proof}

\begin{theorem}[单点紧化]\label{thm:y0-one-point}
设 $X$ 局部紧且 Hausdorff，添加新点 $\infty$，令 $X^*=X\cup\{\infty\}$。规定 $X$ 原有开集仍开，含 $\infty$ 的集合 $W$ 开当且仅当 $X^*\setminus W$ 是 $X$ 的紧子集。则 $X^*$ 为紧 Hausdorff 空间，且诱导于 $X$ 的拓扑不变。若 $X$ 非紧，则 $X$ 在 $X^*$ 中稠密；若 $X$ 紧，则 $\infty$ 为孤立点。
\end{theorem}
\begin{proof}
紧子集闭；含 $\infty$ 的开集取有限交时，补集是有限个紧集的并，仍紧。取任意并时，若其中某成员含 $\infty$，并集的补集是该成员补集的闭子集，故紧。与原有开集取交也保持开性。这证明拓扑公理及子空间拓扑不变。

任意覆盖 $X^*$ 的开族中，取一个含 $\infty$ 的成员，其补集紧，可由有限个其余成员覆盖，故 $X^*$ 紧。$X$ 中两点沿用原来的分离。对 $x\in X$，取紧邻域 $K$ 及开集 $U$ 使 $x\in U\subseteq K$，则 $U$ 和 $X^*\setminus K$ 分离 $x$ 与 $\infty$。若 $X$ 非紧，$\infty$ 的任何邻域都与 $X$ 相交；若 $X$ 紧，$\{\infty\}$ 开。
\end{proof}

\subsection{乘积拓扑与 Tychonoff 定理}
\begin{definition}[乘积拓扑]\label{def:y0-product-topology}
对空间族 $(X_i)_{i\in I}$，乘积 $X=\prod_iX_i$ 的乘积拓扑以集合 $\bigcap_{i\in F}\pi_i^{-1}(U_i)$ 为基，其中 $F\subseteq I$ 有限，各 $U_i\subseteq X_i$ 开，$\pi_i$ 为坐标投影。它是使所有投影连续的最弱拓扑。若允许无限多个坐标同时受限，得到的箱拓扑通常更强。
\end{definition}

以下证明使用第 \ref{thm:y0-zorn} 定理，因而在选择公理下进行。为使任意多个坐标的极限相容，先引入滤子。

\begin{lemma}[超滤子与紧性]\label{lem:y0-ultrafilter}
非空集合 $X$ 上的滤子是子集族 $\mathcal F$，满足 $X\in\mathcal F$、$\varnothing\notin\mathcal F$，对有限交封闭，并且包含成员的所有超集。极大的滤子称为超滤子。每个滤子可扩充为超滤子；超滤子 $\mathcal U$ 对任意 $A\subseteq X$ 恰包含 $A$、$X\setminus A$ 之一。

称滤子收敛于 $x$，若它包含 $x$ 的所有邻域。空间紧当且仅当其上每个超滤子都收敛于至少一点；Hausdorff 空间中的滤子极限至多一个。
\end{lemma}
\begin{proof}
所有包含固定滤子的滤子按包含排序，链的并仍为滤子。Zorn 引理给出极大者。若 $A\notin\mathcal U$，将 $A$ 加入不能生成真滤子，故存在 $F\in\mathcal U$ 使 $F\cap A=\varnothing$；于是 $X\setminus A\in\mathcal U$。两者不能同时属于滤子。

若空间紧，闭集族 $\{\overline F:F\in\mathcal U\}$ 具有有限交性质，故取 $x$ 在其交中。若邻域 $V\ni x$ 不在 $\mathcal U$ 中，取开邻域 $W\subseteq V$，则 $X\setminus W\in\mathcal U$，与 $x\in\overline{X\setminus W}=X\setminus W$ 矛盾。因此超滤子收敛于 $x$。反之，具有有限交性质的闭集族生成滤子，将其扩充为超滤子并取极限 $x$。若某个原闭集 $F$ 不含 $x$，其补集为 $x$ 的邻域，与 $F$ 同时在滤子中，矛盾，所以 $x$ 属于全部原闭集。Hausdorff 情形的唯一性由不交邻域不能同时属于滤子得到。
\end{proof}

\begin{theorem}[Tychonoff 定理]\label{thm:y0-tychonoff}
任意族紧 Hausdorff 空间的乘积，在乘积拓扑下仍为紧 Hausdorff 空间。
\end{theorem}
\begin{proof}
若乘积为空，它紧。否则取乘积上的超滤子 $\mathcal U$。对每个 $i$，族 $\mathcal U_i=\{A\subseteq X_i:\pi_i^{-1}(A)\in\mathcal U\}$ 为 $X_i$ 上的超滤子，由紧性它收敛于 $x_i$。令 $x=(x_i)_{i\in I}$。$x$ 的每个基本邻域是有限个 $\pi_i^{-1}(U_i)$ 的交，各项都在 $\mathcal U$ 中，故交也在其中。这说明 $\mathcal U\to x$，由引理得乘积紧。不同乘积点在某一坐标不同，该坐标中的不交开邻域的原像分离这两点，故乘积 Hausdorff。
\end{proof}

\subsection{Urysohn 引理与连续截断函数}
\begin{theorem}[Urysohn 引理]\label{thm:y0-urysohn}
设 $X$ 为正规 Hausdorff 空间，$A,B\subseteq X$ 为不交闭集。则存在连续函数 $f:X\to[0,1]$，使 $f|_A=0$、$f|_B=1$。
\end{theorem}
\begin{proof}
正规性给出缩小邻域的性质：若闭集 $F\subseteq O$ 且 $O$ 开，则存在开集 $V$ 使 $F\subseteq V\subseteq\overline V\subseteq O$。确实，将 $F$ 与 $X\setminus O$ 用不交开集分离即可。

令 $D$ 为 $[0,1]$ 中的二进有理数。取开集 $U_1=X\setminus B$，再取 $U_0$ 使 $A\subseteq U_0\subseteq\overline U_0\subseteq U_1$。逐层在相邻已定义指标 $r<s$ 的中点插入开集，使
\[
 \overline U_r\subseteq U_{(r+s)/2}
 \subseteq\overline U_{(r+s)/2}\subseteq U_s.
\]
从而对任意 $r<s$ 都有 $\overline U_r\subseteq U_s$。定义
\[
 f(x)=\inf\bigl(\{r\in D:x\in U_r\}\cup\{1\}\bigr).
\]
显然 $f=0$ 于 $A$，$f=1$ 于 $B$。对 $0<a\le1$，有 $\{f<a\}=\bigcup_{r<a}U_r$。对 $0\le a<1$，有
\[
 \{f>a\}=\bigcup_{\substack{r\in D\\r>a}}
                     (X\setminus\overline U_r).
\]
后一等式中，若 $x\notin\overline U_r$，则 $x$ 不属于任何 $U_s$（$s\le r$），所以 $f(x)\ge r>a$；若 $f(x)>a$，在 $a$ 与 $f(x)$ 之间取 $r<s$，则 $x\notin U_s$，从而 $x\notin\overline U_r$。这些原像均开，证明连续性。
\end{proof}

\begin{corollary}[局部紧空间中的截断函数]\label{cor:y0-cutoff}
若 $X$ 局部紧且 Hausdorff，$K\subseteq O\subseteq X$，其中 $K$ 紧、$O$ 开，则有 $\varphi\in C_c(X)$ 满足 $0\le\varphi\le1$、$\varphi|_K=1$、$\operatorname{supp}\varphi\subseteq O$。这里 $C_c(X)$ 表示连续且具有紧支集的函数。
\end{corollary}
\begin{proof}
在紧 Hausdorff 空间 $X^*$ 中，$K$ 闭而 $O$ 开，由正规性取开集 $V$ 使 $K\subseteq V\subseteq\overline V\subseteq O$。对不交闭集 $K$、$X^*\setminus V$ 应用 Urysohn 引理，得到在 $K$ 上为 $1$、在后者上为 $0$ 的函数。限制到 $X$ 后，其支集包含于紧集 $\overline V\subseteq O$，即得结论。
\end{proof}

\subsection{Weierstrass 与 Stone--Weierstrass 逼近}
\begin{theorem}[Bernstein 多项式逼近]\label{thm:y0-bernstein}
对 $f\in C([0,1];\mathbb K)$，其中 $\mathbb K=\R$ 或 $\mathbb C$，多项式
\[
 B_nf(x)=\sum_{k=0}^nf(k/n)\binom nk x^k(1-x)^{n-k}
\]
一致收敛于 $f$。因此任意闭有限区间上的连续函数均可由多项式一致逼近。
\end{theorem}
\begin{proof}
记 $p_{n,k}(x)=\binom nkx^k(1-x)^{n-k}$。二项式公式及其一、二阶导数给出
\[
 \sum_kp_{n,k}=1,\qquad
 \sum_k(k/n-x)^2p_{n,k}=\frac{x(1-x)}n\le\frac1{4n}.
\]
设 $M=\|f\|_\infty$。由紧区间上的一致连续性，给定 $\varepsilon>0$ 可取 $\delta>0$ 使 $|s-t|<\delta$ 时 $|f(s)-f(t)|<\varepsilon$。将求和按 $|k/n-x|<\delta$ 与其补集分开，得
\[
 |B_nf(x)-f(x)|
 \le\varepsilon+2M\delta^{-2}\sum_k(k/n-x)^2p_{n,k}
 \le\varepsilon+\frac{M}{2n\delta^2}.
\]
估计对 $x$ 一致，故结论成立。一般区间由仿射换元得到。这里的一致连续性可由紧性证明：在每点取控制函数振荡的小邻域，有限子覆盖再给出统一的距离界限。
\end{proof}

\begin{theorem}[实 Stone--Weierstrass 定理]\label{thm:y0-stone-real}
设 $X$ 紧 Hausdorff，$\mathcal A\subseteq C(X;\R)$ 为含常数函数且分离点的子代数；分离点指对 $x\ne y$ 存在 $a\in\mathcal A$ 使 $a(x)\ne a(y)$。则 $\mathcal A$ 在一致范数 $\|f\|_\infty=\sup_X|f|$ 下稠密于 $C(X;\R)$。
\end{theorem}
\begin{proof}
设 $\mathcal B$ 为 $\mathcal A$ 的一致闭包，它仍为含常数的代数。对 $g\in\mathcal B$，在包含其值域的有限区间上以多项式逼近绝对值函数，得 $|g|\in\mathcal B$。因此
\[
 \max(g,h)=\tfrac12(g+h+|g-h|),\qquad
 \min(g,h)=\tfrac12(g+h-|g-h|)
\]
也在 $\mathcal B$ 中。

给定 $f\in C(X;\R)$ 与 $\varepsilon>0$。对任意 $x,y\in X$，可找 $a_{x,y}\in\mathcal A$ 满足 $a_{x,y}(x)=f(x)$、$a_{x,y}(y)=f(y)$：若 $x\ne y$，对一个分离它们的函数作仿射变换；若 $x=y$，取常数函数。固定 $x$，开集 $\{t:a_{x,y}(t)>f(t)-\varepsilon\}$ 随 $y$ 覆盖 $X$。取有限子覆盖并令 $g_x$ 为对应函数的最大值，则 $g_x\in\mathcal B$、$g_x>f-\varepsilon$ 于 $X$，且 $g_x(x)=f(x)$。

开集 $\{t:g_x(t)<f(t)+\varepsilon\}$ 随 $x$ 覆盖 $X$。再取有限子覆盖，令 $g$ 为对应 $g_x$ 的最小值，便有 $f-\varepsilon<g<f+\varepsilon$ 于 $X$，且 $g\in\mathcal B$。由闭性 $f\in\mathcal B$，完成证明。
\end{proof}

\begin{corollary}[格形式与复数形式]\label{cor:y0-stone-complex}
实数情形中，可将“子代数”替换为“对有限最大、最小运算封闭的线性子空间”，仍要求含常数及分离点。复数情形中，含常数、分离点且对复共轭封闭的子代数 $\mathcal A\subseteq C(X;\mathbb C)$ 在一致范数下稠密。
\end{corollary}
\begin{proof}
格形式直接使用上一证明的两次有限覆盖构造。复数情形中，实部和虚部都属于 $\mathcal A$，故其中的实值函数构成含常数的实代数，并分离点：一个复值函数在两点不同，其实部或虚部必不同。对任意复值连续函数分别逼近实部和虚部，再组合即可。
\end{proof}

\begin{example}[三角多项式的一致稠密性]
在单位圆上，函数 $z\mapsto z^k$（$k\in\mathbb Z$）的有限复线性组合构成含常数、分离点且对共轭封闭的代数。因此每个连续 $2\pi$ 周期函数都可由 $\sum_{k=-N}^Nc_ke^{ikt}$ 一致逼近。这里只断言存在逼近的三角多项式，不断言原函数的 Fourier 部分和必一致收敛。
\end{example}

\subsection{完备性与 Baire 类别定理}
\begin{definition}[Cauchy 数列、完备性与类别]\label{def:y0-baire}
度量空间中，数列 $(x_n)$ 为 Cauchy 数列，指对任意 $\varepsilon>0$，存在 $N$ 使 $m,n\ge N$ 时 $d(x_m,x_n)<\varepsilon$。每个 Cauchy 数列都在空间内收敛，称空间完备。

若 $(\overline A)^\circ=\varnothing$，称 $A$ 无处稠密；若 $A$ 是可数个无处稠密集的并，称 $A$ 为第一纲集或贫集；否则称为第二纲集。若可数个稠密开集的交仍稠密，称空间为 Baire 空间。
\end{definition}

\begin{theorem}[Baire 类别定理]\label{thm:y0-baire}
每个完备度量空间和每个局部紧 Hausdorff 空间都是 Baire 空间。特别地，非空完备度量空间或非空紧 Hausdorff 空间不能是自身中的第一纲集。
\end{theorem}
\begin{proof}
设 $G_n$ 为稠密开集，$O$ 为非空开集。在完备度量空间中，递归取非空闭球 $K_n=\overline B(x_n,r_n)$，使 $K_1\subseteq O\cap G_1$、$K_{n+1}\subseteq B(x_n,r_n)\cap G_{n+1}$，并使 $0<r_n<2^{-n}$。稠密性保证每一步可进行。中心为 Cauchy 数列，由完备性收敛于 $x$；闭球嵌套给出 $x\in K_n$ 对所有 $n$ 成立，故 $x\in O\cap\bigcap_nG_n$。

在局部紧 Hausdorff 空间中，每个非空开集包含一个非空开集 $V$，使 $\overline V$ 紧且包含于原开集；这由单点紧化空间的正则性得到。于是递归取 $V_n$ 使 $\overline V_1\subseteq O\cap G_1$、$\overline V_{n+1}\subseteq V_n\cap G_{n+1}$，且各闭包紧。闭集 $\overline V_n$ 在 $\overline V_1$ 中具有有限交性质，交非空，交中的点即属于 $O\cap\bigcap_nG_n$。最后，对无处稠密集的闭包取补集，便得到第一纲表述。
\end{proof}

\begin{proposition}[连续函数逐点极限的不连续点]\label{prop:y0-baire-one}
设 $f_n:X\to\R$ 连续，且在每点收敛于有限值 $f$。则 $f$ 的不连续点集为第一纲集。若 $X$ 为 Baire 空间，$f$ 的连续点集稠密。
\end{proposition}
\begin{proof}
固定 $\varepsilon>0$，令
\[
 F_N=\bigcap_{m,n\ge N}\{x:|f_m(x)-f_n(x)|\le\varepsilon/4\}.
\]
各 $F_N$ 闭，且由逐点 Cauchy 性 $X=\bigcup_NF_N$。在 $F_N$ 上取极限得 $|f-f_N|\le\varepsilon/4$。若 $x\in F_N^\circ$，由 $f_N$ 的连续性，可取 $x$ 的开邻域 $U\subseteq F_N^\circ$，使 $f_N$ 在 $U$ 上的振荡小于 $\varepsilon/2$；于是 $f$ 在 $U$ 上的振荡小于 $\varepsilon$。

令 $G_\varepsilon$ 为所有使 $f$ 的振荡小于 $\varepsilon$ 的开集的并。则 $G_\varepsilon$ 开，且
\[
 X\setminus G_\varepsilon\subseteq
          \bigcup_N(F_N\setminus F_N^\circ).
\]
右侧为第一纲集，因为闭集 $F_N$ 的边界无处稠密。$f$ 在 $x$ 连续当且仅当 $x\in\bigcap_{k\ge1}G_{1/k}$，故不连续点集第一纲。在 Baire 空间中，第一纲集不能包含非空开集，因而其补集稠密。
\end{proof}

\subsection{度量空间中的紧性判据}
\begin{definition}[全有界]\label{def:y0-total-bounded}
度量空间 $M$ 若对每个 $\varepsilon>0$ 都能用有限个半径 $\varepsilon$ 的球覆盖，称为全有界。对子集使用限制的度量。全有界蕴含有界，但无限维空间中有界未必全有界。
\end{definition}

\begin{theorem}[度量紧性的等价条件]\label{thm:y0-metric-compact}
度量空间 $M$ 的下列条件等价：紧；每个数列有收敛子列且极限在 $M$ 中；完备且全有界。因此完备度量空间 $X$ 中的子集 $A$ 相对紧，当且仅当 $A$ 全有界。
\end{theorem}
\begin{proof}
紧空间全有界，因为固定半径的全部开球构成覆盖。紧度量空间也完备：对 Cauchy 数列，闭集 $\overline{\{x_n:n\ge N\}}$ 嵌套且非空，紧性给出公共点 $x$，Cauchy 性随即给出 $x_n\to x$。

完备且全有界时，依次用半径 $2^{-k}$ 的有限球覆盖，从当前无限多个数列项中保留落在同一球内的无限子列。取对角子列，得到 Cauchy 子列，再由完备性得到收敛。反之，数列紧性蕴含全有界：若不全有界，递归选取两两距离至少某个 $\varepsilon>0$ 的点，便没有收敛子列。

最后由数列紧性推出开覆盖紧性。给定开覆盖，必有 $\delta>0$ 使每个球 $B(x,\delta)$ 都包含于某一覆盖成员。否则可取 $x_n$ 使 $B(x_n,1/n)$ 不包含于任何成员；取收敛子列 $x_{n_j}\to x$，再取覆盖成员 $U\ni x$ 和 $r>0$ 使 $B(x,r)\subseteq U$，则充分大时 $B(x_{n_j},1/n_j)\subseteq U$，矛盾。用全有界性取有限个半径 $\delta/2$ 的球，再为其中心各选一个包含半径 $\delta$ 球的覆盖成员，即得有限子覆盖。

若 $X$ 完备，$\overline A$ 为完备子空间；$A$ 全有界时，放大覆盖球半径即知其闭包也全有界，故闭包紧。反向由紧集全有界及其子集仍全有界成立。
\end{proof}

\begin{example}[有界不足以推出紧]
在 $\ell^2$ 中，标准单位向量 $e_n$ 均在闭单位球内，但 $\|e_n-e_m\|_2=\sqrt2$（$n\ne m$），所以闭单位球不全有界，也不紧。这与 $\R^n$ 中“闭且有界等价于紧”形成区别；后者由有限网格覆盖给出全有界性，再结合完备性证明。
\end{example}

\begin{exercise}[紧空间上的连续函数空间]
设 $K$ 紧 Hausdorff。证明 $C(K;\mathbb K)$ 在一致范数下完备。
\end{exercise}
\begin{solution}
一致 Cauchy 数列 $(f_n)$ 在每点为标量 Cauchy 数列，令 $f(x)=\lim_nf_n(x)$。将 $\|f_n-f_m\|_\infty<\varepsilon$ 中的 $m$ 取极限，得 $|f_n(x)-f(x)|\le\varepsilon$ 对所有 $x$ 成立，故一致收敛。固定 $x$，先使 $\|f-f_n\|_\infty$ 小，再以 $f_n$ 的连续性控制邻域中的变化，由三角不等式知 $f$ 连续。连续函数在紧空间上的像紧，因而有界，故 $f\in C(K;\mathbb K)$。
\end{solution}

\begin{exercise}[复 Stone--Weierstrass 中的共轭条件]
在单位圆 $\mathbb T$ 上，令 $\mathcal A$ 为仅含非负幂 $z^k$ 的多项式的限制。证明它含常数且分离点，却不能一致逼近 $\overline z$。
\end{exercise}
\begin{solution}
函数 $z$ 本身分离点。对任意 $p\in\mathcal A$，有 $\int_0^{2\pi}p(e^{it})e^{it}\,\dd t=0$，因为每个指数的频率均为正整数。若 $p_n\to\overline z$ 一致，则左侧趋于 $\int_0^{2\pi}e^{-it}e^{it}\,\dd t=2\pi$，矛盾。因此共轭封闭条件不能直接删去。
\end{solution}


\Needspace{8\baselineskip}
\section{测度空间}\label{sec:yosida-00-03}
\noindent\textit{原书位置：中译本印刷第 14--20 页；英文版第 15--20 页。}
% !TeX encoding = UTF-8
本节采用现代的 $\sigma$-代数记号，先从简单函数构造积分，再证明极限交换与乘积积分定理。原书将 $\sigma$-有限性列入主要工作设定；本稿将它单独标出，使读者能看清各结论实际需要什么条件。除涉及乘积测度等处另加说明外，一般测度空间不预设拓扑。

\subsection{可测集、测度与几乎处处}
\begin{definition}[$\sigma$-代数与测度空间]\label{def:y0-measure}
集合 $S$ 上的 $\sigma$-代数 $\Sigma$ 是包含 $S$、对补集和可数并封闭的子集族；其成员称为可测集。映射 $\mu:\Sigma\to[0,\infty]$ 若满足 $\mu(\varnothing)=0$，且对两两不交的 $E_n\in\Sigma$ 有
\[
 \mu\left(\bigcup_{n\ge1}E_n\right)=\sum_{n\ge1}\mu(E_n),
\]
称为测度，$(S,\Sigma,\mu)$ 称为测度空间。若 $\mu(S)<\infty$，称测度有限；若 $S=\bigcup_nS_n$ 且每个 $\mu(S_n)<\infty$，称测度 $\sigma$-有限。

若性质在某个可测零测集之外成立，称它几乎处处成立，简记为 a.e.。若每个可测零测集的所有子集均可测，称测度空间完备。
\end{definition}

原书称含 $S$ 且对补集、可数并封闭的族为“$\sigma$-ring”；按现代通常术语，这是 $\sigma$-代数。现代 $\sigma$-环只要求对差和可数并封闭，不一定包含 $S$。这个区别在非 $\sigma$-紧空间上的紧集生成测度中尤其需要注意。

\begin{proposition}[测度的基本极限性质]\label{prop:y0-measure-continuity}
测度单调且可数次可加。若 $E_n\uparrow E$，则 $\mu(E_n)\uparrow\mu(E)$；若 $E_n\downarrow E$ 且 $\mu(E_1)<\infty$，则 $\mu(E_n)\downarrow\mu(E)$。
\end{proposition}
\begin{proof}
将 $B$ 写为 $A$ 与 $B\setminus A$ 的不交并得单调性。一般可数并可按首次出现的指标分成不交差集，得次可加性。对递增族，将并写为 $E_1$ 与 $E_n\setminus E_{n-1}$ 的不交并，级数部分和即 $\mu(E_n)$。递减情形对 $E_1\setminus E_n$ 使用递增结论，再从有限数 $\mu(E_1)$ 中相减。有限性不能省略，例如计数测度下 $\{n,n+1,\ldots\}\downarrow\varnothing$，每项测度却为无穷。
\end{proof}

\begin{proposition}[测度的完备化]\label{prop:y0-completion-measure}
令 $\widehat\Sigma$ 为所有形如 $E\cup N$ 的集合，其中 $E\in\Sigma$，$N$ 包含于某个 $\Sigma$ 中的零测集。则 $\widehat\Sigma$ 为 $\sigma$-代数，且 $\widehat\mu(E\cup N)=\mu(E)$ 给出唯一延拓 $\mu$ 的完备测度。
\end{proposition}
\begin{proof}
每个新集合 $A$ 都夹在两个原可测集 $E\subseteq A\subseteq F$ 之间，且 $\mu(F\setminus E)=0$。取补集或可数并仍可找到这样的上下夹集，证明 $\widehat\Sigma$ 为 $\sigma$-代数。不同夹集的测度相同，因为它们只相差零测集，所以定义良定。对不交的新集合 $A_n$，选择可测下夹集 $E_n\subseteq A_n$，它们也不交，且 $\bigcup_nA_n$ 与 $\bigcup_nE_n$ 只相差零测集。可数可加性因此继承。完备性及唯一性由零测集的任何子集必须有测度零得到。
\end{proof}

\subsection{可测函数与 Egorov 定理}
\begin{definition}[可测函数与简单函数]\label{def:y0-measurable-function}
函数 $f:S\to[-\infty,\infty]$ 称为可测，若 $\{s:f(s)>a\}\in\Sigma$ 对每个 $a\in\R$ 成立。有限复值函数可测，指其实部与虚部可测，等价于每个复平面开集的原像可测。取有限个值的可测函数称为简单函数，可写作 $\sum_{j=1}^ma_j\mathbf1_{E_j}$，其中 $E_j$ 两两不交且可测。
\end{definition}

可测函数的有限和、乘积、绝对值仍可测；可数上确界、下确界、上极限及下极限也可测。例如 $\{\sup_nf_n>a\}=\bigcup_n\{f_n>a\}$，下确界与极限由类似可数运算表达。每个非负可测函数都有递增的非负简单函数逼近；可取
\[
 s_n(s)=2^{-n}\left\lfloor2^n\min(f(s),n)\right\rfloor,
 \qquad s_n\uparrow f,
\]
其中在 $f=\infty$ 处规定截断值为 $n$。每个取值区间的原像可测，故 $s_n$ 可测。

\begin{theorem}[Egorov 定理]\label{thm:y0-egorov}
设 $E\in\Sigma$、$\mu(E)<\infty$，$f_n$、$f$ 为可测函数，在 $E$ 上几乎处处有限，且 $f_n\to f$ 几乎处处。则对每个 $\eta>0$，存在可测集 $F\subseteq E$，使 $\mu(E\setminus F)<\eta$，且 $f_n\to f$ 在 $F$ 上一致。
\end{theorem}
\begin{proof}
先从 $E$ 中去掉一个可测零测集，使全部函数有限并逐点收敛。对正整数 $k,N$，令
\[
 A_{k,N}=\bigcup_{n\ge N}\{s\in E:|f_n(s)-f(s)|>1/k\}.
\]
固定 $k$ 时，$A_{k,N}\downarrow\varnothing$；由有限测度下的上连续性，可取 $N_k$ 使 $\mu(A_{k,N_k})<\eta2^{-k}$。令 $F=E\setminus\bigcup_kA_{k,N_k}$，则损失测度小于 $\eta$，且 $n\ge N_k$ 时 $\sup_F|f_n-f|\le1/k$。这正是一致收敛。
\end{proof}

\begin{example}[有限测度条件的作用]
在 $\R$ 的 Lebesgue 测度下，$f_n=\mathbf1_{[n,\infty)}$ 逐点趋于零；但从 $\R$ 去掉任意有限测度集后，剩余集合仍与每个 $[n,\infty)$ 相交，故误差上确界始终为 $1$。不能将 Egorov 定理中的有限测度条件直接换成 $\sigma$-有限。
\end{example}

\subsection{积分的构造与单调收敛}
\begin{definition}[Lebesgue 积分]\label{def:y0-integral}
对非负简单函数 $s=\sum_ja_j\mathbf1_{E_j}$，定义 $\int s\,\dd\mu=\sum_ja_j\mu(E_j)$，约定 $0\cdot\infty=0$。对非负可测函数 $f$，定义
\[
 \int f\,\dd\mu=\sup\left\{\int s\,\dd\mu:
                      0\le s\le f,\ s\text{ 为简单函数}\right\}.
\]
实值函数写为 $f=f^+-f^-$，其中 $f^+=\max(f,0)$、$f^-=\max(-f,0)$。若 $\int|f|\,\dd\mu<\infty$，称 $f$ 可积，并定义 $\int f=\int f^+-\int f^-$。复值函数可积也指 $\int|f|<\infty$，积分分别取实部和虚部的积分。局部可积、$L^p$ 等概念将在后续章节使用。
\end{definition}

对简单函数，不同表示可细分到同一不交分割，因此积分良定、单调、可加且正齐次。对于一般函数，先证明下一定理，再从简单函数逼近继承这些性质。

\begin{theorem}[单调收敛定理]\label{thm:y0-monotone}
若 $0\le f_n\uparrow f$ 为可测函数，则 $\int f_n\,\dd\mu\uparrow\int f\,\dd\mu$，允许积分无穷。
\end{theorem}
\begin{proof}
由单调性，左侧极限 $L\le\int f$。取简单函数 $0\le s\le f$ 和 $0<c<1$。令 $E_n=\{f_n\ge cs\}$。这些集合递增，其并包含 $\{s>0\}$：若 $s(x)>0$，则 $f_n(x)\uparrow f(x)\ge s(x)>cs(x)$。所以由测度的下连续性和简单函数公式，$\int s\mathbf1_{E_n}\to\int s$。于是
\[
 \int f_n\ge c\int s\mathbf1_{E_n},\qquad
 L\ge c\int s.
\]
令 $c\uparrow1$，再对 $s$ 取上确界，得 $L\ge\int f$。
\end{proof}

非负函数的积分可加性现在可由各自的递增简单逼近和单调收敛得出；有限实值及复值可积函数的线性性由正负部或实虚部的分解得出。特别地，可积性在加法下保持，且
\[
 \left|\int f\,\dd\mu\right|\le\int|f|\,\dd\mu.
\]
复数情形可选模为 $1$ 的标量 $\alpha$，使 $\alpha\int f$ 为非负实数，再用 $\operatorname{Re}(\alpha f)\le|f|$ 得到估计。若 $f\ge0$，则 $\int f=0$ 当且仅当 $f=0$ a.e.：因为 $\mu(\{f\ge1/n\})\le n\int f$，而 $\{f>0\}$ 是这些集合的并。

\begin{proposition}[与简单函数逼近定义的相容性]\label{prop:y0-simple-l1}
每个可积函数 $f$ 都有可积简单函数 $s_n$ 满足 $\int|s_n-f|\to0$。反之，若可测有限函数 $f$ 是可积简单函数 $s_n$ 的几乎处处极限，且 $(s_n)$ 在积分距离 $\int|s_n-s_m|$ 下 Cauchy，则 $f$ 可积，$\int|s_n-f|\to0$，并且 $\int s_n\to\int f$。
\end{proposition}
\begin{proof}
第一条先对 $f\ge0$ 使用递增简单逼近，由单调收敛及 $\int f<\infty$ 得 $\int(f-s_n)\to0$；再分别处理实虚部和正负部。反向将在下面 Fatou 引理得到后使用它：固定 $n$，
\[
 \int|s_n-f|\le\liminf_{m\to\infty}\int|s_n-s_m|.
\]
右侧在 $n$ 充分大时任意小，故 $f$ 与某个可积 $s_n$ 的差可积，从而 $f$ 可积，并得积分距离趋零。积分估计又给出积分值的收敛。这说明本稿的非负积分构造与原书按简单函数积分距离定义的方式相容。
\end{proof}

\subsection{Fatou 引理、控制收敛与积分的绝对连续性}
\begin{theorem}[Fatou 引理]\label{thm:y0-fatou}
若 $f_n\ge0$ 可测，则
\[
 \int\liminf_nf_n\,\dd\mu\le\liminf_n\int f_n\,\dd\mu.
\]
若实值可积函数 $f_n$ 满足 $f_n\ge h$ a.e.，其中 $h$ 可积，则同样有 $\int\liminf_nf_n\le\liminf_n\int f_n$，左侧按有可积下界的扩展积分解释。若有可积上界，则得到上极限的反向不等式。
\end{theorem}
\begin{proof}
令 $g_N=\inf_{n\ge N}f_n$，则 $0\le g_N\uparrow\liminf_nf_n$，且 $\int g_N\le\inf_{n\ge N}\int f_n$。单调收敛给出结论。有可积下界时对 $f_n-h$ 应用非负版本，再加回 $\int h$；有上界时对 $-f_n$ 应用下界版本。
\end{proof}

\begin{theorem}[控制收敛定理]\label{thm:y0-dominated}
设 $f_n$ 与 $f$ 均可测，$f_n\to f$ a.e.，且存在非负可积函数 $g$ 使 $|f_n|\le g$ a.e.，则 $f$ 可积，$\int|f_n-f|\to0$，并且 $\int f_n\to\int f$。
\end{theorem}
\begin{proof}
有 $|f|\le g$ a.e.，故 $f$ 可积。函数 $h_n=|f_n-f|$ 趋零且 $0\le h_n\le2g$。对非负函数 $2g-h_n$ 使用 Fatou 引理，得
\[
 2\int g\le\liminf_n\left(2\int g-\int h_n\right)
            =2\int g-\limsup_n\int h_n.
\]
所以 $\int h_n\to0$，再用积分估计即可。
\end{proof}
这里把极限函数的可测性单独列出：在未完备的测度空间中，对可测零集上的值作任意修改，可能破坏可测性。若只给出 $f_n$ 几乎处处收敛，则应选取其可测极限代表，或先完备化测度空间，再使用上述结论。

\begin{proposition}[积分的可数可加性与绝对连续性]\label{prop:y0-integral-ac}
若 $f$ 可积，$\nu_f(E)=\int_Ef\,\dd\mu$ 为可数可加的实值或复值集合函数，且 $\mu(E)=0$ 蕴含 $\nu_f(E)=0$。更强地，对任意 $\varepsilon>0$，存在 $\delta>0$，使 $\mu(E)<\delta$ 时 $\int_E|f|\,\dd\mu<\varepsilon$。
\end{proposition}
\begin{proof}
对不交 $E_n$，部分并的示性函数逐点趋于全并的示性函数，乘以 $f$ 后由 $|f|$ 控制，控制收敛给出可数可加性。选 $M\ge1$ 使 $\int_{\{|f|>M\}}|f|<\varepsilon/2$，这由控制收敛得到。于是
\[
 \int_E|f|\le M\mu(E)+\int_{\{|f|>M\}}|f|.
\]
取 $\delta=\varepsilon/(2M)$ 即得结论，零测集的结论也直接成立。
\end{proof}

\subsection{乘积测度与 Fubini--Tonelli 定理}
\begin{lemma}[$\pi$--$\lambda$ 原理]\label{lem:y0-pi-lambda}
集合族 $\mathcal P$ 若对有限交封闭，称为 $\pi$-系。集合族 $\mathcal D$ 若含全集、对补集及可数不交并封闭，称为 $\lambda$-系。若 $\mathcal P\subseteq\mathcal D$，则 $\sigma(\mathcal P)\subseteq\mathcal D$。
\end{lemma}
\begin{proof}
令 $\mathcal L$ 为包含 $\mathcal P$ 的最小 $\lambda$-系。$\lambda$-系对两个嵌套成员的差封闭。固定 $A\in\mathcal P$，集合族 $\{B:A\cap B\in\mathcal L\}$ 为 $\lambda$-系并包含 $\mathcal P$，所以包含 $\mathcal L$。再固定 $B\in\mathcal L$，同样论证 $\{A:A\cap B\in\mathcal L\}$ 包含 $\mathcal L$。于是 $\mathcal L$ 对有限交封闭，结合补集可得有限并封闭，再将可数并分为不交差集，得它对可数并封闭。因此 $\mathcal L$ 是 $\sigma$-代数，必包含 $\sigma(\mathcal P)$，结论成立。
\end{proof}

\begin{theorem}[乘积测度与 Tonelli 定理]\label{thm:y0-tonelli}
设 $(S,\Sigma,\mu)$、$(T,\mathcal T,\nu)$ 均 $\sigma$-有限。令 $\Sigma\otimes\mathcal T$ 为由矩形 $A\times B$ 生成的 $\sigma$-代数。存在唯一测度 $\mu\times\nu$ 满足
\[
 (\mu\times\nu)(A\times B)=\mu(A)\nu(B).
\]
对非负乘积可测函数 $f$，截面可测，两个截面积分函数可测，且
\[
 \int_{S\times T}f\,\dd(\mu\times\nu)
 =\int_S\left(\int_Tf(s,t)\,\dd\nu(t)\right)\dd\mu(s)
 =\int_T\left(\int_Sf(s,t)\,\dd\mu(s)\right)\dd\nu(t),
\]
允许值为无穷，乘积中的 $0\cdot\infty$ 约定为零。
\end{theorem}
\begin{proof}
先设两测度有限。对集合 $E$，记 $E_s=\{t:(s,t)\in E\}$、$E^t=\{s:(s,t)\in E\}$。全部截面可测的集合构成包含矩形的 $\sigma$-代数，故乘积可测集均有可测截面。

考虑满足以下性质的集合 $E$：$s\mapsto\nu(E_s)$、$t\mapsto\mu(E^t)$ 可测，且它们的积分相等。矩形具有这些性质。有限测度保证取补集时可将截面测度及其积分从有限总量中相减；取可数不交并时由单调收敛得到性质。因此此族为包含矩形的 $\lambda$-系，由 $\pi$--$\lambda$ 原理，性质对所有乘积可测集成立。定义
\[
 \lambda(E)=\int_S\nu(E_s)\,\dd\mu(s).
\]
对不交集族，截面不交，测度可加性与单调收敛证明 $\lambda$ 为测度，并有所需矩形值。它也等于反向截面积分。

一般 $\sigma$-有限情形，将 $S$、$T$ 各分为可数个不交的有限测度集 $S_i$、$T_j$。在每个 $S_i\times T_j$ 上用上述有限测度构造，再对 $(i,j)$ 求和，得全空间测度。截面函数的可测性及两向积分相等由非负级数的单调收敛继承。对任意非负简单函数先按示性函数公式计算，再以递增简单函数逼近 $f$，三次使用单调收敛便得 Tonelli 公式。

唯一性在每个有限测度块上由 $\pi$--$\lambda$ 原理得到：两个有限测度相等的集合构成 $\lambda$-系并包含矩形。再对所有块求和，得全空间唯一性。
\end{proof}

\begin{theorem}[Fubini 定理]\label{thm:y0-fubini}
在上一设定下，若乘积可测函数 $f$ 满足 $\int|f|\,\dd(\mu\times\nu)<\infty$，则几乎所有截面可积，截面积分函数可积，且两个累次积分均等于 $\int f\,\dd(\mu\times\nu)$。绝对可积性等价于任意一个 $|f|$ 的累次积分有限。
\end{theorem}
\begin{proof}
对 $|f|$ 应用 Tonelli 定理，截面的绝对积分几乎处处有限，否则其积分不能有限。对实部与虚部的正、负部分分别使用 Tonelli，再作有限量的相减，得到交换次序的公式；由 $|\int f(s,t)\,\dd\nu(t)|\le\int|f(s,t)|\,\dd\nu(t)$ 得截面积分可积。最后的等价性正是 $|f|$ 的 Tonelli 公式。
\end{proof}

对完备化后的乘积测度，可用与原乘积可测函数几乎处处相同的代表元解释这些结论，截面结论仍是几乎处处的。不能只因两个带符号的累次积分各自存在，就跳过绝对可积性的检查。

\subsection{拓扑测度、Borel 集与 Baire 集}
\begin{definition}[Borel、Baire 与 Radon 测度]\label{def:y0-topological-measure}
拓扑空间 $X$ 的 Borel $\sigma$-代数 $\mathcal B(X)$ 由开集生成。现代 Baire $\sigma$-代数由所有连续实值函数的开集原像生成，通常包含于 $\mathcal B(X)$。在局部紧 Hausdorff 空间上，若 Borel 测度在紧集上有限，并满足
\[
 \mu(E)=\inf_{\substack{U\supseteq E\\U\text{ 开}}}\mu(U),\qquad
 \mu(E)=\sup_{\substack{K\subseteq E\\K\text{ 紧}}}\mu(K)
 \quad(E\in\mathcal B(X)),
\]
则本稿称其为 Radon 测度；这一定义明确要求外正则性与内正则性。
\end{definition}

原书的 Baire 集从紧 $G_\delta$ 集生成 $\sigma$-环，其中 $G_\delta$ 表示可数个开集的交；Borel 集也使用紧集生成的表述。这在一般非 $\sigma$-紧空间中与现代全空间 $\sigma$-代数表述须区别对待。在 $\R^n$ 及其闭子空间中，上述区别消失：这些空间可由可数个紧集覆盖，每个闭集为 $G_\delta$，而开集可写为可数个紧集的并，因此原书两种可测集族均等于通常的 Borel 集族。对连续函数的可测性，这些欧氏空间中也有 Baire 与 Borel 一致的结论。

例如，闭集 $F$ 可写为 $\bigcap_m\{x:d(x,F)<1/m\}$；开集 $U\subseteq\R^n$ 可写为紧集
\[
 K_m=\{x:|x|\le m,\ d(x,\R^n\setminus U)\ge1/m\}
\]
的并。更一般拓扑空间中的正则性和 Baire--Borel 延拓，须使用相应的测度表示定理，不能仅由可数可加性推出。

\subsection{Lebesgue 测度与两个逼近性质}
\begin{definition}[Lebesgue 外测度与可测集]\label{def:y0-lebesgue}
对 $E\subseteq\R^n$，令
\[
 m^*(E)=\inf\left\{\sum_j|Q_j|:
               E\subseteq\bigcup_jQ_j,\ Q_j\text{ 为有界开矩形}\right\},
\]
其中 $|Q_j|$ 是边长的乘积。若对每个 $A\subseteq\R^n$ 有
\[
 m^*(A)=m^*(A\cap E)+m^*(A\setminus E),
\]
称 $E$ 为 Lebesgue 可测集；$m^*$ 在这些集合上的限制记为 $m$。
\end{definition}

\begin{proposition}[外测度构造的基本性质]\label{prop:y0-lebesgue-construction}
上述可测集形成完备 $\sigma$-代数，$m$ 为 $\sigma$-有限测度，所有 Borel 集均可测，矩形的测度等于其体积。$m$ 平移不变，且对有限测度可测集 $E$ 有
\[
 m(E)=\inf_{U\supseteq E,\ U\text{ 开}}m(U)
     =\sup_{K\subseteq E,\ K\text{ 紧}}m(K).
\]
\end{proposition}
\begin{proof}
覆盖定义直接给出 $m^*(\varnothing)=0$、单调性及可数次可加性；后一性质通过对各集合选误差小于 $\varepsilon2^{-j}$ 的覆盖，再把覆盖合并得到。

分割等式对补集对称。若 $E,F$ 满足分割等式，对任意测试集先按 $E$、再按 $F$ 分割，得四项分解，从而得 $E\cup F$ 也满足等式。对不交可测集 $E_j$，重复分割给出
\[
 m^*(A)\ge\sum_{j=1}^Nm^*(A\cap E_j)
                   +m^*\left(A\setminus\bigcup_{j\ge1}E_j\right).
\]
令 $N\to\infty$，结合 $m^*(A\cap\bigcup_jE_j)\le\sum_jm^*(A\cap E_j)$，得全并满足分割等式；将一般并化为不交差集，证明 $\sigma$-代数性质。取 $A=\bigcup_jE_j$，同一估计与次可加性给出可数可加性。零外测度集及其子集自动满足分割等式，故测度完备。

坐标半空间可测：对任意矩形覆盖，沿对应坐标超平面切开各矩形，可将交叠处增宽任意小的宽度，使两边仍由开矩形覆盖且总体积增加任意小。于是 $m^*(A\cap H)+m^*(A\setminus H)\le m^*(A)$；反向由次可加性成立。因此矩形可测，开集为可数个有理矩形的并，故 Borel 集可测。

矩形的外测度不超过其体积，可通过稍增宽的开矩形覆盖得到。反向，对矩形内的任意紧闭矩形，用任意可数开矩形覆盖，紧性给出有限子覆盖；按有限个边界坐标细分，有限覆盖的体积和不小于被覆盖闭矩形的体积。令内矩形扩张至原矩形，得反向不等式。平移保持矩形体积并一一对应覆盖，故外测度及可测测度平移不变。可数个有界矩形覆盖全空间，故 $\sigma$-有限。

对有限测度 $E$，取体积总和小于 $m(E)+\varepsilon$ 的开矩形覆盖，并集 $U$ 开，且 $m(U\setminus E)<\varepsilon$，得外正则性。为得内正则性，先取足够大的闭立方体 $Q$，使 $m(E\setminus Q)<\varepsilon$；对 $Q\setminus E$ 外逼近为开集 $V$，使 $m(V\setminus(Q\setminus E))<\varepsilon$。则 $K=Q\setminus V$ 紧、$K\subseteq E$，且 $m(E\setminus K)<2\varepsilon$。
\end{proof}

在实线上，单调不减且右连续函数 $F$ 决定的区间值 $\mu_F((a,b])=F(b)-F(a)$ 称为 Lebesgue--Stieltjes 测度的起点；一般 $F$ 可用区间覆盖的外测度构造完成延拓。\emph{只有 $F(t)=t$ 时才是通常的 Lebesgue 测度}，例如阶跃函数决定点质量。这里不把一般 Stieltjes 测度统称为 Lebesgue 测度。

\begin{theorem}[紧支集连续函数的 $L^1$ 稠密性]\label{thm:y0-cc-density}
若 $G\subseteq\R^n$ 开，$f\in L^1(G)$，则对任意 $\varepsilon>0$，存在 $\varphi\in C_c(G)$ 使 $\int_G|f-\varphi|\,\dd x<\varepsilon$。
\end{theorem}
\begin{proof}
由命题 \ref{prop:y0-simple-l1}，先将 $f$ 逼近为可积简单函数，只需处理 $\mathbf1_E$，其中 $E\subseteq G$ 可测且有限测度。由正则性取紧集 $K\subseteq E$ 和开集 $U\supseteq E$，使 $m(E\setminus K)$、$m(U\setminus E)$ 都任意小。将 $U$ 换为 $U\cap G$，再在其中选含 $K$ 的相对紧开邻域 $V$。由推论 \ref{cor:y0-cutoff} 取 $0\le\varphi\le1$，在 $K$ 上为 $1$ 且支集包含于 $V$。于是
\[
 \int_G|\mathbf1_E-\varphi|
 \le m(E\setminus K)+m(U\setminus E).
\]
将有限个示性函数的逼近按系数组合，得到结论；有限个紧支集的并仍紧且包含于 $G$。
\end{proof}

\begin{theorem}[有限测度集的平移连续性]\label{thm:y0-translation-set}
若 $E\subseteq\R^n$ 可测且 $m(E)<\infty$，则 $m((E+h)\mathbin{\triangle}E)\to0$ 随 $h\to0$，其中 $\triangle$ 为对称差。
\end{theorem}
\begin{proof}
左侧等于 $\|\mathbf1_E(\cdot-h)-\mathbf1_E\|_{L^1}$。取 $\varphi\in C_c(\R^n)$ 使 $\|\mathbf1_E-\varphi\|_1<\varepsilon$。平移不变性给出
\[
 \|\mathbf1_E(\cdot-h)-\mathbf1_E\|_1
 \le2\varepsilon+\|\varphi(\cdot-h)-\varphi\|_1.
\]
当 $|h|\le1$ 时，后者的支集处于一个固定紧集内；紧支集连续函数在全空间一致连续，所以其平移差一致趋零，积分也趋零。最后令 $\varepsilon\downarrow0$。
\end{proof}

\begin{exercise}[逐点收敛与积分收敛]
在 $(0,1)$ 上令 $f_n=n\mathbf1_{(0,1/n)}$。证明它几乎处处趋于零，但积分不趋于零，并说明它不能由一个可积函数统一控制。
\end{exercise}
\begin{solution}
固定 $x>0$，充分大时 $1/n<x$，故 $f_n(x)=0$，而 $\int_0^1f_n=1$。若有可积 $g\ge|f_n|$ a.e.，控制收敛定理将给出积分趋零，矛盾。也可由积分的绝对连续性直接看出：$\int_{(0,1/n)}g\ge1$ 与这些集合测度趋零矛盾。
\end{solution}

\begin{exercise}[非绝对可积时累次积分可以不同]
在 $(0,1)^2$ 上取 $f(x,y)=(x^2-y^2)/(x^2+y^2)^2$。计算两个累次积分，说明为什么不违反 Fubini 定理。
\end{exercise}
\begin{solution}
由于 $\partial_y\bigl(y/(x^2+y^2)\bigr)=f$，对固定 $x>0$ 有 $\int_0^1f(x,y)\,\dd y=1/(1+x^2)$，再积分得 $\pi/4$。由于 $\partial_x\bigl(x/(x^2+y^2)\bigr)=-f$，反向累次积分为 $-\pi/4$。在第一象限极坐标下，$|f|=|\cos(2\theta)|/r^2$，面积元为 $r\,\dd r\,\dd\theta$，原点附近径向积分 $\int_0^\delta\dd r/r$ 发散，故 $f$ 不绝对可积，Fubini 的假设不成立。
\end{solution}


\Needspace{8\baselineskip}
\section{线性空间}\label{sec:yosida-00-04}
\noindent\textit{原书位置：中译本与英文版均为印刷第 20--22 页。}
% !TeX encoding = UTF-8
本节只建立代数结构，不给线性空间预先指定拓扑。后续的范数、半范数、闭子空间及连续算子，都是在这些代数对象上添加分析条件。除特别说明外，标量域 $\mathbb K$ 为 $\R$ 或 $\mathbb C$。

\subsection{线性空间、子空间与生成空间}
\begin{definition}[线性空间]\label{def:y0-vector-space}
域 $\mathbb K$ 上的线性空间是集合 $V$，配有加法与数乘，满足：加法使 $V$ 成为交换群，零元记为 $0$；且对 $x,y\in V$、$\alpha,\beta\in\mathbb K$ 有
\[
 \alpha(x+y)=\alpha x+\alpha y,\qquad
 (\alpha+\beta)x=\alpha x+\beta x,\qquad
 (\alpha\beta)x=\alpha(\beta x),\qquad 1x=x.
\]
其元素称为向量。
\end{definition}

由分配律与加法消去律可得 $0x=0$、$\alpha0=0$、$(-1)x=-x$；若 $\alpha x=0$ 且 $\alpha\ne0$，乘以 $\alpha^{-1}$ 得 $x=0$。标量零与向量零通常共用记号，语境决定其所属空间。

\begin{definition}[子空间与线性张成]\label{def:y0-subspace}
非空子集 $M\subseteq V$ 若对加法与数乘封闭，称为线性子空间。等价地，对任意 $x,y\in M$ 及标量 $\alpha,\beta$，有 $\alpha x+\beta y\in M$。

集合 $A\subseteq V$ 的线性张成 $\operatorname{span}A$ 是 $A$ 中有限个向量的全部线性组合组成的集合，约定 $\operatorname{span}\varnothing=\{0\}$。
\end{definition}

\begin{proposition}[张成与子空间的运算]\label{prop:y0-span}
$\operatorname{span}A$ 是包含 $A$ 的最小子空间。任意族子空间的交仍为子空间，空族的交为 $V$。若 $M,N$ 为子空间，则 $M+N=\{m+n:m\in M,n\in N\}$ 是包含 $M\cup N$ 的最小子空间。
\end{proposition}
\begin{proof}
有限线性组合的和或数乘仍为有限线性组合，故张成是子空间；任何包含 $A$ 的子空间必须包含这些组合，得最小性。交中的元素在各成员中都对线性组合封闭。对 $M+N$，线性组合可分别整理到 $M$ 和 $N$ 中，故它是子空间；包含两者的任一子空间都必须包含它们的和。
\end{proof}

子空间的并通常不是子空间。例如 $\R^2$ 中两条不同的过原点直线，其并不含从两直线各取一个非零向量所得的和。若一族子空间按包含关系成链，则其并为子空间，因为任意两个向量共同属于某个较大的成员。

\begin{example}[后续使用的线性空间]
$\mathbb K^n$、全体数列 $\mathbb K^{\N}$、任意集合上的函数空间 $\mathbb K^S$ 都以逐坐标运算成为线性空间。$C(X;\mathbb K)$、$C_c(X;\mathbb K)$ 是函数空间的子空间；$L^p$ 则把几乎处处相等的可测函数识别为同一个向量。后一识别是一个商构造，它保证范数为零的对象恰为零向量。
\end{example}

\subsection{线性无关、基与维数}
\begin{definition}[线性无关与 Hamel 基]\label{def:y0-linear-independence}
有限向量族 $v_1,\ldots,v_n$ 线性无关，指 $\sum_j\alpha_jv_j=0$ 蕴含全部 $\alpha_j=0$。任意子集 $B\subseteq V$ 线性无关，指其中每个有限的互异向量族都线性无关。若 $B$ 同时张成 $V$，称为 Hamel 基。

若 $V$ 有 $n$ 个向量组成的基，称 $V$ 为 $n$ 维，记 $\dim V=n$；若没有有限基，称为无限维。零空间的空基给出维数 $0$。
\end{definition}

\begin{lemma}[有限交换引理]\label{lem:y0-exchange}
若 $u_1,\ldots,u_m$ 线性无关，而 $v_1,\ldots,v_n$ 张成 $V$，则 $m\le n$；且可用 $u_1,\ldots,u_m$ 替换生成族中的 $m$ 个向量而仍张成 $V$。
\end{lemma}
\begin{proof}
把 $u_1$ 表示为 $v_j$ 的线性组合。因为 $u_1\ne0$，有一个系数非零，可解出对应的 $v_j$，从而用 $u_1$ 替换它，仍保持张成。若已将 $u_1,\ldots,u_k$ 换入，把 $u_{k+1}$ 表示为当前生成族的线性组合；由于 $u_{k+1}$ 不属于前 $k$ 个 $u$ 的张成，至少有一个尚未替换的 $v_j$ 的系数非零，再解出它并替换。若 $m>n$，到第 $n$ 步时已无剩余 $v_j$，$u_{n+1}$ 将属于前 $n$ 个 $u$ 的张成，矛盾。
\end{proof}

\begin{corollary}[有限维空间的基定理]\label{cor:y0-finite-basis}
有限维空间中，任意两个基的元素个数相同，任意线性无关族可扩充为基，任意有限生成族可删减为基。在 $n$ 维空间中，任意 $n$ 个线性无关向量都是基，任意 $n$ 个生成向量也都是基。
\end{corollary}
\begin{proof}
对两个基分别用交换引理，得其元素个数互相不超过，从而相同。从有限生成族中逐一删去可由其余向量表示者，最终得到基；交换引理先将无关族换入有限生成族，再只删去多余的其余向量，即得到包含原无关族的基。若 $n$ 个无关向量未张成，可再加入一个不在其张成中的向量，与交换引理矛盾。若 $n$ 个生成向量相关，可删去至少一个，导致由少于 $n$ 个向量张成空间，也与交换引理矛盾。
\end{proof}

任意维空间中基的存在性已在定理 \ref{thm:y0-hamel-basis} 中用 Zorn 引理证明。这里有限维的结论只用有限交换，不需要一般的选择公理。具有基时，每个向量的有限坐标表达唯一，因两个表达相减给出有限线性关系。

\subsection{线性算子、核、值域与延拓}
\begin{definition}[线性算子与线性泛函]\label{def:y0-linear-operator}
设 $X,Y$ 为同一标量域上的线性空间，$D(T)\subseteq X$ 为子空间。映射 $T:D(T)\to Y$ 若满足
\[
 T(\alpha x+\beta y)=\alpha Tx+\beta Ty,
\]
称为线性算子。定义域、值域与核分别记为
\[
 D(T),\qquad R(T)=T[D(T)],\qquad N(T)=\ker T=\{x\in D(T):Tx=0\}.
\]
若 $Y=\mathbb K$，称 $T$ 为线性泛函。双射的线性算子称为线性同构。
\end{definition}

\begin{proposition}[核、值域与逆算子]\label{prop:y0-operator}
核及值域均为子空间。$T$ 单射当且仅当 $N(T)=\{0\}$；此时 $T^{-1}:R(T)\to D(T)$ 线性。映射 $T$ 若要在整个 $Y$ 上具有逆，还必须满足 $R(T)=Y$。
\end{proposition}
\begin{proof}
线性性给出 $T0=0$，以及核和值域对线性组合封闭。因为 $Tx=Ty$ 当且仅当 $T(x-y)=0$，单射性等价于零核。单射时，对 $u=Tx$、$v=Ty$，有 $\alpha u+\beta v=T(\alpha x+\beta y)$，故 $T^{-1}(\alpha u+\beta v)=\alpha x+\beta y$，证明逆算子线性。整个 $Y$ 上的逆需对每个 $y$ 有原像，即满射。
\end{proof}

两个算子相等，要求定义域相同且逐点取值相同。若 $D(T_1)\subseteq D(T_2)$ 且 $T_2|_{D(T_1)}=T_1$，称 $T_2$ 为 $T_1$ 的延拓，记 $T_1\subseteq T_2$；这也是将算子看作图时的包含关系。线性泛函的配对在本稿中对向量线性，例如 $\langle T,x\rangle=T(x)$；若使用内积实现泛函，则须遵守第一变量线性的既定内积约定。

\begin{proposition}[由基上的取值定义算子]\label{prop:y0-basis-map}
设 $B$ 为 $X$ 的 Hamel 基。任意映射 $a:B\to Y$ 唯一延拓为线性算子 $T:X\to Y$，由
\[
 T\left(\sum_{b\in F}\alpha_bb\right)
          =\sum_{b\in F}\alpha_ba(b)
 \qquad(F\subseteq B\text{ 有限})
\]
定义。
\end{proposition}
\begin{proof}
基的有限坐标表达唯一，故定义不依赖表达；将两向量的表达统一到一个有限基向量集合上，即可验证线性性。任何延拓都须满足同一公式，故唯一。
\end{proof}

这是代数构造；没有连续性结论。后续分析中必须另行检查定义的算子是否有界或连续。

\subsection{商空间与第一同构定理}
\begin{theorem}[商线性空间]\label{thm:y0-quotient}
设 $M\subseteq X$ 为子空间。关系 $x\sim y$ 当且仅当 $x-y\in M$ 是等价关系，其等价类为 $[x]=x+M$。在集合 $X/M=\{x+M:x\in X\}$ 上定义
\[
 [x]+[y]=[x+y],\qquad\alpha[x]=[\alpha x].
\]
这些运算良定，使 $X/M$ 成为线性空间。商映射 $q:X\to X/M$、$q(x)=[x]$ 为满射线性算子，且 $\ker q=M$。
\end{theorem}
\begin{proof}
$x-x=0\in M$ 给出自反性；$x-y\in M$ 蕴含其负向量 $y-x\in M$，给出对称性；两个差相加给出传递性。若 $x'=x+m$、$y'=y+n$，则 $(x'+y')-(x+y)=m+n\in M$，且 $\alpha x'-\alpha x=\alpha m\in M$，所以运算不依赖代表元。线性空间公理由原空间逐条传递，零元为 $M=[0]$，负元为 $[-x]$。商映射的线性性和满射性直接由定义成立；$[x]=[0]$ 当且仅当 $x\in M$，故核为 $M$。
\end{proof}

\begin{theorem}[商空间的因子分解与第一同构定理]\label{thm:y0-first-isomorphism}
设 $T:X\to Y$ 线性，$M\subseteq\ker T$。则存在唯一线性算子 $\widetilde T:X/M\to Y$ 使 $T=\widetilde T\circ q$，其公式为 $\widetilde T([x])=Tx$。特别地，
\[
 X/\ker T\ \cong\ R(T).
\]
对定义在子空间上的算子，相应结论为 $D(T)/N(T)\cong R(T)$。
\end{theorem}
\begin{proof}
若 $[x]=[y]$，则 $x-y\in M\subseteq\ker T$，所以 $Tx=Ty$，定义良定。线性性由商运算继承，唯一性由 $q$ 满射得到。取 $M=\ker T$ 并将目标限制为 $R(T)$，所得算子满射，且 $\widetilde T([x])=0$ 蕴含 $x\in\ker T$、$[x]=0$，所以为同构。
\end{proof}

\begin{corollary}[有限维的秩--零度公式]\label{cor:y0-rank-nullity}
若 $X$ 有限维且 $T:X\to Y$ 线性，则 $\dim X=\dim\ker T+\dim R(T)$。
\end{corollary}
\begin{proof}
取核的一组基 $e_1,\ldots,e_k$ 并扩充为 $X$ 的基 $e_1,\ldots,e_k,u_1,\ldots,u_r$。向量 $Tu_1,\ldots,Tu_r$ 张成值域；若它们有线性关系，则对应的 $u_j$ 组合属于核，再由原基的线性无关性得系数全零。因此它们为值域的基，维数公式成立。
\end{proof}

\begin{example}[商空间消除积分中的代表元歧义]
在可积函数构成的线性空间中，令 $M$ 为几乎处处为零的函数组成的子空间。商空间把 $f$ 与 $g$ 识别为同一元素，恰当且仅当 $f=g$ a.e.。积分距离 $\int|f-g|$ 对代表元不敏感，并且在商空间中只有零元素的距离为零，这就是 $L^1$ 空间的代数起点。抽象商空间无需子空间闭；要得到 Hausdorff 的拓扑商或真正的商范数，则需后续的闭性条件。
\end{example}

\begin{exercise}[代数补空间与投影]
假设选择公理。证明每个子空间 $M\subseteq X$ 都存在子空间 $N$，使 $X=M\oplus N$，即每个 $x$ 唯一写成 $m+n$。由此构造投影 $P:X\to M$，并说明它未必连续。
\end{exercise}
\begin{solution}
取 $M$ 的 Hamel 基并用定理 \ref{thm:y0-hamel-basis} 扩充为 $X$ 的基。令 $N$ 为新增基向量的张成。基的唯一有限表达给出 $X=M+N$ 及 $M\cap N=\{0\}$。定义 $P(m+n)=m$，则 $P$ 线性、$P^2=P$、$R(P)=M$、$\ker P=N$。例如在 $\ell^2$ 中，取包含全部单位向量 $e_n$ 的 Hamel 基，按命题 \ref{prop:y0-basis-map} 定义泛函 $f(e_n)=n$，其余基向量取值为零。它在单位向量上无界，因而不连续；对应的投影 $P(x)=f(x)e_1$ 满足 $P^2=P$，同样不连续。
\end{solution}


% !TeX encoding = UTF-8
\section*{本章参考文献与阅读衔接}
\phantomsection
\addcontentsline{toc}{section}{本章参考文献与阅读衔接}

本章主线依据 \cite{yosida-zh,yosida-en}。原书第零章末的参考条目已核对：中译本与英文版均位于印刷第 22 页，即两份本地 PDF 的第 37 页。它按主题列出拓扑空间的三组参考作者和测度空间的两组参考作者，没有另列集合论与线性空间书目。原书的作者后缀“[1]”指各作者在全书参考书目中的第一条，\emph{不是本稿的参考文献编号}。对应关系如下。

\begin{enumerate}
 \item \textbf{P. Alexandroff--H. Hopf [1]}：\emph{Topologie}, Vol.~I, Springer, 1935，见 \cite{alexandroff-hopf}。英文原书总书目印刷第 469 页（PDF 第 484 页）使用拼写 Alexandrov；章末及中译本使用 Alexandroff，均指同一作者。
 \item \textbf{N. Bourbaki [1]}：\emph{Topologie G\'en\'erale}，Hermann，原书记录年份 1940--1942，见 \cite{bourbaki-topology}。核对位置为英文原书总书目印刷第 470 页（PDF 第 485 页）。这是分册出版的书目记录；不可把现代合订本的重印年份当作原书所引年份。
 \item \textbf{J. L. Kelley [1]}：\emph{General Topology}, D.~Van Nostrand, 1955，见 \cite{kelley-topology}。核对位置为英文原书总书目印刷第 475 页（PDF 第 490 页）。Springer 的 GTM 27 版为后来的再版，书目中另作说明。
 \item \textbf{P. R. Halmos [1]}：\emph{Measure Theory}, D.~Van Nostrand, 1950，见 \cite{halmos-measure}。核对位置为英文原书总书目印刷第 473 页（PDF 第 488 页）。
 \item \textbf{S. Saks [1]}：\emph{Theory of the Integral}, Warszawa, 1937，见 \cite{saks-integral}。核对位置为英文原书总书目印刷第 481 页（PDF 第 496 页）。
\end{enumerate}

以上五条均已并入书后的统一参考文献表，正文不再维护另一套编号。这里记录的是吉田原书的参考指引；本次实际核对了原书章末与总书目的对应条目，未把这五部参考书的全文标为已经逐页查阅。

阅读衔接上，第 \ref{sec:yosida-00-01} 节的 Zorn 引理用于代数基及延拓构造；第 \ref{sec:yosida-00-02} 节的紧性、类别定理与连续函数逼近分别服务于弱拓扑、Baire 方法和表示定理；第 \ref{sec:yosida-00-03} 节的积分工具为 $L^p$ 空间、卷积及积分算子准备基础；第 \ref{sec:yosida-00-04} 节的核、值域和商空间则在后续算子论中反复出现。第一章从这些基础出发，以半范数建立拓扑向量空间。


% !TeX encoding = UTF-8
\SetChapterInfo{subtitle={局部凸拓扑、函数空间与分布的共同基础},keywords={半范数; 拟范数; 完备化; 分布},prereq={线性代数、度量空间、Lebesgue 积分},goals={用半范数组织拓扑; 掌握连续性与完备性; 建立分布的基本运算}}
\chapter{半范数}\label{ch:yosida-01}
\ChapterSource{本章以吉田耕作第六版中译本第一章（印刷页 23--70）为主线，以英文第六版第一章（印刷页 23--67）辅助核对，保留原书全部十四节的主题与顺序。以下是统一记号后重新组织的讲述，定理编号属于本稿；标为“补充”的论证、辨析及习题用于展开原书略写的步骤。参见 \cite{yosida-zh,yosida-en}。}

半范数的作用不只是给向量赋予一个“长度”。在函数空间中，函数大小、各阶导数大小以及不同区域上的行为，往往需要分别测量。有限个或无限个这样的测量共同确定收敛，连续算子则把一种测量控制在另一种测量之内。本章从这一观点出发，依次建立局部凸空间、赋范与拟赋范空间、完备化和商空间，再把这些工具用于广义函数。

除特别说明外，标量域 $\mathbb K$ 是 $\R$ 或 $\mathbb{C}$，所有拓扑向量空间均为 Hausdorff 空间。$\Omega\subset\R^n$ 是开集，$K\Subset\Omega$ 表示 $K$ 是包含于 $\Omega$ 的紧集。多重指标 $\alpha=(\alpha_1,\ldots,\alpha_n)$ 满足 $|\alpha|=\sum\alpha_j$、$\alpha!=\prod\alpha_j!$，并记 $\partial^\alpha=\prod\partial_{x_j}^{\alpha_j}$。分布配对对测试函数是线性的，不加入复共轭；内积则约定对第一个变量线性。

\section{半范数与局部凸线性拓扑空间}\label{sec:yosida-01-01}
\noindent\textit{原书位置：中译本第 23--30 页；英文版第 23--30 页。}

\subsection{半范数及其几何}
\begin{definition}[半范数、均衡集与吸收集]
线性空间 $X$ 上的半范数是满足
\[
 p(x+y)\le p(x)+p(y),\qquad p(ax)=|a|p(x)
\]
的函数 $p:X\to[0,\infty)$。集合 $U\subset X$ 称为均衡的，若 $|a|\le1$ 蕴含 $aU\subset U$；称为吸收的，若对每个 $x\in X$，存在 $t>0$ 使 $x\in tU$。均衡凸集亦称绝对凸集。
\end{definition}
半范数可以在非零向量处为零。例如在 $\R^2$ 上，$p(x_1,x_2)=|x_1|$ 不测量第二个坐标。它的单位球是一条无限长的竖直带，而不是有界圆盘。

\begin{proposition}[基本不等式与半范数球]\label{prop:y1-seminorm-basic}
半范数满足 $p(0)=0$、$p(-x)=p(x)$ 以及
\[
 |p(x)-p(y)|\le p(x-y).
\]
若 $c>0$，则 $U_c=\{x:p(x)\le c\}$ 是吸收的绝对凸集，且
\[
 p(x)=c\inf\{t>0:x\in tU_c\}.
\]
半范数的零空间 $\ker p$ 是线性子空间。
\end{proposition}
\begin{proof}
齐次性给出前两个等式。由 $p(x)\le p(x-y)+p(y)$ 及交换 $x,y$ 后的不等式得到第三式。若 $x,y\in U_c$、$0\le s\le1$，则 $p(sx+(1-s)y)\le sc+(1-s)c=c$；均衡性由齐次性得到。选择 $t>p(x)/c$ 即得吸收性。最后 $x\in tU_c$ 等价于 $p(x)\le tc$，对 $t$ 取下确界即得公式。零空间的封闭于加法和数乘也直接来自两条公理。
\end{proof}

\begin{definition}[Minkowski 泛函]
若 $U$ 是吸收的绝对凸集，定义
\[
 p_U(x)=\inf\{t>0:x\in tU\}.
\]
这称为 $U$ 的 Minkowski 泛函或规范泛函。
\end{definition}
\begin{theorem}[Minkowski 泛函是半范数]\label{thm:y1-gauge}
$p_U$ 是有限值半范数，且
\[
 \{p_U<1\}\subset U\subset\{p_U\le1\}.
\]
若 $U$ 是拓扑向量空间中零点的绝对凸邻域，则 $p_U$ 连续。若 $U$ 还开，则 $U=\{p_U<1\}$；若 $U$ 还闭，则 $U=\{p_U\le1\}$。
\end{theorem}
\begin{proof}
吸收性保证下确界有限。对 $a\ne0$，均衡性蕴含 $aU=|a|U$，故 $p_U(ax)=|a|p_U(x)$。给定 $r>p_U(x)$、$s>p_U(y)$，均衡性保证 $x/r,y/s\in U$。凸性给出
\[
 \frac{x+y}{r+s}=\frac r{r+s}\frac xr+\frac s{r+s}\frac ys\in U.
\]
令 $r\downarrow p_U(x)$、$s\downarrow p_U(y)$ 得次可加性。若 $p_U(x)<1$，选 $r<1$ 使 $x\in rU\subset U$；另一包含关系显然。$U$ 为邻域时，$(\varepsilon/2)U\subset\{p_U<\varepsilon\}$，所以 $p_U$ 在零点连续，反三角不等式保证处处连续。

若 $U$ 开且 $x\in U$，数乘连续性给出某个 $r<1$ 使 $x/r\in U$，于是 $p_U(x)<1$。若 $U$ 闭且 $p_U(x)\le1$，则 $x/(1+1/j)\in U$，取极限得 $x\in U$。这也说明一般不能把两个包含关系不加条件地写成等式。
\end{proof}

\subsection{由半范数族生成拓扑}
\begin{theorem}[局部凸空间的半范数刻画]\label{thm:y1-lc}
设 $\mathcal P$ 是 $X$ 上一族半范数，且满足分离条件：每个 $x\ne0$ 都有 $p\in\mathcal P$ 使 $p(x)>0$。以
\[
 U(p_1,\ldots,p_r;\varepsilon_1,\ldots,\varepsilon_r)
 =\{x:p_j(x)<\varepsilon_j,\ 1\le j\le r\}
\]
为零点的邻域基，以平移得到各点邻域，则 $X$ 成为 Hausdorff 局部凸拓扑向量空间，各 $p\in\mathcal P$ 连续。反过来，每个 Hausdorff 局部凸拓扑向量空间的拓扑均由它的连续半范数生成。
\end{theorem}
\begin{proof}
上述集合的有限交仍具有同样形式，并且均为吸收的绝对凸集。把各阈值缩小一半得到 $V+V\subset U$，因而加法连续。在 $(a_0,x_0)$ 处，利用
\[
 p_j(ax-a_0x_0)\le |a|p_j(x-x_0)+|a-a_0|p_j(x_0)
\]
同时使 $a$ 接近 $a_0$、$x$ 接近 $x_0$，即可验证数乘的联合连续性。各半范数连续来自命题 \ref{prop:y1-seminorm-basic}。若 $x\ne y$，选 $p(x-y)=c>0$，则 $x+\{p<c/3\}$ 和 $y+\{p<c/3\}$ 不相交，故拓扑 Hausdorff。

反向证明中，每个零邻域包含绝对凸零邻域 $U$。定理 \ref{thm:y1-gauge} 给出连续半范数 $p_U$，且 $\{p_U<1\}\subset U$，所以这些半范数生成原拓扑。Hausdorff 性保证它们分离点，否则一个非零向量将落在所有零邻域中。
\end{proof}

图 \ref{fig:y1-seminorm} 比较单个半范数与分离点的半范数族；带状邻域的无界方向正是被忽略的方向。
\begin{figure}[htbp]
\centering
\begin{tikzpicture}[fa diagram]
\begin{scope}
  \fill[LogicColor!12] (-.8,-1.7) rectangle (.8,1.7);
  \draw[fa main,dashed] (-.8,-1.7)--(-.8,1.7) (.8,-1.7)--(.8,1.7);
  \draw[fa axis] (-2,0)--(2,0) node[right] {$x_1$};
  \draw[fa axis] (0,-1.9)--(0,2) node[above] {$x_2$};
  \draw[fa accent] (0,-1.7)--(0,1.7);
  \node[fa note] at (0,-2.5) {$p_1(x)=|x_1|<1$\\竖直方向是 $\ker p_1$};
\end{scope}
\begin{scope}[xshift=6.5cm]
  \fill[LogicColor!12] (-.8,-.8) rectangle (.8,.8);
  \draw[fa main,dashed] (-.8,-.8) rectangle (.8,.8);
  \draw[fa axis] (-2,0)--(2,0) node[right] {$x_1$};
  \draw[fa axis] (0,-1.9)--(0,2) node[above] {$x_2$};
  \node[fa note] at (0,-2.5) {$p_1(x)<1,\ p_2(x)=|x_2|<1$\\$\ker p_1\cap\ker p_2=\{0\}$};
\end{scope}
\end{tikzpicture}
\caption{一个半范数可以忽略非零方向；有限个半范数球的交形成基本零邻域。本图采用 $\R^2$ 的坐标例子，虚线边界不属于开球；一般局部凸空间不一定存在有界零邻域。}
\label{fig:y1-seminorm}
\end{figure}


把半范数族扩充为一切有限最大值及正倍数，不改变拓扑。于是许多估计可写成 $q(Tx)\le Cp(x)$；若原来的族没有经过这种扩充，则右边必须允许有限个半范数的最大值。

\subsection{光滑函数空间与测试函数空间}
函数的支集定义为其非零点集在定义域中的闭包。记
\[
 \mathcal E^k(\Omega)=C^k(\Omega),\qquad
 p_{K,m}(f)=\max_{|\alpha|\le m}\sup_{x\in K}|\partial^\alpha f(x)|,
\]
其中 $m\le k$；$k=\infty$ 时允许任意有限的 $m$。这些半范数定义局部 $C^k$ 拓扑：$f_j\to f$ 就是允许的各阶导数在每个紧集上一致收敛。取紧集穷竭 $K_j\subset\operatorname{int}K_{j+1}$，$\bigcup K_j=\Omega$，所有紧集都包含于某个 $K_j$，于是可用可数个 $p_{K_j,m}$ 定义同一拓扑。度量化将在下一节证明。

光滑紧支集函数并非只有零函数。例如
\[
 b(x)=\begin{cases}\exp\bigl(-1/(1-|x|^2)\bigr),&|x|<1,\\0,&|x|\ge1\end{cases}
\]
属于 $C_c^\infty(\R^n)$。其边界光滑性来自：任意幂次的 $t^{-1}$ 乘 $e^{-1/t}$，在 $t\downarrow0$ 时都趋于零。

令
\[
 \Test_K(\Omega)=\{\varphi\in C_c^\infty(\Omega):\supp\varphi\subset K\},\qquad
 \Test(\Omega)=C_c^\infty(\Omega).
\]
$\Test_K$ 由 $p_{K,m}$ 赋拓扑；$\Test$ 取使所有包含映射 $\Test_K\to\Test$ 连续的最强局部凸拓扑，称为局部凸归纳极限拓扑。利用穷竭可写为严格的可数归纳极限 $\Test=\bigcup_j\Test_{K_j}$。其中 $\Test_{K_j}$ 在 $\Test_{K_{j+1}}$ 中是闭子空间，诱导拓扑等于原拓扑。

\begin{proposition}[归纳极限的半范数判别]\label{prop:y1-inductive}
$\Test$ 上的半范数 $q$ 连续，当且仅当其在每个 $\Test_K$ 上的限制连续。更一般地，线性映射 $T:\Test\to Y$（$Y$ 局部凸）连续，当且仅当所有 $T|_{\Test_K}$ 连续。
\end{proposition}
\begin{proof}
一个半范数 $q$ 的各限制连续时，把 $q$ 加入生成 $\Test$ 拓扑的半范数族，仍使各包含映射连续；最强性表明 $q$ 本来就连续。必要性由复合映射的连续性得到。对一般线性映射，只须把此结论用于 $q\circ T$，其中 $q$ 遍历 $Y$ 的连续半范数。
\end{proof}

\begin{theorem}[测试函数的序列收敛]\label{thm:y1-D-sequence}
$\varphi_j\to\varphi$ 在 $\Test(\Omega)$ 中成立，当且仅当存在共同紧集 $K\Subset\Omega$ 包含所有 $\supp\varphi_j$ 及 $\supp\varphi$，且对每个 $m$，$p_{K,m}(\varphi_j-\varphi)\to0$。
\end{theorem}
\begin{proof}
充分性由包含映射连续得到。证明必要性时先令极限为零。若不存在共同紧集，可逐次选子列 $\varphi_{j_l}$ 和无域内聚点的 $x_l\in\Omega$，使 $\varphi_{j_l}(x_l)\ne0$；这是因为其支集不断逃出紧集穷竭，非零点可在逃出的支集附近选取。定义
\[
 q(\psi)=\sum_{l=1}^\infty
 \frac{|\psi(x_l)|}{|\varphi_{j_l}(x_l)|}.
\]
对任意紧支集 $\psi$ 此和只有有限项；在每个 $\Test_K$ 上，它是有限个点评价绝对值之和，所以由命题 \ref{prop:y1-inductive} 连续。然而 $q(\varphi_{j_l})\ge1$，矛盾。共同紧集存在后，全局半范数 $\psi\mapsto\sup_\Omega|\partial^\alpha\psi|$ 的各层限制连续，故在 $\Test$ 中连续，得到全部导数的一致收敛。一般极限通过相减化为零极限。
\end{proof}

图 \ref{fig:y1-test-support} 解释为什么局部光滑收敛不能替代共同紧支集条件。
\begin{figure}[htbp]
\centering
\begin{tikzpicture}[fa diagram,x=1cm,y=1cm]
  \fill[LogicColor!8] (-1.2,-.2) rectangle (1.2,1.55);
  \node at (0,1.85) {固定紧集 $K$};
  \draw[fa axis] (-2,0)--(8,0) node[right] {$x$};
  \foreach \a/\h/\c in {0/1/LogicColor,3/.65/ExColor,6/.35/DefColor}{
    \draw[thick,draw=\c] plot[domain=-1:1,samples=40] ({\a+\x},{\h*exp(1-1/(1.001-\x*\x))});
  }
  \node[fa note] at (3,-.65) {$\eta(x)$\hspace{1.2cm}$\eta(x-j)$\hspace{1cm}$a_j\eta(x-j)$};
  \draw[->,fa accent] (2.1,1.2)--(6.8,1.2) node[midway,above] {支集逃向无穷远};
\end{tikzpicture}
\caption{测试函数收敛需要共同紧支集。对非零 $\eta\in\Test(\R)$，平移在每个固定紧集上最终为零，故在 $\mathcal E$ 中趋零；即使 $a_j>0$ 趋零，$a_j\eta(\cdot-j)$ 仍不在 $\Test$ 中趋零。图中衰减高度只作示意，不给出统一导数估计。}
\label{fig:y1-test-support}
\end{figure}


\begin{pitfall}[局部收敛与测试函数收敛]
若 $\eta\ne0$ 是紧支集光滑函数，则 $\eta(x-j)$ 在 $\mathcal E(\R)$ 中趋于零，却不在 $\Test(\R)$ 中趋于零；即使乘上 $e^{-j^2}$ 也一样。共同紧支集是额外条件。$\Test$ 通常不可度量，不能把序列判别未经证明地用于任意非线性映射的连续性。
\end{pitfall}

\subsection{正则化与逼近}
\begin{proposition}[紧支集连续函数的光滑逼近]\label{prop:y1-mollify}
取非负 $\rho\in\Test(\R^n)$ 满足 $\int\rho=1$、$\supp\rho\subset\overline B(0,1)$，并令 $\rho_\varepsilon(x)=\varepsilon^{-n}\rho(x/\varepsilon)$。若 $f\in C_c(\R^n)$，则
\[
 f_\varepsilon(x)=\int f(y)\rho_\varepsilon(x-y)\,dy
\]
属于 $C_c^\infty$，支集包含于 $\supp f+\overline B(0,\varepsilon)$，且 $\|f_\varepsilon-f\|_\infty\to0$。
\end{proposition}
\begin{proof}
对光滑核在积分号下求导，得 $\partial^\alpha f_\varepsilon=f*\partial^\alpha\rho_\varepsilon$；积分区间可固定在 $\supp f$，故求导合法。支集结论来自卷积核仅在距离不超过 $\varepsilon$ 时非零。紧支集连续函数在整个 $\R^n$ 上一致连续，故
\[
 |f_\varepsilon(x)-f(x)|\le
 \sup_{|h|\le\varepsilon}|f(x-h)-f(x)|\longrightarrow0
\]
一致成立。若 $\supp f\Subset\Omega$，充分小的 $\varepsilon$ 保证 $f_\varepsilon\in\Test(\Omega)$。
\end{proof}

\begin{exercise}[补充：半范数的零空间]
证明 $X/\ker p$ 上 $\|x+\ker p\|=p(x)$ 是范数；再说明一个半范数族不分离点时，怎样先取公共零空间的商而得到 Hausdorff 空间。
\end{exercise}
\begin{solution}
若 $x-y\in\ker p$，反三角不等式给出 $p(x)=p(y)$，所以定义良好，正定性在商空间上成立，其余公理由 $p$ 继承。对半范数族令 $N=\bigcap_p\ker p$，所有 $p$ 都下降为 $X/N$ 上的半范数，且下降后的族分离商空间中的点。
\end{solution}

\section{范数和拟范数}\label{sec:yosida-01-02}
\noindent\textit{原书位置：中译本第 31--33 页；英文版第 30--32 页。}

\subsection{范数与强收敛}
\begin{definition}[赋范线性空间]
半范数 $\|\cdot\|$ 若满足 $\|x\|=0\Rightarrow x=0$，称为范数。它诱导平移不变距离 $d(x,y)=\|x-y\|$。在此距离中的收敛称为范数收敛或强收敛，记作 $x_j\to x$，原书也用 $\mathrm{s}\!\operatorname{-lim}x_j=x$。
\end{definition}
范数满足反三角不等式；若 $x_j\to x$、$y_j\to y$、$a_j\to a$，则
\[
 \|x_j\|\to\|x\|,\qquad x_j+y_j\to x+y,\qquad a_jx_j\to ax.
\]
前两条由三角不等式得到，最后一条由
\[
 \|a_jx_j-ax\|\le |a_j|\|x_j-x\|+|a_j-a|\|x\|
\]
及收敛标量序列有界得到。这是在最简单情形下的拓扑与线性运算相容性。

\subsection{吉田意义的拟范数}
\begin{definition}[拟范数：本章的历史用法]\label{def:y1-Fnorm}
函数 $q:X\to[0,\infty)$ 若满足
\[
 q(x)=0\Longleftrightarrow x=0,\quad q(-x)=q(x),\quad
 q(x+y)\le q(x)+q(y),
\]
并满足两条数乘连续性条件
\[
 a_j\to0\ \Longrightarrow\ q(a_jx)\to0\quad(x\text{ 固定}),\qquad
 q(x_j)\to0\ \Longrightarrow\ q(ax_j)\to0\quad(a\text{ 固定}),
\]
称为吉田意义的拟范数。它定义距离 $d(x,y)=q(x-y)$。
\end{definition}
\begin{remark}[补充：与现代术语的区别]
原书的 quasi-norm 在这里不要求 $q(ax)=|a|q(x)$，但要求真正的三角不等式。现代文献常以 quasi-norm 指满足绝对齐次性和放宽三角不等式 $\|x+y\|\le C(\|x\|+\|y\|)$ 的函数。为避免混用，本章在谈完备度量线性空间时使用定义 \ref{def:y1-Fnorm}，也可称之为一种 F-范数。
\end{remark}

\begin{lemma}[Kakutani 的一致数乘结论]\label{lem:y1-uniform-scalars}
若 $q(x_j)\to0$，则对每个 $M<\infty$，
\[
 \sup_{|a|\le M}q(ax_j)\longrightarrow0.
\]
因此定义 \ref{def:y1-Fnorm} 中数乘实际上联合连续。
\end{lemma}
\begin{proof}
这是原书借助 Egorov 定理给出的关键步骤，现将其展开。令 $r_j(a)=q(ax_j)$。由反三角不等式与第一条数乘条件，各 $r_j$ 在标量域 $\mathbb K$ 上连续；第二条条件给出对每个 $a$，$r_j(a)\to0$。在标量域的单位球上应用 Egorov 定理，可取有限正 Lebesgue 测度的可测集 $A$，使 $r_j\to0$ 在 $A$ 上一致。复标量时把 $\mathbb{C}$ 看作 $\R^2$。

可测集示性函数的 $L^1$ 平移连续性给出 $|A\mathbin\triangle(A+h)|\to0$。这里这一事实可由有限测度集在测度上用有限个矩形逼近、再利用矩形的平移连续性证明。因此存在 $\delta>0$，使 $|h|<\delta$ 时 $A\cap(A+h)\ne\varnothing$，即 $h=a-b$，其中 $a,b\in A$。于是
\[
 r_j(h)\le r_j(a)+r_j(-b)=r_j(a)+r_j(b)
 \le2\sup_{c\in A}r_j(c).
\]
这给出 $r_j\to0$ 在 $|h|<\delta$ 上一致。给定 $M$，选整数 $N>M/\delta$，对 $|a|\le M$ 写 $a=N(a/N)$，次可加性给出 $r_j(a)\le Nr_j(a/N)$，即得结论。

若 $a_j\to a$、$x_j\to x$，则
\[
 q(a_jx_j-ax)\le q(a_j(x_j-x))+q((a_j-a)x)\to0.
\]
第一项使用刚证的一致结论，第二项使用固定向量的数乘条件。加法与 $q$ 自身的连续性则来自三角和反三角不等式。
\end{proof}

\subsection{可数半范数族的度量化}
\begin{theorem}[可数半范数族与平移不变距离]\label{thm:y1-countable-metric}
设分离点的半范数族为 $p_1,p_2,\ldots$。则
\[
 q(x)=\sum_{j=1}^\infty2^{-j}\frac{p_j(x)}{1+p_j(x)},\qquad d(x,y)=q(x-y)
\]
定义吉田意义的拟范数与距离，生成同一局部凸拓扑。并且 $d$-Cauchy 等价于对每个 $j$ 在 $p_j$ 下 Cauchy。
\end{theorem}
\begin{proof}
函数 $h(t)=t/(1+t)$ 增加、$h(0)=0$，且
\[
 h(s+t)=\frac{s}{1+s+t}+\frac{t}{1+s+t}\le h(s)+h(t).
\]
因此各求和项满足三角不等式，级数由 $\sum2^{-j}$ 控制。分离性给出正定性。若 $q(x_l)\to0$，则每个固定项趋于零，从而 $p_j(x_l)\to0$。反向结论先固定足够小的级数尾，再控制有限个首项得到。此“有限首项加一致小尾项”论证也证明两个邻域系统一致，以及两种 Cauchy 条件等价。半范数的齐次性与同一论证验证两条数乘条件。
\end{proof}
于是 $\mathcal E^k(\Omega)$、$\Test_K(\Omega)$ 可度量。不同的穷竭或等价半范数族给出不同公式的距离，但拓扑与 Cauchy 概念相同。

\begin{proposition}[补充：可赋范性的判别]\label{prop:y1-normable}
Hausdorff 局部凸空间可由一个范数生成拓扑，当且仅当它有一个有界的零邻域。
\end{proposition}
\begin{proof}
范数单位球显然是有界零邻域。反向取有界零邻域中的绝对凸零邻域 $U$。其规范泛函 $p_U$ 连续。若 $p_U(x)=0$，则 $x\in tU$ 对每个 $t>0$ 成立；有界性使 $tU$ 最终包含于任意给定零邻域，故 $x=0$。对任意零邻域 $V$，有界性给出 $U\subset cV$，于是 $\{p_U<1/c\}\subset V$（可先取均衡 $V$）。所以 $p_U$ 的范数拓扑等于原拓扑。
\end{proof}
有界集的正式定义见第 \ref{sec:yosida-01-07} 节。此判别说明可度量性远弱于可赋范性：多个局部测量不能总由一种全局测量替代。

\section{几例赋范线性空间}\label{sec:yosida-01-03}
\noindent\textit{原书位置：中译本第 33--39 页；英文版第 32--38 页。}

\subsection{有界连续函数与可积函数}
拓扑空间 $S$ 上的有界连续函数组成 $C_b(S)$，赋予范数 $\|f\|_\infty=\sup_S|f|$。范数收敛就是全域一致收敛。原书此处的 $C(S)$ 限于有界连续函数；若 $S$ 紧，所有连续函数都自动有界，记号 $C(S)$ 无歧义。

对测度空间 $(S,\Sigma,\mu)$ 与 $1\le p<\infty$，令
\[
 L^p(S,\mu)=\{f\text{ 可测}:\int_S|f|^p\,d\mu<\infty\}/\text{几乎处处相等},
 \qquad\|f\|_p=\left(\int_S|f|^p\,d\mu\right)^{1/p}.
\]
商掉几乎处处相等是正定性的必要条件。$L^\infty$ 由本质有界函数的等价类组成，范数为
\[
 \|f\|_\infty=\operatorname*{ess\,sup}|f|
 =\inf\{M\ge0:|f|\le M\text{ 几乎处处}\}.
\]
本质上确界与逐点上确界不同：零测集上的任意修改都不改变前者。

\begin{lemma}[Young 不等式]
若 $1<p<\infty$、$1/p+1/p'=1$，则对 $a,b\ge0$，
\[
 ab\le\frac{a^p}{p}+\frac{b^{p'}}{p'}.
\]
等号成立当且仅当 $a^p=b^{p'}$。
\end{lemma}
\begin{proof}
固定 $b$，函数 $a\mapsto a^p/p-ab+b^{p'}/p'$ 的导数为 $a^{p-1}-b$，唯一极小值点为 $a=b^{1/(p-1)}$，极小值为零。$b=0$ 的情形直接验证。
\end{proof}
\begin{theorem}[Hölder 与 Minkowski 不等式]\label{thm:y1-holder}
对共轭指数 $p,p'$（包含 $p=1,p'=\infty$），
\[
 \int|fg|\,d\mu\le\|f\|_p\|g\|_{p'}.
\]
对 $1\le p\le\infty$，$\|f+g\|_p\le\|f\|_p+\|g\|_p$，因而 $L^p$ 是赋范线性空间。
\end{theorem}
\begin{proof}
端点 Hölder 由 $|g|\le\|g\|_\infty$ 几乎处处得到。内部指数时，若两范数不为零，将 Young 不等式用于 $|f|/\|f\|_p$ 和 $|g|/\|g\|_{p'}$ 后积分；任一范数为零则直接成立。

先由 $|f+g|^p\le2^{p-1}(|f|^p+|g|^p)$ 确保 $f+g\in L^p$。当 $1<p<\infty$，利用 Hölder 和 $(p-1)p'=p$，
\[
 \|f+g\|_p^p
 \le\int |f+g|^{p-1}(|f|+|g|)\,d\mu
 \le\|f+g\|_p^{p-1}(\|f\|_p+\|g\|_p).
\]
若左侧非零，约去相应因子即得 Minkowski；否则结论显然。$p=1,\infty$ 直接由逐点三角不等式得到。
\end{proof}
当 $1<p<\infty$ 且 $f,g$ 非零时，Minkowski 等号成立恰在 $f=cg$ 几乎处处、$c>0$ 的情形。实值时是同号正比例，复值时还要求相同相位；证明由 Hölder 的模值比例条件与逐点三角等号条件合并得到。$p=1$ 只须几乎处处无相消，不要求固定比例；$p=\infty$ 的等号条件也不同。

\begin{proposition}[有限测度空间上的极限]
若 $\mu(S)<\infty$ 且 $f\in L^\infty$，则 $\lim_{p\to\infty}\|f\|_p=\|f\|_\infty$。
\end{proposition}
\begin{proof}
总测度为零或 $\|f\|_\infty=0$ 时直接成立。记 $M=\|f\|_\infty>0$，上界为 $\|f\|_p\le\mu(S)^{1/p}M$。对 $0<\varepsilon<M$，集合 $A_\varepsilon=\{|f|>M-\varepsilon\}$ 有正测度，所以
\[
 \|f\|_p\ge\mu(A_\varepsilon)^{1/p}(M-\varepsilon).
\]
先令 $p\to\infty$，再令 $\varepsilon\downarrow0$。
\end{proof}

\subsection{数列空间}
在 $\N$ 上取计数测度，即得
\[
 \ell^p=\{x=(x_j):\sum|x_j|^p<\infty\},\qquad
 \ell^\infty=\{x:\sup_j|x_j|<\infty\}.
\]
另外 $c_0=\{x:x_j\to0\}$、$c=\{x:x_j\text{ 有有限极限}\}$ 均取上确界范数。它们是 $\ell^\infty$ 的闭子空间：一致极限保留趋零或存在极限的条件。有限支集数列在 $c_0$ 与 $\ell^p$（$p<\infty$）中稠密，却不在 $\ell^\infty$ 中稠密，常数数列 $(1,1,\ldots)$ 与任何有限支集数列距离都是 $1$。

\subsection{带符号测度、Jordan 分解与 Hahn 分解}
以下先取 $\sigma$-代数 $\Sigma$。原书用 $\sigma$-环表述的情形可对可测集合分别限制处理；为使全空间上的公式直接有意义，此处假定 $S\in\Sigma$。一个有限值带符号测度 $\nu:\Sigma\to\R$ 是可数可加的集合函数。有限值指每个可测集的值都是有限实数，随后将证明其全变差有限。

\begin{lemma}[有限值带符号测度的值域有界]\label{lem:y1-measure-bounded}
$\sup_{A\in\Sigma}|\nu(A)|<\infty$。
\end{lemma}
\begin{proof}
记 $V(A)=\sup_{B\subset A}|\nu(B)|$。若 $V(S)=\infty$，归纳构造递减 $E_j$，使 $V(E_j)=\infty$ 且 $|\nu(E_j)|\ge j$。给定 $E_j$，选 $B\subset E_j$，$|\nu(B)|>|\nu(E_j)|+j+1$。$B$ 与 $E_j\setminus B$ 的测度绝对值都大于 $j+1$，并且至少一者有无穷大的 $V$；选该者为 $E_{j+1}$。可是可数可加性应用于不交分解
\[
 E_1=\left(\bigcap_jE_j\right)\mathbin{\dot\cup}
 \bigcup_j(E_j\setminus E_{j+1})
\]
给出 $\nu(E_j)\to\nu(\bigcap_jE_j)$，与绝对值无界矛盾。
\end{proof}

\begin{theorem}[Jordan 与 Hahn 分解]\label{thm:y1-jordan}
定义
\[
 \nu^+(A)=\sup_{B\subset A}\nu(B),\qquad
 \nu^-(A)=-\inf_{B\subset A}\nu(B).
\]
则 $\nu^+,\nu^-$ 是有限正测度，$\nu=\nu^+-\nu^-$。存在可测分解 $S=P\mathbin{\dot\cup}N$，使 $P$ 的每个可测子集测度非负，$N$ 的每个可测子集测度非正；而且 $\nu^-(P)=0$、$\nu^+(N)=0$。这种分解称为 Hahn 分解。
\end{theorem}
\begin{proof}
上一引理保证上下确界有限。对不交 $A_j$，任意 $B\subset\bigcup A_j$ 有
\[
 \nu(B)=\sum_j\nu(B\cap A_j)\le\sum_j\nu^+(A_j),
\]
故 $\nu^+(\bigcup A_j)\le\sum\nu^+(A_j)$。反向对前 $m$ 个 $A_j$ 各取逼近上确界的子集，将它们合并，误差取小于任意 $\varepsilon$；再令 $\varepsilon\downarrow0$、$m\to\infty$。这证明可数可加性。对 $-\nu$ 同理。由 $\nu(B)=\nu(A)-\nu(A\setminus B)$ 对 $B\subset A$ 取上确界，得 $\nu^+(A)=\nu(A)+\nu^-(A)$。

选 $B_j$ 使 $\nu(B_j)>\nu^+(S)-2^{-j}$。由分解公式得 $\nu^-(B_j)<2^{-j}$、$\nu^+(S\setminus B_j)<2^{-j}$。令 $P=\liminf B_j$、$N=S\setminus P$。对每个 $m$，
\[
 N\subset\bigcup_{j\ge m}(S\setminus B_j),\qquad
 \nu^+(N)\le\sum_{j\ge m}2^{-j}\to0.
\]
对 $P_m=\bigcap_{j\ge m}B_j$，有 $\nu^-(P_m)\le\inf_{j\ge m}2^{-j}=0$；$P=\bigcup_mP_m$，故 $\nu^-(P)=0$。这正是所需的符号性质。
\end{proof}
原书以负值的“负变差”记下确界；此处 $\nu^-$ 取非负值，因此 Jordan 公式使用减号。这个记号改变不改变数学内容。

\begin{proposition}[全变差范数]\label{prop:y1-variation}
令 $|\nu|=\nu^++\nu^-$，则
\[
 \|\nu\|_{\mathrm{TV}}=|\nu|(S)
 =\sup_{\substack{h\text{ 可测}\\|h|\le1}}
 \left|\int h\,d\nu\right|
 =\sup_{\{A_j\}}\sum_j|\nu(A_j)|,
\]
最后的上确界取遍有限可测分割。有限带符号测度组成赋范实线性空间。有限复测度也组成赋范复线性空间，其范数由同一分割公式给出，积分上确界中允许复值 $h$。
\end{proposition}
\begin{proof}
实情形中 $|\int h\,d\nu|\le\int|h|\,d|\nu|\le|\nu|(S)$；取 $h=1_P-1_N$ 达到上界。每个分割的和不超过全变差，而分割 $P,N$ 达到上界。复情形中 $\nu=\nu_1+i\nu_2$，由前述实测度结论知分割和不超过 $|\nu_1|(S)+|\nu_2|(S)<\infty$。对一个有限分割，取简单函数在 $A_j$ 上为 $\overline{\nu(A_j)}/|\nu(A_j)|$（零值时取零），积分等于该分割和。反向对简单 $h$ 逐项估计，再用一致简单函数逼近一般有界可测 $h$，得积分公式。两种情形下范数的三角性及齐次性来自积分公式；若范数为零，取 $h=1_A$ 得 $\nu(A)=0$ 对所有 $A$ 成立。
\end{proof}

\begin{exercise}[补充：全变差不是总质量的绝对值]
在两点空间取 $\nu=\delta_a-\delta_b$。计算 $\nu(S)$ 和 $\|\nu\|_{\mathrm{TV}}$，并用积分上确界解释两者为什么不同。
\end{exercise}
\begin{solution}
总质量为零而全变差为二。选择 $h(a)=1,h(b)=-1$，则 $\int h\,d\nu=2$；全变差可以分别选择不同位置的符号，避免正负质量抵消。
\end{solution}

\section{几例拟赋范线性空间}\label{sec:yosida-01-04}
\noindent\textit{原书位置：中译本第 39--40 页；英文版第 38--39 页。}

\subsection{局部光滑拓扑}
按定理 \ref{thm:y1-countable-metric}，$\mathcal E^k(\Omega)$ 和固定紧集上的 $\Test_K$ 都可取距离
\[
 d(f,g)=\sum_{j=1}^\infty2^{-j}\frac{p_j(f-g)}{1+p_j(f-g)},
\]
其中 $p_j$ 枚举相应的半范数。此拟范数可有统一上界 $1$，但这不意味着全空间是拓扑有界集：拓扑有界性涉及标量缩放与所有邻域，而不是某个距离数值的上界。

\subsection{有限测度空间上的测度收敛}
\begin{definition}[可测函数空间]
设 $\mu(S)<\infty$。用 $L^0(S,\mu)$（原书记 $M(S,\Sigma,\mu)$）表示几乎处处有限的可测函数按几乎处处相等组成的空间。定义
\[
 q(f)=\int_S\frac{|f|}{1+|f|}\,d\mu,\qquad d(f,g)=q(f-g).
\]
\end{definition}
\begin{proposition}[积分距离与测度收敛]\label{prop:y1-measure-convergence}
上述 $q$ 是吉田意义的拟范数，并且 $d(f_j,f)\to0$ 当且仅当对每个 $\varepsilon>0$，
\[
 \mu\{|f_j-f|\ge\varepsilon\}\longrightarrow0.
\]
\end{proposition}
\begin{proof}
正定性在等价类上成立。三角不等式使用 $t/(1+t)$ 的单调性与次可加性。固定 $f$ 时，$q(af)\to0$ 随 $a\to0$ 由控制收敛得到，控制函数是可积的常数 $1$。

令 $A_\varepsilon=\{|u|\ge\varepsilon\}$，则
\[
 \frac{\varepsilon}{1+\varepsilon}\mu(A_\varepsilon)
 \le q(u)\le\mu(A_\varepsilon)
 +\frac{\varepsilon}{1+\varepsilon}\mu(S).
\]
第一式证明距离收敛推出测度收敛；第二式先令序列指标趋于无穷，再令 $\varepsilon\downarrow0$，证明反向。固定 $a\ne0$ 时，$\{|au|\ge\varepsilon\}=\{|u|\ge\varepsilon/|a|\}$，故第二条数乘条件成立。联合连续性也可引用引理 \ref{lem:y1-uniform-scalars}。
\end{proof}
原书还给出同一拓扑的另一种拟范数。令
\[
 r(f)=\inf_{\varepsilon>0}\bigl(\varepsilon+\mu\{|f|\ge\varepsilon\}\bigr),
 \qquad \widetilde q(f)=\arctan r(f).
\]
由 $\{|f+g|\ge a+b\}\subset\{|f|\ge a\}\cup\{|g|\ge b\}$ 得 $r(f+g)\le r(f)+r(g)$；$\arctan$ 在非负轴上增加且次可加，故 $\widetilde q$ 满足三角不等式。$r(f_j)\to0$ 与测度收敛等价：一向取固定阈值，另一向选择接近下确界的小阈值，再比较较大阈值的异常集。于是它也满足数乘条件，并定义相同拓扑。

\begin{example}[测度收敛与几乎处处收敛]
在 $[0,1)$ 上依次枚举第 $m$ 层的二进区间
\[
 I_{m,k}=[k2^{-m},(k+1)2^{-m}),\quad 0\le k<2^m,
\]
取相应示性函数作为一个序列。其积分距离趋于零，却在每个点上都无极限，因为每层恰有一次取值为 $1$，其余大量次取值为 $0$。相反，有限测度空间上的几乎处处收敛必推出测度收敛，由控制收敛应用于异常集示性函数即可。
\end{example}
测度收敛序列可以抽出几乎处处收敛的子列：选 $j_k$ 使 $\mu\{|f_{j_k}-f|>2^{-k}\}<2^{-k}$，可数次可加性给出异常事件的上极限为零测集。这一选子列技术也用于后面的完备性证明。

\begin{remark}[补充：非局部凸性需要测度条件]
在 $[0,1]$ 的 Lebesgue 测度上，$L^0$ 的测度拓扑并非局部凸。事实上其连续线性泛函只能为零：连续性给出某个 $\delta>0$，使 $q(f)<\delta$ 时 $|L(f)|<1$。任何支集测度小于 $\delta$ 的 $f$，其任意标量倍 $tf$ 都有 $q(tf)\le\mu(\supp f)<\delta$，故 $L(f)=0$。把 $[0,1]$ 分成有限个长度小于 $\delta$ 的区间，将任意 $f$ 分解于这些区间上，即得 $L(f)=0$。若此非零 Hausdorff 空间局部凸，Hahn--Banach 定理将给出分离点的连续线性泛函，矛盾。有限原子空间则可局部凸，所以不能对任意测度空间一概作此断言。
\end{remark}

\section{准 Hilbert 空间}\label{sec:yosida-01-05}
\noindent\textit{原书位置：中译本第 40--44 页；英文版第 39--42 页。}

\subsection{平行四边形恒等式与极化}
\begin{definition}[准 Hilbert 空间]
赋范空间若满足
\[
 \|x+y\|^2+\|x-y\|^2=2\|x\|^2+2\|y\|^2,
\]
称为准 Hilbert 空间，亦称预 Hilbert 空间。这里尚不要求完备。
\end{definition}
\begin{theorem}[Jordan--von Neumann 极化定理]\label{thm:y1-polarization}
实赋范空间满足平行四边形恒等式，当且仅当其范数由唯一实内积生成，此内积为
\[
 B(x,y)=\frac14\bigl(\|x+y\|^2-\|x-y\|^2\bigr).
\]
复情形中，以第一变量线性为约定，唯一内积为
\[
 (x,y)=\frac14\left(\|x+y\|^2-\|x-y\|^2
 +i\|x+iy\|^2-i\|x-iy\|^2\right).
\]
\end{theorem}
\begin{proof}
实情形中对称性及 $B(x,x)=\|x\|^2$ 直接成立。平行四边形恒等式分别用于 $x+z,y$ 和 $x-z,y$，相减得到
\[
 B(x+y,z)+B(x-y,z)=2B(x,z).
\]
取 $y=x$ 得 $B(2x,z)=2B(x,z)$。再令 $x=(u+v)/2$、$y=(u-v)/2$，得到 $B(u+v,z)=B(u,z)+B(v,z)$。于是整数及有理数齐次性成立；$B$ 的定义与范数连续性使有理数逼近推广为实齐次性。这完整证明了双线性。

复空间先视为实空间，得到 $B$。由 $\|ix\|=\|x\|$ 可知 $B(ix,iy)=B(x,y)$，所以 $B(ix,y)=-B(x,iy)$，且 $B(x,ix)=0$。定义 $(x,y)=B(x,y)+iB(x,iy)$，便有 $(ix,y)=i(x,y)$、$(y,x)=\overline{(x,y)}$ 及 $(x,x)=\|x\|^2$。这证明复内积性质。内积展开 $\|x\pm y\|^2$ 和 $\|x\pm iy\|^2$ 反过来给出极化公式，故唯一性与必要性都成立。
\end{proof}

\begin{theorem}[内积生成范数与 Cauchy--Schwarz 不等式]\label{thm:y1-schwarz}
若 $(x,y)$ 是正定的实或复内积，则 $\|x\|=(x,x)^{1/2}$ 是范数，满足平行四边形恒等式，并且
\[
 |(x,y)|\le\|x\|\|y\|.
\]
等号成立当且仅当 $x,y$ 线性相关。
\end{theorem}
\begin{proof}
$y=0$ 时显然。否则令 $c=(x,y)/\|y\|^2$，利用第一变量线性展开
\[
 0\le\|x-cy\|^2=\|x\|^2-\frac{|(x,y)|^2}{\|y\|^2}.
\]
等号恰在 $x=cy$ 时成立。因而
\[
 \|x+y\|^2=\|x\|^2+2\operatorname{Re}(x,y)+\|y\|^2
 \le(\|x\|+\|y\|)^2.
\]
这是三角不等式；其余范数公理及平行四边形恒等式直接由内积展开得到。
\end{proof}
\begin{remark}[补充：完备性与共轭位置]
仅有内积公理不能推出 Hilbert 空间，还须证明完备性；原书中译本这一处的表述应按英文版的 pre-Hilbert 理解。极化公式中虚部的符号取决于哪一个变量线性。本稿内积第一变量线性，而分布对测试函数线性，二者不应混用。
\end{remark}

\subsection{基本例子与 Sobolev 型内积}
$L^2(S,\mu)$ 和 $\ell^2$ 分别取
\[
 (f,g)=\int f\overline g\,d\mu,\qquad
 (x,y)=\sum_jx_j\overline{y_j}.
\]
积分与级数的绝对收敛由 Hölder 保证。对 $k\in\N_0$，令 $\widehat H^k(\Omega)$ 为满足以下有限范数的 $C^k$ 函数：
\[
 \|f\|_{k,2}^2=\sum_{|\alpha|\le k}\int_\Omega|\partial^\alpha f|^2\,dx,
 \qquad (f,g)_{k,2}=\sum_{|\alpha|\le k}\int_\Omega
 \partial^\alpha f\,\overline{\partial^\alpha g}\,dx.
\]
它是准 Hilbert 空间，$C_c^\infty(\Omega)$ 也可用此内积赋予准 Hilbert 结构。它们一般尚不完备；后面将用弱导数解释其完备化。原书也考虑 $C_c^k$，其完备化相同：先用正则化在全部有限阶导数范数中逼近，即可将 $C_c^k$ 换成 $C_c^\infty$。

\subsection{Bergman 空间与 Hardy--Lebesgue 类}
平面开集 $G$ 上的 Bergman 空间为
\[
 A^2(G)=\{f\text{ 在 }G\text{ 全纯}:\int_G|f|^2\,dA<\infty\},
 \qquad(f,g)=\int_Gf\overline g\,dA.
\]
它继承 $L^2$ 内积。第 \ref{sec:yosida-01-09} 节将证明全纯性在 $L^2$ 极限中保留，从而得到完备性。

单位圆盘 $\mathbb D$ 上的 Hardy--Lebesgue 类（现代记作 $H^2(\mathbb D)$）由满足
\[
 \|f\|_{H^2}^2=\sup_{0<r<1}\frac1{2\pi}
 \int_0^{2\pi}|f(re^{i\theta})|^2\,d\theta<\infty
\]
的全纯函数组成。
\begin{theorem}[Hardy 类与系数空间的等距对应]\label{thm:y1-hardy}
若 $f(z)=\sum_{j=0}^\infty c_jz^j$，则 $f\in H^2$ 当且仅当 $(c_j)\in\ell^2$，且 $\|f\|_{H^2}^2=\sum|c_j|^2$。系数映射是线性等距双射，因此内积为 $(f,g)=\sum c_j\overline{d_j}$。
\end{theorem}
\begin{proof}
在每个固定 $r<1$ 上，Taylor 级数一致绝对收敛，可逐项相乘积分。圆周指数函数的正交性给出
\[
 \frac1{2\pi}\int_0^{2\pi}|f(re^{i\theta})|^2\,d\theta
 =\sum_{j=0}^\infty|c_j|^2r^{2j}.
\]
令 $r\uparrow1$，单调收敛即得范数公式。反过来，给定 $\ell^2$ 系数，Cauchy--Schwarz 给出对 $|z|\le r<1$，
\[
 \sum_{j\ge N}|c_jz^j|
 \le\left(\sum_{j\ge N}|c_j|^2\right)^{1/2}
 \frac{r^N}{\sqrt{1-r^2}}.
\]
因此级数在紧子圆盘上一致收敛，定义全纯函数，前一公式保证其属于 $H^2$。这也证明满射性。
\end{proof}

\begin{exercise}[补充：哪些序列范数来自内积]
在 $\mathbb K^2$ 上取 $\ell^p$ 范数。用 $(1,0)$ 和 $(0,1)$ 验证：$1\le p\le\infty$ 中只有 $p=2$ 满足平行四边形恒等式。
\end{exercise}
\begin{solution}
有限 $p$ 时，左边为 $2\cdot2^{2/p}$，右边为 $4$，故 $p=2$。$p=\infty$ 时左边为 $2$，也不等于 $4$。$p=2$ 的范数由标准内积生成，故对所有向量成立。
\end{solution}

\section{线性算子的连续性}\label{sec:yosida-01-06}
\noindent\textit{原书位置：中译本第 44--46 页；英文版第 42--44 页。}

线性算子的定义域 $D(T)$ 是线性子空间，取原空间的诱导拓扑。由 $T(x+h)-Tx=Th$，在零点连续等价于在定义域每一点连续。

\begin{theorem}[连续线性算子的半范数估计]\label{thm:y1-operator-seminorm}
设 $X,Y$ 分别由半范数族 $\mathcal P,\mathcal Q$ 定义拓扑。线性映射 $T:D(T)\subset X\to Y$ 连续，当且仅当对每个 $q\in\mathcal Q$，存在 $p_1,\ldots,p_r\in\mathcal P$ 与 $C>0$，使
\[
 q(Tx)\le C\max_{1\le j\le r}p_j(x)\quad(x\in D(T)).
\]
若 $\mathcal P$ 对有限最大值封闭，则可简写为 $q(Tx)\le Cp(x)$。
\end{theorem}
\begin{proof}
连续性给出有限半范数邻域，其与定义域的交映入 $\{q<1\}$。令 $p(x)=\max_jp_j(x)/\varepsilon_j$。若 $p(x)>0$，对每个 $t>p(x)$，$x/t$ 落在该邻域中，于是 $q(Tx)<t$，令 $t\downarrow p(x)$。若 $p(x)=0$，所有倍数 $ax$ 都在该邻域中，$|a|q(Tx)<1$ 对任意 $a$ 成立，故 $q(Tx)=0$。反向估计把每个有限的目标邻域条件转成有限的源邻域条件，因此在零点连续。
\end{proof}
\begin{corollary}[线性泛函、范数估计与逆算子]
局部凸空间上的线性泛函连续，当且仅当其绝对值受有限个连续半范数的最大值控制。赋范空间之间的线性算子连续，当且仅当 $\|Tx\|\le C\|x\|$。而逆映射 $T^{-1}:R(T)\to D(T)$ 存在且连续，当且仅当存在 $c>0$，使
\[
 \|Tx\|\ge c\|x\|\quad(x\in D(T)).
\]
\end{corollary}
\begin{proof}
前两项分别以标量绝对值及单一范数作为目标半范数。后一项的下界首先使 $T$ 单射；把它改写成 $\|T^{-1}y\|\le c^{-1}\|y\|$ 即得连续性，反向同理。这里逆算子只定义于值域，尚未断言值域是整个目标空间。
\end{proof}

\begin{definition}[算子范数]
全定义的有界线性算子 $T:X\to Y$ 的算子范数为
\[
 \|T\|=\sup_{\|x\|\le1}\|Tx\|
 =\inf\{C\ge0:\|Tx\|\le C\|x\|\ \forall x\}.
\]
当 $X\ne\{0\}$ 时，还等于 $\sup_{\|x\|=1}\|Tx\|$。此处“有界”指把单位球映成有界集，等价于连续；不指在整个 $X$ 上函数值有统一上界。
\end{definition}
若 $S,T$ 定义域不同，则 $D(S+T)=D(S)\cap D(T)$，$D(ST)=\{x\in D(T):Tx\in D(S)\}$。对全定义有界算子，
\[
 \|S+T\|\le\|S\|+\|T\|,\quad
 \|aT\|=|a|\|T\|,\quad \|ST\|\le\|S\|\|T\|.
\]
这些公式逐点估计后在单位球取上确界即可；例如 $\|STx\|\le\|S\|\|Tx\|\le\|S\|\|T\|\|x\|$。于是 $\|T^m\|\le\|T\|^m$，$T^0=I$。

\begin{example}[算子的非交换性与定义域]
在 $L^2(\R)$ 上，乘法算子 $Mu(t)=tu(t)$ 的定义域为 $\{u:tu\in L^2\}$，微分算子 $Du=iu'$ 可取定义域 $H^1(\R)$。在共同的不变子空间 $\Test(\R)$ 上直接算得
\[
 (DM-MD)u=i(tu)'-tiu'=iu.
\]
这两个算子在 $L^2$ 范数下都不是全定义有界算子。给出共同定义域后，交换关系才有明确含义。
\end{example}

\begin{proposition}[补充：有界下界与闭值域]
若 $X$ 是 Banach 空间，$T:X\to Y$ 有界且 $\|Tx\|\ge c\|x\|$，则 $R(T)$ 在赋范空间 $Y$ 中闭。
\end{proposition}
\begin{proof}
若 $Tx_j\to y$，下界使 $x_j$ 是 Cauchy 序列，故趋于 $x\in X$。连续性给出 $y=Tx$。源空间的完备性在此不可省去，例如未完备空间到其完备化的等距包含有稠密但不闭的值域。
\end{proof}

\section{有界集和囿空间}\label{sec:yosida-01-07}
\noindent\textit{原书位置：中译本第 46--48 页；英文版第 44--46 页。}

\begin{definition}[拓扑有界集与囿空间]
拓扑向量空间中的 $B$ 称为有界集，若每个零邻域 $U$ 都吸收 $B$：存在 $c>0$ 使 $B\subset cU$（可先把 $U$ 换成均衡零邻域）。局部凸空间称为囿空间或 bornological 空间，若每个吸收全部有界集的绝对凸集都是零邻域。
\end{definition}
\begin{proposition}[有界集的半范数判别]\label{prop:y1-bounded}
局部凸空间中，$B$ 有界当且仅当对每个连续半范数 $p$，$\sup_{x\in B}p(x)<\infty$。连续线性映射把有界集映为有界集。紧集以及收敛序列连同其极限均有界。
\end{proposition}
\begin{proof}
有界性应用于 $\{p<1\}$ 得一向；反向对一个含有限个半范数条件的零邻域，选择足够大的共同缩放常数即可。若 $T$ 连续，给定目标零邻域 $V$，取源零邻域 $U$ 使 $T(U)\subset V$，则 $B\subset cU$ 推出 $T(B)\subset cV$。紧集的有界性由零邻域的吸收性和有限子覆盖证明；收敛序列加极限是紧集。这最后一条对一般拓扑向量空间也成立。
\end{proof}

\begin{theorem}[囿空间的等价刻画]\label{thm:y1-bornological}
局部凸空间 $X$ 是囿空间，当且仅当每个在全部有界集上有界的半范数都连续。因而从囿空间到局部凸空间的线性映射连续，当且仅当它把有界集映为有界集。
\end{theorem}
\begin{proof}
设 $X$ 囿，半范数 $p$ 在各有界集上有界。则 $\{p\le1\}$ 吸收所有有界集，故是零邻域；由半范数的缩放得 $p$ 连续。反向设每个这种半范数连续。若绝对凸集 $U$ 吸收全部有界集，它尤其吸收每个单点，因此 $p_U$ 有定义；吸收有界集保证 $p_U$ 在有界集上有界，于是连续，$\{p_U<1\}\subset U$ 证明 $U$ 是邻域。

对线性映射 $T:X\to Y$，若它保有界集，则每个连续目标半范数 $q$ 的复合 $q\circ T$ 在源有界集上有界，故连续；半范数刻画推出 $T$ 连续。反向已由上一命题证明。
\end{proof}
赋范空间都是囿空间：若半范数 $p$ 在范数单位球上有界于 $C$，则 $p(x)\le C\|x\|$，所以连续。不要求赋范空间完备。

\begin{theorem}[拟赋范空间上保有界集的线性映射]
设 $X,Y$ 均为定义 \ref{def:y1-Fnorm} 意义下的拟赋范空间。线性映射 $T:X\to Y$ 连续，当且仅当它保拓扑有界集。
\end{theorem}
\begin{proof}
只证明充分性。若 $x_j\to0$，可选正整数 $n_j\to\infty$，使 $n_jx_j\to0$：对每个整数 $k$，固定数乘的连续性给出 $kx_j\to0$，选择递增的 $N_k$ 使 $j\ge N_k$ 时 $q_X(kx_j)<1/k$，再令 $n_j=k$ 对 $N_k\le j<N_{k+1}$。于是 $\{n_jx_j\}$ 有界，故 $\{n_jTx_j\}$ 有界。拓扑有界性保证有界集乘以趋零标量时一致趋于零，因此 $Tx_j=n_j^{-1}(n_jTx_j)\to0$。度量空间上线性映射的序列连续等价于连续，故结论成立。
\end{proof}
\begin{corollary}[补充：可度量局部凸空间是囿空间]
每个可度量的局部凸空间都是囿空间。
\end{corollary}
\begin{proof}
对在有界集上有界的半范数 $p$，将它看作线性映射 $X\to X/\ker p$ 的目标范数。此映射保有界集，上一定理（目标赋范）保证连续，故 $p$ 连续，再引用定理 \ref{thm:y1-bornological}。
\end{proof}
囿空间关注“有界性足以检测连续性”；它与以后出现的桶空间不同。后者的定义要求所有闭的吸收绝对凸集都是邻域。两者的假设和用途需要分别检查。

\section{广义函数和广义导数}\label{sec:yosida-01-08}
\noindent\textit{原书位置：中译本第 48--54 页；英文版第 46--52 页。}

\subsection{测试函数空间的有界集与囿性}
\begin{theorem}[测试空间有界集的完整刻画]\label{thm:y1-D-bounded}
$B\subset\Test(\Omega)$ 有界，当且仅当存在共同紧集 $K\Subset\Omega$，使 $B\subset\Test_K$，且对每个 $m$，
\[
 \sup_{\varphi\in B}p_{K,m}(\varphi)<\infty.
\]
$\Test(\Omega)$ 是囿空间。
\end{theorem}
\begin{proof}
若无共同紧集，与定理 \ref{thm:y1-D-sequence} 同样选出 $\varphi_j\in B$ 及无域内聚点的非零取值点 $x_j$。半范数
\[
 q(\psi)=\sum_{j=1}^\infty j\frac{|\psi(x_j)|}{|\varphi_j(x_j)|}
\]
各层限制为有限和，故连续；却有 $q(\varphi_j)\ge j$，与有界性矛盾。各阶导数的一致有界由全局连续半范数 $\sup_\Omega|\partial^\alpha\psi|$ 得到。反向给定任意连续半范数 $q$，其在 $\Test_K$ 上由有限个 $p_{K,m}$ 控制，所以在 $B$ 上有界，命题 \ref{prop:y1-bounded} 即给出充分性。

最后，若半范数 $q$ 在 $\Test$ 的有界集上有界，其在每个 $\Test_K$ 上也保有界性，因为包含映射连续。$\Test_K$ 可度量且局部凸，故囿，$q|_{\Test_K}$ 连续。归纳极限判别给出 $q$ 连续，于是 $\Test$ 囿。这个证明无需把不可度量的 $\Test$ 当作度量空间。
\end{proof}

\subsection{定义、局部有限阶估计与典型分布}
\begin{definition}[广义函数或分布]
$\Omega$ 上的分布是连续线性泛函 $T:\Test(\Omega)\to\mathbb K$；全体分布记作 $\Dist(\Omega)$。配对记为 $T(\varphi)=\Pair{T}{\varphi}$。
\end{definition}
\begin{theorem}[分布的局部有限阶判别]\label{thm:y1-distribution-estimate}
线性泛函 $T$ 是分布，当且仅当对每个 $K\Subset\Omega$，存在 $m\in\N_0$ 与 $C_K>0$，使
\[
 |\Pair{T}{\varphi}|\le C_Kp_{K,m}(\varphi)\quad(\varphi\in\Test_K).
\]
它也等价于 $T$ 在 $\Test$ 的每个有界集上有界。
\end{theorem}
\begin{proof}
连续性限制到 $\Test_K$，用第 \ref{sec:yosida-01-06} 节的半范数估计；$p_{K,m}$ 随 $m$ 增大，所以有限最大值可用一个较高阶半范数替代。反向由归纳极限判别得到。最后一条由 $\Test$ 的囿性以及半范数 $|T(\cdot)|$ 的判别得到。
\end{proof}
阶数 $m$ 可随紧集变化；一般分布未必有一个对整个 $\Omega$ 通用的有限阶数。

\begin{example}[局部可积函数、测度和 Dirac 分布]
对 $f\in L^1_{\Loc}(\Omega)$，令 $T_f(\varphi)=\int_\Omega f\varphi\,dx$。其估计为 $|T_f(\varphi)|\le\|f\|_{L^1(K)}p_{K,0}(\varphi)$。若复 Radon 测度 $\nu$ 的全变差在紧集上有限，则 $T_\nu(\varphi)=\int\varphi\,d\nu$ 同样定义分布。特别地，$\delta_a(\varphi)=\varphi(a)$ 是集中于 $a$ 的 Dirac 分布。仅有 $\sigma$-有限性不能代替局部有限全变差条件。
\end{example}

\begin{proposition}[局部可积函数到分布的嵌入是单射]\label{prop:y1-regular-injective}
若 $T_f=0$，则 $f=0$ 几乎处处。因而 $T_f=T_g$ 等价于 $f=g$ 几乎处处。
\end{proposition}
\begin{proof}
给定任意 $V\Subset\Omega$ 的开集，在包含 $\overline V$ 的紧邻域内取足够小的 $\varepsilon$，则对 $x\in V$，$y\mapsto\rho_\varepsilon(x-y)$ 是测试函数，所以 $f*\rho_\varepsilon(x)=0$。局部可积函数的正则化在 $L^1(V)$ 中趋于 $f$。这一逼近可由局部截断化为 $L^1(\R^n)$，再用
\[
 \|f*\rho_\varepsilon-f\|_1
 \le\int\rho_\varepsilon(h)\|f(\cdot-h)-f\|_1\,dh
\]
及 $L^1$ 平移连续性证明；后者由简单函数的矩形逼近得到。因此 $f=0$ 在 $V$ 中几乎处处，取可数穷竭完成证明。这避免将局部可积的 $f$ 在整个无界域上当成有限测度。
\end{proof}

\subsection{广义导数与跳跃项}
\begin{definition}[分布导数]
对任意多重指标 $\alpha$，定义
\[
 \Pair{\partial^\alpha T}{\varphi}
 =(-1)^{|\alpha|}\Pair{T}{\partial^\alpha\varphi}.
\]
\end{definition}
这是分布，因为测试函数的微分保持紧支集，且提高有限阶估计中的阶数。各分布都能任意次求导，偏导相互交换。若 $f\in C^1(\Omega)$，分部积分给出 $\partial_jT_f=T_{\partial_jf}$；测试函数在边界附近为零，所以不产生边界项。

\begin{example}[Heaviside 函数与分段光滑函数]
对 $H=1_{(0,\infty)}$，
\[
 \Pair{H'}{\varphi}=-\int_0^\infty\varphi'(x)\,dx=\varphi(0),\qquad H'=\delta_0.
\]
若 $f$ 在有限个点 $a_j$ 外分段 $C^1$，具有有限单侧极限，且分段导数局部可积，则
\[
 \partial T_f=T_{f'_{\mathrm{pw}}}
 +\sum_j\bigl(f(a_j+)-f(a_j-)\bigr)\delta_{a_j}.
\]
证明是在各光滑区间分部积分，相邻端点项合并为右值减左值。点值本身任意更改不影响公式。特别地，$(|x|)''=2\delta_0$。
\end{example}

\begin{proposition}[边界与 Green 公式]\label{prop:y1-boundary-distribution}
设 $U\Subset\R^n$ 有光滑边界，$f\in C^2(\overline U)$，并把 $f$ 延零到 $U$ 外。取外法向 $n$，则
\[
 \partial_jT_{f1_U}=T_{(\partial_jf)1_U}
 -f n_j\,dS\big|_{\partial U},
\]
且对测试函数 $\varphi$，
\[
 \Pair{\Delta T_{f1_U}}{\varphi}
 =\int_U\Delta f\,\varphi\,dx
 +\int_{\partial U}\bigl(f\partial_n\varphi-(\partial_nf)\varphi\bigr)\,dS.
\]
\end{proposition}
\begin{proof}
散度定理给出 $-\int_U f\partial_j\varphi=\int_U\partial_jf\,\varphi-\int_{\partial U}fn_j\varphi$。Green 第二恒等式给出 $\int_U f\Delta\varphi=\int_U\Delta f\,\varphi+\int_{\partial U}(f\partial_n\varphi-\partial_nf\,\varphi)$。原书取内法向时两条边界符号相应反转；必须先固定法向约定。
\end{proof}

\subsection{光滑乘子与广义 Leibniz 公式}
对 $a\in C^\infty(\Omega)$ 定义 $\Pair{aT}{\varphi}=\Pair{T}{a\varphi}$。Leibniz 公式给出 $p_{K,m}(a\varphi)\le C_{a,K,m}p_{K,m}(\varphi)$，所以定义的是分布，且
\[
 \partial_j(aT)=(\partial_ja)T+a\partial_jT.
\]
两边作用于 $\varphi$ 后，只需将 $\partial_j(a\varphi)$ 展开即得。

\begin{theorem}[Hörmander 的广义 Leibniz 公式]\label{thm:y1-general-leibniz}
为与后面的 Fourier 记号衔接，此处临时记 $D_j=(1/i)\partial_j$。若 $P(\xi)$ 是多项式，$a$ 光滑，$T$ 是分布，则
\[
 P(D)(aT)=\sum_\beta\frac{D^\beta a}{\beta!}
 (\partial_\xi^\beta P)(D)T.
\]
右边是有限和。
\end{theorem}
\begin{proof}
先对单项式 $P(\xi)=\xi^\alpha$ 逐次应用乘积求导，得
\[
 D^\alpha(aT)=\sum_{\beta\le\alpha}\binom\alpha\beta
 (D^\beta a)D^{\alpha-\beta}T.
\]
而 $\partial_\xi^\beta\xi^\alpha/\beta!=\binom\alpha\beta\xi^{\alpha-\beta}$，所以单项式情形成立。再以线性组合得到一般多项式。这补足了原书利用指数函数和多项式 Taylor 恒等式辨认系数的论证。
\end{proof}
一般两个分布之间没有自然乘法：例如 $\delta_0^2$ 不由以上定义产生。对光滑乘子的要求是保证乘积仍把每个测试函数送入测试空间。

\begin{exercise}[补充：Dirac 分布与乘子]
证明在 $\R$ 上，$x\delta_0=0$、$x\delta_0'=-\delta_0$，并求 $a\delta_0'$ 的一般表达式。
\end{exercise}
\begin{solution}
分别作用于 $\varphi$。$x\delta_0$ 给出零；$x\delta_0'$ 给出 $-(x\varphi)'(0)=-\varphi(0)$；一般有 $a\delta_0'=a(0)\delta_0'-a'(0)\delta_0$。由此可见不能将 Dirac 分布误当成只在一点有非零数值的普通函数。
\end{solution}
进一步的拓扑、限制、运算与例子参见第 \ref{ch:distributions}、\ref{ch:operations} 章。

\section{B-空间和 F-空间}\label{sec:yosida-01-09}
\noindent\textit{原书位置：中译本第 54--58 页；英文版第 52--56 页。}

\subsection{完备性与术语}
\begin{definition}[B-空间、F-空间与 Hilbert 空间]
在平移不变距离下，每个 Cauchy 序列均收敛的空间称为完备。完备赋范空间称为 Banach 空间，原书简称 B-空间；完备拟赋范空间称为 F-空间。完备准 Hilbert 空间称为 Hilbert 空间。现代通常把“Fréchet 空间”限定为完备、可度量的局部凸空间；一般 F-空间不必局部凸。
\end{definition}
收敛序列总是 Cauchy；极限唯一，因为 $d(x,y)\le d(x,x_j)+d(x_j,y)$。完备性要求极限落在空间内，而不只落在某个更大的函数集合内。

\begin{theorem}[光滑函数空间与固定支集测试空间完备]\label{thm:y1-smooth-complete}
$\mathcal E^k(\Omega)$（$k$ 有限或无穷）以及 $\Test_K(\Omega)$ 都是 Fréchet 空间。
\end{theorem}
\begin{proof}
对 $\mathcal E^k$ 中的 Cauchy 序列 $f_j$，每个允许的 $\partial^\alpha f_j$ 在每个紧集上一致 Cauchy，故有连续局部极限 $g_\alpha$，不同紧集上的极限相容。需要证明 $g_\alpha$ 真是 $g_0$ 的导数。在域内任意小闭长方体上，对 $|\alpha|<k$ 用微积分基本定理写
\[
 \partial^\alpha f_j(x+te_l)-\partial^\alpha f_j(x)
 =\int_0^t\partial^{\alpha+e_l}f_j(x+se_l)\,ds.
\]
一致极限允许交换积分与极限，得同样的 $g$ 公式；再由连续性求导，得到 $\partial_lg_\alpha=g_{\alpha+e_l}$。归纳即有 $g_0\in C^k$，并在所有半范数下收敛。无限阶情形对每个有限阶重复即可。若原序列均支集于 $K$，域外 $\Omega\setminus K$ 上所有函数及其极限为零，故极限仍属于 $\Test_K$。这补足了“导数各有极限，所以极限光滑”中的关键相容性步骤。
\end{proof}

\subsection{可积函数空间的完备性}
\begin{theorem}[Riesz--Fischer 完备性]\label{thm:y1-Lp-complete}
对任意测度空间，$L^p$（$1\le p\le\infty$）都是 Banach 空间。若 $1\le p<\infty$，每个 $L^p$-Cauchy 序列都有一个几乎处处收敛于其范数极限的子列。
\end{theorem}
\begin{proof}
先设 $p<\infty$。从 Cauchy 序列抽出 $f_{j_k}$，使
\[
 \|f_{j_{k+1}}-f_{j_k}\|_p\le2^{-k}.
\]
令 $g_N=|f_{j_1}|+\sum_{k=1}^N|f_{j_{k+1}}-f_{j_k}|$。Minkowski 保证 $\|g_N\|_p\le\|f_{j_1}\|_p+1$；单调收敛或 Fatou 引理给出极限 $g\in L^p$。因此 $g<\infty$ 几乎处处，差分级数绝对收敛，定义几乎处处有限的
\[
 f=f_{j_1}+\sum_{k=1}^\infty(f_{j_{k+1}}-f_{j_k}).
\]
$|f|\le g$ 保证 $f\in L^p$。再次用 Fatou 与 Minkowski，
\[
 \|f-f_{j_k}\|_p\le\sum_{l\ge k}2^{-l}\longrightarrow0.
\]
原序列由 Cauchy 性和这一收敛子列也趋于 $f$。

对于 $p=\infty$，同样选择子列使差分范数不超过 $2^{-k}$。去掉可数个零测集后，每个差分逐点不超过该上界，故差分级数在余集上一致收敛。极限本质有界，且 $L^\infty$ 距离的尾项不超过 $\sum_{l\ge k}2^{-l}$；再由 Cauchy 性处理原序列。甚至原序列在去掉一个共同零测集后可取一致收敛的代表元。
\end{proof}
$C_b(S)$ 也完备：一致 Cauchy 序列逐点取极限，一致收敛保留连续性和有界性。由此 $c_0,c$ 是 Banach 空间；$\ell^p$ 是 $L^p$ 的计数测度特例。因而 $L^2,\ell^2$ 是 Hilbert 空间，定理 \ref{thm:y1-hardy} 又使 $H^2$ 成为 Hilbert 空间。

\begin{proposition}[补充：有界算子空间的完备性]
若 $Y$ 是 Banach 空间，则 $\mathcal L(X,Y)$ 在算子范数下是 Banach 空间，不要求 $X$ 完备。
\end{proposition}
\begin{proof}
算子范数 Cauchy 序列 $T_j$ 对每个 $x$ 使 $T_jx$ 在 $Y$ 中 Cauchy，定义 $Tx=\lim T_jx$。逐点极限保留线性；$\|T_j\|$ 统一有界，所以 $T$ 有界。给定 $\varepsilon$，$j,l\ge N$ 时 $\|(T_j-T_l)x\|\le\varepsilon\|x\|$，令 $l\to\infty$，再在单位球取上确界，得到 $\|T_j-T\|\le\varepsilon$。
\end{proof}

\subsection{Bergman 空间、测度收敛空间与全体数列}
\begin{proposition}[Bergman 空间完备]
$A^2(G)$ 是 $L^2(G)$ 的闭子空间，因而是 Hilbert 空间。
\end{proposition}
\begin{proof}
全纯函数在 $\overline B(z,r)\subset G$ 上的圆周平均或 Taylor 展开给出
\[
 |f(z)|^2\le\frac1{\pi r^2}\int_{B(z,r)}|f(w)|^2\,dA(w).
\]
例如展开 $f(z+\rho e^{i\theta})=\sum a_j\rho^je^{ij\theta}$，角积分消掉异次项，再对 $\rho\,d\rho$ 积分，常数项就给出右边至少 $|a_0|^2$。在任意紧集 $K\Subset G$ 上可取统一 $r>0$，所以 $L^2$-Cauchy 序列的全纯代表元在 $K$ 上一致 Cauchy。其局部一致极限 $f$ 全纯，可由 Cauchy 积分公式取极限证明。与此同时 $L^2$ 完备性给出范数极限 $g$，两种极限在每个紧子域上几乎处处一致，故 $g$ 有全纯代表 $f$。所以闭性成立。
\end{proof}

\begin{proposition}[测度收敛空间完备]
若 $\mu(S)<\infty$，则第 \ref{sec:yosida-01-04} 节的 $L^0(S,\mu)$ 是 F-空间。
\end{proposition}
\begin{proof}
从距离 Cauchy 序列抽出子列 $f_{j_k}$，使
\[
 \mu\{|f_{j_{k+1}}-f_{j_k}|>2^{-k}\}<2^{-k}.
\]
尾部异常集之并的测度不超过 $\sum_{l\ge k}2^{-l}$，趋于零；因此几乎处处只有有限次异常，差分级数的尾部绝对值受 $\sum2^{-l}$ 控制，得到有限极限 $f$。有限测度上的几乎处处收敛推出测度收敛，故子列在距离中趋于 $f$；原序列由 Cauchy 性得到同样结论。
\end{proof}
原书的空间 $(s)$，即所有实或复数列 $\mathbb K^{\N}$，在
\[
 q(x)=\sum_{j=1}^\infty2^{-j}\frac{|x_j|}{1+|x_j|}
\]
下是 F-空间且局部凸：Cauchy 性等价于各坐标 Cauchy，逐坐标取极限后，有限首项与小尾项论证给出距离收敛。其拓扑就是逐坐标收敛的乘积拓扑。
\begin{pitfall}[补充核查：原书的另一计数公式]
原书还列出 $\arctan\inf_{\varepsilon>0}(\varepsilon+\#\{j:|x_j|>\varepsilon\})$ 作为替代公式。若它按未加权的普通计数理解，则不能在\emph{全体}数列上作为拟范数：取 $x_j=j$，任意有限阈值以上的坐标数都是无穷，所以任意 $a\ne0$ 时该表达式在 $ax$ 上恒为 $\pi/2$，不满足 $a\to0$ 的数乘连续性。本稿因此采用上面的加权求和公式来处理全体数列；这里没有把有缺陷的替代式作为定理使用。
\end{pitfall}

\subsection{Sobolev 空间及其完备性}
\begin{definition}[弱导数与 Sobolev 空间]
对整数 $k\ge0$ 和 $1\le p<\infty$，定义
\[
 W^{k,p}(\Omega)=\{f\in L^p(\Omega):\partial^\alpha T_f=T_{g_\alpha},\
 g_\alpha\in L^p(\Omega),\ |\alpha|\le k\}.
\]
由分布嵌入单射，$g_\alpha$ 唯一，记为弱导数 $\partial^\alpha f$。范数为
\[
 \|f\|_{k,p}=\left(\sum_{|\alpha|\le k}\|\partial^\alpha f\|_p^p\right)^{1/p}.
\]
当 $p=2$，记 $H^k(\Omega)=W^{k,2}(\Omega)$，取第 \ref{sec:yosida-01-05} 节的导数内积。
\end{definition}
\begin{theorem}[Sobolev 空间完备]\label{thm:y1-sobolev-complete}
$W^{k,p}(\Omega)$ 是 Banach 空间；$H^k(\Omega)$ 是 Hilbert 空间。
\end{theorem}
\begin{proof}
若 $f_j$ 在 Sobolev 范数下 Cauchy，则各 $\partial^\alpha f_j$ 在 $L^p$ 中 Cauchy，设极限为 $g_\alpha$，特别地 $f_j\to g_0=f$。对任意 $\varphi\in\Test$，Hölder（$p=1$ 时用上确界）允许在
\[
 \int\partial^\alpha f_j\,\varphi
 =(-1)^{|\alpha|}\int f_j\partial^\alpha\varphi
\]
两边取极限，得到 $\int g_\alpha\varphi=(-1)^{|\alpha|}\int f\partial^\alpha\varphi$。所以 $g_\alpha$ 正是 $f$ 的弱导数，$f\in W^{k,p}$，全部有限项范数同时收敛。这是“函数及其导数独立取极限”后必须验证的相容性。$p=2$ 时范数由内积生成，结合完备性即得 Hilbert 性。
\end{proof}

\begin{exercise}[补充：不完备的自然函数空间]
令 $c_{00}$ 为有限支集数列空间，取 $\ell^2$ 范数。构造一个不在 $c_{00}$ 中收敛的 Cauchy 序列，并确定其在 $\ell^2$ 中的极限。
\end{exercise}
\begin{solution}
取 $x^{(N)}=(1,1/2,\ldots,1/N,0,\ldots)$。差的平方范数为相应的 $\sum j^{-2}$ 尾部，故 Cauchy。极限 $(1/j)_{j\ge1}$ 属于 $\ell^2$，但不属于 $c_{00}$；若在 $c_{00}$ 内有极限，坐标连续性会迫使它等于该无限支集数列，矛盾。
\end{solution}

\section{完备化}\label{sec:yosida-01-10}
\noindent\textit{原书位置：中译本第 58--61 页；英文版第 56--59 页。}

\subsection{由 Cauchy 序列构造完备化}
\begin{theorem}[拟赋范空间的完备化]\label{thm:y1-completion}
每个拟赋范空间 $X$ 都能线性等距嵌入某个 F-空间 $\widehat X$，且像稠密；该完备化在固定原空间嵌入的线性等距同构意义下唯一。若原拟范数是范数，则 $\widehat X$ 是 Banach 空间；若原范数来自内积，则其完备化是 Hilbert 空间。
\end{theorem}
\begin{proof}
令 $\mathcal C$ 为 $X$ 的全部 Cauchy 序列，定义 $x_j\sim y_j$ 当且仅当 $q(x_j-y_j)\to0$。三角不等式证明它是等价关系。定义
\[
 [x_j]+[y_j]=[x_j+y_j],\qquad a[x_j]=[ax_j],\qquad
 \widehat q([x_j])=\lim_jq(x_j).
\]
反三角不等式保证实数序列 $q(x_j)$ Cauchy，故极限存在。加法与拟范数不依赖代表元；固定数乘的连续性保证 $[ax_j]$ 仍为 Cauchy 序列且不依赖代表元。正定性、对称性及三角不等式直接继承。

\emph{补充检查数乘条件。} 一个 Cauchy 序列 $x_j$ 满足 $q(ax_j)\to0$ 随 $a\to0$ 对 $j$ 一致。确实，数乘在 $(0,0)$ 联合连续，所以给定误差可先选小的向量邻域和标量邻域，使两者乘积小；再把 $x_j$ 的尾部写成某个固定 $x_N$ 加上该小向量，有限个首项另行控制。因此 $\widehat q(a[x_j])\to0$。对固定 $a$，源空间的数乘连续性给出：$q(x)<\delta$ 时 $q(ax)<\varepsilon/2$。若 $\widehat q(\xi)<\delta/2$，其代表元最终满足 $q(x_j)<\delta$，故 $\widehat q(a\xi)\le\varepsilon/2<\varepsilon$。于是 $\widehat q$ 满足两条数乘条件。这避免未经论证交换两重极限。

嵌入 $Jx=[x,x,\ldots]$ 线性等距。对 $\xi=[x_j]$，有 $\widehat q(\xi-Jx_N)=\lim_jq(x_j-x_N)\to0$，所以像稠密。若 $\xi_k$ 在 $\widehat X$ 中 Cauchy，由稠密性取 $z_k\in X$ 使 $\widehat q(Jz_k-\xi_k)<2^{-k}$。三角不等式使 $z_k$ 在 $X$ 中 Cauchy，令 $\xi=[z_k]$；由刚证明的稠密逼近及 Cauchy 性，$Jz_k\to\xi$，从而 $\xi_k\to\xi$，得完备性。

若另有稠密等距嵌入 $J':X\to Z$，$Z$ 完备，则映射 $[x_j]\mapsto\lim_jJ'x_j$ 良定义、线性等距且满射：任何 $z\in Z$ 可由原空间的像逼近。它固定 $X$，所以唯一。范数齐次性以及平行四边形恒等式均可在极限中继承。
\end{proof}

\begin{proposition}[补充：算子的延拓性质]\label{prop:y1-extension}
设 $Y$ 是 Banach 空间，$T:X\to Y$ 为赋范空间上的有界线性算子，则存在唯一有界线性延拓 $\widehat T:\widehat X\to Y$，且 $\|\widehat T\|=\|T\|$。
\end{proposition}
\begin{proof}
定义 $\widehat T([x_j])=\lim_jTx_j$。估计 $\|Tx_j-Tx_l\|\le\|T\|\|x_j-x_l\|$ 保证极限存在及代表元无关，并给出 $\|\widehat T\|\le\|T\|$。限制于原空间得到反向不等式。稠密性保证唯一性。
\end{proof}

\subsection{Sobolev 空间的完备化解释}
令 $H_0^k(\Omega)$ 为 $C_c^\infty(\Omega)$ 在 $H^k$ 范数下的完备化。第 \ref{sec:yosida-01-09} 节的完备性与弱导数相容性证明说明，它可等距识别为
\[
 H_0^k(\Omega)=\overline{C_c^\infty(\Omega)}^{\,H^k(\Omega)}.
\]
因此它是 $H^k(\Omega)$ 的闭子空间。原书还从 $\widehat H^k$ 出发作抽象完备化；本稿将现代弱导数定义的空间直接记为 $H^k$，以免“经典导数空间”和“已完备的 Sobolev 空间”混淆。对任意开集，$H_0^k$ 的定义不需要边界光滑；“零边界值”的迹解释则需要另加相应条件。

\begin{lemma}[补充：$L^p$ 平移连续与正则化逼近]\label{lem:y1-Lp-translation}
若 $1\le p<\infty$，$f\in L^p(\R^n)$，则 $\|f(\cdot-h)-f\|_p\to0$ 随 $h\to0$。因而 $f*\rho_\varepsilon\to f$ 在 $L^p$ 中。
\end{lemma}
\begin{proof}
先对有界矩形的示性函数证明：平移后对称差的体积趋零，$L^p$ 距离的 $p$ 次方正是该体积。有限个这样的简单函数亦成立。一般 $L^p$ 函数可先截断大小和远方尾部，再用简单函数及有限测度集的矩形逼近，在 $L^p$ 中逼近为上述函数。平移等距，于是
\[
 \|f(\cdot-h)-f\|_p
 \le2\|f-g\|_p+\|g(\cdot-h)-g\|_p
\]
完成一般情形。卷积逼近由 Jensen 不等式和 Fubini 得
\[
 \|f*\rho_\varepsilon-f\|_p^p
 \le\int\rho_\varepsilon(h)\|f(\cdot-h)-f\|_p^p\,dh.
\]
核支集收缩至零，右边趋零。
\end{proof}

\begin{theorem}[全空间上的光滑紧支集稠密性]\label{thm:y1-sobolev-density}
对每个整数 $k\ge0$ 与 $1\le p<\infty$，$C_c^\infty(\R^n)$ 在 $W^{k,p}(\R^n)$ 中稠密。特别地
\[
 H_0^k(\R^n)=H^k(\R^n)=W^{k,2}(\R^n).
\]
\end{theorem}
\begin{proof}
取 $\chi\in C_c^\infty$ 在单位球上为 $1$，令 $\chi_R(x)=\chi(x/R)$。Leibniz 公式给出
\[
 \partial^\alpha(\chi_Rf)-\partial^\alpha f
 =(\chi_R-1)\partial^\alpha f
 +\sum_{0<\beta\le\alpha}\binom\alpha\beta
 (\partial^\beta\chi_R)\partial^{\alpha-\beta}f.
\]
第一项的 $L^p$ 范数由 $\partial^\alpha f$ 的可积尾部趋零；其余项的系数不超过 $C_\beta R^{-|\beta|}$，且支集位于远方环带，故也趋零。因此 $\chi_Rf\to f$ 在 Sobolev 范数下，先把问题化为紧支集。

固定 $R$ 后令 $f_{R,\varepsilon}=(\chi_Rf)*\rho_\varepsilon$。它光滑且紧支集。利用分布导数的定义，将微分移入卷积核，再移到 $\chi_Rf$ 的弱导数，得到
\[
 \partial^\alpha f_{R,\varepsilon}
 =(\partial^\alpha(\chi_Rf))*\rho_\varepsilon.
\]
引理 \ref{lem:y1-Lp-translation} 保证每个有限阶导数在 $L^p$ 下逼近；先选大 $R$，再选小 $\varepsilon$，使总误差任意小。
\end{proof}

\begin{example}[补充：在有边界的域中两空间可以不同]
在 $(0,1)$ 上，常数 $1\in H^1(0,1)$，却不属于 $H_0^1(0,1)$。若 $\varphi_j\in C_c^\infty(0,1)$ 在 $H^1$ 中趋于 $1$，令 $v_j=\varphi_j-1$。对任意 $v\in C^1([0,1])$，由 $v(0)=v(x)-\int_0^xv'(t)\,dt$ 并对 $x$ 积分，得
\[
 |v(0)|\le\|v\|_{L^2(0,1)}+\|v'\|_{L^2(0,1)}.
\]
应用于 $v_j$，右侧趋零，左侧恒为 $1$，矛盾。这也说明将全空间的密度结论直接搬到任意有边界域会出错。
\end{example}

\section{B-空间的商空间}\label{sec:yosida-01-11}
\noindent\textit{原书位置：中译本第 61--62 页；英文版第 59--60 页。}

\begin{definition}[商范数]
设 $X$ 赋范，$M\subset X$ 为闭线性子空间。对商空间中的陪集 $x+M$ 定义
\[
 \|x+M\|_{X/M}=\inf_{m\in M}\|x+m\|=\operatorname{dist}(x,M).
\]
规范映射记为 $Q:X\to X/M$，$Qx=x+M$。
\end{definition}
\begin{proposition}[商范数、商拓扑与开映射]\label{prop:y1-quotient-norm}
该定义不依赖代表元，是商空间上的范数；其拓扑正是商拓扑。$Q$ 线性连续，$\|Q\|\le1$，且 $Q$ 是开映射。
\end{proposition}
\begin{proof}
代表元变化只是把 $m$ 平移，故下确界不变。齐次性显然，三角性由对两个下确界选取任意精度的近似代表再相加得到。范数为零时 $x$ 在 $M$ 的闭包中，闭性给出 $x\in M$，故正定。

有 $\|Qx\|\le\|x\|$。更准确地，
\[
 Q(B_X(x,r))=B_{X/M}(Qx,r).
\]
一向由压缩估计；反向若 $\|\xi-Qx\|<r$，下确界的定义允许选取该陪集中的代表 $z$，使 $\|z\|<r$，于是 $\xi=Q(x+z)$。这证明 $Q$ 开，也证明范数拓扑等于商拓扑。
\end{proof}
若 $M$ 不闭，同一公式仍是半范数，其零空间为 $\overline M/M$；所以闭性不是一个仅为方便而附加的条件。

\begin{theorem}[Banach 空间的闭子空间商仍完备]\label{thm:y1-quotient-complete}
若 $X$ 为 Banach 空间且 $M$ 闭，则 $X/M$ 为 Banach 空间。
\end{theorem}
\begin{proof}
从商空间的 Cauchy 序列 $\xi_j$ 抽出子列，使 $\|\xi_{j_{k+1}}-\xi_{j_k}\|<2^{-k-1}$。对每个差陪集选代表 $u_k\in X$，使
\[
 Qu_k=\xi_{j_{k+1}}-\xi_{j_k},\qquad\|u_k\|<2^{-k}.
\]
选 $x_1$ 满足 $Qx_1=\xi_{j_1}$。级数 $x=x_1+\sum_ku_k$ 在 Banach 空间中绝对收敛。部分和 $s_k=x_1+\sum_{l<k}u_l$ 满足 $Qs_k=\xi_{j_k}$，因此 $\xi_{j_k}\to Qx$。原序列由 Cauchy 性也趋于 $Qx$。这里仅需要近似达到下确界的代表，完全不需要最近点存在。
\end{proof}

图 \ref{fig:y1-quotient} 给出商范数的有限维模型；证明中的代表元只须近似达到下确界。
\begin{figure}[htbp]
\centering
\begin{tikzpicture}[fa diagram]
  \draw[fa axis] (-.7,0)--(5.6,0) node[right] {$M$};
  \draw[fa main] (-.5,2)--(5,2) node[right] {$x+M$};
  \node[fa dot,label=above:$x$] at (3.7,2) {};
  \node[fa dot,label=below:$0$] at (1,0) {};
  \draw[fa accent] (1,0)--(1,2) node[midway,left] {$d$};
  \draw[fa main,dashed] (1,0)--(3.7,2);
  \draw[fa main,dashed] (1,0)--(4.7,2);
  \node[fa note] at (3.2,-1) {$\|Qx\|=\inf_{m\in M}\|x+m\|$};
  \node[fa box] at (8.5,1) {商空间把整条陪集\\视为同一个点 $Qx$};
\end{tikzpicture}
\caption{商范数在 $\R^2$ 欧氏范数中的模型：在同一陪集中改变代表元，再取距原点的距离下确界。本例中 $d$ 由垂足达到；一般 Banach 空间的闭子空间未必有最近点，正文完备性证明只选择近似代表。}
\label{fig:y1-quotient}
\end{figure}



\begin{proposition}[补充：经商空间分解算子]
若有界线性算子 $T:X\to Y$ 满足 $M\subset\ker T$，则存在唯一有界线性算子 $\widetilde T:X/M\to Y$，满足 $T=\widetilde TQ$，并且 $\|\widetilde T\|=\|T\|$。
\end{proposition}
\begin{proof}
定义 $\widetilde T(x+M)=Tx$。$M\subset\ker T$ 保证代表元无关；由 $\|Tx\|=\|T(x+m)\|\le\|T\|\|x+m\|$ 对 $m$ 取下确界，得到 $\|\widetilde T\|\le\|T\|$。反向来自 $T=\widetilde TQ$ 与 $\|Q\|\le1$。满射性保证唯一性。
\end{proof}

\begin{exercise}[补充：商空间的具体计算]
在 $C([0,1])$ 上取上确界范数，令 $M=\{f:f(0)=0\}$。证明 $C([0,1])/M$ 与标量域等距同构。
\end{exercise}
\begin{solution}
评价映射 $f\mapsto f(0)$ 的核为 $M$，故诱导双射。所有代表的范数至少为 $|f(0)|$；而陪集中有常函数 $f(0)$，其范数恰为 $|f(0)|$，故商范数就是 $|f(0)|$。
\end{solution}

\section{单位分解}\label{sec:yosida-01-12}
\noindent\textit{原书位置：中译本第 63--64 页；英文版第 60--62 页。}

\subsection{局部有限覆盖的作用}
一族集合称为局部有限，若每个点都有邻域只与其中有限多个集合相交。在开集 $\Omega$ 中，对闭支集组成的局部有限族，每个紧集只与有限多个支集相交；因此相应的函数和在局部是有限和，可以逐项求导。这是把局部定义拼成全局对象的基本机制。

\begin{theorem}[从属任意开覆盖的光滑单位分解]\label{thm:y1-partition}
设 $\Omega\subset\R^n$ 开，$\{G_i\}_{i\in I}$ 是其开覆盖。存在可数函数族 $\{\theta_j\}_{j\ge1}\subset C_c^\infty(\Omega)$，满足
\[
 0\le\theta_j\le1,\qquad \supp\theta_j\subset G_{i(j)},\qquad
 \sum_j\theta_j(x)=1\quad(x\in\Omega),
\]
且支集族在 $\Omega$ 中局部有限。这里允许多个 $j$ 从属于同一个覆盖元素。
\end{theorem}
\begin{proof}
取紧集穷竭 $K_m\subset\operatorname{int}K_{m+1}$，置 $K_0=K_{-1}=\varnothing$。紧环层
\[
 A_m=K_m\setminus\operatorname{int}K_{m-1}
\]
包含于开环层 $W_m=\operatorname{int}K_{m+1}\setminus K_{m-2}$。对每个 $x\in A_m$，选择覆盖元素 $G_i$ 及小球 $B(x,2r)$，使闭球包含于 $W_m\cap G_i$。有限个小球 $B(x,r)$ 覆盖 $A_m$。在每个大球中取非负光滑紧支集函数 $b_j$，在对应小球上严格为正。

所有环层的这些函数组成可数族，且支集在 $\Omega$ 中局部有限：固定紧集包含于某个 $K_N$，而 $m\ge N+2$ 的支集在 $K_N$ 外，剩下层各只有有限个函数。函数
\[
 b(x)=\sum_jb_j(x)
\]
因局部有限而光滑，因小球覆盖各环层而处处为正。令 $\theta_j=b_j/b$，则正性、求和为一和支集条件全部成立。每个 $b_j$ 的支集紧包含于 $\Omega$，所以 $\theta_j$ 也可延零成 $\R^n$ 上的光滑紧支集函数；分母可能在域边界变小不影响这一点，因为每个单独支集与边界保持正距离。
\end{proof}
这给出了原书“散开覆盖”的标准局部有限含义，同时确保局部有限的是\emph{支集}而非仅非零开集，消除了归一化与求导中的潜在缺口。

\begin{corollary}[光滑截断函数]\label{cor:y1-cutoff}
若 $K\Subset\Omega$、$O\subset\Omega$ 开且含 $K$，则存在 $\chi\in\Test(\Omega)$，$0\le\chi\le1$，在 $K$ 的某个邻域上为 $1$，且 $\supp\chi\subset O$。
\end{corollary}
\begin{proof}
用有限个闭球包含于 $O$ 的小球覆盖 $K$，取在小球上为正的紧支集光滑函数，相加得 $b$，其支集紧包含于 $O$，且 $\min_K b=c>0$。取光滑实函数 $h$，在 $t\le c/2$ 上为零，在 $t\ge3c/4$ 上为一，并介于零与一之间。$\chi=h\circ b$ 即满足要求。
\end{proof}
特别地，可以取 $\chi_m$ 在 $K_m$ 附近为一，于是 $\chi_mf\to f$ 在 $\mathcal E(\Omega)$ 中成立，说明 $\Test$ 在 $\mathcal E$ 中稠密。这个稠密性不是说 $\Test$ 与 $\mathcal E$ 的拓扑相同。

\subsection{补充：分布的限制与粘合}
若 $U\subset\Omega$ 开，$\Test(U)$ 的函数延零到 $\Omega$ 后仍是测试函数，因此可定义 $T|_U$。下面的局部原则使后面支集的定义成立。
\begin{proposition}[分布的局部唯一性与粘合]\label{prop:y1-gluing}
若分布 $T$ 在开集族 $U_i$ 上均为零，则它在 $\bigcup_iU_i$ 上为零。更一般地，若 $T_i\in\Dist(U_i)$ 在交集上相容，则存在唯一 $T\in\Dist(\bigcup_iU_i)$，其各限制等于 $T_i$。
\end{proposition}
\begin{proof}
先取从属于 $U_i$ 的单位分解 $\theta_j$。对支集于并集的 $\varphi$，只有有限个 $\theta_j\varphi$ 非零，因此 $\varphi=\sum_j\theta_j\varphi$。若各局部限制为零，对有限和应用 $T$ 即得零。

一般情形定义
\[
 T(\varphi)=\sum_jT_{i(j)}(\theta_j\varphi).
\]
对固定 $\Test_K$ 只有有限项；每项由光滑乘子连续性及局部有限阶估计控制，所以 $T$ 是分布。若 $\supp\varphi\subset U_i$，各 $\theta_j\varphi$ 的紧支集在 $U_i\cap U_{i(j)}$ 中，相容性使和等于 $\sum_jT_i(\theta_j\varphi)=T_i(\varphi)$。局部零结论保证唯一性，因此也保证结果与单位分解选择无关。
\end{proof}
专题部分第 \ref{ch:distributions} 章给出限制与粘合的进一步应用。

\section{具有紧支集的广义函数}\label{sec:yosida-01-13}
\noindent\textit{原书位置：中译本第 64--67 页；英文版第 62--65 页。}

\subsection{支集及光滑函数对偶}
\begin{definition}[分布支集]
称 $T$ 在开集 $U$ 上为零，若它对所有支集于 $U$ 的测试函数取零。由命题 \ref{prop:y1-gluing}，所有这种开集的并仍是零区域。其补集称为 $\supp T$，即 $T$ 的支集。
\end{definition}
支集是相对闭集。若 $\varphi$ 的紧支集与 $\supp T$ 不相交，则 $T(\varphi)=0$。由定义可直接得
\[
 \supp(\partial^\alpha T)\subset\supp T,\qquad
 \supp(aT)\subset\supp a\cap\supp T.
\]
例如 $\supp\delta_a=\{a\}$，而非零光滑函数生成的分布支集就是其函数支集；后一事实利用局部可积嵌入单射。

$\mathcal E(\Omega)=C^\infty(\Omega)$ 的有界集由每个紧集上的每阶导数一致有界刻画，直接来自半范数有界判别。它是可度量局部凸空间，故囿，因此其线性泛函连续等价于在所有有界集上有界。

\begin{theorem}[紧支集分布等同于光滑函数的连续对偶]\label{thm:y1-Eprime}
每个紧支集分布 $T$ 唯一延拓为 $\mathcal E(\Omega)$ 上的连续线性泛函。反过来，每个 $L\in\mathcal E(\Omega)'$ 限制于 $\Test$ 后都是紧支集分布。这给出自然的线性对应
\[
 \CompactDist(\Omega)=\{T\in\Dist(\Omega):\supp T\Subset\Omega\}.
\]
连续性还等价于存在一个 $K\Subset\Omega$、整数 $m\ge0$ 和常数 $C$，使
\[
 |L(f)|\le C\max_{|\alpha|\le m}\sup_K|\partial^\alpha f|
 \quad(f\in C^\infty(\Omega)).
\]
\end{theorem}
\begin{proof}
设 $\supp T=F$ 紧，取 $\chi\in\Test$ 在 $F$ 附近为一，定义 $L(f)=T(\chi f)$。若换用另一个 $\chi_1$，$(\chi-\chi_1)f$ 是紧支集测试函数且在 $F$ 附近为零，故 $T$ 对它取零，定义与截断选择无关。由 $T$ 在固定紧集 $K=\supp\chi$ 上的有限阶估计及 Leibniz 公式，得到上述对 $f$ 的估计，因此 $L$ 连续。它限制到 $\Test$ 等于 $T$；$\Test$ 在 $\mathcal E$ 中稠密，故连续延拓唯一。

反向，连续性给出有限个 $p_{K_j,m_j}$ 控制 $|L(f)|$，取有限个紧集的并及最大的阶数得所列单一估计。于是 $L|_{\Test}$ 是分布，并在 $\Omega\setminus K$ 上为零，故支集包含于 $K$。估计的充分性显然。注意这里“一种延拓”指的是延拓\emph{泛函}唯一；实现它的截断函数可以有很多个。
\end{proof}
每个紧支集分布因此具有全局有限阶估计，但估计中的紧集可以稍大于其真实支集。Dirac 分布及其导数可作用于任何光滑函数；一般非紧支集分布则不能未经条件直接作用于常函数 $1$。

\subsection{点支集分布的完整结构}
\begin{theorem}[点支集分布是 Dirac 导数的有限和]\label{thm:y1-point-support}
若 $\supp T\subset\{a\}$，则存在 $m<\infty$ 及唯一系数 $c_\alpha$，使
\[
 T=\sum_{|\alpha|\le m}c_\alpha\partial^\alpha\delta_a.
\]
只在需要时选取足够大的 $m$；补零后各系数唯一。
\end{theorem}
\begin{proof}
平移后设 $a=0$。由定理 \ref{thm:y1-Eprime}，存在固定紧邻域 $K$ 上的 $m$ 阶估计。若 $\partial^\alpha f(0)=0$ 对 $|\alpha|\le m$ 成立，取 $\chi\in\Test$ 在零附近为一，令 $\chi_\varepsilon(x)=\chi(x/\varepsilon)$，其支集在充分小时落于 $K$。支集性质给出 $T(f)=T(\chi_\varepsilon f)$。

Taylor 余项保证对 $|\gamma|\le m$，在 $|x|\le C\varepsilon$ 上，
\[
 |\partial^\gamma f(x)|\le C_f\varepsilon^{m+1-|\gamma|}.
\]
Leibniz 公式中的一项为 $\varepsilon^{-|\beta|}(\partial^\beta\chi)(x/\varepsilon)\partial^{\alpha-\beta}f(x)$，因此对 $|\alpha|\le m$，
\[
 \sup_K|\partial^\alpha(\chi_\varepsilon f)|
 \le C_f'\varepsilon^{m+1-|\alpha|}\le C_f''\varepsilon
 \quad(0<\varepsilon\le1).
\]
有限阶估计使 $T(f)=0$。这展开了截断函数求导产生负幂后，Taylor 余项仍能抵消它们的理由。

一般 $f$ 减去 $m$ 阶 Taylor 多项式 $P_mf=\sum_{|\alpha|\le m}\partial^\alpha f(0)x^\alpha/\alpha!$ 后满足上述零喷射条件，所以 $T(f)=T(P_mf)$。这里 $T$ 已延拓到 $\mathcal E$，可以合法作用于多项式。由于
\[
 (\partial^\alpha\delta_0)(f)=(-1)^{|\alpha|}\partial^\alpha f(0),
\]
取 $c_\alpha=(-1)^{|\alpha|}T(x^\alpha/\alpha!)$ 即得表示。最后以在零附近等于单项式的紧支集测试函数分别检验各喷射，得到系数唯一。
\end{proof}
有限多个点支集的分布也是各点 Dirac 导数的有限和：取互不相交的小邻域和相应截断，将分布分解成各点支集分布后应用定理。点支集结论不能把“有限阶”换成任意无穷导数级数，连续性正是截断为有限喷射的原因。

\begin{exercise}[补充：点支集分布的阶数]
证明在 $\R$ 上 $\delta_0^{(m)}$ 的阶数恰为 $m$：它受 $m$ 阶测试半范数控制，但 $m\ge1$ 时不能由 $m-1$ 阶半范数控制。
\end{exercise}
\begin{solution}
上界由 $|\delta_0^{(m)}(\varphi)|=|\varphi^{(m)}(0)|$ 给出。选 $\eta\in\Test$ 使 $\eta^{(m)}(0)\ne0$，令 $\varphi_\varepsilon(x)=\varepsilon^{m-1}\eta(x/\varepsilon)$。支集保持在一个固定紧邻域中，所有至 $m-1$ 阶导数统一有界，而 $|\varphi_\varepsilon^{(m)}(0)|=\varepsilon^{-1}|\eta^{(m)}(0)|\to\infty$。因此不存在低阶估计。
\end{solution}

\section{广义函数的直接积}\label{sec:yosida-01-14}
\noindent\textit{原书位置：中译本第 67--70 页；英文版第 65--67 页。}

\subsection{分离变量函数的稠密性}
记 $u\otimes v$ 为函数 $(x,y)\mapsto u(x)v(y)$。有限个这样的函数的线性组合构成测试空间的代数张量积；它不是集合的笛卡尔积，也不是把同一变量中的两个分布相乘。

\begin{theorem}[分离变量的测试函数逼近]\label{thm:y1-tensor-density}
给定 $\varphi\in\Test(\R^{n+m})$，存在
\[
 \varphi_j(x,y)=\sum_{l=1}^{N_j}u_{jl}(x)v_{jl}(y),\qquad
 u_{jl}\in\Test(\R^n),\ v_{jl}\in\Test(\R^m),
\]
使 $\varphi_j\to\varphi$ 在 $\Test(\R^{n+m})$ 中成立。即有限分离变量函数的和在完整测试拓扑下稠密。
\end{theorem}
\begin{proof}
展开原书的高斯正则化证明。令 $d=n+m$、$z=(x,y)$，定义
\[
 \Phi_t(z)=(4\pi t)^{-d/2}\int_{\R^d}
 \varphi(w)\exp\left(-\frac{|z-w|^2}{4t}\right)\,dw,
 \qquad t>0.
\]
对每个导数 $\partial^\alpha$，分部积分将微分从核移到 $\varphi$，得到 $\partial^\alpha\Phi_t=G_t*\partial^\alpha\varphi$。换元 $w=z+2\sqrt t\,s$ 后，积分为归一化高斯下 $\partial^\alpha\varphi(z+2\sqrt t\,s)$ 的平均。分成 $|s|\le R$ 与 $|s|>R$：前一部分因紧支集光滑函数的全域一致连续性而随 $t\downarrow0$ 一致逼近，后一部分由可积高斯尾部先取大 $R$ 控制。因此所有有限阶导数一致趋于 $\partial^\alpha\varphi$。

积分中的 $|z-w|^2$ 在复变量延拓时应理解为 $\sum(z_l-w_l)^2$，不用复共轭。$w$ 的积分区间固定紧集，所以核对复 $z$ 是整函数并在复紧集上有统一可积控制；故 $\Phi_t$ 是 $d$ 变量整函数。其多元 Taylor 部分和为多项式，在每个复紧集上一致收敛；由 Cauchy 积分公式，各有限阶导数也一致收敛。选择 $t_j\downarrow0$ 及足够高阶多项式 $P_j$，使在 $|z|\le j$ 上、$|\alpha|\le j$ 时，$\partial^\alpha(P_j-\Phi_{t_j})$ 的绝对值小于 $1/j$。于是 $P_j\to\varphi$ 在 $\mathcal E(\R^d)$ 中。

选固定 $a\in\Test(\R^n)$、$b\in\Test(\R^m)$，分别在 $\supp\varphi$ 的两坐标投影附近为一。令 $\varphi_j=a(x)b(y)P_j(x,y)$。它们的支集全部在固定紧集 $\supp a\times\supp b$ 中，且 Leibniz 公式给出所有导数一致趋于 $a(x)b(y)\varphi=\varphi$，所以是 $\Test$ 收敛。每个多项式是有限个 $x^\alpha y^\beta$ 的和，乘上 $a,b$ 后即是所需分离变量和。固定截断而非随 $j$ 逃向无穷的截断，是最后一步的要点。
\end{proof}
同一结论适用于 $\Test(\Omega\times\Lambda)$：先用有限个紧包含于乘积域的矩形邻域及单位分解分解一个给定测试函数，再对每片选择支集留在对应坐标域中的截断。因支集紧，只需有限片，每片的上述逼近可以合并。

\subsection{带参数作用、直积及 Fubini 性质}
\begin{lemma}[分布作用后的参数函数仍是测试函数]\label{lem:y1-parameter}
设 $S\in\Dist(\Lambda)$，$\varphi\in\Test(\Omega\times\Lambda)$。则
\[
 \psi(x)=\Pair{S}{\varphi(x,\cdot)}
\]
属于 $\Test(\Omega)$，且 $\partial_x^\alpha\psi(x)=\Pair{S}{\partial_x^\alpha\varphi(x,\cdot)}$。若 $\supp\varphi\subset K\times L$ 且 $S$ 在 $\Test_L$ 上受 $s$ 阶估计控制，则
\[
 \max_{|\alpha|\le r}\sup_K|\partial_x^\alpha\psi|
 \le C_S\max_{\substack{|\alpha|\le r\\|\beta|\le s}}
 \sup_{K\times L}|\partial_x^\alpha\partial_y^\beta\varphi|.
\]
\end{lemma}
\begin{proof}
$x$ 在坐标投影支集外时，$\varphi(x,\cdot)=0$，所以 $\psi$ 有紧支集。对一阶差商，Taylor 公式给出
\[
 \frac{\varphi(x+he_j,y)-\varphi(x,y)}h-\partial_{x_j}\varphi(x,y)
 \longrightarrow0
\]
在任意有限阶的 $y$ 导数半范数下一致成立，且 $y$ 支集固定于 $L$。$S$ 连续允许取此极限，得到导数公式；逐次重复证明所有阶。对导数公式应用 $S$ 的有限阶估计即得最后的上界。它还证明带参数作用把测试空间的有界集送入有界集。
\end{proof}

\begin{theorem}[分布直积的存在、唯一性与交换作用顺序]\label{thm:y1-tensor}
对 $T\in\Dist(\Omega)$、$S\in\Dist(\Lambda)$，存在唯一 $T\otimes S\in\Dist(\Omega\times\Lambda)$，使
\[
 \Pair{T\otimes S}{u(x)v(y)}=\Pair{T}{u}\Pair{S}{v}.
\]
对一般测试函数，具有分布版 Fubini 公式
\[
 \Pair{T\otimes S}{\varphi}
 =\Pair{T_x}{\Pair{S_y}{\varphi(x,y)}}
 =\Pair{S_y}{\Pair{T_x}{\varphi(x,y)}}.
\]
\end{theorem}
\begin{proof}
由上一引理，先作用 $S$ 再作用 $T$ 有意义，并且在固定 $K\times L$ 上有
\[
 |T_x(S_y\varphi)|\le C_TC_S
 \max_{\substack{|\alpha|\le r\\|\beta|\le s}}
 \sup_{K\times L}|\partial_x^\alpha\partial_y^\beta\varphi|.
\]
每个一般紧集都包含在两个坐标投影紧集的乘积内，故局部有限阶判别证明它是分布。反向作用顺序同样定义分布。两者对分离变量函数都取值 $T(u)S(v)$，故对其有限线性组合相同。定理 \ref{thm:y1-tensor-density} 的稠密性及两个分布的连续性证明它们在全部测试函数上相同，并同时证明满足所列乘积条件的分布唯一。
\end{proof}
交换公式是“交换作用于不同变量的顺序”；若把 $T\otimes S$ 与 $S\otimes T$ 看作不同次序乘积域上的分布，还须配合坐标交换映射。原书以 $\times$ 记直接积，此处用 $\otimes$ 突出独立变量与张量含义。

\begin{proposition}[补充：支集、微分与普通函数的对应]
分布直积满足
\[
 \supp(T\otimes S)=\supp T\times\supp S,
\]
以及
\[
 \partial_x^\alpha\partial_y^\beta(T\otimes S)
 =(\partial^\alpha T)\otimes(\partial^\beta S).
\]
若 $f\in L^1_{\Loc}(\Omega)$、$g\in L^1_{\Loc}(\Lambda)$，则 $T_f\otimes T_g=T_{f(x)g(y)}$。
\end{proposition}
\begin{proof}
若一点在 $\supp T\times\supp S$ 外，可取一个矩形邻域，其中某个因子分布为零，直积在此邻域为零，得支集包含关系。反向若 $a\in\supp T$、$b\in\supp S$，任意包含 $(a,b)$ 的邻域都有较小矩形 $U\times V$。由支集定义存在支集于 $U,V$ 的 $u,v$，使 $T(u)\ne0$、$S(v)\ne0$，于是直积在 $uv$ 上非零，故 $(a,b)$ 属于支集。微分公式逐次使用导数定义和 Fubini 公式。普通函数情形中，固定乘积紧集上的绝对可积性由 $\int_K|f|\int_L|g|<\infty$ 保证，经典 Fubini 给出对应公式。
\end{proof}
例如 $\delta_a\otimes\delta_b=\delta_{(a,b)}$，且 $\delta_a'\otimes\delta_b$ 作用于 $\varphi$ 给出 $-\partial_x\varphi(a,b)$。这都合法，仍没有定义同一变量中的 $\delta_a^2$。更深入的卷积与张量应用见第 \ref{ch:convolution} 章。

\begin{exercise}[补充：直积与微分的检验]
在 $\R^2$ 上令 $W=H(x)\otimes\delta_0(y)$。写出 $W(\varphi)$，计算 $\partial_xW$ 与 $\partial_yW$，并确定其支集。
\end{exercise}
\begin{solution}
$W(\varphi)=\int_0^\infty\varphi(x,0)\,dx$。因此 $\partial_xW=\delta_0(x)\otimes\delta_0(y)=\delta_{(0,0)}$，$\partial_yW=H(x)\otimes\delta_0'(y)$，后者作用于测试函数为 $-\int_0^\infty\partial_y\varphi(x,0)\,dx$。支集为 $[0,\infty)\times\{0\}$。
\end{solution}


\section*{本章参考文献与阅读衔接}
\phantomsection
\addcontentsline{toc}{section}{本章参考文献与阅读衔接}
本章十四节的基本材料来源于 \cite{yosida-zh,yosida-en} 第一章。原书章末还指引读者阅读 Bourbaki、Grothendieck、Köthe、Banach、Dunford--Schwartz 及 Hille--Phillips 关于局部凸空间与 Banach 空间的著作，以及 Schwartz、Gelfand--Shilov、Hörmander 与 Friedman 关于分布的著作。本稿实际使用的分布参考书见 \cite{schwartz,gelfand}；上述其余著作仅保留为原书的阅读线索，未作为已核查的直接引文。

第 \ref{ch:yosida-02} 章将从完备性推出一致有界、开映射与闭图像定理；第 \ref{ch:yosida-03} 章继续发展 Hilbert 空间的投影与表示。广义函数的进一步例子、卷积、Fourier 变换及偏微分方程应用，可接续阅读第 \ref{ch:test}--\ref{ch:weighted} 章。本章中分布的基本定义与证明仍独立保留，以保证主线阅读不依赖专题部分。

% !TeX encoding = UTF-8
\SetChapterInfo{subtitle={由完备性获得统一控制与连续性},keywords={一致有界; 测度极限; 开映射; 闭图像; 亚椭圆性},prereq={Baire 定理、第一章函数空间与分布},goals={掌握三大基本定理; 区分闭性与有界性; 理解定性正则性如何产生估计}}
\chapter{Baire--Hausdorff 定理的应用}\label{ch:yosida-02}
\ChapterSource{本章依照 \cite{yosida-zh,yosida-en} 第二章七节的顺序重新组织。实际核对范围：中译本印刷第 71--83 页（PDF 第 86--98 页），英文第六版印刷第 68--81 页（PDF 第 83--96 页）。中译本第 84 页为空白页；英文第 81 页下部已进入第三章。以下采用本稿记号与编号，补充步骤另作说明，并非逐字翻译。}

第一章提供了大量完备函数空间。本章解释完备性如何把逐点信息变成统一控制：一致有界性把各点上的有界性变为邻域上的一致估计；开映射把满射变为可控地求解方程；闭图像把极限相容性变为连续性。三者都以第零章的 Baire 定理为基础。

除特别说明外，标量域为 $\R$ 或 $\mathbb C$，所有空间均为 Hausdorff 拓扑向量空间。沿用第一章术语：$F$ 空间的拓扑由完备平移不变度量给出，不要求局部凸；Fr\'echet 空间另要求局部凸。Banach 空间是完备赋范空间。第一纲集是可数个无处稠密集之并，其补集称为余第一纲集。非空完备度量空间中的余第一纲集稠密。

\section{一致有界性定理及共鸣定理}\label{sec:yosida-02-01}
\noindent\textit{原书位置：中译本第 71--72 页；英文版第 68--70 页。}

\begin{theorem}[一致有界性：原书的次可加形式]\label{thm:y2-uniform-general}
设 $X$ 是非第一纲的 Hausdorff 线性拓扑空间，$Y$ 具有第一章意义的拟范数 $q$，$T_a:X\to Y$ 连续，并满足
\[
 q(T_a(x+y))\le q(T_ax)+q(T_ay),\qquad
 q(T_a(tx))=q(tT_ax)\quad(t\ge0).
\]
若对每个 $x\in X$，集合 $\{T_ax:a\in A\}$ 拓扑有界，则对任意 $\varepsilon>0$，存在零邻域 $V$，使 $q(T_ax)\le\varepsilon$ 对所有 $x\in V,a\in A$ 成立。
\end{theorem}
\begin{proof}
固定 $\varepsilon>0$，令
\[
 E_m=\left\{x:\sup_a\bigl[q(m^{-1}T_ax)+q(m^{-1}T_a(-x))\bigr]\le\varepsilon/2\right\}.
\]
各 $E_m$ 闭。拓扑有界性意味着对固定 $x$，$m^{-1}T_a(\pm x)$ 对 $a$ 一致趋零，故 $X=\bigcup_mE_m$。于是某 $E_m$ 含有 $x_0+W$，其中 $W$ 是零邻域且 $0\in W$。对 $h\in W$，次可加性和正数缩放条件给出
\[
 q(T_a(h/m))\le q(m^{-1}T_a(x_0+h))+q(m^{-1}T_a(-x_0))\le\varepsilon.
\]
取 $V=m^{-1}W$ 即可。拟范数不要求绝对齐次，不能把拓扑有界误换成拟范数数值有界。
\end{proof}

\begin{theorem}[Banach--Steinhaus 与共鸣定理]\label{thm:y2-banach-steinhaus}
设 $X$ 为 Banach 空间，$Y$ 为赋范空间，$T_a:X\to Y$ 为有界线性算子。若 $\sup_a\|T_ax\|<\infty$ 对每个 $x\in X$ 成立，则 $\sup_a\|T_a\|<\infty$。更一般地，若 $X$ 是 Baire 局部凸空间、$Y$ 是局部凸空间，则逐点有界的连续线性算子族等度连续。
\end{theorem}
\begin{proof}
令 $E_m=\{x:\sup_a\|T_ax\|\le m\}$。这些闭集覆盖 $X$，故某 $E_m$ 包含 $B(x_0,r)$。当 $\|h\|<r$ 时，$x_0+h,x_0\in E_m$，故 $\sup_a\|T_ah\|\le2m$。对单位向量取 $h=(r/2)x$，得 $\sup_a\|T_a\|\le4m/r$。

一般版本对每个目标连续半范数 $q$ 使用 $E_m=\{x:\sup_aq(T_ax)\le m\}$。同样的平移相减给出某零邻域上的统一界，再缩放得到任意小的 $q$ 界。对有限个目标半范数条件取源邻域的交，即得等度连续性。特别地，Fr\'echet 空间上逐点有界的连续线性泛函族满足 $\sup_a|T_ax|\le C\max_{j\le r}p_j(x)$，其中 $p_j$ 是有限多个源半范数。
\end{proof}

图 \ref{fig:y2-uniform} 对应 Banach 版本中从非零中心的内球转到零附近估计的步骤。
\begin{figure}[htbp]
\centering
\begin{tikzpicture}[fa diagram]
  \draw[fa region] (-2,-1.7) rectangle (2.2,1.8);
  \node at (-1.45,1.4) {$E_m$};
  \draw[fa main] (.2,0) circle (1);
  \node[fa dot,label=below:$x_0$] at (.2,0) {};
  \node[fa dot,label=above:$x_0+h$] at (.8,.45) {};
  \draw[->,fa accent] (.2,0)--(.8,.45);
  \draw[->,fa main] (2.6,.1)--(4.3,.1) node[midway,above] {相减};
  \draw[fa main,fill=LogicColor!8] (6,0) circle (1);
  \node[fa dot,label=below:$0$] at (6,0) {};
  \node[fa dot,label=above:$h$] at (6.6,.45) {};
  \draw[->,fa accent] (6,0)--(6.6,.45);
  \node[fa note] at (2.9,-2.4) {$x_0+h,x_0\in E_m\quad\Longrightarrow\quad\sup_a\|T_ah\|\le2m$};
\end{tikzpicture}
\caption{一致有界性证明的转折：Baire 定理在某个 $E_m$ 中找到非空内球，线性相减把它变为零附近的统一估计，再缩放至单位球。仅取 $x_0$ 与 $x_0+h$ 相减已足够，不要求 $E_m$ 本身是子空间。}
\label{fig:y2-uniform}
\end{figure}



\begin{corollary}[强极限的连续性]\label{cor:y2-strong-limit}
若 $T_n:X\to Y$ 如上且 $T_nx\to Tx$ 对每个 $x$ 成立，则 $T$ 有界线性，且 $\|T\|\le\liminf_n\|T_n\|$。
\end{corollary}
\begin{proof}
收敛序列逐点有界，故 $M=\sup_n\|T_n\|<\infty$。在和与数乘的恒等式中取极限得线性，并有 $\|Tx\|\le M\|x\|$。又由 $\|T_nx\|\le\|T_n\|\|x\|$ 取下极限，再对单位球取上确界即得结论。这里不要求 $Y$ 完备，因为已经假定极限在 $Y$ 中存在。
\end{proof}

\begin{theorem}[Neumann 级数]\label{thm:y2-neumann}
若 $X$ 为 Banach 空间，$A\in\mathcal L(X)$ 且 $\|A\|<1$，则
\[
 (I-A)^{-1}=\sum_{j=0}^\infty A^j,\qquad
 \|(I-A)^{-1}\|\le(1-\|A\|)^{-1}.
\]
级数在算子范数下收敛。原书以 $A=I-T$ 表述此结果。
\end{theorem}
\begin{proof}
第一章已证明 $\mathcal L(X)$ 完备，而 $\sum_j\|A^j\|\le\sum_j\|A\|^j<\infty$。若 $S_N=\sum_{j=0}^NA^j$，则 $(I-A)S_N=S_N(I-A)=I-A^{N+1}$。按算子范数取极限得双侧逆，范数界来自几何级数。
\end{proof}

\begin{example}[源空间完备性的作用]
给有限支集序列空间 $c_{00}$ 配上 $\|x\|_\infty$，令 $T_nx=nx_n$。每个固定 $x$ 只在有限多个坐标非零，所以 $\sup_n|T_nx|<\infty$，但 $\|T_n\|=n$。$c_{00}$ 不完备，其完备化为 $c_0$。
\end{example}

\begin{exercise}
设 $X$ 为 Banach 空间，$T_n\in\mathcal L(X,Y)$，$Y$ 赋范。证明：$T_nx\to0$ 对所有 $x$ 成立，当且仅当 $\sup_n\|T_n\|<\infty$ 且在某个稠密子集 $D\subset X$ 上 $T_nd\to0$。
\end{exercise}
\begin{solution}
必要性由共鸣定理，且可取 $D=X$。反向令 $M=\sup_n\|T_n\|$。给定 $x,\varepsilon$，取 $d\in D$ 使 $M\|x-d\|<\varepsilon/2$（$M=0$ 时显然）。当 $n$ 足够大时 $\|T_nd\|<\varepsilon/2$，于是 $\|T_nx\|\le M\|x-d\|+\|T_nd\|<\varepsilon$。无需 $D$ 为子空间。
\end{solution}

\begin{exercise}
设 $T:X\to Y$ 为 Banach 空间之间的有界可逆算子，$E\in\mathcal L(X,Y)$ 且 $\|T^{-1}E\|<1$。证明 $T+E$ 可逆，给出逆的级数及范数界。
\end{exercise}
\begin{solution}
由 $T+E=T(I+T^{-1}E)$，对 $A=-T^{-1}E$ 使用 Neumann 级数得
\[
 (T+E)^{-1}=\sum_{j=0}^\infty(-T^{-1}E)^jT^{-1},\qquad
 \|(T+E)^{-1}\|\le\frac{\|T^{-1}\|}{1-\|T^{-1}E\|}.
\]
级数在 $\mathcal L(Y,X)$ 中按范数收敛，并给出双侧逆。
\end{solution}

\section{Vitali--Hahn--Saks 定理}\label{sec:yosida-02-02}
\noindent\textit{原书位置：中译本第 73--74 页；英文版第 70--72 页。}

本节测度均为可数可加的有限值实或复测度，$|\nu|$ 表示全变差。对正测度 $m$，这里的 $m$-绝对连续明确指
\[
 \forall\varepsilon>0\ \exists\delta>0:\quad m(E)<\delta\Longrightarrow|\nu(E)|<\varepsilon.
\]
对有限值复测度，它等价于在 $m$-零集上为零。非平凡方向可用反证法选 $m(E_j)<2^{-j}$ 而 $|\nu(E_j)|\ge\varepsilon$，令 $F_k=\bigcup_{j\ge k}E_j$。则 $m(\bigcap_kF_k)=0$，由 $|\nu|$ 的从上连续性得 $|\nu|(F_k)\to0$，矛盾。此论证使用了第一章已证明的复测度全变差有限。

\begin{lemma}[集合空间的完备性]\label{lem:y2-measure-algebra}
在正测度空间 $(S,\Sigma,m)$ 上，令 $\Sigma_0=\{E\in\Sigma:m(E)<\infty\}$，把相差 $m$-零集的集合视为同一点。则 $d([E],[F])=m(E\triangle F)$ 是完备度量。
\end{lemma}
\begin{proof}
由 $\|\mathbf1_E-\mathbf1_F\|_1=m(E\triangle F)$ 得到向 $L^1$ 的等距嵌入。若 $\mathbf1_{E_j}$ 为 Cauchy 序列，$L^1$ 完备性给出极限 $f$，并有几乎处处收敛的子列。子列只取 $0,1$，所以 $f=\mathbf1_E$ 几乎处处；用子列的可测上极限可选择可测 $E$。又 $m(E)=\|f\|_1<\infty$，所以极限仍在像中。
\end{proof}

\begin{theorem}[Vitali--Hahn--Saks]\label{thm:y2-vhs}
设 $m(S)<\infty$，$\nu_n$ 为 $m$-绝对连续的复测度。若对每个 $E\in\Sigma$，有限极限 $\nu(E)=\lim_n\nu_n(E)$ 存在，则
\[
 \forall\varepsilon>0\ \exists\delta>0:\quad
 m(E)<\delta\Longrightarrow\sup_n|\nu_n(E)|<\varepsilon.
\]
而且 $\nu$ 是 $m$-绝对连续的复测度。
\end{theorem}
\begin{proof}
每个 $[E]\mapsto\nu_n(E)$ 在引理的完备空间上良定义且连续，因为
\[
 |\nu_n(E)-\nu_n(F)|\le|\nu_n(E\setminus F)|+|\nu_n(F\setminus E)|
\]
在 $m(E\triangle F)\to0$ 时趋零。固定 $\varepsilon>0$，令
\[
 F_k=\{[E]:|\nu_n(E)-\nu_l(E)|\le\varepsilon/4\ \text{对所有 }n,l\ge k\}.
\]
这些闭集由逐集合 Cauchy 性覆盖整个空间。Baire 定理给出 $k,E_0,\eta>0$，使 $d([E],[E_0])<\eta$ 蕴含 $[E]\in F_k$。若 $m(B)<\eta$，令 $E_1=E_0\cup B$、$E_2=E_0\setminus B$。两者均在该球内，且 $\nu_n(B)=\nu_n(E_1)-\nu_n(E_2)$，故对 $n\ge k$ 有
\[
 |\nu_n(B)-\nu_k(B)|\le\varepsilon/2.
\]
再用有限多个 $\nu_1,\ldots,\nu_k$ 的绝对连续性缩小 $\delta<\eta$，使它们在 $m(B)<\delta$ 时均小于 $\varepsilon/4$。便得全部 $n$ 的统一界小于 $\varepsilon$。

取极限得 $\nu$ 的绝对连续性，逐集合取极限也给出有限可加性。若 $E_j$ 两两不交，令 $E=\bigcup_jE_j$、$R_N=\bigcup_{j>N}E_j$。因 $m(S)<\infty$，$m(R_N)\to0$，故
\[
 \nu(E)-\sum_{j=1}^N\nu(E_j)=\nu(R_N)\longrightarrow0.
\]
这证明可数可加性。有限值复测度的全变差自动有限；此处无需额外假定变差对 $n$ 一致有界，这也对应中译本第 73 页的校注。
\end{proof}

\begin{corollary}[一致可数可加性与无限参考测度]\label{cor:y2-nikodym}
若 $\nu_n(E)$ 对每个可测集收敛，无需预给参考测度，则对 $E_j\downarrow\varnothing$ 有 $\sup_n|\nu_n(E_j)|\to0$，逐集合极限是复测度。若各 $\nu_n$ 对任意正测度 $m$ 满足上述绝对连续性，则极限也对 $m$ 绝对连续，即使 $m(S)=\infty$。
\end{corollary}
\begin{proof}
定义有限正控制测度
\[
 \mu(E)=\sum_{n=1}^\infty2^{-n}\frac{|\nu_n|(E)}{1+|\nu_n|(S)}.
\]
非负求和可交换，故 $\mu$ 可数可加且 $\mu(S)\le1$。每个 $|\nu_n|\le2^n(1+|\nu_n|(S))\mu$，从而各 $\nu_n$ 都对 $\mu$ 绝对连续。应用定理，$\mu(E_j)\to0$ 导出一致从上连续性和极限的可数可加性。分母的 $1+$ 处理了零测度。

若另有 $m$，在 $\{[E]:m(E)<\infty\}$ 上重复定理的闭集和 Baire 论证，便得对 $m$ 的统一绝对连续性；所有 $m(B)<\delta$ 的集合本来就在该空间内。因此极限也对 $m$ 绝对连续。可数可加性已由有限控制测度获得，不需要在无限 $m$ 下错误地断言所有尾集的 $m$ 测度趋零。
\end{proof}

\begin{example}[密度几乎处处收敛不能替代逐集合收敛]
在 $[0,1]$ 上令 $f_n=n\mathbf1_{(0,1/n)}$、$\nu_n(E)=\int_Ef_n\,dx$。虽然 $f_n\to0$ 几乎处处，取 $E_n=(0,1/n)$ 却有 $m(E_n)\to0$、$\nu_n(E_n)=1$。所以该列不一致绝对连续，由定理它不可能对每个可测集都有极限。
\end{example}

\begin{exercise}
证明定理的假设还蕴含：对每个 $\varepsilon>0$ 存在 $\delta>0$，使 $m(E)<\delta$ 时 $\sup_n|\nu_n|(E)<\varepsilon$。
\end{exercise}
\begin{solution}
先证一般估计 $|\rho|(E)\le4\sup_{B\subset E}|\rho(B)|$，其中 $B$ 可测。实带符号测度 $a$ 的正负分解给出 $|a|(E)\le2\sup_{B\subset E}|a(B)|$。写 $\rho=a+ib$，由分割定义 $|\rho|(E)\le|a|(E)+|b|(E)$，即得估计。若 $m(E)<\delta$，所有 $B\subset E$ 也满足此条件。将定理的误差取为 $\varepsilon/8$，得 $\sup_n|\nu_n|(E)\le\varepsilon/2<\varepsilon$。
\end{solution}

\section{广义函数序列的逐项可微性}\label{sec:yosida-02-03}
\noindent\textit{原书位置：中译本第 75 页；英文版第 72 页。}

\begin{theorem}[分布的逐测试函数极限]\label{thm:y2-distribution-limit}
设 $T_n\in\Dist(\Omega)$，对每个 $\varphi\in\Test(\Omega)$，$T(\varphi)=\lim_n\Pair{T_n}{\varphi}$ 存在且有限。则 $T\in\Dist(\Omega)$。对每个 $K\Subset\Omega$，存在 $C_K>0$ 和整数 $m_K\ge0$，使
\[
 \sup_n|\Pair{T_n}{\varphi}|\le C_K
 \max_{|\alpha|\le m_K}\sup_K|\partial^\alpha\varphi|
 \quad(\varphi\in\Test_K).
\]
\end{theorem}
\begin{proof}
第一章已证明 $\Test_K$ 是 Fr\'echet 空间。各 $T_n|_{\Test_K}$ 为连续线性泛函，逐点收敛使它们逐点有界。由定理 \ref{thm:y2-banach-steinhaus} 得等度连续性，再把有限多个导数半范数合并，即得所述估计。取极限得到 $T$ 的同一估计；线性直接由极限得到。对每个紧集的有限阶估计正是第一章的分布连续性判别，所以 $T$ 为分布。这里在各 $\Test_K$ 上使用 Baire 定理，并未假定整个 LF 空间 $\Test(\Omega)$ 是 Fr\'echet 空间。
\end{proof}

\begin{corollary}[逐项求导与光滑乘法]\label{cor:y2-termwise}
若 $T_n\to T$ 是逐测试函数的分布收敛，则对任意多重指标 $\alpha$ 及 $a\in C^\infty(\Omega)$，$\partial^\alpha T_n\to\partial^\alpha T$，且 $aT_n\to aT$。
\end{corollary}
\begin{proof}
对固定 $\varphi$，$\partial^\alpha\varphi,a\varphi\in\Test(\Omega)$，且
\[
 \Pair{\partial^\alpha T_n}{\varphi}
 =(-1)^{|\alpha|}\Pair{T_n}{\partial^\alpha\varphi}
 \longrightarrow(-1)^{|\alpha|}\Pair{T}{\partial^\alpha\varphi}.
\]
乘法同理。该收敛不等于导数作为普通函数逐点或一致收敛。
\end{proof}

\begin{example}[正则化阶跃与 Dirac 质量]
取 $\rho\in\Test(\R)$、$\rho\ge0$、$\int\rho=1$，令 $\rho_\varepsilon(x)=\varepsilon^{-1}\rho(x/\varepsilon)$，$H_\varepsilon(x)=\int_{-\infty}^x\rho_\varepsilon(t)\,dt$。由有界收敛，$H_\varepsilon\to H=\mathbf1_{(0,\infty)}$ 于分布意义；故 $H_\varepsilon'=\rho_\varepsilon\to H'=\delta_0$。求导后极限已不是普通局部可积函数。
\end{example}

\begin{exercise}
令 $u_n(x)=\sin(nx)$。证明 $u_n\to0$、$u_n'\to0$ 于分布意义，而 $u_n'$ 不逐点收敛到零。说明分布级数的逐项求导需要什么假设。
\end{exercise}
\begin{solution}
分部积分得 $|\int\sin(nx)\varphi(x)\,dx|\le n^{-1}\|\varphi'\|_1\to0$。推论给出 $u_n'\to0$，但 $u_n'(0)=n$。若级数部分和 $S_N=\sum_{j=1}^NT_j$ 在每个测试函数上收敛到有限值，则定理保证极限为分布，并且 $\partial^\alpha S_N\to\partial^\alpha S$。无需级数以函数范数绝对收敛；也不能只凭各项可微便断言级数可逐项求导。
\end{solution}

\section{奇点凝聚原理}\label{sec:yosida-02-04}
\noindent\textit{原书位置：中译本第 75--77 页；英文版第 72--75 页。}

\begin{theorem}[有界点集合的二择性]\label{thm:y2-bounded-points}
设 $X$ 为 Banach 空间，$T_j:X\to Y_j$ 为有界线性算子，$Y_j$ 赋范。则
\[
 B=\{x:\sup_j\|T_jx\|<\infty\}
\]
或者等于 $X$，或者是第一纲集。
\end{theorem}
\begin{proof}
写 $B=\bigcup_m E_m$，其中 $E_m=\{x:\sup_j\|T_jx\|\le m\}$ 闭。若 $B$ 不是第一纲集，某 $E_m$ 有内点。与共鸣定理相同的平移相减和缩放给出 $\sup_j\|T_j\|<\infty$，故 $B=X$。反之，各 $E_m$ 均闭且无内点，便是无处稠密集。
\end{proof}

\begin{corollary}[Banach--Steinhaus 奇点凝聚原理]\label{cor:y2-condensation}
设 $T_{p,q}:X\to Y_p$ 为有界线性算子，$p,q\ge1$。若对每个 $p$ 存在 $x_p$ 使 $\sup_q\|T_{p,q}x_p\|=\infty$，则
\[
 G=\{x:\sup_q\|T_{p,q}x\|=\infty\ \text{对所有 }p\}
\]
是稠密的余第一纲集。这里可等价地把上确界无穷写为 $\limsup_{q\to\infty}\|T_{p,q}x\|=\infty$，不能换成普通极限。
\end{corollary}
\begin{proof}
对每个 $p$，定理中的有界点集合 $B_p$ 不等于 $X$，所以是第一纲集。$G=X\setminus\bigcup_pB_p$ 的补集仍是第一纲集，由 Baire 定理得稠密性。
\end{proof}

\begin{theorem}[处处无有限导数的连续函数：补全锯齿论证]\label{thm:y2-nowhere-differentiable}
在 $C([0,1];\R)$ 中，处处无有限导数的函数构成一个包含稠密余第一纲集的集合；端点的导数按单侧理解。
\end{theorem}
\begin{proof}
对 $m\ge1$ 令
\[
 M_m=\left\{f:\exists t\in[0,1/2],\quad
 |f(t+h)-f(t)|\le mh\ \text{对所有 }0<h<1/2\right\}.
\]
若 $f_j\to f$ 一致且 $f_j\in M_m$，由相应 $t_j$ 的紧性可取 $t_j\to t$。固定 $h\in(0,1/2)$，连续性和一致收敛使不等式传到 $t$，故 $f\in M_m$；所以 $M_m$ 闭。

现证 $M_m$ 无内点。给定 $f$ 和误差 $\varepsilon$，先用分段线性 $g$ 一致逼近 $f$ 到 $\varepsilon/2$，记其有限个斜率绝对值的最大值为 $L$。在 $g$ 的每个线性区间上细分，叠加端点为零、振幅小于 $\varepsilon/2$ 的连续三角波 $w$，使每段斜率为 $\pm M$ 且 $M>L+m$；只需把三角波周期取得足够短，并使每个原区间包含整数个周期。于是 $g+w$ 的每个线性小段斜率的绝对值都大于 $m$。对每个 $t\in[0,1/2]$，即使 $t$ 是顶点，也有一个充分短的右侧线性区间，其差商绝对值大于 $m$。所以 $g+w\notin M_m$，而 $\|g+w-f\|_\infty<\varepsilon$。

因此 $G_+=C([0,1])\setminus\bigcup_mM_m$ 是稠密余第一纲集。若 $f$ 在 $t\in[0,1/2]$ 有有限右导数，则小 $h$ 的差商有界；较大 $h$ 的差商由 $2\|f\|_\infty/h$ 控制，故 $f\in M_m$ 对某个 $m$ 成立。由此 $G_+$ 中函数在左半区间无有限右导数。反射 $Rf(t)=f(1-t)$ 是同胚，$G_+\cap R(G_+)$ 仍为稠密余第一纲集，其函数在整个区间无有限导数。
\end{proof}

\begin{example}[连续周期函数的 Fourier 部分和可以无界]\label{ex:y2-fourier-divergence}
在实 Banach 空间 $C(\mathbb T)$ 上使用标准 Dirichlet 核
\[
 D_N(s)=\sum_{k=-N}^Ne^{iks}
 =\frac{\sin((N+1/2)s)}{\sin(s/2)},\qquad
 S_Nf(t)=\frac1{2\pi}\int_{-\pi}^{\pi}f(t-s)D_N(s)\,ds.
\]
对固定 $t$，$f\mapsto S_Nf(t)$ 为有界实线性泛函，其范数为 $(2\pi)^{-1}\|D_N\|_1$。上界来自积分估计；下界可取连续周期函数近似 $\operatorname{sgn}D_N$，例如 $D_N/\sqrt{D_N^2+\delta^2}$，再令 $\delta\downarrow0$，并相应平移反射。

这些范数无界。事实上 $\sin(s/2)\le s/2$ 给出
\[
 \int_0^\pi|D_N(s)|\,ds
 \ge2\int_0^{(N+1/2)\pi}\frac{|\sin v|}{v}\,dv
 \ge c\sum_{j=1}^{N-1}\frac1{j+1},
\]
最后一步在 $[j\pi+\pi/6,j\pi+5\pi/6]$ 上用 $|\sin v|\ge1/2$；$c>0$ 与 $N$ 无关。因此由共鸣定理的逆否命题，对每个固定 $t$ 存在 $f_t$ 使 $\sup_N|S_Nf_t(t)|=\infty$。

给定任意可数点集，加入一个可数稠密集得到 $\{t_p\}$。奇点凝聚原理给出一个连续 $f$，使上述无界性在全部 $t_p$ 同时发生。对这个 $f$，令
\[
 F_m=\bigcap_N\{t:|S_Nf(t)|\le m\}.
\]
各 $F_m$ 闭，且不可能含非空开弧，否则与 $\{t_p\}$ 稠密且部分和无界矛盾。所以 $P=\mathbb T\setminus\bigcup_mF_m$ 是余第一纲集，并包含所有预先指定的点。在 $P$ 上有 $\limsup_N|S_Nf(t)|=\infty$。

补充说明：$P$ 的势为连续统，无需连续统假设。写 $P=\bigcap_mO_m$ 为稠密开集之交，在任意开弧内递归选择两两分离的闭子弧，令第 $m$ 层均落在 $O_1\cap\cdots\cap O_m$，直径趋零。每个二进制无限序列确定一个不同的交点，因而向 $P$ 嵌入 $\{0,1\}^{\mathbb N}$；反向上界由 $P\subset\mathbb T$ 给出。这展开了原书关于无须连续统假设的注记。
\end{example}

\begin{exercise}
设 $X$ 为 Banach 空间，$T_n:X\to Y_n$ 有界线性且 $\sup_n\|T_n\|=\infty$。证明存在稠密余第一纲集 $G$，使每个 $x\in G$ 满足 $\limsup_n\|T_nx\|=\infty$。该结论是否保证 $\|T_nx\|\to\infty$？
\end{exercise}
\begin{solution}
若有界点集合等于 $X$，共鸣定理给出算子范数一致有界，矛盾。因此有界点集合为第一纲集，其补集就是所求 $G$。不保证普通极限：在 $X=Y_n=\R$ 上取 $T_{2k}=0$、$T_{2k-1}x=kx$。对 $x\ne0$ 上极限无穷，而偶数子列恒为零。
\end{solution}

\section{开映射定理}\label{sec:yosida-02-05}
\noindent\textit{原书位置：中译本第 77--79 页；英文版第 75--77 页。}

\begin{lemma}[像的闭包含零邻域]\label{lem:y2-almost-open}
设 $T:X\to Y$ 连续线性，且其值域 $R(T)$ 在 $Y$ 中不是第一纲集。那么对每个 $X$ 的零邻域 $U$，$\overline{T(U)}$ 包含 $Y$ 的某个零邻域。
\end{lemma}
\begin{proof}
取对称零邻域 $W$ 使 $W-W\subset U$。零邻域吸收每个向量，所以 $X=\bigcup_{n\ge1}nW$，$R(T)\subset\bigcup_n n\overline{T(W)}$。若所有这些闭集无内点，值域便是第一纲集，矛盾。故 $\overline{T(W)}$ 有内点。取非空开集 $O\subset\overline{T(W)}$，则 $O-O$ 为零邻域，且
\[
 O-O\subset\overline{T(W)}-\overline{T(W)}
 \subset\overline{T(W-W)}\subset\overline{T(U)}.
\]
中间包含关系由减法连续性得到。
\end{proof}

\begin{theorem}[Banach 开映射定理]\label{thm:banach-open}
设 $X,Y$ 为 $F$ 空间，$T:X\to Y$ 为连续线性满射。则 $T$ 是开映射。特别地，它适用于 Banach 空间和 Fr\'echet 空间。
\end{theorem}
\begin{proof}
取给出拓扑的完备平移不变度量 $d_X,d_Y$，以 $B_X(r),B_Y(r)$ 表示零点开球。$Y$ 完备，故不是第一纲集。由引理，对每个 $r>0$，$\overline{T(B_X(r))}$ 含某零球。

固定 $\varepsilon>0$，令 $r_j=\varepsilon2^{-j-2}$，$j\ge0$。选择 $\eta_j>0$ 使
\[
 B_Y(\eta_j)\subset\overline{T(B_X(r_j))},\qquad
 \eta_j\le2^{-j}.
\]
任取 $y\in B_Y(\eta_0)$。由闭包定义可取 $x_0\in B_X(r_0)$，使残差 $y-Tx_0\in B_Y(\eta_1)$；再取 $x_1\in B_X(r_1)$，使 $y-T(x_0+x_1)\in B_Y(\eta_2)$。如此继续得到
\[
 d_X(x_j,0)<r_j,\qquad
 d_Y\left(y-T\sum_{j=0}^Nx_j,0\right)<\eta_{N+1}.
\]
平移不变度量的三角不等式使部分和为 Cauchy 序列，故在 $X$ 中趋于 $x$，且 $d_X(x,0)\le\sum_jr_j=\varepsilon/2<\varepsilon$。连续性和 $\eta_j\to0$ 给出 $Tx=y$。因此 $B_Y(\eta_0)\subset T(B_X(\varepsilon))$。平移到任意点即得开性。此证明只用度量的三角不等式，未假定度量对标量绝对齐次，也未用局部凸性。
\end{proof}

\begin{corollary}[有界逆与商空间]\label{cor:y2-inverse}
连续线性双射 $T:X\to Y$ 在 $F$ 空间之间具有连续逆。若 $X,Y$ 为 Banach 空间，则逆有界。若 $T$ 只是 Banach 空间之间的有界满射，则存在 $C>0$，对每个 $y$ 可选 $Tx=y$ 且 $\|x\|\le C\|y\|$；诱导映射 $X/\ker T\to Y$ 是 Banach 空间同构。
\end{corollary}
\begin{proof}
逆像开集的条件正是 $T$ 的开性。对满射，由开映射得 $B_Y(\delta)\subset T(B_X(1))$。当 $y\ne0$ 时，将 $y$ 缩放至范数 $\delta/2$ 再选原像，得 $\|x\|<2\|y\|/\delta$；$y=0$ 取 $x=0$。核闭，商空间由第一章知为 Banach，诱导映射为有界双射，再用有界逆结论。这里选取原像的估计并不自动产生一个线性的右逆。
\end{proof}

\begin{example}[等价范数与不完备反例]
同一线性空间上的两个 Banach 范数若满足 $\|x\|_1\le C\|x\|_2$，则反向估计也成立：对恒等双射应用有界逆定理。完备性缺失时可失败。在 $C^1[0,1]$ 上，$\|f\|_{C^1}=\|f\|_\infty+\|f'\|_\infty$ 完备，而 $\|f\|_\infty$ 不完备；恒等映射从前者到后者连续双射，逆不连续，因为 $f_n(t)=\sin(nt)/n$ 一致趋零，导数却不一致趋零。
\end{example}

\begin{exercise}
设 $T:X\to Y$ 为 Banach 空间之间的有界单射。证明：$R(T)$ 闭当且仅当存在 $c>0$，使 $\|Tx\|\ge c\|x\|$ 对所有 $x$ 成立。
\end{exercise}
\begin{solution}
若值域闭，$R(T)$ 按 $Y$ 的范数完备，有界逆定理用于 $T:X\to R(T)$，得所需下界。反之，若 $Tx_n\to y$，下界使 $x_n$ 为 Cauchy 序列，源空间完备性给出 $x_n\to x$。连续性使 $Tx=y$，故值域闭。单射本身不足，例如 $T:\ell^2\to\ell^2$，$(Tx)_j=x_j/j$；有限支集序列在值域内，但 $(1/j)_j$ 是值域的极限点而不在值域内。
\end{solution}

\section{闭图像定理}\label{sec:yosida-02-06}
\noindent\textit{原书位置：中译本第 79--81 页；英文版第 77--79 页。}

\begin{definition}[闭算子与可闭算子]
线性算子 $T:D(T)\subset X\to Y$ 的图像是 $G(T)=\{(x,Tx):x\in D(T)\}$。若它在 $X\times Y$ 中闭，称 $T$ 闭；若 $\overline{G(T)}$ 仍为某线性算子的图像，称 $T$ 可闭，该算子记为 $\overline T$。对可度量空间，闭性等价于：$x_n\in D(T)$、$x_n\to x$、$Tx_n\to y$ 蕴含 $x\in D(T)$ 且 $Tx=y$。
\end{definition}

\begin{proposition}[可闭判别与逆算子的闭性]\label{prop:y2-closable}
在可度量拓扑向量空间之间，$T$ 可闭当且仅当 $x_n\in D(T)$、$x_n\to0$、$Tx_n\to y$ 蕴含 $y=0$。闭单射的逆 $T^{-1}:R(T)\to X$ 也是闭算子。
\end{proposition}
\begin{proof}
若图像闭包为图像，其中 $(0,y)$ 只能有 $y=0$。反向，$\overline{G(T)}$ 是闭线性子空间，若同时含 $(x,y_1),(x,y_2)$，相减得 $(0,y_1-y_2)$。可度量性允许用图像中的序列逼近它，假设即给出 $y_1=y_2$。所以闭包确为某线性算子的图像，且是包含 $G(T)$ 的最小闭图像。逆算子的图像由同胚 $(x,y)\mapsto(y,x)$ 得到，所以仍闭。
\end{proof}

\begin{theorem}[Banach 闭图像定理]\label{thm:banach-closed-graph}
设 $X,Y$ 为 $F$ 空间，全定义线性算子 $T:X\to Y$ 若闭，则连续。对 Banach 空间，结论等价于存在 $C$ 使 $\|Tx\|\le C\|x\|$。
\end{theorem}
\begin{proof}
积空间 $X\times Y$ 由乘积度量完备，$G(T)$ 是其闭线性子空间，因而也是 $F$ 空间。投影 $\pi_X:G(T)\to X$ 连续线性双射，因为 $D(T)=X$。由开映射定理其逆连续，而 $T=\pi_Y\pi_X^{-1}$ 连续。
\end{proof}

\begin{proposition}[图像范数与算子比较]\label{prop:y2-graph-comparison}
若 $X,Y_1,Y_2$ 为 Banach 空间，$T_1:D(T_1)\subset X\to Y_1$ 闭，则
\[
 \|x\|_{T_1}=\|x\|_X+\|T_1x\|_{Y_1}
\]
使 $D(T_1)$ 成为 Banach 空间。若 $T_2:D(T_2)\subset X\to Y_2$ 可闭且 $D(T_1)\subset D(T_2)$，则存在 $C>0$，使
\[
 \|T_2x\|_{Y_2}\le C(\|x\|_X+\|T_1x\|_{Y_1})\quad(x\in D(T_1)).
\]
\end{proposition}
\begin{proof}
图像范数的 Cauchy 序列同时使 $x_n$ 和 $T_1x_n$ 收敛；闭性保证两个极限仍构成图像中的点，所以范数完备。对 $T_2:D(T_1)\to Y_2$，若 $x_n\to x$ 于图像范数且 $T_2x_n\to y$，则 $x\in D(T_2)$，并有 $x_n-x\to0$、$T_2(x_n-x)\to y-T_2x$。可闭判别给出 $y=T_2x$。故该映射图像闭，由闭图像定理有界。这是原书归于 H\"ormander 的算子比较定理；和范数与平方和范数等价。
\end{proof}

\begin{example}[闭而无界的微分算子]
在 $X=Y=C[0,1]$ 上取 $D(T)=C^1[0,1]$、$Tf=f'$。若 $f_n\to f$、$f_n'\to g$ 都一致，恒等式
\[
 f_n(t)-f_n(0)=\int_0^tf_n'(s)\,ds
\]
传到极限得 $f(t)-f(0)=\int_0^tg(s)\,ds$，故 $f\in C^1$ 且 $f'=g$，所以 $T$ 闭。另一方面 $f_n(t)=t^n$ 满足 $\|f_n\|_\infty=1$、$\|Tf_n\|_\infty=n$，故按 $C$ 范数无界。它的定义域不是整个 $C[0,1]$；换用图像范数以后它是有界的。
\end{example}

\begin{example}[微分算子的可闭性与形式转置]
设 $P=\sum_{|\alpha|\le k}a_\alpha(x)\partial^\alpha$，$a_\alpha\in C^k(\Omega)$。在 $L^2(\Omega)$ 中取
\[
 D(P)=\{u\in C^k(\Omega)\cap L^2(\Omega):Pu\in L^2(\Omega)\}.
\]
若 $u_n\to0$、$Pu_n\to g$ 于 $L^2$，对 $\varphi\in\Test(\Omega)$ 分部积分得
\[
 \int g\varphi=\lim_n\int(Pu_n)\varphi
 =\lim_n\int u_nP^t\varphi=0,\qquad
 P^t\varphi=\sum_{|\alpha|\le k}(-1)^{|\alpha|}\partial^\alpha(a_\alpha\varphi).
\]
紧支集消去边界项，$P^t\varphi\in L^2$ 允许用 Schwarz 不等式取极限。因此 $g=0$，$P$ 可闭。这里的双线性配对不加复共轭，$P^t$ 是形式转置；Hilbert 伴随的系数则须取共轭。
\end{example}

\begin{exercise}
在 $L^2(\R)$ 上定义 $D(T)=\{u\in L^2:xu\in L^2\}$、$(Tu)(x)=xu(x)$。证明 $T$ 闭而无界，并说明其图像范数空间为什么完备。
\end{exercise}
\begin{solution}
若 $u_n\to u$、$xu_n\to v$ 于 $L^2$，可逐次选取同一子列使两者都几乎处处收敛，于是 $v=xu$ 几乎处处。故 $u\in D(T)$ 且 $Tu=v$，图像闭。取 $u_n=\mathbf1_{[n,n+1]}$，有 $\|u_n\|_2=1$、$\|Tu_n\|_2\ge n$，所以无界。图像范数完备由命题 \ref{prop:y2-graph-comparison}，也可对两部分的 Cauchy 序列用刚才的闭性论证。
\end{solution}

\section{闭图像定理的一个应用（H\"ormander 定理）}\label{sec:yosida-02-07}
\noindent\textit{原书位置：中译本第 82--83 页；英文版第 80--81 页。}

本节展示一个有代表性的思路：先由定性光滑性保证导数算子处处有定义，再以闭图像定理得到统一估计，最后代入指数函数，把估计变成多项式零点的限制。

\begin{definition}[本节使用的亚椭圆性]
令 $D_j=(1/i)\partial_{x_j}$，$P(D)$ 为常系数偏微分算子，$P(\zeta)$ 为对应复多项式。本节沿用原书的 $L^2_{\Loc}$ 表述：若在任意开子域上，$u\in L^2_{\Loc}$ 且 $P(D)u$ 在某开集光滑，必有 $u$ 在该开集具有光滑代表，则称 $P(D)$ 亚椭圆。现代通常要求同一结论对所有分布 $u$ 成立；该较强定义当然蕴含本节所用性质。下面的必要条件仅需 $L^2_{\Loc}$ 形式。
\end{definition}

分布等式使用形式转置 $P^t(D)=P(-D)$：$P(D)u=0$ 指 $\int uP(-D)\varphi=0$ 对所有测试函数成立；不加入 Hilbert 配对的复共轭。

\begin{theorem}[H\"ormander 的复特征根必要条件]\label{thm:y2-hormander}
若 $P(D)$ 亚椭圆，则对每个 $A>0$ 存在 $B_A<\infty$，使
\[
 P(\zeta)=0,\quad\zeta=\xi+i\eta\in\mathbb C^n,\quad
 |\eta|\le A\quad\Longrightarrow\quad|\zeta|\le B_A.
\]
等价地，零点集上的 $|\zeta|\to\infty$ 必迫使 $|\operatorname{Im}\zeta|\to\infty$。这采用中译本第 82 页校注中修正的常数次序。
\end{theorem}
\begin{proof}
选非空有界开集 $V\Subset W\subset\R^n$，并使 $W$ 有界。令
\[
 U=\{u\in L^2(W):P(D)u=0\text{ 于分布意义}\}.
\]
对每个测试函数，$u\mapsto\int_WuP(-D)\varphi$ 在 $L^2(W)$ 上连续，故其核闭；$U$ 为这些核的交，所以是 Banach 空间。亚椭圆性保证每个 $u\in U$ 在 $W$ 光滑，从而 $\partial_ju|_V\in L^2(V)$。

算子 $A_j:U\to L^2(V)$、$A_ju=\partial_ju|_V$ 图像闭：若 $u_l\to u$ 于 $L^2(W)$、$A_ju_l\to v$ 于 $L^2(V)$，对 $\varphi\in\Test(V)$ 分部积分取极限得
\[
 \int_Vv\varphi=-\int_Vu\partial_j\varphi,
\]
所以 $v=\partial_ju$。闭图像定理给出常数 $C$，使
\[
 \sum_{j=1}^n\|\partial_ju\|_{L^2(V)}^2
 \le C\|u\|_{L^2(W)}^2\quad(u\in U).
\]
若 $P(\zeta)=0$，取 $u(x)=e^{ix\cdot\zeta}$。因 $W$ 有界，$u\in L^2(W)$ 且 $P(D)u=P(\zeta)u=0$，代入得
\[
 |\zeta|^2\int_Ve^{-2x\cdot\eta}\,dx
 \le C\int_We^{-2x\cdot\eta}\,dx.
\]
选择 $R$ 使 $W\subset B(0,R)$。若 $|\eta|\le A$，左边积分不小于 $|V|e^{-2RA}$，右边积分不大于 $|W|e^{2RA}$，所以
\[
 |\zeta|^2\le C\frac{|W|}{|V|}e^{4RA}.
\]
右边与具体零点无关，取平方根即得 $B_A$。这补全了原书“虚部有界则根有界”的积分比值步骤。
\end{proof}

\begin{example}[Laplace 算子与遗漏方向的导数]
对 Laplace 算子取 $P(\zeta)=-\sum_j\zeta_j^2$。若 $P(\xi+i\eta)=0$，比较实虚部得 $|\xi|^2=|\eta|^2$、$\xi\cdot\eta=0$，所以 $|\zeta|^2=2|\eta|^2$，满足必要条件。但必要条件本身不构成 Weyl 引理的证明；其充分性留待后续理论。

在 $\R^2$ 上，$P(D)=D_1$ 的零点包含 $(0,t)$，$t\in\R$，虚部恒为零而模无界，因而不亚椭圆。也可直接取 $u(x_1,x_2)=|x_2|$：它局部平方可积且 $D_1u=0$，却不光滑。这里的缺陷来自算子没有控制第二个方向。
\end{example}

\begin{exercise}
对常系数算子 $P(D)$，假定每个局部平方可积的齐次解均光滑。证明对任意多重指标 $\alpha$ 及非空有界开集 $V\Subset W$，存在 $C_{\alpha,V,W}$，使
\[
 \|\partial^\alpha u\|_{L^2(V)}\le C_{\alpha,V,W}\|u\|_{L^2(W)}
 \quad(u\in L^2(W),\ P(D)u=0).
\]
这些常数能否由本章论证直接算出？
\end{exercise}
\begin{solution}
齐次解空间 $U$ 在 $L^2(W)$ 中闭。光滑性使 $u\mapsto\partial^\alpha u|_V$ 全定义。若 $u_l\to u$ 于 $U$ 且 $\partial^\alpha u_l\to v$ 于 $L^2(V)$，对测试函数配对得 $\int v\varphi=(-1)^{|\alpha|}\int u\partial^\alpha\varphi$，故 $v=\partial^\alpha u$。闭图像定理给出估计。该论证保证常数存在，并未提供可计算的最佳常数或其随 $\alpha$、区域变化的定量规律。
\end{solution}


\section*{本章参考文献与阅读衔接}
\phantomsection
\addcontentsline{toc}{section}{本章参考文献与阅读衔接}
本章直接核对的资料是两版吉田第二章 \cite{yosida-zh,yosida-en}。原书章末列出 Banach、Bourbaki、Dunford--Schwartz、Hille--Phillips 与 H\"ormander 的著作；它们是原书的进一步阅读线索，本稿未将其列作已直接查阅的文献。中译本第 73 页关于复测度及绝对连续性的校注、第 82 页关于常数次序的校注，已在相应节说明。

本章八道习题均为补充学习题，解答统一输出于书后。第三章转向 Hilbert 空间的投影与表示；分布极限与闭微分算子的例子可与第 \ref{ch:distributions}、\ref{ch:operations}、\ref{ch:pde} 章对读。H\"ormander 特征根条件的充分性和一般 Weyl 引理不在本章证明范围内。

% !TeX encoding = UTF-8
% 吉田第六版：中译本第 3 章；此文件当前仅为目录与正文占位。
\SetChapterInfo{subtitle={吉田耕作第六版主线 · 正文待撰写}}
\chapter{正交投影及 F. Riesz 表示定理}\label{ch:yosida-03}
\YosidaChapterPlaceholder{3}{The Orthogonal Projection and F. Riesz' Representation Theorem}{85--106}{81--101}

\section{正交投影}\label{sec:yosida-03-01}
\YosidaSectionPlaceholder{85}{81}
% TODO: 定义与记号；定理与完整证明；例子；习题；中英文核对。

\section{“殆正交”元}\label{sec:yosida-03-02}
\YosidaSectionPlaceholder{88}{84}
% TODO: 定义与记号；定理与完整证明；例子；习题；中英文核对。

\section{Ascoli--Arzelà 定理}\label{sec:yosida-03-03}
\YosidaSectionPlaceholder{89}{85}
% TODO: 定义与记号；定理与完整证明；例子；习题；中英文核对。

\section{正交基、Bessel 不等式与 Parseval 等式}\label{sec:yosida-03-04}
\YosidaSectionPlaceholder{90}{86}
% TODO: 定义与记号；定理与完整证明；例子；习题；中英文核对。

\section{E. Schmidt 正交化}\label{sec:yosida-03-05}
\YosidaSectionPlaceholder{92}{88}
% TODO: 定义与记号；定理与完整证明；例子；习题；中英文核对。

\section{F. Riesz 表示定理}\label{sec:yosida-03-06}
\YosidaSectionPlaceholder{93}{90}
% TODO: 定义与记号；定理与完整证明；例子；习题；中英文核对。

\section{Lax--Milgram 定理}\label{sec:yosida-03-07}
\YosidaSectionPlaceholder{95}{92}
% TODO: 定义与记号；定理与完整证明；例子；习题；中英文核对。

\section{Lebesgue--Nikodym 定理的一个证明}\label{sec:yosida-03-08}
\YosidaSectionPlaceholder{97}{93}
% TODO: 定义与记号；定理与完整证明；例子；习题；中英文核对。

\section{Aronszajn--Bergman 再生核}\label{sec:yosida-03-09}
\YosidaSectionPlaceholder{99}{95}
% TODO: 定义与记号；定理与完整证明；例子；习题；中英文核对。

\section{P. Lax 的负指数范数}\label{sec:yosida-03-10}
\YosidaSectionPlaceholder{102}{98}
% TODO: 定义与记号；定理与完整证明；例子；习题；中英文核对。

\section{广义函数的局部结构}\label{sec:yosida-03-11}
\YosidaSectionPlaceholder{104}{100}
% TODO: 定义与记号；定理与完整证明；例子；习题；中英文核对。
扩展阅读：第 \ref{ch:distributions} 章的局部结构。\par

\section*{本章参考文献整理（占位）}
\phantomsection
\addcontentsline{toc}{section}{本章参考文献整理（占位）}
待核对原书本章参考文献。中译本印刷第 105 页起；后续按实际引用合并至本书参考文献表。

% !TeX encoding = UTF-8
% 吉田第六版：中译本第 4 章；此文件当前仅为目录与正文占位。
\SetChapterInfo{subtitle={吉田耕作第六版主线 · 正文待撰写}}
\chapter{Hahn--Banach 定理}\label{ch:yosida-04}
\YosidaChapterPlaceholder{4}{The Hahn--Banach Theorems}{107--124}{102--118}

\section{实线性空间中的 Hahn--Banach 扩张定理}\label{sec:yosida-04-01}
\YosidaSectionPlaceholder{107}{102}
% TODO: 定义与记号；定理与完整证明；例子；习题；中英文核对。

\section{广义极限}\label{sec:yosida-04-02}
\YosidaSectionPlaceholder{109}{103}
% TODO: 定义与记号；定理与完整证明；例子；习题；中英文核对。

\section{局部凸的完备线性拓扑空间}\label{sec:yosida-04-03}
\YosidaSectionPlaceholder{110}{104}
% TODO: 定义与记号；定理与完整证明；例子；习题；中英文核对。

\section{复线性空间中的 Hahn--Banach 扩张定理}\label{sec:yosida-04-04}
\YosidaSectionPlaceholder{111}{105}
% TODO: 定义与记号；定理与完整证明；例子；习题；中英文核对。

\section{赋范线性空间中的 Hahn--Banach 扩张定理}\label{sec:yosida-04-05}
\YosidaSectionPlaceholder{112}{106}
% TODO: 定义与记号；定理与完整证明；例子；习题；中英文核对。

\section{非平凡连续线性泛函的存在性}\label{sec:yosida-04-06}
\YosidaSectionPlaceholder{113}{107}
% TODO: 定义与记号；定理与完整证明；例子；习题；中英文核对。

\section{线性映射的拓扑}\label{sec:yosida-04-07}
\YosidaSectionPlaceholder{115}{110}
% TODO: 定义与记号；定理与完整证明；例子；习题；中英文核对。

\section{嵌 \texorpdfstring{$X$}{X} 入其二次对偶 \texorpdfstring{$X''$}{X''}}\label{sec:yosida-04-08}
\YosidaSectionPlaceholder{117}{112}
% TODO: 定义与记号；定理与完整证明；例子；习题；中英文核对。

\section{几例对偶空间}\label{sec:yosida-04-09}
\YosidaSectionPlaceholder{119}{114}
% TODO: 定义与记号；定理与完整证明；例子；习题；中英文核对。

\section*{本章参考文献整理（占位）}
\phantomsection
\addcontentsline{toc}{section}{本章参考文献整理（占位）}
待核对原书本章参考文献。中译本印刷第 124 页起；后续按实际引用合并至本书参考文献表。

% !TeX encoding = UTF-8
% 吉田第六版：中译本第 5 章；此文件当前仅为目录与正文占位。
\SetChapterInfo{subtitle={吉田耕作第六版主线 · 正文待撰写}}
\chapter{强收敛和弱收敛}\label{ch:yosida-05}
\YosidaChapterPlaceholder{5}{Strong Convergence and Weak Convergence}{125--150}{119--144}

\section{弱收敛和弱\texorpdfstring{$^*$}{*}收敛}\label{sec:yosida-05-01}
\YosidaSectionPlaceholder{125}{120}
% TODO: 定义与记号；定理与完整证明；例子；习题；中英文核对。

\section{自反 B-空间的局部弱列紧性、一致凸性}\label{sec:yosida-05-02}
\YosidaSectionPlaceholder{131}{126}
% TODO: 定义与记号；定理与完整证明；例子；习题；中英文核对。

\section{Dunford 定理和 Gelfand--Mazur 定理}\label{sec:yosida-05-03}
\YosidaSectionPlaceholder{133}{128}
% TODO: 定义与记号；定理与完整证明；例子；习题；中英文核对。

\section{弱可测性和强可测性、Pettis 定理}\label{sec:yosida-05-04}
\YosidaSectionPlaceholder{135}{130}
% TODO: 定义与记号；定理与完整证明；例子；习题；中英文核对。

\section{Bochner 积分}\label{sec:yosida-05-05}
\YosidaSectionPlaceholder{137}{132}
% TODO: 定义与记号；定理与完整证明；例子；习题；中英文核对。

\section*{本章参考文献整理（占位）}
\phantomsection
\addcontentsline{toc}{section}{本章参考文献整理（占位）}
待核对原书本章参考文献。中译本印刷第 141 页起；后续按实际引用合并至本书参考文献表。

\section*{本章附录：局部凸空间中的弱拓扑和对偶性}
\phantomsection
\addcontentsline{toc}{section}{本章附录：局部凸空间中的弱拓扑和对偶性}
\label{app:yosida-05}
中译本印刷第 141--150 页；英文版印刷第 136--144 页。正文待撰写。
% 中译本总目录未展开以下子目；依据英文版附录目录辅助补足。
\subsection*{极集}
\phantomsection
\addcontentsline{toc}{subsection}{极集}
\label{sec:yosida-05-app-01}
正文待撰写；子目中文名称依据英文版辅助整理。

\subsection*{桶空间}
\phantomsection
\addcontentsline{toc}{subsection}{桶空间}
\label{sec:yosida-05-app-02}
正文待撰写；子目中文名称依据英文版辅助整理。

\subsection*{半自反性与自反性}
\phantomsection
\addcontentsline{toc}{subsection}{半自反性与自反性}
\label{sec:yosida-05-app-03}
正文待撰写；子目中文名称依据英文版辅助整理。

\subsection*{Eberlein--Shmulyan 定理}
\phantomsection
\addcontentsline{toc}{subsection}{Eberlein--Shmulyan 定理}
\label{sec:yosida-05-app-04}
正文待撰写；子目中文名称依据英文版辅助整理。

% !TeX encoding = UTF-8
% 吉田第六版：中译本第 6 章；此文件当前仅为目录与正文占位。
\SetChapterInfo{subtitle={吉田耕作第六版主线 · 正文待撰写}}
\chapter{Fourier 变换和微分方程}\label{ch:yosida-06}
\YosidaChapterPlaceholder{6}{Fourier Transform and Differential Equations}{151--200}{145--192}

\section{速降函数的 Fourier 变换}\label{sec:yosida-06-01}
\YosidaSectionPlaceholder{151}{146}
% TODO: 定义与记号；定理与完整证明；例子；习题；中英文核对。
扩展阅读：第 \ref{ch:fourier} 章。\par

\section{缓增分布的 Fourier 变换}\label{sec:yosida-06-02}
\YosidaSectionPlaceholder{155}{149}
% TODO: 定义与记号；定理与完整证明；例子；习题；中英文核对。
扩展阅读：第 \ref{ch:fourier} 章的温和分布。\par

\section{卷积}\label{sec:yosida-06-03}
\YosidaSectionPlaceholder{161}{156}
% TODO: 定义与记号；定理与完整证明；例子；习题；中英文核对。
扩展阅读：第 \ref{ch:convolution} 章。\par

\section{Paley--Wiener 定理、单边 Laplace 变换}\label{sec:yosida-06-04}
\YosidaSectionPlaceholder{167}{161}
% TODO: 定义与记号；定理与完整证明；例子；习题；中英文核对。

\section{Titchmarsh 定理}\label{sec:yosida-06-05}
\YosidaSectionPlaceholder{172}{166}
% TODO: 定义与记号；定理与完整证明；例子；习题；中英文核对。

\section{Mikusiński 的算子演算}\label{sec:yosida-06-06}
\YosidaSectionPlaceholder{175}{169}
% TODO: 定义与记号；定理与完整证明；例子；习题；中英文核对。

\section{Sobolev 引理}\label{sec:yosida-06-07}
\YosidaSectionPlaceholder{179}{173}
% TODO: 定义与记号；定理与完整证明；例子；习题；中英文核对。
扩展阅读：第 \ref{ch:fourier} 章的 Sobolev 空间。\par

\section{Gårding 不等式}\label{sec:yosida-06-08}
\YosidaSectionPlaceholder{181}{175}
% TODO: 定义与记号；定理与完整证明；例子；习题；中英文核对。

\section{Friedrichs 定理}\label{sec:yosida-06-09}
\YosidaSectionPlaceholder{183}{177}
% TODO: 定义与记号；定理与完整证明；例子；习题；中英文核对。

\section{Malgrange--Ehrenpreis 定理}\label{sec:yosida-06-10}
\YosidaSectionPlaceholder{188}{182}
% TODO: 定义与记号；定理与完整证明；例子；习题；中英文核对。
相关模型：第 \ref{ch:pde} 章的基本解；一般定理仍待本节撰写。\par

\section{具有一致强度的微分算子}\label{sec:yosida-06-11}
\YosidaSectionPlaceholder{193}{188}
% TODO: 定义与记号；定理与完整证明；例子；习题；中英文核对。

\section{亚椭圆性（Hörmander 定理）}\label{sec:yosida-06-12}
\YosidaSectionPlaceholder{195}{189}
% TODO: 定义与记号；定理与完整证明；例子；习题；中英文核对。

% !TeX encoding = UTF-8
\SetChapterInfo{subtitle={由对偶配对、定义域与闭图像理解无界算子},keywords={对偶; 伴随; 自伴; Cayley 变换; 亏空间; 闭值域},prereq={第三章正交投影与 Riesz 表示; 第四章 Hahn--Banach; 闭图像与开映射},goals={区分转置与伴随; 把定义域纳入算子等式; 证明闭值域及可解性判据}}
\chapter{对偶算子}\label{ch:yosida-07}
\ChapterSource{按 \cite{yosida-zh,yosida-en} 第七章五节组织，2026 年 10 月 26 日补写首轮学习正文。原书位置为中译本印刷第 201--218 页、英文版第 193--208 页。实际查看英文全章扫描页及中文五节入口页；范围详见本章 COVERAGE。本稿用 $T'$ 表示线性对偶算子、$T^*$ 表示 Hilbert 伴随，内积对第一变量线性。闭值域证明的图空间步骤使用限制映射与商拓扑，回应中译本第 214 页校注指出的对偶空间识别问题。}

无界算子的公式只有连同定义域才构成算子。本章从连续泛函的延拓定义出发，经图像正交补研究闭性，再由边界条件、Cayley 变换和闭值域刻画自伴性及方程可解性。五节均包含证明、模型与附解答习题；这里的首轮写作不替代独立数学审校。

\Needspace{12\baselineskip}
\section{对偶算子}\label{sec:yosida-07-01}
\noindent\textit{原书起始页：中译本 201；英文版 193。}

本节 $X,Y$ 为 Hausdorff 局部凸空间，$X',Y'$ 表示连续线性对偶，配对写作 $\langle x,f\rangle=f(x)$，对两个变量均线性。若无特别说明，$T:D(T)\subset X\to Y$ 为线性算子，且 $D(T)$ 在 $X$ 中稠密。

\begin{definition}[对偶算子及其定义域]
令
\[
 D(T')=\{g\in Y':\text{存在 }f\in X',\quad g(Tx)=f(x)\ (x\in D(T))\},
 \qquad T'g=f.
\]
稠密性保证 $f$ 唯一：两个连续泛函在稠密子空间相等，便在全空间相等。若 $D(T)$ 不稠密，Hahn--Banach 给出非零泛函在其上为零，此时以上关系一般只能定义多值关系。
\end{definition}

\begin{proposition}[连续算子的转置]\label{prop:y7-transpose}
若 $T:X\to Y$ 连续，则 $D(T')=Y'$，且 $T'g=g\circ T$。$T'$ 对弱星拓扑 $\sigma(Y',Y)\to\sigma(X',X)$ 连续，也对强对偶拓扑 $\beta(Y',Y)\to\beta(X',X)$ 连续。若 $X,Y$ 为赋范空间，则
\[
 \|T'\|=\|T\|.
\]
\end{proposition}
\begin{proof}
复合泛函连续，故定义域为全空间。弱星连续性由 $(T'g)(x)=g(Tx)$ 直接得出。强对偶半范数为 $p_B(f)=\sup_{x\in B}|f(x)|$，其中 $B$ 遍历有界集；由 $T(B)$ 有界及 $p_B(T'g)=p_{T(B)}(g)$ 得连续性。
赋范情形有 $\|T'g\|\le\|T\|\|g\|$。反向由 Hahn--Banach 的对偶范数公式
$\|Tx\|=\sup_{\|g\|\le1}|g(Tx)|\le\|T'\|\|x\|$ 得到。
\end{proof}

\begin{proposition}[代数运算的定义域]
有界算子的线性组合满足 $(aS+bT)'=aS'+bT'$，复数标量也不取共轭。在赋范空间间，若 $T:D(T)\subset X\to Y$ 稠密定义，$S:Y\to Z$ 有界，则
\[
 (ST)'=T'S',\qquad D(T'S')=\{h\in Z':S'h\in D(T')\}.
\]
若 $S,T$ 均稠密定义、可能无界，且 $ST$ 也稠密定义，则一般只有 $T'S'\subset(ST)'$；此处复合的定义域必须逐项保留。
\end{proposition}
\begin{proof}
线性组合代入配对即可。$S$ 有界时 $D(ST)=D(T)$，故 $h(STx)=(S'h)(Tx)$ 能延拓为 $X$ 上连续泛函，当且仅当 $S'h\in D(T')$。两者无界时，对 $h\in D(T'S')$ 仍可作同样计算，但反向未必成立，因为 $D(ST)$ 可能严格小于 $D(T)$。
\end{proof}

\begin{example}[转置不保持算子强收敛]
在 $\ell^2$ 上令 $L(x_1,x_2,\ldots)=(x_2,x_3,\ldots)$，$R(x_1,x_2,\ldots)=(0,x_1,x_2,\ldots)$。有 $L^nx\to0$ 对每个 $x$ 成立。以线性对偶配对 $f_a(x)=\sum_jx_ja_j$ 识别 $\ell^2$ 与其对偶，则 $(L^n)'f_a=f_{R^na}$。非零 $a$ 满足 $\|R^na\|=\|a\|$，所以对偶算子未强收敛到零；但对每个 $x$，$((L^n)'f_a)(x)\to0$，即仍弱星收敛。
\end{example}

\begin{exercise}[矩阵与标量]\label{ex:y7-matrix}
在 $\mathbb C^n$ 上分别用双线性配对 $f_a(x)=\sum_jx_ja_j$ 和第一变量线性的标准内积表示矩阵 $A$ 的对偶算子、Hilbert 伴随。计算 $(iA)'$ 与 $(iA)^*$。
\end{exercise}
\begin{solution}
$A'$ 在 $a$ 坐标下为 $A^{\mathsf T}$；由 $(Ax,y)=(x,A^*y)$ 得 $A^*=\overline A^{\mathsf T}$。所以 $(iA)'=iA'$，$(iA)^*=-iA^*$。共轭来自 Riesz 识别的共轭线性，而不是线性对偶配对。
\end{solution}

\begin{exercise}[分布求导的转置]\label{ex:y7-distribution}
令 $D=\partial_j:\Test(\Omega)\to\Test(\Omega)$。求其对偶算子，并说明它与分布导数的关系。
\end{exercise}
\begin{solution}
测试空间上求导连续，故转置在全部 $\Dist(\Omega)$ 定义，$\langle\varphi,D'u\rangle=\langle\partial_j\varphi,u\rangle$。按分布导数定义，这等于 $-\langle\varphi,\partial_ju\rangle$，故 $D'u=-\partial_ju$。这里使用双线性分布配对，不插入复共轭。
\end{solution}

\Needspace{12\baselineskip}
\section{伴随算子}\label{sec:yosida-07-02}
\noindent\textit{原书起始页：中译本 204；英文版 195。}

本节 $H,K$ 为 Hilbert 空间，内积对第一变量线性。令 $J_H:H\to H'$ 为共轭线性等距映射 $J_Hy=(\,\cdot\,,y)$。对稠密定义 $T:D(T)\subset H\to K$，定义
\[
 D(T^*)=\{y\in K:\exists z\in H,\ (Tx,y)_K=(x,z)_H\quad\forall x\in D(T)\},
 \qquad T^*y=z.
\]
于是 $T^*=J_H^{-1}T'J_K$，其定义域也按此复合确定。等价地，$y\in D(T^*)$ 当且仅当存在 $C_y$ 使 $|(Tx,y)|\le C_y\|x\|$ 对全部 $x\in D(T)$ 成立。这里控制的是 $\|x\|$，并非图范数。

\begin{theorem}[图像、闭性和可闭性]\label{thm:y7-closure}
写 $G(T)=\{(x,Tx):x\in D(T)\}$。在 $K\oplus H$ 中有
\[
 G(T^*)=\{(-Tx,x):x\in D(T)\}^{\perp}.
\]
因此 $T^*$ 总为闭算子。$T$ 可闭当且仅当 $D(T^*)$ 在 $K$ 中稠密，此时 $\overline T=T^{**}$。特别地，稠密定义的 $T$ 闭当且仅当 $T=T^{**}$。
\end{theorem}
\begin{proof}
$(y,z)$ 与所有 $(-Tx,x)$ 正交，等价于 $-(Tx,y)+(x,z)=0$，这正是伴随定义。正交补闭，故 $T^*$ 闭。

再看图像闭包 $M=\overline{G(T)}\subset H\oplus K$。有
\[
 G(T)^\perp=\{(-T^*y,y):y\in D(T^*)\}.
\]
故 $(0,v)\in M=(G(T)^\perp)^\perp$ 当且仅当 $v\perp D(T^*)$。图像闭包为单值算子的图像，当且仅当其竖直截面只有零，即 $D(T^*)$ 稠密。这也是可闭性的充要条件。此时用同一正交补公式再取一次伴随，得到 $G(T^{**})=M$，从而 $T^{**}=\overline T$。
\end{proof}

可闭性的另一表述是：若 $x_n\in D(T)$，$x_n\to0$ 且 $Tx_n\to y$，则 $y=0$。闭图像定理只在 $D(T)=H$ 时推出有界性；稠密定义、闭、无界可以同时成立。若 $T\in\mathcal L(H,K)$，由 Riesz 表示可知 $D(T^*)=K$，且 $\|T^*\|=\|T\|$。算子包含 $S\subset T$ 总导致 $T^*\subset S^*$。

\begin{example}[一个不可闭的稠密定义算子]
在 $H=\ell^2$ 的有限支集序列空间 $c_{00}$ 上令
\[
 Tx=\Big(\sum_{n\ge1}n x_n\Big)e_1.
\]
取 $x^{(n)}=e_n/n$，则 $x^{(n)}\to0$ 而 $Tx^{(n)}=e_1$，故不可闭。由 $(Tx,y)=(\sum n x_n)\overline{y_1}$，该泛函对 $\ell^2$ 范数有界当且仅当 $y_1=0$。于是 $D(T^*)=e_1^\perp$，$T^*y=0$，其定义域果然不稠密。
\end{example}

\begin{exercise}[积分核的伴随]\label{ex:y7-kernel}
设 $k\in L^2((a,b)^2)$，$(Kf)(s)=\int_a^bk(s,t)f(t)\,\dd t$。证明 $K$ 有界并求 $K^*$。
\end{exercise}
\begin{solution}
由积分的 Cauchy--Schwarz 不等式得到 $\|Kf\|_2\le\|k\|_2\|f\|_2$。交换积分的依据为
\[
 \iint|k(s,t)f(t)\overline{g(s)}|\,\dd s\,\dd t
 \le\|k\|_2\|f\|_2\|g\|_2<\infty.
\]
因此 Fubini 定理给出
\[
 (K^*g)(t)=\int_a^b\overline{k(s,t)}g(s)\,\dd s.
\]
若 $k(s,t)=\overline{k(t,s)}$ 几乎处处，则 $K=K^*$。
\end{solution}

\begin{exercise}[图范数与闭算子]\label{ex:y7-graph}
证明闭算子 $T:D(T)\subset H\to K$ 的定义域在 $\|x\|_T=(\|x\|^2+\|Tx\|^2)^{1/2}$ 下完备。说明这为何不推出 $T$ 对原范数有界。
\end{exercise}
\begin{solution}
图范数 Cauchy 列使 $(x_n,Tx_n)$ 在 $H\oplus K$ 中 Cauchy，极限因图像闭而属于图像，故图范数完备。闭图像定理此时给出的是从图范数定义域到 $K$ 的有界性，这已由 $\|Tx\|\le\|x\|_T$ 直接成立。原范数定义域可能不完备，不能偷换两个范数。
\end{solution}

\Needspace{12\baselineskip}
\section{对称算子和自伴算子}\label{sec:yosida-07-03}
\noindent\textit{原书起始页：中译本 205；英文版 197。}

\begin{definition}[三个不同的概念]
稠密定义 $A:D(A)\subset H\to H$ 称为对称，若 $(Ax,y)=(x,Ay)$ 对 $x,y\in D(A)$ 成立，即 $A\subset A^*$；称为自伴，若 $A=A^*$，这同时要求定义域相等；称为本质自伴，若可闭且 $\overline A$ 自伴。对称算子可闭，闭包仍对称，因为 $A\subset A^*$ 给出 $\overline A=A^{**}\subset A^*$，再用 $(\overline A)^*=A^*$。
\end{definition}

\begin{example}[实乘法算子]
在 $L^2(\mathbb R)$ 上令 $M_mf=mf$，$m$ 为几乎处处有限的实可测函数，取最大定义域 $D(M_m)=\{f\in L^2:mf\in L^2\}$。这是自伴算子。
\end{example}
\begin{proof}
截断 $E_n=\{|t|\le n,\ |m(t)|\le n\}$，则 $\mathbf1_{E_n}f\to f$ 于 $L^2$，证明定义域稠密。$m$ 实保证对称。若 $g\in D(M_m^*)$ 且 $M_m^*g=h$，在 $E_n$ 上用任意 $L^2$ 探针得 $h=mg$ 几乎处处于 $E_n$。取并集得 $mg=h\in L^2$，故 $g\in D(M_m)$。于是两个定义域相等。取 $m(t)=t$ 得坐标算子，其无界性由支于 $[n,n+1]$ 的单位向量得到。
\end{proof}

\begin{example}[同一微分表达式，不同定义域]\label{exm:y7-boundary}
在 $L^2(0,1)$ 上，令 $Pf=-if'$，$D(P)=H_0^1(0,1)$。这里 $H^1$ 元素取绝对连续代表且弱导数属于 $L^2$，$H_0^1$ 等于两端迹为零的子空间。则
\[
 P^*g=-ig',\qquad D(P^*)=H^1(0,1).
\]
所以 $P$ 闭且对称，却不自伴。对 $\theta\in\mathbb R$，同一表达式在
\[
 D(P_\theta)=\{f\in H^1(0,1):f(1)=e^{i\theta}f(0)\}
\]
上定义自伴算子。
\end{example}
\begin{proof}
一维 Sobolev 事实可由积分直接看出：若 $f'=h\in L^2$（分布意义），则 $f(t)-\int_0^th(s)\,\dd s$ 的分布导数为零，故为常数。迹估计 $|f(0)|+|f(1)|\le C\|f\|_{H^1}$ 来自绝对连续公式和 Cauchy--Schwarz。零迹函数的零延拓仍在 $H^1(\mathbb R)$，经向内部缩放和光滑化可由 $C_c^\infty(0,1)$ 在 $H^1$ 中逼近，说明零迹刻画。

分部积分给出边界型
\[
 (Pf,g)-(f,-ig')=-i[f\overline g]_0^1.
\]
若 $g\in D(P^*)$，先用紧支集探针推出 $g$ 的弱导数属于 $L^2$；反向 $g\in H^1$ 时，$f$ 的零迹消去边界项。导数及迹的闭性给出 $P$ 闭。对 $P_\theta$，任意允许的端点值均能由仿射函数实现，边界型对全部 $f\in D(P_\theta)$ 消失当且仅当 $e^{i\theta}\overline{g(1)}=\overline{g(0)}$，即 $g(1)=e^{i\theta}g(0)$。于是伴随定义域恰好相等。
\end{proof}

\begin{proposition}[逆算子与满射对称算子]
若 $A$ 自伴且单射，则 $A^{-1}:R(A)\to H$ 稠密定义且自伴，逆算子未必有界。若稠密定义的对称算子 $B$ 满射，则 $B$ 自伴。
\end{proposition}
\begin{proof}
$R(A)^\perp=\ker A^*=0$ 保证逆算子稠密定义。$z\in D((A^{-1})^*)$ 且 $(A^{-1})^*z=w$，等价于 $(x,z)=(Ax,w)$ 对全部 $x\in D(A)$ 成立，亦即 $w\in D(A^*)=D(A)$ 且 $Aw=z$。故 $z\in R(A)$，$w=A^{-1}z$；反向同理，所以逆算子自伴。
对满射对称算子 $B$，有 $\ker B^*=R(B)^\perp=0$。若 $z\in D(B^*)$，选 $x\in D(B)$ 使 $Bx=B^*z$，则 $B^*(z-x)=0$，故 $z=x\in D(B)$，得到自伴。
\end{proof}

\begin{theorem}[从闭算子构造非负自伴算子]\label{thm:y7-tstar-t}
若 $T:D(T)\subset H\to K$ 闭且稠密定义，则 $T^*T$ 与 $TT^*$ 自伴且非负，且
\[
 (I+T^*T)^{-1}\in\mathcal L(H),\qquad\|(I+T^*T)^{-1}\|\le1.
\]
这里 $D(T^*T)=\{x\in D(T):Tx\in D(T^*)\}$。
\end{theorem}
\begin{proof}
闭图像的正交分解给出，对任意 $h\in H$，唯一存在 $x\in D(T),y\in D(T^*)$ 使
\[
 (h,0)=(x,Tx)+(-T^*y,y).
\]
第二分量给出 $y=-Tx$，第一分量即 $h=x+T^*Tx$。故 $I+T^*T$ 双射，且 $\|x\|\le\|h\|$，因为 $(h,x)=\|x\|^2+\|Tx\|^2$。
令 $B=(I+T^*T)^{-1}$。对 $h=(I+T^*T)x,k=(I+T^*T)z$，有 $(Bh,k)=(x,z)+(Tx,Tz)=(h,Bk)$，所以 $B$ 有界自伴且单射；$R(B)^\perp=\ker B^*=0$，说明 $D(T^*T)$ 稠密。记 $A=T^*T$。它对称且非负。若 $u\in D(A^*)$，取 $v\in D(A)$ 使 $(I+A)v=(I+A^*)u$，则 $u-v\in\ker(I+A^*)=R(I+A)^\perp=0$。故 $A=A^*$。$T^*$ 也闭且稠密定义，且 $T^{**}=T$，对 $T^*$ 重复论证即得 $TT^*$。
\end{proof}

\begin{proposition}[有界自伴算子的范数]
若 $A\in\mathcal L(H)$ 自伴，则 $\|A\|=\sup_{\|x\|=1}|(Ax,x)|$。
\end{proposition}
\begin{proof}
令右侧为 $c\le\|A\|$，$q(x)=(Ax,x)$ 为实数。由
$4t\operatorname{Re}(Ax,y)=q(x+ty)-q(x-ty)$，$t>0$，得
$|\operatorname{Re}(Ax,y)|\le c(\|x\|^2+t^2\|y\|^2)/(2t)$。取 $t=\|x\|/\|y\|$，并在复情形旋转 $y$ 的相位，使实部等于模，得到 $|(Ax,y)|\le c\|x\|\|y\|$。再对单位 $y$ 取上确界即得 $\|A\|\le c$；零向量情形显然。
\end{proof}

\begin{exercise}[亏空间检验]\label{ex:y7-deficiency}
对例 \ref{exm:y7-boundary} 的 $P$，求 $\ker(P^*-iI)$ 和 $\ker(P^*+iI)$。说明为何不能从形式表达式实对称推出自伴。
\end{exercise}
\begin{solution}
$-ig'=ig$ 给出 $g=ce^{-t}$，$-ig'=-ig$ 给出 $g=ce^t$；两者均在 $H^1(0,1)$，故两个亏空间各一维。$P$ 要求两个零端点而 $P^*$ 不要求，二者虽有相同微分表达式却不是同一算子。
\end{solution}

\begin{exercise}[实轴动量算子]\label{ex:y7-momentum}
借助第六章 Plancherel 定理，证明 $-i\,\dd/\dd t$ 在 $L^2(\mathbb R)$ 上取定义域 $H^1(\mathbb R)$ 时自伴。
\end{exercise}
\begin{solution}
令 $\mathcal U=(2\pi)^{-1/2}\FT$，则 $\mathcal U$ 酉，分布恒等式给出 $\mathcal U(-if')=\xi\mathcal Uf$。$f\in H^1$ 恰等价于 $\xi\mathcal Uf\in L^2$。故该算子为 $\mathcal U^{-1}M_\xi\mathcal U$，定义域也完整对应。实乘法算子自伴，酉变换保持伴随关系和定义域，结论成立。
\end{solution}

\Needspace{12\baselineskip}
\section{酉算子、Cayley 变换}\label{sec:yosida-07-04}
\noindent\textit{原书起始页：中译本 210；英文版 202。}

以下 $H$ 为复 Hilbert 空间。全空间上线性等距算子满足 $U^*U=I$，由极化恒等式可见保范数等价于保内积。若还满射，称为酉算子，此时 $U^{-1}=U^*$ 且 $UU^*=I$。$\ell^2$ 上右移保范数但不满射。平移 $f(t)\mapsto f(t+a)$ 及归一化 Fourier 变换均为酉算子。有界算子称为正规，若 $A^*A=AA^*$；无界情形同一等式还必须包含两个乘积的定义域相等，不能仅比较形式表达式。

\begin{theorem}[Cayley 变换及其逆]\label{thm:y7-cayley}
设 $A$ 闭、对称、稠密定义。则 $R(A\pm iI)$ 闭，且
\[
 U_A=(A-iI)(A+iI)^{-1}:R(A+iI)\longrightarrow R(A-iI)
\]
是两个闭子空间之间的满射等距映射。$I-U_A$ 单射，$R(I-U_A)=D(A)$，且
\[
 A=i(I+U_A)(I-U_A)^{-1}.
\]
反向，若 $U:M\to H$ 是定义于闭子空间 $M$ 的线性等距映射，$R(I-U)$ 稠密，则 $I-U$ 单射，以上逆公式在 $D(A)=R(I-U)$ 上定义唯一的闭对称算子，其 Cayley 变换为 $U$。
\end{theorem}
\begin{proof}
对称性使 $(Ax,x)$ 为实，故
\begin{equation}\label{eq:y7-shift-estimate}
 \|(A\pm iI)x\|^2=\|Ax\|^2+\|x\|^2.
\end{equation}
这保证单射及逆映射的界；若 $(A+iI)x_n$ 收敛，则 $x_n,Ax_n$ 均为 Cauchy，由闭性得极限仍在图像，故值域闭，负号同理。式 \eqref{eq:y7-shift-estimate} 也给出 $U_A$ 等距。写 $y=(A+iI)x$，则 $(I-U_A)y=2ix$，$(I+U_A)y=2Ax$，立即得到单射、值域与逆公式。

反向若 $Uy=y$，则对任意 $z\in M$，$(y,(I-U)z)=(y,z)-(Uy,Uz)=0$，由稠密性得 $y=0$。故可定义
$A((I-U)y)=i(I+U)y$。利用 $(Uy,Uz)=(y,z)$ 展开，得
$(i(I+U)y,(I-U)z)=((I-U)y,i(I+U)z)$，所以对称。
若 $x_n=(I-U)y_n\to x$，$Ax_n=i(I+U)y_n\to v$，则
$y_n=(x_n-iAx_n)/2$ 收敛。$M$ 闭及 $U$ 连续给出 $x=(I-U)y$、$v=i(I+U)y$，故 $A$ 闭。最后 $(A+iI)(I-U)y=2iy$、$(A-iI)(I-U)y=2iUy$，说明变换正好恢复 $U$，也证明唯一性。
\end{proof}

\begin{theorem}[自伴判据与亏空间分解]\label{thm:y7-deficiency}
对闭稠密定义的对称算子 $A$，下列等价：$A$ 自伴；$R(A+iI)=R(A-iI)=H$；$U_A$ 是 $H$ 上的酉算子。置
\[
 N_+=\ker(A^*-iI)=R(A+iI)^\perp,\qquad
 N_-=\ker(A^*+iI)=R(A-iI)^\perp.
\]
则有代数直和（不声称为 $H$ 内积正交直和）
\[
 D(A^*)=D(A)\mathbin{\dot+}N_+\mathbin{\dot+}N_-,
 \quad A^*(x_0+x_++x_-)=Ax_0+ix_+-ix_-.
\]
\end{theorem}
\begin{proof}
一般恒等式 $R(B)^\perp=\ker B^*$ 直接来自伴随定义，注意 $(A+iI)^*=A^*-iI$。若 $A=A^*$，则由式 \eqref{eq:y7-shift-estimate} 两个核均为零；两个值域既闭又稠密，故等于 $H$。反之若两者满射，任取 $y\in D(A^*)$，取 $x\in D(A)$ 使 $(A+iI)x=(A^*+iI)y$，于是 $y-x\in\ker(A^*+iI)=0$，得 $D(A^*)=D(A)$。酉判据随后由前一定理成立。

对分解，任取 $y\in D(A^*)$，在 $H=R(A+iI)\oplus N_+$ 中分解
$(A^*+iI)y=(A+iI)x_0+w_+$。令 $x_+=w_+/(2i)$，则 $x_-=y-x_0-x_+$ 属于 $N_-$，给出存在性。若 $x_0+x_++x_-=0$，施以 $A^*+iI$ 得 $(A+iI)x_0+2ix_+=0$；两项在正交子空间中，故均为零。由式 \eqref{eq:y7-shift-estimate} 得 $x_0=0$，继而 $x_+=x_-=0$。作用公式来自各子空间的定义。
\end{proof}

\begin{exercise}[乘法算子的 Cayley 变换]\label{ex:y7-cayley-mult}
求 $M_t$ 在 $L^2(\mathbb R)$ 上的 Cayley 变换，并检验逆公式的定义域。
\end{exercise}
\begin{solution}
$U$ 为乘以 $u(t)=(t-i)/(t+i)$ 的算子，$|u(t)|=1$，故为酉算子。$I-U$ 乘以 $2i/(t+i)$。其值域正是 $\{f:tf\in L^2\}$：写 $f=2ig/(t+i)$ 时 $tf\in L^2$；反之 $g=(t+i)f/(2i)\in L^2$。逆公式乘数为 $i(1+u)/(1-u)=t$，定义域也相等。
\end{solution}

\begin{exercise}[为什么两个符号都要检查]\label{ex:y7-halfline}
在 $L^2(0,\infty)$ 上取 $Af=-if'$，$D(A)=H_0^1(0,\infty)$。已知同样的弱导数论证给出 $D(A^*)=H^1(0,\infty)$。求两个亏空间，判断 $A$ 是否自伴。
\end{exercise}
\begin{solution}
$\ker(A^*-iI)$ 由 $e^{-t}$ 张成，$\ker(A^*+iI)=\{0\}$，因为 $e^t\notin L^2(0,\infty)$。$A$ 闭且对称，故 $R(A-iI)=H$，而 $R(A+iI)$ 的正交补一维。一个符号满射并不足以推出自伴；这里 $A\ne A^*$。
\end{solution}

\Needspace{12\baselineskip}
\section{闭值域定理}\label{sec:yosida-07-05}
\noindent\textit{原书起始页：中译本 214；英文版 205。}

此节 $X,Y$ 为 Banach 空间，$T:D(T)\subset X\to Y$ 闭且稠密定义。记 $N(T)=\ker T$，$R(T)=T(D(T))$。对 $M\subset X$ 写 $M^\perp=\{f\in X':f|_M=0\}$；对 $F\subset X'$ 写 ${}^\perp F=\{x\in X:f(x)=0\ \forall f\in F\}$。在 Hilbert 空间中通过 Riesz 识别后可用通常的正交补。

\begin{lemma}[有界算子的闭值域]\label{lem:y7-bounded-range}
若 $S:E\to F$ 为 Banach 空间间有界算子，则 $R(S)$ 范数闭当且仅当 $R(S')$ 范数闭。
\end{lemma}
\begin{proof}
若 $R(S)$ 闭，则 $E/N(S)\to R(S)$ 为 Banach 空间间有界双射，开映射定理给出 $\operatorname{dist}(x,N(S))\le C\|Sx\|$。对 $f\in N(S)^\perp$，令 $g_0(Sx)=f(x)$，则良定义且 $|g_0(Sx)|\le C\|f\|\|Sx\|$。Hahn--Banach 延拓给出 $g\in F'$ 且 $S'g=f$，因此 $R(S')=N(S)^\perp$ 闭。

反向设 $R(S')$ 闭。置 $Z=\overline{R(S)}$，将 $S$ 看成 $S_0:E\to Z$。限制映射 $F'\to Z'$ 满射，故 $R(S_0')=R(S')$；$S_0'$ 单射且有闭值域，逆映射定理给出 $\|g\|\le C\|S_0'g\|$（$g\in Z'$）。令 $B_E$ 为闭单位球，$K=\overline{S_0(B_E)}$。分离定理说明 $C^{-1}B_Z\subset K$：否则存在 $g$ 使某个 $\|z\|\le C^{-1}$ 满足
$\operatorname{Re}g(z)>\sup_{w\in K}\operatorname{Re}g(w)=\|S_0'g\|$，与 $|g(z)|\le C^{-1}\|g\|\le\|S_0'g\|$ 矛盾。
因此对任意余项 $r\in Z$，可选 $x$ 使 $\|x\|\le C\|r\|$ 且 $\|r-S_0x\|\le\|r\|/2$。从 $r=z$ 迭代，得到 $\sum\|x_n\|\le2C\|z\|$，余项趋零。$E$ 完备及 $S_0$ 连续保证 $S_0(\sum x_n)=z$，故 $R(S)=Z$ 闭。
\end{proof}

\begin{theorem}[Banach 闭值域定理]\label{thm:y7-closed-range}
以下四个条件等价：
\[
 \begin{aligned}
 &R(T)\text{ 在 }Y\text{ 中范数闭};\qquad
 R(T')\text{ 在 }X'\text{ 中范数闭};\\
 &R(T)={}^\perp N(T');\qquad R(T')=N(T)^\perp.
 \end{aligned}
\]
不附闭值域条件时，总有 $\overline{R(T)}={}^\perp N(T')$，但不能一般去掉闭包。
\end{theorem}
\begin{proof}
第一步将无界算子约化为有界算子，同时明确图空间的对偶。$G=G(T)$ 在 $\|(x,Tx)\|=\|x\|+\|Tx\|$ 下为 Banach 空间，令 $S:G\to Y$，$S(x,Tx)=Tx$。显然 $R(S)=R(T)$。

令 $q:X'\times Y'\to G'$ 为限制映射，$q(f,g)(x,Tx)=f(x)+g(Tx)$。Hahn--Banach 保证 $q$ 满射，它是 Banach 空间间有界满射，因而是商映射。由定义
\[
 \ker q=\{(-T'g,g):g\in D(T')\},\qquad S'g=q(0,g).
\]
于是
\[
 q^{-1}(R(S'))=\ker q+(\{0\}\times Y')=R(T')\times Y'.
\]
商拓扑说明 $R(S')$ 闭当且仅当 $R(T')$ 闭。用引理 \ref{lem:y7-bounded-range}，即得前两个条件等价。这一步把 $G'$ 处理为商空间，不能把它直接当成 $X'\times Y'$ 的子空间。

第二步，$N(T')$ 恰为 $R(T)$ 的湮灭子：若 $g$ 在 $R(T)$ 上为零，则 $g\in D(T')$ 且 $T'g=0$。Hahn--Banach 分离给出 $\overline{R(T)}={}^\perp N(T')$，所以第一、第三条件等价。

第三步，假设 $R(T)$ 闭，证明 $R(T')=N(T)^\perp$。包含 $R(T')\subset N(T)^\perp$ 直接来自配对。对反向，$N(T)$ 在 $X$ 中闭，由 $G/N(S)\to R(T)$ 的开映射定理有
\[
 \operatorname{dist}(x,N(T))+\|Tx\|\le C\|Tx\|\quad(x\in D(T)).
\]
故对 $f\in N(T)^\perp$，$g_0(Tx)=f(x)$ 良定义，且 $|g_0(Tx)|\le C\|f\|\|Tx\|$。将 $g_0$ 延拓到 $Y$ 得 $g\in Y'$，定义保证 $g\in D(T')$ 且 $T'g=f$。最后若第四条件成立，$N(T)^\perp$ 本来就是范数闭子空间，故第二条件成立，完成等价性证明。
\end{proof}

\begin{corollary}[Hilbert 空间中的可解性与估计]
若 $T:D(T)\subset H\to K$ 闭且稠密定义，则值域闭当且仅当存在 $C$ 使
\[
 \|x\|\le C\|Tx\|\quad(x\in D(T)\cap N(T)^\perp).
\]
此时 $R(T)=N(T^*)^\perp$，$R(T^*)=N(T)^\perp$。故 $Tx=f$ 可解当且仅当 $f\perp N(T^*)$，且存在唯一的 $x\perp N(T)$ 的解。
\end{corollary}
\begin{proof}
闭值域时，上述图空间开映射估计及 $\operatorname{dist}(x,N(T))=\|x\|$ 给出不等式。反向若 $Tx_n\to y$，减去 $x_n$ 的核投影不改变值域，且所得向量仍在定义域；估计使其为 Cauchy，闭性推出 $y\in R(T)$。余下等式由 Riesz 识别转译闭值域定理。
\end{proof}

当 $H=K$ 且 $\operatorname{Re}(Tx,x)\ge c\|x\|^2$（$c>0$）时，Cauchy--Schwarz 给出 $\|Tx\|\ge c\|x\|$。故 $T$ 核为零且值域闭，闭值域定理推出 $T^*$ 满射。该不等式本身不保证 $T$ 满射。

\begin{example}[稠密值域却不闭]
在 $\ell^2$ 上令 $(Tx)_n=x_n/n$，则 $T=T^*$ 为有界单射，有限支集序列均在其值域中，故值域稠密。$y=(1/n)_{n\ge1}\in\ell^2$ 不在值域，因为其唯一形式原像为 $(1,1,\ldots)\notin\ell^2$。截断 $y$ 却都在值域中。因此 $N(T^*)^\perp=\ell^2$ 严格大于 $R(T)$；也有 $\|Te_n\|=1/n$，直接否定闭值域估计。
\end{example}

\begin{exercise}[满射与下界]\label{ex:y7-surjectivity}
设 $T:D(T)\subset H\to K$ 闭且稠密定义。证明 $T$ 满射当且仅当存在 $c>0$ 使 $\|T^*y\|\ge c\|y\|$（$y\in D(T^*)$）。
\end{exercise}
\begin{solution}
满射时 $R(T)$ 闭且 $N(T^*)=R(T)^\perp=0$；闭值域定理给出 $R(T^*)$ 闭，再用前一推论的估计。反向下界使 $T^*$ 单射、值域闭（用其闭性），故 $R(T)$ 闭，又因 $N(T^*)=0$ 得 $\overline{R(T)}=K$，所以 $R(T)=K$。
\end{solution}

\begin{exercise}[周期微分方程的相容条件]\label{ex:y7-periodic}
对例 \ref{exm:y7-boundary} 的 $P_0$，证明 $R(P_0)=\{g\in L^2(0,1):\int_0^1g=0\}$，并求 $P_0f=g$ 的全部解。
\end{exercise}
\begin{solution}
积分 $-if'=g$ 得 $f(t)=c+i\int_0^tg(s)\,\dd s$。周期端点条件等价于 $\int_0^1g=0$。反向该条件下公式定义 $H^1$ 周期解；$c$ 任意，选择 $c=-i\int_0^1\int_0^tg(s)\,\dd s\,\dd t$ 得唯一零均值解。右侧是连续泛函的核，因而值域闭，也等于常数函数空间 $N(P_0^*)$ 的正交补。
\end{solution}


\Needspace{10\baselineskip}
\section*{本章参考文献与阅读衔接}
\phantomsection\addcontentsline{toc}{section}{本章参考文献与阅读衔接}
直接参考为两版吉田第七章 \cite{yosida-zh,yosida-en}。原书章末列出的 Stone、Akhiezer--Glazman、Dunford--Schwartz、Banach 以及节内 von Neumann 原始文献本轮未直接查阅。上述论证按本稿记号重新组织；Banach 闭值域定理已给出有界情形的分离与逐次逼近证明，再借助图空间商映射推广到闭算子。

第六章的形式转置决定微分表达式，本章边界条件决定 Hilbert 伴随。下一章讨论预解式和谱；式 \eqref{eq:y7-shift-estimate} 与 $T^*T$ 构造是其基础。自伴延拓分类和无界谱理论留待后续，十道习题的答案集中在书后。

% !TeX encoding = UTF-8
% 吉田第六版：中译本第 8 章；此文件当前仅为目录与正文占位。
\SetChapterInfo{subtitle={吉田耕作第六版主线 · 正文待撰写}}
\chapter{预解式和谱}\label{ch:yosida-08}
\YosidaChapterPlaceholder{8}{Resolvent and Spectrum}{219--242}{209--230}

\section{预解式和谱}\label{sec:yosida-08-01}
\YosidaSectionPlaceholder{219}{209}
% TODO: 定义与记号；定理与完整证明；例子；习题；中英文核对。

\section{预解方程和谱半径}\label{sec:yosida-08-02}
\YosidaSectionPlaceholder{221}{211}
% TODO: 定义与记号；定理与完整证明；例子；习题；中英文核对。

\section{平均遍历定理}\label{sec:yosida-08-03}
\YosidaSectionPlaceholder{223}{213}
% TODO: 定义与记号；定理与完整证明；例子；习题；中英文核对。

\section{关于伪预解式的 Hille 型遍历定理}\label{sec:yosida-08-04}
\YosidaSectionPlaceholder{226}{215}
% TODO: 定义与记号；定理与完整证明；例子；习题；中英文核对。

\section{殆周期函数的平均值}\label{sec:yosida-08-05}
\YosidaSectionPlaceholder{229}{218}
% TODO: 定义与记号；定理与完整证明；例子；习题；中英文核对。

\section{对偶算子的预解式}\label{sec:yosida-08-06}
\YosidaSectionPlaceholder{235}{224}
% TODO: 定义与记号；定理与完整证明；例子；习题；中英文核对。

\section{Dunford 积分}\label{sec:yosida-08-07}
\YosidaSectionPlaceholder{236}{225}
% TODO: 定义与记号；定理与完整证明；例子；习题；中英文核对。

\section{预解式的孤立奇点}\label{sec:yosida-08-08}
\YosidaSectionPlaceholder{238}{228}
% TODO: 定义与记号；定理与完整证明；例子；习题；中英文核对。

\section*{本章评注和参考文献整理（占位）}
\phantomsection
\addcontentsline{toc}{section}{本章评注和参考文献整理（占位）}
待核对原书本章评注和参考文献（中译本印刷第 242 页），并整理相应论述。

% !TeX encoding = UTF-8
% 吉田第六版：中译本第 9 章；此文件当前仅为目录与正文占位。
\SetChapterInfo{subtitle={吉田耕作第六版主线 · 正文待撰写}}
\chapter{半群的分析理论}\label{ch:yosida-09}
\YosidaChapterPlaceholder{9}{Analytical Theory of Semi-groups}{243--286}{231--273}

\section{\texorpdfstring{$(C_0)$}{(C0)} 类半群}\label{sec:yosida-09-01}
\YosidaSectionPlaceholder{243}{232}
% TODO: 定义与记号；定理与完整证明；例子；习题；中英文核对。

\section{局部凸空间中 \texorpdfstring{$(C_0)$}{(C0)} 类等度连续半群、几例半群}\label{sec:yosida-09-02}
\YosidaSectionPlaceholder{246}{234}
% TODO: 定义与记号；定理与完整证明；例子；习题；中英文核对。

\section{\texorpdfstring{$(C_0)$}{(C0)} 类等度连续半群的无穷小生成元}\label{sec:yosida-09-03}
\YosidaSectionPlaceholder{248}{237}
% TODO: 定义与记号；定理与完整证明；例子；习题；中英文核对。

\section{无穷小生成元 \texorpdfstring{$A$}{A} 的预解式}\label{sec:yosida-09-04}
\YosidaSectionPlaceholder{252}{240}
% TODO: 定义与记号；定理与完整证明；例子；习题；中英文核对。

\section{几例无穷小生成元}\label{sec:yosida-09-05}
\YosidaSectionPlaceholder{254}{242}
% TODO: 定义与记号；定理与完整证明；例子；习题；中英文核对。

\section{具等度连续幂的连续线性算子的指数函数}\label{sec:yosida-09-06}
\YosidaSectionPlaceholder{256}{244}
% TODO: 定义与记号；定理与完整证明；例子；习题；中英文核对。

\section{\texorpdfstring{$(C_0)$}{(C0)} 类等度连续半群用其无穷小生成元的表示和刻画}\label{sec:yosida-09-07}
\YosidaSectionPlaceholder{258}{246}
% TODO: 定义与记号；定理与完整证明；例子；习题；中英文核对。

\section{压缩半群和耗散算子}\label{sec:yosida-09-08}
\YosidaSectionPlaceholder{261}{250}
% TODO: 定义与记号；定理与完整证明；例子；习题；中英文核对。

\section{\texorpdfstring{$(C_0)$}{(C0)} 类等度连续群、Stone 定理}\label{sec:yosida-09-09}
\YosidaSectionPlaceholder{263}{251}
% TODO: 定义与记号；定理与完整证明；例子；习题；中英文核对。

\section{全纯半群}\label{sec:yosida-09-10}
\YosidaSectionPlaceholder{266}{254}
% TODO: 定义与记号；定理与完整证明；例子；习题；中英文核对。

\section{闭算子的分数幂}\label{sec:yosida-09-11}
\YosidaSectionPlaceholder{271}{259}
% TODO: 定义与记号；定理与完整证明；例子；习题；中英文核对。

\section{半群的收敛性、Trotter--Kato 定理}\label{sec:yosida-09-12}
\YosidaSectionPlaceholder{281}{269}
% TODO: 定义与记号；定理与完整证明；例子；习题；中英文核对。

\section{对偶半群、Phillips 定理}\label{sec:yosida-09-13}
\YosidaSectionPlaceholder{284}{272}
% TODO: 定义与记号；定理与完整证明；例子；习题；中英文核对。

% !TeX encoding = UTF-8
\SetChapterInfo{subtitle={从紧性判别到 Fredholm 择一},keywords={全有界; 核算子; Rellich; Schauder; Fredholm; Dirichlet; 核空间},prereq={Ascoli 定理; Hahn--Banach; Sobolev 空间; 闭值域; Riesz 投影},goals={证明紧算子的谱结构; 识别相容条件; 用紧嵌入处理椭圆边值问题}}
\chapter{紧算子}\label{ch:yosida-10}
\ChapterSource{按 \cite{yosida-zh,yosida-en} 第十章六节及附录组织，2026 年 10 月 27 日补写为学习讲义。直接查看英文印刷 274--293 页全部扫描图，以及中译本印刷 287、290、294--296、299、302、306 页。两版参考范围分别为 287--306、274--293；英文第 274 页上部仍属第九章。查阅范围与未展开结果详见本章 COVERAGE。}

有界集在无限维空间中未必有收敛子列，紧算子却把每个有界序列送到具有收敛子列的序列。本章利用这一性质证明：紧算子在非零谱点附近保留有限维线性代数的结构，方程的障碍由有限多个对偶条件描述。应用中，椭圆方程的估计先给出有界解算子，Sobolev 紧嵌入再把它变为紧算子。

除另有说明，$X,Y$ 为实或复 Banach 空间，讨论谱时取复数域；$B_X$ 为闭单位球，$T'$ 为线性对偶算子。Hilbert 内积对第一变量线性，Hilbert 伴随写作 $T^*$。本章的核算子指 nuclear operator，不能与泛指由积分核定义的算子混同。

\Needspace{10\baselineskip}
\section{B-空间中的紧集}\label{sec:yosida-10-01}
\noindent\textit{原书起始页：中译本 287；英文版 274。}

\begin{definition}[全有界与相对紧]
集合 $M\subset X$ 称为全有界，若对每个 $\varepsilon>0$ 可用有限多个半径 $\varepsilon$ 的球覆盖；称为相对紧，若其闭包紧。在 Banach 空间中，两者等价。闭且有界本身不够，例如 $\ell^2$ 的单位向量满足 $\|e_n-e_m\|=\sqrt2$。
\end{definition}
\begin{proof}
紧性蕴含有限网。反之，逐次用半径 $2^{-j}$ 的有限覆盖作对角选取，得到 Cauchy 子列；完备性保证其收敛。因此闭包序列紧，因而紧。
\end{proof}

\begin{theorem}[Fréchet--Kolmogorov 判别]\label{thm:y10-kolmogorov}
设 $1\le p<\infty$。$M\subset L^p(\R^d)$ 相对紧，当且仅当以下三项成立：
\begin{align*}
 &\sup_{f\in M}\|f\|_p<\infty,\\
 &\lim_{h\to0}\sup_{f\in M}\|f(\cdot+h)-f\|_p=0,\\
 &\lim_{R\to\infty}\sup_{f\in M}\|1_{\{|x|>R\}}f\|_p=0.
\end{align*}
第二项排除越来越快的振荡，第三项排除质量逃向无穷远。
\end{theorem}
\begin{proof}
必要性：紧闭包有界。对单个 $f\in L^p$，平移连续性由 $C_c^\infty$ 的稠密性和移位等距性得到，尾部范数则由积分的绝对连续性趋于零。在 $M$ 中取有限 $\varepsilon$ 网，将这两个性质从有限个网点传给全体元素；平移差的误差至多 $2\varepsilon$，尾部误差至多 $\varepsilon$，故为一致极限。

充分性：取非负 $\rho\in C_c^\infty(B(0,1))$，$\int\rho=1$，令 $\rho_\delta(x)=\delta^{-d}\rho(x/\delta)$。Minkowski 不等式给出
\[
 \sup_{f\in M}\|\rho_\delta*f-f\|_p
 \le\sup_{|h|\le\delta,\,f\in M}\|f(\cdot-h)-f\|_p\longrightarrow0.
\]
固定 $\delta$，由 Hölder 不等式，函数 $g_f=\rho_\delta*f$ 及其一阶导数在全空间一致有界，常数可依赖 $\delta$。这对 $p=1$ 同样成立，因为卷积核及其导数有界。故 $\{g_f\}$ 在每个闭球上一致有界且等度连续。Ascoli 定理使其限制在该球上的集合在一致范数、从而在 $L^p$ 范数下全有界。

又当 $R>\delta$ 时，由卷积核的支集和 Minkowski 不等式，
\[
 \|1_{\{|x|>R\}}g_f\|_p
 \le\|1_{\{|x|>R-\delta\}}f\|_p.
\]
先选 $\delta$ 使平滑误差小，再选 $R$ 使尾部一致小，最后选球上的有限网并将网点在球外延零，得到 $M$ 的任意精度有限网。$L^p$ 完备，故 $M$ 相对紧。
\end{proof}

\begin{example}[紧性的两种损失]
固定非零 $f\in C_c^\infty(\R)$，$f_n(x)=f(x-n)$ 的平移连续性对 $n$ 一致，但尾部条件失败。另一方面，在 $L^2(\R)$ 中取 $g_n(x)=1_{[0,2\pi]}(x)e^{inx}$，它们有共同紧支集且范数相同，却彼此正交，故不相对紧；此时失败的是一致平移连续性。
\end{example}

\begin{exercise}[紧闭包的凸化]\label{ex:y10-convex}
证明 Banach 空间中全有界集的凸包仍全有界，其闭凸包因而紧。
\end{exercise}
\begin{solution}
取有限 $\varepsilon$ 网 $F$。每个有限凸组合与 $\operatorname{conv}F$ 中某点的距离至多 $\varepsilon$。$\operatorname{conv}F$ 位于有限维子空间且有界，故有有限 $\varepsilon$ 网；合成即得 $2\varepsilon$ 网。闭包的紧性使用环境空间完备性。
\end{solution}

\begin{exercise}[去掉尾部条件]\label{ex:y10-tail}
对上述 $f_n$，证明一致平移连续性，并用支集分离说明没有收敛子列。
\end{exercise}
\begin{solution}
$\|f_n(\cdot+h)-f_n\|_p=\|f(\cdot+h)-f\|_p$。若两个平移量相差大于支集直径，则支集不交，$\|f_n-f_m\|_p^p=2\|f\|_p^p$。任意无限子列都含任意远的两项，故不可能是 Cauchy 列。
\end{solution}

\Needspace{10\baselineskip}
\section{紧算子与核算子}\label{sec:yosida-10-02}
\noindent\textit{原书起始页：中译本 290；英文版 277。}

\begin{definition}[紧算子]
$T\in\mathcal L(X,Y)$ 称为紧算子，若 $T(B_X)$ 相对紧，记 $T\in\mathcal K(X,Y)$。等价地，任何有界列 $x_n$ 的像 $Tx_n$ 都有范数收敛子列。
\end{definition}

\begin{proposition}[闭理想性质]\label{prop:y10-ideal}
紧算子的线性组合仍紧；与有界算子在任意一侧复合仍紧。若紧算子 $T_n$ 在算子范数下收敛于 $T$，则 $T$ 紧。特别地，有限秩算子及其范数极限紧。
\end{proposition}
\begin{proof}
线性组合的像包含在两个相对紧集的和中；右复合将单位球送入一个有界集，左复合则保持紧闭包的连续像紧。给定 $\varepsilon>0$，选 $n$ 使 $\|T-T_n\|<\varepsilon$，再取 $T_n(B_X)$ 的有限 $\varepsilon$ 网，便得到 $T(B_X)$ 的有限 $2\varepsilon$ 网。有限维有界集全有界，故有限秩情形成立。
\end{proof}

本命题只推出有限秩范数极限必紧；一般 Banach 空间上，反方向涉及逼近性质，不在此作无条件断言。

\begin{example}[连续核与平方可积核]
若 $k\in C([a,b]^2)$，则 $Tf(x)=\int_a^b k(x,y)f(y)\,\dd y$ 在 $C[a,b]$ 上紧。事实上 $\|Tf\|_\infty\le(b-a)\|k\|_\infty\|f\|_\infty$，且对 $\|f\|_\infty\le1$，
\[
 |Tf(x)-Tf(z)|\le(b-a)\sup_y|k(x,y)-k(z,y)|.
\]
核的一致连续性和 Ascoli 定理给出紧性。

若 $(S,\mu)$ 为 $\sigma$-有限测度空间，$k\in L^2(S\times S)$，则对应积分算子在 $L^2(S)$ 上紧，且 $\|T\|\le\|k\|_2$。Cauchy--Schwarz 后积分给出此界。矩形乘积函数的有限线性组合在 $L^2(S\times S)$ 中稠密：可先用 $\sigma$-有限分解截到有限测度，再用生成乘积 $\sigma$-代数的矩形逼近简单函数。它们的积分算子为有限秩，且核的 $L^2$ 逼近给出算子范数逼近。这类算子称 Hilbert--Schmidt 积分算子。
\end{example}

\begin{definition}[核算子与核范数]\label{def:y10-nuclear}
若存在 $f_n\in X'$、$y_n\in Y$ 满足
\[
 Tx=\sum_{n=1}^\infty f_n(x)y_n,\qquad
 \sum_{n=1}^\infty\|f_n\|\,\|y_n\|<\infty,
\]
则 $T$ 为核算子。其核范数 $\nu(T)$ 为所有此类表示的上式总和之下确界。这等价于原书使用的 $\ell^1$ 系数乘以两个有界序列的表示。
\end{definition}
\begin{proposition}[核算子必紧]
上式在单位球上一致绝对收敛，部分和 $T_N$ 为有限秩，且
\[
 \|T-T_N\|\le\sum_{n>N}\|f_n\|\,\|y_n\|.
\]
因此核算子紧，$\|T\|\le\nu(T)$，且 $\nu(ATB)\le\|A\|\nu(T)\|B\|$。
\end{proposition}
\begin{proof}
对 $\|x\|\le1$ 逐项估计即得尾和界，紧性由命题 \ref{prop:y10-ideal}。对复合把表示改成 $(f_n\circ B)(x)Ay_n$，再取下确界。
\end{proof}

\begin{proposition}[Sobolev 阶差足够大时的核嵌入]\label{prop:y10-nuclear-sobolev}
设 $G\subset\R^d$ 有界开，$k>j\ge0$ 为整数且 $k-j>d$。则自然嵌入 $H_0^k(G)\to H_0^j(G)$ 为核算子。
\end{proposition}
\begin{proof}
把 $\overline G$ 放在一个周期立方体 $Q$ 的内部。$H_0^s(G)$ 的零延拓是 $H^s(Q)$ 中的闭子空间，因其由测试函数的闭包定义；导数范数与 Fourier 范数等价。周期空间的正交模 $e_m$，$m\in\mathbb Z^d$，在 $H^s(Q)$ 的范数等价于 $(1+|m|^2)^{s/2}$。于是嵌入 $H^k(Q)\to H^j(Q)$ 有秩一展开，归一化后的系数为
\[
 (1+|m|^2)^{-(k-j)/2},\qquad
 \sum_{m\in\mathbb Z^d}(1+|m|^2)^{-(k-j)/2}<\infty.
\]
最后一个和由格点球壳计数 $O(r^{d-1})$ 收敛。将零延拓、此嵌入、$H^j(Q)$ 到零延拓闭子空间的正交投影以及逆零延拓复合，得到所求映射。等价范数不改变核性。
\end{proof}

原书随后联系迹类。这里可直接证明：两个 Hilbert--Schmidt 算子 $A:H\to K$、$B:K\to H$ 的乘积 $BA$ 为核算子。取 $K$ 的正交基 $e_n$，有
\[
 BAx=\sum_n(x,A^*e_n)Be_n,\quad
 \sum_n\|A^*e_n\|\,\|Be_n\|\le\|A\|_{\rm HS}\|B\|_{\rm HS}.
\]
一般迹类与奇异值理论留待谱表示以后讨论。

\begin{exercise}[对角紧算子]\label{ex:y10-diagonal}
设 $D_a:\ell^2\to\ell^2$，$D_ax=(a_nx_n)$，$a\in\ell^\infty$。证明 $D_a$ 紧当且仅当 $a_n\to0$，并证明 $a\in\ell^1$ 时它为核算子。
\end{exercise}
\begin{solution}
截断前 $N$ 项的有限秩算子与 $D_a$ 的范数差为 $\sup_{n>N}|a_n|$。若 $a_n$ 不趋于零，存在子列 $|a_{n_k}|\ge\varepsilon$，其像 $a_{n_k}e_{n_k}$ 两两距离至少 $\sqrt2\varepsilon$，故不紧。核表示是 $D_ax=\sum_n a_nx_ne_n$，绝对范数和为 $\sum_n|a_n|$。
\end{solution}

\begin{exercise}[紧算子改善弱收敛]\label{ex:y10-weak-strong}
若 $T$ 紧且 $x_n\rightharpoonup x$，证明 $Tx_n\to Tx$。
\end{exercise}
\begin{solution}
弱收敛列由一致有界原理有界。若范数收敛失败，可选子列距 $Tx$ 至少 $\varepsilon$；紧性再给出范数收敛子列。它同时弱收敛于 $Tx$，因 $f(Tx_n)=T'f(x_n)$，故范数极限只能为 $Tx$，矛盾。
\end{solution}

\Needspace{10\baselineskip}
\section{Rellich--Gårding 定理}\label{sec:yosida-10-03}
\noindent\textit{原书起始页：中译本 294；英文版 281。}

\begin{theorem}[零边界 Sobolev 紧嵌入]\label{thm:y10-rellich}
设 $G\subset\R^d$ 为有界开集，$k>j\ge0$ 为整数。则 $H_0^k(G)\hookrightarrow H_0^j(G)$ 紧。更一般地，若
\[
 T\in\mathcal L(H_0^k(G),H_0^k(G)),\qquad
 \|T\varphi\|_{H^k}\le C\|\varphi\|_{H^j}
 \quad(\varphi\in C_c^\infty(G)),
\]
则 $T$ 在 $H_0^k(G)$ 上紧。
\end{theorem}
\begin{proof}
令 $u_n$ 在 $H_0^k(G)$ 中有界。由测试函数闭包定义，零延拓 $\widetilde u_n\in H^k(\R^d)$，其导数也是相应导数的零延拓，且范数保持。它们具有共同有界支集，因而 $L^2$ 范数有界且 $L^1$ 范数一致有界。

Fourier 变换满足 $|\widehat {\widetilde u_n}(\xi)|\le C$，且
\[
 |\widehat {\widetilde u_n}(\xi)-\widehat {\widetilde u_n}(\eta)|
 \le |\xi-\eta|\int|x|\,|\widetilde u_n(x)|\,\dd x\le C_1|\xi-\eta|.
\]
在频率闭球上应用 Ascoli，再取对角子列，得到每个频率球上的一致收敛。由 Plancherel 及整数 Sobolev 范数的 Fourier 等价式，频率高于 $R$ 的差满足
\[
 \int_{|\xi|>R}(1+|\xi|^2)^j
 |\widehat {\widetilde u_n}-\widehat {\widetilde u_m}|^2\,\dd\xi
 \le C_2(1+R^2)^{j-k}.
\]
先令 $R$ 大，再令 $m,n$ 大，得到该子列在 $H^j(\R^d)$ 中 Cauchy。限制回 $G$ 后在 $H^j(G)$ 中收敛；每个 $u_n\in H_0^j(G)$，而后者闭，极限仍属于它。这证明嵌入紧。

对第二项，先由 $C_c^\infty$ 在 $H_0^k$ 中稠密，将不等式延至全空间。对任意有界列选出在 $H_0^j$ 中 Cauchy 的子列，不等式使其 $T$ 像在 $H_0^k$ 中 Cauchy，即得紧性。
\end{proof}

这里不要求 $G$ 的边界光滑，因为使用的是零边界闭包空间。若改成 $H^k(G)$，则需另行提供延拓条件，例如有界 Lipschitz 域。全空间上的同类嵌入通常不紧：固定函数的远距离平移仍保持所有 Sobolev 范数。

\begin{example}[正则性如何产生紧性]
设 $S:L^2(G)\to H_0^1(G)$ 有界，则作为 $L^2(G)\to L^2(G)$ 的算子 $JS$ 紧，其中 $J$ 为自然嵌入。紧性来自解算子与紧嵌入的复合；不要求 $S$ 在目标 $H_0^1$ 范数下紧。
\end{example}

\begin{exercise}[平移估计的另一条证明]\label{ex:y10-translation-estimate}
证明 $u\in H^1(\R^d)$ 时，$\|u(\cdot+h)-u\|_2\le |h|\|\nabla u\|_2$。用它及定理 \ref{thm:y10-kolmogorov} 重证 $H_0^1(G)\hookrightarrow L^2(G)$ 紧。
\end{exercise}
\begin{solution}
对光滑函数写差为 $\int_0^1h\cdot\nabla u(x+th)\,\dd t$，由 Minkowski 与平移等距性得界，再由稠密性推广。零延拓的有界族满足此一致估计，且共同有界支集给出尾部条件，故相对紧。
\end{solution}

\begin{exercise}[紧性与核性的阶差]\label{ex:y10-orders}
比较 $H_0^k(G)\to H_0^j(G)$ 的本节结论与命题 \ref{prop:y10-nuclear-sobolev}；为何不应将两者的条件混用？
\end{exercise}
\begin{solution}
紧性只需 $k-j>0$，Fourier 高频因子即趋于零。所给核性充分条件需 $k-j>d$，因为要求格点上的系数绝对可和。系数趋于零与系数可和是不同要求；核性证明提供更强结论，也需要更强估计。
\end{solution}

\Needspace{10\baselineskip}
\section{Schauder 定理}\label{sec:yosida-10-04}
\noindent\textit{原书起始页：中译本 295；英文版 282。}

\begin{theorem}[紧性的对偶刻画]\label{thm:y10-schauder}
有界线性算子 $T:X\to Y$ 紧，当且仅当 $T':Y'\to X'$ 紧。
\end{theorem}
\begin{proof}
若 $T$ 紧，令 $K=\overline{T(B_X)}$，它为紧度量空间。对任意 $f_n\in B_{Y'}$，将其限制于 $K$，得到函数 $F_n(y)=f_n(y)$。它们一致有界且满足
\[
 |F_n(y)-F_n(z)|\le\|y-z\|.
\]
Ascoli 给出在 $K$ 上一致收敛的子列。因此
\[
 \|T'f_n-T'f_m\|
 =\sup_{x\in B_X}|(f_n-f_m)(Tx)|\longrightarrow0.
\]
$X'$ 完备，故 $T'$ 紧。

反之，已证方向应用于 $T'$，得到 $T'':X''\to Y''$ 紧。自然等距嵌入 $J_X,J_Y$ 满足 $T''J_X=J_YT$。故 $J_YT(B_X)$ 相对紧于 $Y''$。$J_Y(Y)$ 是闭子空间，因为 $Y$ 完备；其闭包因此仍在 $J_Y(Y)$ 内。由等距性拉回，$T(B_X)$ 相对紧于 $Y$。
\end{proof}

此证明使用自然嵌入的等距性与闭像，不要求 $X$ 或 $Y$ 自反。Hilbert 空间中经 Riesz 对应可写成 $T$ 紧当且仅当 $T^*$ 紧，但 Banach 空间中的 $T'$ 始终作用在对偶空间上。

\begin{exercise}[对偶方向的估计]\label{ex:y10-adjoint-net}
若 $\{y_1,\ldots,y_N\}$ 是 $T(B_X)$ 的 $\varepsilon$ 网，证明对 $f,g\in B_{Y'}$ 有
\[
 \|T'f-T'g\|\le\max_{1\le i\le N}|(f-g)(y_i)|+2\varepsilon.
\]
\end{exercise}
\begin{solution}
对每个 $x\in B_X$ 选 $i$ 使 $\|Tx-y_i\|<\varepsilon$。因 $\|f-g\|\le2$，有 $|(f-g)(Tx)|\le|(f-g)(y_i)|+2\varepsilon$，对 $x$ 取上确界即得。
\end{solution}

\begin{exercise}[恒等算子]\label{ex:y10-identity}
用 Schauder 定理说明：$I_X$ 紧当且仅当 $I_{X'}$ 紧，并给出这两个条件的维数含义。
\end{exercise}
\begin{solution}
$I_X'=I_{X'}$，故二者等价。由 Riesz 引理，无限维赋范空间的单位球不全有界，所以 $I_X$ 紧当且仅当 $X$ 有限维；这时 $X'$ 也有限维。逆向也可由自然等距嵌入 $X\hookrightarrow X''$ 得到。
\end{solution}

\Needspace{10\baselineskip}
\section{Riesz--Schauder 理论}\label{sec:yosida-10-05}
\noindent\textit{原书起始页：中译本 296；英文版 283。}

\begin{lemma}[闭值域与商空间估计]\label{lem:y10-closed-range}
设 $K\in\mathcal K(X)$，$\lambda\ne0$，$A=\lambda I-K$。则 $N=\ker A$ 有限维，$R(A)$ 闭，且存在 $c>0$ 使
\[
 \|Ax\|\ge c\,\operatorname{dist}(x,N)\quad(x\in X).
\]
\end{lemma}
\begin{proof}
在 $N$ 上 $K=\lambda I$，所以 $N$ 的单位球相对紧；$N$ 闭，故由 Riesz 引理它有限维。

若估计不成立，可按距离归一化并减去核向量，选 $x_n$ 使 $\operatorname{dist}(x_n,N)=1$、$\|x_n\|\le2$、$Ax_n\to0$。取 $Kx_n$ 的收敛子列，利用 $\lambda x_n=Ax_n+Kx_n$ 得 $x_n\to x$，且 $Ax=0$，与距离为 $1$ 矛盾。

若 $Ax_n\to y$，估计使 $\operatorname{dist}(x_n,N)$ 有界。选核向量 $z_n$ 使 $x_n-z_n$ 有界，取 $K(x_n-z_n)$ 收敛子列。同一恒等式给出 $x_n-z_n\to x$，于是 $Ax=y$，所以值域闭。
\end{proof}

\begin{theorem}[Fredholm 择一的可逆性部分]\label{thm:y10-alternative}
对紧算子 $K$ 及 $\lambda\ne0$，$A=\lambda I-K$ 单射当且仅当满射，此时有有界逆。因此紧算子的每个非零谱点都是特征值。
\end{theorem}
\begin{proof}
设 $A$ 单射但不满射。令 $R_n=A^nX$，$R_0=X$。因为 $A^n=\lambda^nI-K_n$，其中 $K_n$ 是紧算子的多项式且常数项为零，上一引理说明每个 $R_n$ 闭。所有包含 $R_{n+1}\subset R_n$ 均严格：若某处相等，$A^nx=A^{n+1}y$，由单射性消去 $A^n$ 得 $x=Ay$，将推出满射。

由 Riesz 引理选 $x_n\in R_n$，$\|x_n\|=1$，$\operatorname{dist}(x_n,R_{n+1})>1/2$。若 $m>n$，则
\[
 \lambda^{-1}(Kx_n-Kx_m)=x_n-x_m-\lambda^{-1}(Ax_n-Ax_m),
\]
右侧除 $x_n$ 外均在 $R_{n+1}$，所以 $\|Kx_n-Kx_m\|>|\lambda|/2$，与紧性矛盾。故单射推出满射。

若 $A$ 满射，$A'$ 单射；$K'$ 由 Schauder 定理紧，已证方向使 $A'$ 满射。若 $Ax=0$，则对任意 $f\in X'$ 可写 $f=A'g$，从而 $f(x)=g(Ax)=0$。Hahn--Banach 分离点，故 $x=0$。有界逆由开映射定理得到。
\end{proof}

\begin{theorem}[非零谱的离散性]\label{thm:y10-discrete-spectrum}
紧算子 $K$ 的非零谱点至多可数，每个均孤立且特征空间有限维。对每个 $\varepsilon>0$，满足 $|\lambda|\ge\varepsilon$ 的谱点只有有限个；故非零谱的唯一可能聚点为零。
\end{theorem}
\begin{proof}
有限维特征空间已证。若存在无穷多个不同特征值 $\lambda_n$ 且 $|\lambda_n|\ge\varepsilon$，由 $|\lambda_n|\le\|K\|$ 取子列使 $\lambda_n\to\lambda\ne0$。取对应特征向量 $v_n$，令 $E_n=\operatorname{span}(v_1,\ldots,v_n)$。不同特征值的特征向量线性无关，故可选 $y_n\in E_n$，$\|y_n\|=1$、$\operatorname{dist}(y_n,E_{n-1})>1/2$。

因 $(K-\lambda_nI)E_n\subset E_{n-1}$，对于 $m<n$，$\lambda_n^{-1}Ky_n-\lambda_m^{-1}Ky_m$ 是 $y_n$ 与 $E_{n-1}$ 中向量之和，其范数大于 $1/2$。但 $Ky_n$ 有收敛子列，且 $\lambda_n\to\lambda\ne0$，这迫使 $\lambda_n^{-1}Ky_n$ 也有收敛子列，矛盾。于是每个远离零的区域只有有限谱点，按 $\varepsilon=1/r$ 取并即得可数性与孤立性。
\end{proof}

\begin{theorem}[有限秩谱投影与相容条件]\label{thm:y10-fredholm}
设 $\lambda\ne0$ 是 $K$ 的谱点，$P$ 为其 Riesz 投影。则 $P$ 有限秩，$\lambda$ 是预解式的有限阶极点；对充分大 $r$，
\[
 X=\ker(\lambda I-K)^r\oplus R((\lambda I-K)^r).
\]
此外，对所有 $\lambda\ne0$，
\begin{align*}
 \dim\ker(\lambda I-K)&=\dim\ker(\lambda I-K'),\\
 (\lambda I-K)x=y\text{ 可解}
 &\Longleftrightarrow f(y)=0\quad\text{对所有 }f\in\ker(\lambda I-K').
\end{align*}
有解时全体解为一个特解加上 $\ker(\lambda I-K)$。
\end{theorem}
\begin{proof}
取只包围 $\lambda$ 且不包围零的小圆 $\Gamma$。由 $R(z,K)=z^{-1}I+z^{-1}R(z,K)K$，得到
\[
 P=\frac1{2\pi i}\int_\Gamma z^{-1}R(z,K)K\,\dd z.
\]
被积算子紧，积分是有限和的算子范数极限，故 $P$ 紧。投影在闭空间 $PX$ 上是恒等算子，所以 $PX$ 有限维。定理 \ref{thm:y8-riesz-projection} 和 \ref{thm:y8-pole} 给出极点、幂零部分和直和分解。

令 $E=PX$、$F=(I-P)X$。$A=\lambda I-K$ 在 $F$ 上可逆，在 $E$ 上为有限维算子。因此 $\dim\ker A=\operatorname{codim}R(A)$，后者由闭值域的湮灭子对应等于 $\dim\ker A'$。若 $\lambda\in\rho(K)$，等式两侧同为零。又 Hahn--Banach 给出闭子空间 $R(A)={}^{\perp}\ker A'$，这恰是相容条件。解集描述由线性性直接得到。
\end{proof}

这里的维数等式是普通特征空间的维数相等；$\operatorname{rank}P$ 则是代数重数，可能严格大于普通特征空间维数。对 Hilbert 伴随，相容条件中的核写成 $\ker(\overline\lambda I-K^*)$，共轭来自 Riesz 对应。

\begin{proposition}[幂紧推广]\label{prop:y10-power-compact}
若 $V\in\mathcal L(X)$ 且某个 $V^q$ 紧，则上述非零谱结构、闭值域与 Fredholm 相容条件仍对 $V$ 成立。
\end{proposition}
\begin{proof}
对 $\lambda\ne0$，若 $\lambda^q\notin\sigma(V^q)$，由
\[
 (\lambda I-V)(\lambda^{q-1}I+\lambda^{q-2}V+\cdots+V^{q-1})
 =\lambda^qI-V^q
\]
及各因子交换可知 $\lambda I-V$ 可逆。若 $\lambda^q\in\sigma(V^q)$，取 $V^q$ 在该谱点的有限秩投影 $Q$。它与 $V$ 交换。$QX$ 有限维，而在 $(I-Q)X$ 上 $\lambda^qI-V^q$ 可逆，故 $\lambda I-V$ 也可逆。问题完全归结为 $QX$ 上的矩阵。因此非零谱均由有限维块产生，闭值域、零指数与极点性质随之成立。远离零的谱点有限，由其 $q$ 次幂的有限性得到。
\end{proof}

\begin{exercise}[共振的相容条件]\label{ex:y10-resonance}
在 $\ell^2$ 上取 $Ke_n=n^{-1}e_n$。解 $(I-K)x=y$ 并给出可解条件及全体解。
\end{exercise}
\begin{solution}
第一坐标要求 $y_1=0$；对 $n\ge2$，$x_n=y_n/(1-1/n)$，乘子至多为 $2$，故所得尾部属于 $\ell^2$。$x_1$ 任意。这正是 $y\perp\ker(I-K)$ 的条件，其中 $\ker(I-K)=\operatorname{span}\{e_1\}$。
\end{solution}

\begin{exercise}[零谱点的例外]\label{ex:y10-zero}
证明无限维 Banach 空间上的紧算子必有 $0\in\sigma(K)$，并说明零不一定是特征值。
\end{exercise}
\begin{solution}
若 $K$ 有有界逆，则 $I=K^{-1}K$ 紧，迫使空间有限维。$\ell^2$ 上 $Ke_n=n^{-1}e_n$ 单射，故其零谱点不是特征值；同理第八章的 Volterra 算子给出另一个例子。
\end{solution}

\begin{exercise}[代数重数与几何重数]\label{ex:y10-jordan}
取 $K=\left(\begin{smallmatrix}a&1\\0&a\end{smallmatrix}\right)\oplus0$ 作用于 $\mathbb C^2\oplus\ell^2$，$a\ne0$。求 $a$ 的 Riesz 投影、两种重数和预解式极点阶数。
\end{exercise}
\begin{solution}
$K$ 有限秩，故紧。$P=I_{\mathbb C^2}\oplus0$，代数重数为 $2$，几何重数为 $1$。左上块逆为 $(z-a)^{-1}I_2+(z-a)^{-2}E_{12}$，故是二阶极点；无限维补空间在 $a$ 附近的逆 $z^{-1}I$ 全纯。
\end{solution}

\Needspace{10\baselineskip}
\section{Dirichlet 问题}\label{sec:yosida-10-06}
\noindent\textit{原书起始页：中译本 299；英文版 286。}

原书用高阶强椭圆算子说明紧性与可解性的联系。这里先在变分形式中给出完整论证，再说明它与分布方程的对应。设 $G\subset\R^d$ 有界开，$V=H_0^m(G)$，$H=L^2(G)$，$J:V\to H$ 为自然嵌入。令 $V^\dagger$ 表示连续共轭线性泛函空间，使 $u\mapsto a(u,\cdot)$ 线性。

\begin{theorem}[Gårding 形式的 Fredholm 择一]\label{thm:y10-dirichlet}
设 $a:V\times V\to\mathbb C$ 是有界半双线性形式，对第一变量线性，且存在 $c>0$、$\beta\ge0$ 使
\[
 \operatorname{Re}a(u,u)+\beta\|Ju\|_H^2\ge c\|u\|_V^2.
\]
则 $Au=a(u,\cdot):V\to V^\dagger$ 的值域闭，核与余核有限维且维数相同。方程 $a(u,v)=F(v)$ 对所有 $v\in V$ 可解，当且仅当
\[
 F(v)=0\quad\text{对所有满足 }a(w,v)=0\ (w\in V)\text{ 的 }v.
\]
特别地，或者齐次问题有非零解，或者对每个 $F\in V^\dagger$ 都有唯一解，且解连续依赖 $F$。
\end{theorem}
\begin{proof}
取实数 $\gamma>\beta$，令 $b(u,v)=a(u,v)+\gamma(Ju,Jv)_H$。由 Lax--Milgram 定理 \ref{thm:y3-lax-milgram}，$B:u\mapsto b(u,\cdot)$ 是 $V$ 到 $V^\dagger$ 的有界同构。定义 $M:V\to V^\dagger$ 为 $Mu(v)=(Ju,Jv)_H$。因为 $J$ 紧而 $H\to V^\dagger$，$h\mapsto(h,J\cdot)_H$ 有界，$M$ 紧。于是
\[
 A=B(I-C),\qquad C=\gamma B^{-1}M\in\mathcal K(V).
\]
Riesz--Schauder 理论说明 $I-C$ 以及 $A$ 值域闭，核与余核维数相同且有限。

为得到具体相容条件，给 $V$ 任选其标准 Hilbert 内积，写 $a(u,v)=(\mathcal Au,v)_V$，$F(v)=(g,v)_V$。因 $\mathcal A$ 值域闭，$g\in R(\mathcal A)$ 等价于 $g\perp\ker\mathcal A^*$。而 $v\in\ker\mathcal A^*$ 恰指 $a(w,v)=0$ 对所有 $w$ 成立，故得到所写条件。若核为零，则余核为零，$A$ 双射，有界逆定理给出稳定估计。
\end{proof}

\begin{example}[高阶算子与非齐次边界数据]
设
\[
 Lu=\sum_{|\alpha|,|\eta|\le m}(-1)^{|\alpha|}
 \partial^\alpha\bigl(a_{\alpha\eta}(x)\partial^\eta u\bigr),\quad
 a(u,v)=\sum_{\alpha,\eta}\int_G a_{\alpha\eta}
 \partial^\eta u\,\overline{\partial^\alpha v}\,\dd x.
\]
假定系数有界，并明确假定此形式满足上述 Gårding 不等式；第 \ref{sec:yosida-06-08} 节提供相应强椭圆、系数正则条件下的来源。给定 $f\in L^2(G)$ 与一个边界提升 $u_1\in H^m(G)$，寻求 $u-u_1\in V$。写 $u=u_1+w$，则
\[
 a(w,v)=(f,Jv)_H-a(u_1,v)=F(v),\qquad v\in V.
\]
右侧属于 $V^\dagger$，故上一定理给出精确可解条件。在形式对称时，相容条件直接对所有齐次解 $v$ 检验；若齐次核为零，每对 $f,u_1$ 都产生唯一弱解，并有
\[
 \|w\|_{H^m}\le C\bigl(\|f\|_2+\|u_1\|_{H^m}\bigr).
\]
测试 $v\in C_c^\infty(G)$ 并分部积分可知这是分布意义的 $Lu=f$。反之，在 $u-u_1\in H_0^m(G)$ 的前提下，分布等式由测试函数稠密性延至 $V$。这里没有把弱解自动提升为 $H^{2m}$ 解；那需要额外正则性与边界假设。任意开集上的 $H_0^m$ 条件也不直接写成逐点边界导数条件。
\end{example}

原书采取 $I+\alpha L$ 的可逆平移，再构造紧算子 $S_1T$；本节用 $A+\gamma M$ 的变分同构实现同一机制。两种写法均先取得强制性，再借助紧嵌入撤去平移。非齐次问题只有在齐次核为零时才对所有数据唯一可解。

\begin{exercise}[区间上的共振]\label{ex:y10-dirichlet-resonance}
在区间 $(0,\pi)$ 上考察零边界问题
\[
 -u''-u=f,\qquad u\in H_0^1(0,\pi).
\]
证明它可解当且仅当 $\int_0^\pi f(x)\sin x\,\dd x=0$，并描述全体解。
\end{exercise}
\begin{solution}
形式为 $a(u,v)=\int u'\overline{v'}-\int u\overline v$，加 $2\|u\|_2^2$ 后强制。齐次弱解满足分布方程 $u''=-u$，在一维可积分或由常系数方程得 $u=A\sin x+B\cos x$；零端点使 $B=0$。形式对称，故相容条件正为对 $\sin x$ 正交。有解时任一特解加 $A\sin x$ 为全体解。
\end{solution}

\begin{exercise}[强制情形没有共振]\label{ex:y10-coercive}
设 $a(u,u)=\|\nabla u\|_2^2+\|u\|_2^2$，$u\in H_0^1(G)$。证明每个 $F\in V^\dagger$ 都有唯一解并给出估计。
\end{exercise}
\begin{solution}
用 $\|u\|_V^2=\|u\|_2^2+\|\nabla u\|_2^2$，形式强制常数为 $1$。Lax--Milgram 给出唯一解；取 $v=u$ 得 $\|u\|_V^2\le\|F\|\|u\|_V$，故 $\|u\|_V\le\|F\|$。此时不必再用紧性撤去平移。
\end{solution}

\Needspace{10\baselineskip}
\section*{本章附录：A. Grothendieck 的核空间}
\phantomsection\addcontentsline{toc}{section}{本章附录：A. Grothendieck 的核空间}
\label{app:yosida-10}
\noindent\textit{原书页码：中译本 302--306；英文版 289--293。}

核空间把一个算子的可和分解提升为局部凸拓扑的性质。以下设 $E$ 为 Hausdorff 局部凸空间。对连续半范数 $p$，记 $E_p=E/\ker p$，范数为 $p$，其 Banach 完备化记为 $\widehat E_p$。若 $q\ge p$，则有范数至多为 $1$ 的自然映射 $i_{q,p}:\widehat E_q\to\widehat E_p$；它未必单射。

\begin{definition}[局部凸空间到 Banach 空间的核映射]
连续线性映射 $T:E\to Y$ 称为核映射，若存在等度连续的 $f_n\in E'$、有界的 $y_n\in Y$ 及 $c\in\ell^1$，使
\[
 Tx=\sum_n c_nf_n(x)y_n.
\]
等度连续意味着存在连续半范数 $p$ 使 $|f_n(x)|\le p(x)$。级数在每个 $p$-有界集上一致绝对收敛，且
\[
 \Bigl\|Tx-\sum_{n\le N}c_nf_n(x)y_n\Bigr\|
 \le p(x)\sup_n\|y_n\|\sum_{n>N}|c_n|.
\]
因此 $T$ 连续，并将某个零邻域送入相对紧集。
\end{definition}

\begin{definition}[核空间]\label{def:y10-nuclear-space}
$E$ 称为核空间，若对每个连续半范数 $p$，存在连续半范数 $q\ge p$，使 $i_{q,p}:\widehat E_q\to\widehat E_p$ 为 Banach 空间间的核算子。这与原书用凸平衡零邻域及其 Minkowski 泛函的定义一致。
\end{definition}

\begin{theorem}[到 Banach 空间的映射判别]\label{thm:y10-nuclear-characterization}
$E$ 为核空间，当且仅当从 $E$ 到任意 Banach 空间的每个连续线性映射都为核映射。
\end{theorem}
\begin{proof}
若 $E$ 核且 $T:E\to Y$ 连续，选连续半范数 $p$ 使 $\|Tx\|\le p(x)$。$T$ 因子分解为 $E\to\widehat E_p\to Y$。再选核映射 $i_{q,p}$，便有分解 $E\to\widehat E_q\to\widehat E_p\to Y$，其可和表示拉回 $E$ 后是核表示。

反之，对自然映射 $j_p:E\to\widehat E_p$ 使用假设，写 $j_px=\sum c_nf_n(x)y_n$。取连续半范数 $q_0$ 使 $|f_n(x)|\le q_0(x)$，令 $q=p+q_0$。每个 $f_n$ 都在 $\ker q$ 上为零，故延为 $\widehat E_q$ 上范数至多为 $1$ 的泛函。级数因此定义核算子 $\widehat E_q\to\widehat E_p$，它在稠密的 $E/\ker q$ 上恰是自然映射，故等于 $i_{q,p}$。
\end{proof}

\begin{proposition}[Hilbert 半范数的来源]\label{prop:y10-hilbert-seminorm}
若 $E$ 核，对任意连续半范数 $p$，都存在由内积诱导的连续半范数 $r\ge p$。因而核空间的拓扑可由 Hilbert 半范数生成。
\end{proposition}
\begin{proof}
将自然映射写成 $j_px=\sum c_nf_n(x)y_n$，吸收相位后可取 $c_n\ge0$，并设 $|f_n(x)|\le q(x)$。定义
\[
 \alpha x=(\sqrt{c_n}f_n(x))_n\in\ell^2,\qquad
 \beta t=\sum_n\sqrt{c_n}t_ny_n\in\widehat E_p.
\]
Cauchy--Schwarz 给出 $\|\alpha x\|_2\le(\sum c_n)^{1/2}q(x)$，且 $\|\beta\|\le(\sum c_n\|y_n\|^2)^{1/2}$。于是 $j_p=\beta\alpha$，从而 $p(x)\le\|\beta\|\|\alpha x\|_2$。令 $r(x)=\max(1,\|\beta\|)\|\alpha x\|_2$ 即可。商掉 $\ker r$ 后是预 Hilbert 空间，其完备化为 Hilbert 空间。
\end{proof}

\begin{example}[乘积空间与快速下降序列]
对任意指标集 $I$，赋乘积拓扑的 $\mathbb C^I$ 为核空间。每个连续半范数 $p$ 都满足 $p(x)\le C\max_{i\in F}|x_i|$，其中 $F$ 有限：这是因为一个控制 $p$ 的基本零邻域仅限制有限坐标，其他坐标可任意伸缩。选右侧为 $q$，则 $\widehat E_q$ 有限维，$i_{q,p}$ 有限秩，因而核。

令 $s$ 为快速下降序列空间，使用 Hilbert 范数
\[
 p_k(x)^2=\sum_{n\ge1}(1+n)^{2k}|x_n|^2,\qquad k=0,1,\ldots.
\]
这些范数与通常的加权上确界半范数定义相同拓扑：上确界由同阶 $p_k$ 控制，而 $p_k$ 由更高一阶的加权上确界控制。$\widehat s_{p_k}$ 是加权 $\ell^2$，其间自然嵌入 $p_{k+2}\to p_k$ 的归一化对角系数为 $(1+n)^{-2}$，绝对可和。任意连续半范数都被某个 $Cp_k$ 控制，再与该核嵌入复合，得到 $s$ 为核空间。
\end{example}

\begin{proposition}[Banach 核空间只能有限维]
若 Banach 空间 $E$ 带其范数拓扑且为核空间，则 $\dim E<\infty$。
\end{proposition}
\begin{proof}
由定理 \ref{thm:y10-nuclear-characterization}，恒等映射 $I_E$ 核，因而紧，Riesz 引理迫使 $E$ 有限维。反之有限维恒等映射有有限秩表示，所以有限维空间核。
\end{proof}

原书还列出核性的子空间、Hausdorff 商、乘积及适当归纳极限稳定性，进而联系 $\Test$、$\Smooth$、$\Sch$ 及其强对偶的核性。这些一般稳定性定理的证明本附录未展开，也不从上述两个例子直接推出。原书末尾的 Minlos 定理把核空间上的适当连续柱分布延拓为测度；仅有有限维相容性并不足以省去连续性条件。本讲义不在此陈述其全部版本或证明，后续专题如需使用应另列准确假设。

\begin{exercise}[等度连续不可省略]\label{ex:y10-equicontinuity}
在 $E=\mathbb C^{\mathbb N}$ 的乘积拓扑上，坐标泛函 $f_n(x)=x_n$ 各自连续。证明 $\{f_n\}$ 不等度连续。
\end{exercise}
\begin{solution}
每个零邻域都包含只限制有限坐标的基本零邻域。任选未受限制的坐标 $n$，该坐标可任意大，故不存在零邻域使所有 $|f_n(x)|\le1$。因此逐个连续不等于等度连续；甚至 $\sum2^{-n}x_n$ 也不对所有乘积空间元素收敛。
\end{solution}

\begin{exercise}[自然映射不必单射]\label{ex:y10-canonical-kernel}
在 $E=\mathbb C^2$ 上取 $q(x)=|x_1|+|x_2|$、$p(x)=|x_1|$。计算 $i_{q,p}$ 的核，并说明这为何不妨碍核空间定义。
\end{exercise}
\begin{solution}
$\widehat E_q=\mathbb C^2$，$\widehat E_p\cong\mathbb C$，自然映射为 $(x_1,x_2)\mapsto x_1$，核为 $\{0\}\times\mathbb C$。核空间的条件要求映射具有可和秩一展开，并不要求其单射；此例为有限秩映射，当然满足要求。
\end{solution}


\Needspace{6\baselineskip}
\section*{本章参考说明与阅读衔接}
\phantomsection\addcontentsline{toc}{section}{本章参考说明与阅读衔接}
直接来源为吉田第十章及章内附录 \cite{yosida-zh,yosida-en}；原书所引 Grothendieck、Gårding、Minlos 等原始文献未计为直接阅读。讲义补充有限网、商空间估计与变分形式的推导。核空间的一般稳定性和 Minlos 定理仅说明位置，未在此证明。下一章将从有限秩谱投影进一步走向正规与自伴算子的谱表示。

% !TeX encoding = UTF-8
\SetChapterInfo{subtitle={从极大理想到谱积分与调和分析},keywords={Gelfand 表示; 谱测度; 自伴算子; Friedrichs; 紧群; Tannaka; Wiener},prereq={预解式; Cayley 变换; 紧算子; Riesz--Markov; Stone--Weierstrass},goals={构造谱表示; 控制无界演算定义域; 连接群表示与 Fourier 分析}}
\chapter[赋范环和谱表示]{赋范环\\和谱表示}\label{ch:yosida-11}
\ChapterSource{按 \cite{yosida-zh,yosida-en} 第十一章十六节组织，2026 年 10 月 27 日补写。直接查看英文印刷 294--362 页扫描图，第 362 页上部为本章结尾、下部开始第十二章，故更正本章英文范围。中译本参考范围为 307--376 页，选页及证明边界详见本章 COVERAGE。}

第十章的紧性使非零谱问题归结为有限维空间，本章转向一般正规与自伴算子。第一条线索是把交换 Banach 代数的元素表示为函数；第二条线索是用投影值测度把这些函数重新解释为算子；第三条线索则把谱分解用于群表示、正定函数和 Tauber 型逼近。

本章 $H$ 为非零复 Hilbert 空间，内积对第一变量线性。$\mathcal A$ 表示代数，$A$ 或 $S$ 表示算子，$E$ 表示投影值测度。实轴谱族统一采用右连续约定 $E_t=E((-\infty,t])$。Fourier 变换仍取 $\widehat f(\xi)=\int e^{-ix\xi}f(x)\,\dd x$，原书约定在对应公式中换算。测度表示的外部前提在使用处注明。

\Needspace{10\baselineskip}
\section{赋范环的极大理想}\label{sec:yosida-11-01}
\noindent\textit{原书起始页：中译本 309；英文版 295。}

\begin{definition}[Banach 代数与特征]
复 Banach 空间 $\mathcal A$ 若具有结合、双线性的乘法且 $\|ab\|\le\|a\|\|b\|$，称为 Banach 代数。以下先假定其交换、有单位 $1$ 且 $\|1\|=1$。特征指非零复线性乘法泛函 $\chi:\mathcal A\to\mathbb C$，全体记为 $\Delta(\mathcal A)$。真理想中的极大元称为极大理想。
\end{definition}

本节的理想是满足 $\mathcal A M\subset M$ 的线性子空间 $M$；“真”表示 $M\ne\mathcal A$。有单位时，这等价于 $1\notin M$。

\begin{theorem}[极大理想与特征对应]\label{thm:y11-characters}
每个真理想包含于某极大理想。每个极大理想 $M$ 闭，$\mathcal A/M\cong\mathbb C$，且唯一对应 $\chi\in\Delta(\mathcal A)$ 使 $M=\ker\chi$。每个特征自动连续，$\|\chi\|=\chi(1)=1$。
\end{theorem}
\begin{proof}
用包含给定真理想的真理想族及 Zorn 引理取极大元；链的并仍真，因为它不含 $1$。若 $\overline M=\mathcal A$，则有 $a\in M$ 使 $\|1-a\|<1$；Neumann 级数使 $a$ 可逆，导致 $1\in M$，矛盾。闭包为理想，由极大性 $M=\overline M$。

商代数是复 Banach 除代数。对其元素 $b$，谱非空：若预解式处处存在，则在无穷远趋零，逐连续泛函应用 Liouville 定理会使预解式恒为零，矛盾。取 $\lambda\in\sigma(b)$，除代数中不可逆元素只能为零，故 $b=\lambda1$。这给出 Gelfand--Mazur 结论及商特征。反之，特征的核的商是一维域，所以核极大。

对任意特征，$\chi(1)=1$；$a-\chi(a)1$ 不可逆，否则对其逆应用 $\chi$ 得矛盾。又 $|\lambda|>\|a\|$ 时 Neumann 级数给出逆，故 $|\chi(a)|\le\|a\|$。结合单位值，得到范数为 $1$，无需预先假定连续性。
\end{proof}

\begin{example}[连续函数代数]
若 $K$ 紧 Hausdorff，则 $C(K)$ 的特征恰为点值 $f\mapsto f(t)$。若极大理想 $M$ 的函数无公共零点，紧性给出 $f_1,\ldots,f_n\in M$ 使 $h=\sum|f_j|^2>0$；$h\in M$ 且可逆，矛盾。因此 $M\subset\{f:f(t)=0\}$，极大性使它们相等。
\end{example}

\begin{exercise}[单位化]\label{ex:y11-unitization}
对无单位 Banach 代数 $\mathcal B$，在 $\mathbb C\oplus\mathcal B$ 上写出乘法与范数，使 $(1,0)$ 成为单位，并验证乘法估计。
\end{exercise}
\begin{solution}
定义 $(\lambda,a)(\mu,b)=(\lambda\mu,\lambda b+\mu a+ab)$，范数 $|\lambda|+\|a\|$。三角不等式和 $\|ab\|\le\|a\|\|b\|$ 给出积的范数至多两范数之积；完备性来自直和完备性。
\end{solution}

\begin{exercise}[交换性的位置]\label{ex:y11-noncommutative}
说明矩阵代数 $M_2(\mathbb C)$ 不能套用本节极大理想到复数特征的对应。
\end{exercise}
\begin{solution}
$M_2(\mathbb C)$ 的双边理想只有零与全代数，故零理想极大，但商仍为矩阵代数而非域。若有复特征，其核只能为零，得到从四维复空间到一维空间的单射，矛盾。交换性用于将极大理想商变成除代数。
\end{solution}

\Needspace{10\baselineskip}
\section{根、半单性}\label{sec:yosida-11-02}
\noindent\textit{原书起始页：中译本 312；英文版 298。}

\begin{definition}[Gelfand 变换与根]
对 $a\in\mathcal A$ 定义 $\widehat a(\chi)=\chi(a)$。赋 $\Delta(\mathcal A)$ 逐点收敛拓扑，称为特征空间。根为 $\operatorname{rad}\mathcal A=\bigcap_{\chi}\ker\chi$；根为零时称半单。
\end{definition}

\begin{theorem}[交换 Banach 代数的函数表示]\label{thm:y11-gelfand}
$\Delta(\mathcal A)$ 为紧 Hausdorff 空间；Gelfand 变换是到 $C(\Delta(\mathcal A))$ 的连续代数同态，并且
\[
 \sigma_{\mathcal A}(a)=\{\chi(a):\chi\in\Delta(\mathcal A)\},\qquad
 \|\widehat a\|_\infty=r(a)=\lim_n\|a^n\|^{1/n}.
\]
其核为根。因此交换代数中，根由所有拟幂零元素组成；半单性等价于 Gelfand 变换单射，不自动等价于等距。
\end{theorem}
\begin{proof}
把特征嵌入 $\prod_{a\in\mathcal A}\{z:|z|\le\|a\|\}$。线性、乘法与 $\chi(1)=1$ 均为闭条件，故特征空间是紧乘积的闭子集，且坐标函数分离点。若 $a-\lambda1$ 不可逆，其生成理想真，故包含于某极大理想；相应特征满足 $\chi(a)=\lambda$。反向已在上一节证明。谱半径公式的证明与第八章相同：Neumann 展开与预解式的 Cauchy 系数估计只用代数乘法及完备性，因此适用于 $\mathcal A$。由谱等式即得范数等式，核和根的断言随定义得到。
\end{proof}

\begin{theorem}[交换算子 C 星代数]\label{thm:y11-commutative-cstar}
设 $\mathcal A\subset\mathcal L(H)$ 为含单位、范数闭、对伴随封闭的交换子代数。则 Gelfand 变换是到 $C(\Delta(\mathcal A))$ 的等距满 $*$ 同构。
\end{theorem}
\begin{proof}
每个元素正规。由 $\|T^*T\|=\|T\|^2$ 及交换性，反复平方得 $\|T^{2^n}\|=\|T\|^{2^n}$，故 $r(T)=\|T\|$。于是 Gelfand 变换等距，像闭。若 $S=S^*$，$e^{itS}$ 酉，所以 $|e^{it\chi(S)}|\le1$ 对正负 $t$ 均成立，迫使 $\chi(S)$ 实。分解 $T$ 为实部与虚部，得 $\chi(T^*)=\overline{\chi(T)}$。像因而含常数、分离点且对共轭封闭，Stone--Weierstrass 定理使其一致稠密；闭性给出满射。
\end{proof}

\begin{exercise}[非零根]\label{ex:y11-radical}
在代数 $\mathbb C[\varepsilon]/(\varepsilon^2)$ 上取范数 $\|a+b\varepsilon\|=|a|+|b|$。求特征、根和谱。
\end{exercise}
\begin{solution}
特征使 $\chi(\varepsilon)^2=0$，故唯一特征为 $a+b\varepsilon\mapsto a$。根为 $\mathbb C\varepsilon$，$\sigma(a+b\varepsilon)=\{a\}$。非零元素 $\varepsilon$ 幂零，说明一般交换 Banach 代数不一定半单。
\end{solution}

\begin{exercise}[Wiener 逆定理的离散版本]\label{ex:y11-wiener-discrete}
证明 $\ell^1(\mathbb Z)$ 卷积代数的特征为 $a\mapsto\sum_n a_nz^n$，$|z|=1$；从而绝对收敛 Fourier 级数处处非零时，其倒数也有绝对收敛 Fourier 级数。
\end{exercise}
\begin{solution}
令 $z=\chi(\delta_1)$，因 $\delta_1\delta_{-1}=\delta_0$，有 $|z|\le1$ 且 $|z^{-1}|\le1$，故 $|z|=1$。由有限支集序列稠密及特征连续性得到级数公式。反向这些公式显然为特征。处处非零即在每个极大理想商中非零，所以元素可逆；其逆属于 $\ell^1$，给出所求倒数级数。
\end{solution}

\Needspace{10\baselineskip}
\section{有界正规算子的谱分解}\label{sec:yosida-11-03}
\noindent\textit{原书起始页：中译本 315；英文版 302。}

\begin{theorem}[连续函数演算]\label{thm:y11-continuous-calculus}
若 $N\in\mathcal L(H)$ 正规，则存在唯一等距含单位 $*$ 同态
\[
 C(\sigma(N))\longrightarrow\mathcal L(H),\qquad f\longmapsto f(N),
\]
把坐标函数 $z$ 送到 $N$，并满足 $\|f(N)\|=\|f\|_\infty$。
\end{theorem}
\begin{proof}
令 $\mathcal A=C^*(I,N)$，即 $N,N^*$ 的多项式的范数闭包。它交换，故应用定理 \ref{thm:y11-commutative-cstar}。映射 $\chi\mapsto\chi(N)$ 单射，因为这些多项式稠密。还需验证 $\sigma_{\mathcal A}(N)=\sigma_{\mathcal L(H)}(N)$。

若 $T\in\mathcal A$ 在 $\mathcal L(H)$ 中可逆，则 $S=T^*T$ 有 $cI\le S\le MI$，$c>0$。对自伴算子的二次型范数公式给出 $\|I-S/M\|\le1-c/M<1$，故 $S^{-1}$ 的 Neumann 级数属于 $\mathcal A$。于是 $T^{-1}=S^{-1}T^*\in\mathcal A$。取 $T=N-\lambda I$ 便得谱相等，故特征空间与 $\sigma(N)$ 同胚。逆 Gelfand 变换给出演算，唯一性由 $z,\bar z$ 多项式在 $C(\sigma(N))$ 中稠密得到。
\end{proof}

\begin{theorem}[正规算子的投影值测度]\label{thm:y11-normal-pvm}
存在唯一 Borel 投影值测度 $E$，使 $E(\sigma(N))=I$，
\[
 E(B)E(C)=E(B\cap C),\qquad
 N=\int_{\sigma(N)}z\,\dd E(z).
\]
$E$ 的可列可加性采用强算子拓扑。对任意有界 Borel 函数 $f$，
\[
 \|f(N)x\|^2=\int|f|^2\,\dd\mu_x,\qquad
 \mu_x(B)=(E(B)x,x).
\]
\end{theorem}
\begin{proof}
先固定 $x$。泛函 $f\mapsto(f(N)x,x)$ 为正，因为非负连续函数有连续平方根。用外部前提 Riesz--Markov 定理 \ref{thm:y4-riesz-markov} 得到有限正测度 $\mu_x$。映射 $f\mapsto f(N)x$ 满足
\[
 (f(N)x,g(N)x)=\int f\overline g\,\dd\mu_x.
\]
连续函数在 $L^2(\mu_x)$ 中稠密，因此该映射延为 $L^2(\mu_x)$ 到循环空间 $H_x=\overline{C(\sigma(N))x}$ 的酉映射。在此映射下 $N$、$N^*$ 分别为乘 $z$、乘 $\bar z$，所以 $H_x$ 为约化子空间。

以 Zorn 引理取极大正交循环空间族；若其正交和不为 $H$，补空间仍约化，可再取一个非零循环空间，矛盾。因此 $H=\bigoplus_\alpha H_{x_\alpha}$。在每个 $L^2(\mu_{x_\alpha})$ 上以乘 $1_B$ 定义 $E(B)$，再作正交直和。乘法关系和强可加性由指示函数及控制收敛得到。即使直和不可数，每个向量也只有至多可数个非零分量，故此步骤无可分性要求。相同构造定义 $f(N)$ 并给出范数公式。

若另有这样的投影值测度，积分多项式先给出同一连续演算。各向量的标量测度由连续积分唯一决定；极化后得到所有 $(E(B)x,y)$，故投影测度唯一。其支集恰为 $\sigma(N)$，否则一个非零连续函数被演算为零，会违背等距性。
\end{proof}

本节将原书的交换正规族结论完整证明到单个正规算子的演算及谱测度。有限个两两交换且连同伴随一起交换的正规算子也可对其生成的交换 C 星代数作同样构造，得到联合谱上的投影值测度；任意族的原书双交换子与范畴集构造不在此重写。

\begin{exercise}[正规性不可省略]\label{ex:y11-jordan-normal}
证明 $J=\left(\begin{smallmatrix}0&1\\0&0\end{smallmatrix}\right)$ 没有使 $J=\int z\,\dd E(z)$ 的正交投影值谱表示。
\end{exercise}
\begin{solution}
投影值积分与其伴随分别为乘 $z$ 与乘 $\bar z$，必交换，故必正规。但 $JJ^*=\operatorname{diag}(1,0)$，$J^*J=\operatorname{diag}(0,1)$，不相等。
\end{solution}

\begin{exercise}[投影的谱测度]\label{ex:y11-projection-pvm}
设 $P=P^*=P^2$。写出其谱测度和 $f(P)$。
\end{exercise}
\begin{solution}
$E(B)=1_B(0)(I-P)+1_B(1)P$，故 $f(P)=f(0)(I-P)+f(1)P$。若某个子空间为零，相应谱点可从支集中删去。
\end{solution}

\Needspace{10\baselineskip}
\section{酉算子的谱分解}\label{sec:yosida-11-04}
\noindent\textit{原书起始页：中译本 319；英文版 306。}

\begin{theorem}[单位圆上的谱积分]\label{thm:y11-unitary-pvm}
若 $U$ 酉，则其谱测度支于单位圆 $\mathbb T$，且对所有 $n\in\mathbb Z$，
\[
 U^n=\int_{\mathbb T}z^n\,\dd E(z),\qquad U^*=\int_{\mathbb T}\bar z\,\dd E(z).
\]
反之，单位圆上总质量为 $I$ 的投影值测度的坐标积分是酉算子。
\end{theorem}
\begin{proof}
$U$ 正规，$U^*U=I$ 使连续演算中 $|z|^2=1$ 于谱上；也可对 $|\lambda|>1$ 用 $U$ 的 Neumann 级数，对 $|\lambda|<1$ 用 $U^{-1}$ 的级数得谱包含。各幂由演算乘法得到，反向由积分 $z\bar z=1$ 得 $U^*U=UU^*=I$。
\end{proof}

选定 Borel 辐角 $\operatorname{Arg}:\mathbb T\to[0,2\pi)$，定义 $B=\operatorname{Arg}(U)$。则 $B$ 为有界自伴算子，$U=e^{iB}$，但这个对数依赖所选分支，不能据此宣称唯一。若用实参数谱族 $F_t=E(\{z:\operatorname{Arg}z\le t\})$，应把 $z=1$ 的质量计入 $t=0$；原书使用另一端点记法时须整体换算，不能重复计算该原子。

\begin{example}[平移的谱表示]
在 $\ell^2(\mathbb Z)$ 上取双边移位 $(Sx)_n=x_{n-1}$。圆周取归一化测度 $\dd\theta/(2\pi)$，则酉 Fourier 级数映射 $\mathcal Fx(\theta)=\sum_nx_ne^{in\theta}$ 将它变为乘 $e^{i\theta}$。因此谱为 $\mathbb T$，谱投影为 $\mathcal F^{-1}1_B(e^{i\theta})\mathcal F$。单点测度为零，所以没有特征值。
\end{example}

\begin{exercise}[谱投影的原子]\label{ex:y11-unitary-atom}
证明 $E(\{z_0\})H=\ker(U-z_0I)$。
\end{exercise}
\begin{solution}
由 $\|(U-z_0)x\|^2=\int|z-z_0|^2\,\dd\mu_x$，该范数为零当且仅当 $\mu_x$ 集中于 $z_0$。这又等价于 $E(\{z_0\})x=x$。
\end{solution}

\begin{exercise}[酉对数不唯一]\label{ex:y11-unitary-log}
给出两个不同的有界自伴算子 $B_1,B_2$，使 $e^{iB_1}=e^{iB_2}=I$。
\end{exercise}
\begin{solution}
取 $B_1=0$，$B_2=2\pi P$，其中 $P$ 为任意非零正交投影。连续演算给出 $e^{2\pi iP}=(I-P)+P=I$。
\end{solution}

\Needspace{10\baselineskip}
\section{单位分解}\label{sec:yosida-11-05}
\noindent\textit{原书起始页：中译本 322；英文版 309。}

\begin{definition}[实轴单位分解]
一族正交投影 $E_t$ 称为右连续单位分解，若
\[
 E_sE_t=E_{\min(s,t)},\quad
 \mathop{\mathrm{s\!-!lim}}_{t\to-\infty}E_t=0,\quad
 \mathop{\mathrm{s\!-!lim}}_{t\to+\infty}E_t=I,\quad
 \mathop{\mathrm{s\!-!lim}}_{h\downarrow0}E_{t+h}=E_t.
\]
它对应 $E((a,b])=E_b-E_a$ 的唯一投影值测度；反过来每个实轴投影值概率测度给出此谱族。这里的“单位分解”是 resolution of identity，与光滑单位分解不是同一概念。
\end{definition}

\begin{proposition}[从谱族到测度]
对 $x,y\in H$，$t\mapsto(E_tx,y)$ 有界变差，其全变差至多 $\|x\|\|y\|$。相应标量 Stieltjes 测度 $\mu_{x,y}$ 满足 $\mu_x=\mu_{x,x}\ge0$，并定义上述唯一投影值测度。
\end{proposition}
\begin{proof}
对递增分割 $t_j$，令 $P_j=E_{t_j}-E_{t_{j-1}}$，这些投影两两正交。因此
\[
 \sum_j|(P_jx,y)|=\sum_j|(P_jx,P_jy)|
 \le\Bigl(\sum_j\|P_jx\|^2\Bigr)^{1/2}
     \Bigl(\sum_j\|P_jy\|^2\Bigr)^{1/2}\le\|x\|\|y\|.
\]
标量 Stieltjes 延拓和极化得到所有 $\mu_{x,y}$。对每个 Borel 集 $B$，由此界及 Riesz 表示定义有界算子 $E(B)$，使 $(E(B)x,y)=\mu_{x,y}(B)$。区间上的投影乘法关系，依次固定一个变量再用有界单调类延拓到所有 Borel 集，得到 $E(B)E(C)=E(B\cap C)$ 及自伴性。标量可加性先给弱可加性；不交投影的正交性使余项的范数平方等于相应正测度余项，故为强可加性。
\end{proof}

\begin{theorem}[谱积分的定义域]\label{thm:y11-spectral-integral}
对有限值 Borel 函数 $f:\R\to\mathbb C$，定义
\[
 D(T_f)=\{x:\int|f|^2\,\dd\mu_x<\infty\},\qquad
 T_fx=\lim_{n\to\infty}\int f1_{\{|f|\le n\}}\,\dd E\,x.
\]
此极限存在，定义稠密闭算子，且 $\|T_fx\|^2=\int|f|^2\,\dd\mu_x$，$T_f^*=T_{\bar f}$。实值 $f$ 给出自伴算子。
\end{theorem}
\begin{proof}
对简单函数由正交性直接得到范数公式。用有界简单函数一致逼近有界 Borel 函数，公式传到有界演算。截断之差的平方范数恰为 $|f|^2$ 在相应层集的积分，故其 Cauchy 性等价于定义域条件。投影 $P_n=E(\{|f|\le n\})$ 强趋于 $I$，且 $P_nH\subset D(T_f)$，故定义域稠密。

若 $x_j\to x$、$T_fx_j\to y$，则 $f1_{\{|f|\le n\}}(E)x=P_ny$。于是截断平方积分有界于 $\|y\|^2$，单调收敛使 $x\in D(T_f)$，再令 $n\to\infty$ 得 $T_fx=y$。这证明闭性。显然 $T_{\bar f}\subset T_f^*$。若 $y\in D(T_f^*)$，在 $P_nH$ 上比较伴随得 $T_{\bar f}P_ny=P_nT_f^*y$，同样的积分估计使 $y\in D(T_{\bar f})$。故伴随恰相等。
\end{proof}

\begin{exercise}[强收敛与范数收敛]\label{ex:y11-projection-limits}
在 $\ell^2$ 上令 $P_n$ 为前 $n$ 个坐标的投影。证明 $P_n\to I$ 强收敛，但不按算子范数收敛。
\end{exercise}
\begin{solution}
$\|(I-P_n)x\|^2=\sum_{j>n}|x_j|^2\to0$，而 $\|(I-P_n)e_{n+1}\|=1$，所以 $\|I-P_n\|=1$。
\end{solution}

\begin{exercise}[无界积分]\label{ex:y11-integral-domain}
在 $L^2(\R)$ 的坐标谱测度上取 $f(t)=t^2$。求 $D(T_f)$，并说明 $D(T_f)$ 与 $D(T_t)$ 的关系。
\end{exercise}
\begin{solution}
$D(T_f)=\{u:\int t^4|u(t)|^2\,\dd t<\infty\}$，而 $D(T_t)$ 只要求二阶加权平方积分有限。前者真包含于后者；例如 $u(t)=1_{(1,\infty)}t^{-2}$ 属于后者而不属于前者。
\end{solution}

\Needspace{10\baselineskip}
\section{自伴算子的谱分解}\label{sec:yosida-11-06}
\noindent\textit{原书起始页：中译本 327；英文版 313。}

\begin{theorem}[无界自伴谱定理]\label{thm:y11-self-adjoint}
每个自伴算子 $A$ 都有唯一的实轴投影值测度 $E$，使
\[
 A=\int_\R t\,\dd E(t),\qquad
 D(A)=\{x:\int_\R t^2\,\dd\mu_x(t)<\infty\}.
\]
对于 $z\notin\R$，
\[
 (A-zI)^{-1}=\int_\R(t-z)^{-1}\,\dd E(t).
\]
\end{theorem}
\begin{proof}
由定理 \ref{thm:y7-cayley}，$U=(A-iI)(A+iI)^{-1}=I-2i(A+iI)^{-1}$ 酉，且 $\ker(I-U)=0$。取其单位圆谱测度 $P$，则 $P(\{1\})=0$。以下映射为 Borel 同构：
\[
 c(t)=\frac{t-i}{t+i}:\R\longrightarrow\mathbb T\setminus\{1\},\qquad
 c^{-1}(z)=i\frac{1+z}{1-z}
\]
\par\Needspace{6\baselineskip}
定义 $E(B)=P(c(B))$，按上一节构造自伴算子 $A_0=\int t\,\dd E(t)$，则
\[
 (A_0+iI)^{-1}=\frac{I-U}{2i}=(A+iI)^{-1}.
\]
相同逆算子的值域给出相同定义域，再取逆得 $A_0=A$。一般非实预解式由有界乘子 $(t-z)^{-1}$ 与 $t-z$ 相乘检验双侧逆。若另有谱测度表示 $A$，经 $c$ 推回后表示同一 $U$；酉谱测度唯一，故 $E$ 唯一。
\end{proof}

若 $z=e^{i\theta}$，直接计算得 $c^{-1}(z)=-\cot(\theta/2)$。中译本此处有半角校注；不能写成 $-\cot\theta$。端点 $z=1$ 的投影质量为零，而不要求 $1$ 不属于 $U$ 的谱：无界 $A$ 常使它成为连续谱点。

\begin{example}[动量算子的定义域]
令 $\mathcal F_0=(2\pi)^{-1/2}\mathcal F$ 为 $L^2(\R)$ 上的酉 Fourier 变换。对 $P=-i\,\dd/\dd x$，$D(P)=H^1(\R)$，有
\[
 \mathcal F_0P\mathcal F_0^{-1}=M_\xi,\qquad
 E_P(B)=\mathcal F_0^{-1}1_B(\xi)\mathcal F_0.
\]
先对 Schwartz 函数分部积分，再用 Fourier 加权范数闭包得到定义域相等。$C_c^\infty$ 上同一微分表达式只是其稠密限制，不能直接把该限制称作自伴算子。
\end{example}

\begin{exercise}[Cayley 的端点]\label{ex:y11-cayley-endpoint}
对 $A=M_t$ 作用于 $L^2(\R)$，证明 $1\in\sigma(c(A))$，但 $P(\{1\})=0$。
\end{exercise}
\begin{solution}
$c(t)$ 对有限实数从不等于 $1$，故单点投影为零；但 $c(t)\to1$ 当 $|t|\to\infty$。取支于 $(n,n+1)$ 的单位向量 $u_n$，则 $\|(c(A)-I)u_n\|\to0$，故不存在有界逆，$1$ 属于谱。
\end{solution}

\begin{exercise}[谱截断是图范数逼近]\label{ex:y11-spectral-core}
设 $P_n=E([-n,n])$。证明 $P_nx\to x$ 对每个 $x\in H$ 成立，且对 $x\in D(A)$ 有 $AP_nx\to Ax$。
\end{exercise}
\begin{solution}
两项误差的平方分别为 $\mu_x(\{|t|>n\})$ 与 $\int_{|t|>n}t^2\,\dd\mu_x$，均趋于零。因此 $\bigcup_nP_nH$ 在 $D(A)$ 的图范数中稠密。
\end{solution}

\Needspace{10\baselineskip}
\section{实算子和半有界算子、Friedrichs 定理}\label{sec:yosida-11-07}
\noindent\textit{原书起始页：中译本 329；英文版 316。}

\begin{definition}[实性与半有界性]
共轭算子 $C$ 是满足 $C^2=I$ 的反线性等距映射，例如 $L^2$ 上的复共轭。若 $CD(S)=D(S)$ 且 $SC=CS$，称 $S$ 对 $C$ 为实算子。稠密定义的对称算子 $S$ 称下半有界，若存在 $m\in\R$ 使 $(Su,u)\ge m\|u\|^2$；上半有界可化为 $-S$ 下半有界。
\end{definition}

\begin{theorem}[Friedrichs 扩张]\label{thm:y11-friedrichs}
每个稠密定义、下半有界的对称算子 $S$ 都有保持同一下界的自伴扩张 $S_F$。该扩张由 $S$ 的二次型闭包唯一确定，不意味着所有自伴扩张唯一。
\end{theorem}
\begin{proof}
平移后先设 $S\ge I$，令 $q(u,v)=(Su,v)$，$u,v\in D(S)$。这是内积，且 $\|u\|\le\|u\|_q$。以 $V$ 表示其完备化，映入 $H$ 的连续延拓记为 $j$。关键是 $j$ 单射：若 $u_n$ 在 $q$ 范数中收敛于 $u\in V$，且 $u_n\to0$ 于 $H$，则对 $v\in D(S)$，
\[
 q(u_n,v)=(u_n,Sv)\longrightarrow0.
\]
故 $u$ 与 $V$ 中的稠密子空间 $D(S)$ 正交，$u=0$。这也证明原二次型可闭。因而可把 $V$ 看成 $H$ 中的稠密子空间，并把 $q$ 延为闭形式。

对 $f\in H$，Riesz 表示给出唯一 $Rf\in V$ 使 $q(Rf,v)=(f,v)$。由 $q\ge\|\cdot\|^2$ 得 $R:H\to H$ 有界、正、自伴、单射：单射使用 $V$ 在 $H$ 中稠密；自伴由 $q$ 的对称性得到。故 $R$ 值域稠密。定义 $S_F=R^{-1}$，$D(S_F)=R(H)$。它对称，且 $S_F\pm iI$ 满射，因为 $I\pm iR$ 满射：其范数下界为 $1$，伴随有同一下界，故闭值域且稠密。自伴判据给出 $S_F=S_F^*$。

若 $u\in D(S)$，则 $q(u,v)=(Su,v)$ 先对 $v\in D(S)$ 成立，再延到 $V$，故 $u=RSu$，从而 $S_Fu=Su$。又 $(S_Fu,u)=q(u,u)\ge\|u\|^2$，保持下界。一般 $m$ 情形对 $S+(1-m)I$ 应用此构造再移回。闭形式决定 $R$，故决定此扩张。
\end{proof}

\begin{proposition}[实算子的扩张]
稠密定义的实对称算子有自伴扩张。若还半有界，则 Friedrichs 扩张本身为实算子。
\end{proposition}
\begin{proof}
对闭包仍有实性。共轭将 $\ker(S^*-iI)$ 反线性等距映到 $\ker(S^*+iI)$，故亏指数相等；第 \ref{sec:yosida-11-15} 节的 Cayley 扩张判据给出自伴扩张。半有界时，$C$ 保持型范数和闭型空间 $V$，Riesz 定义的 $R$ 与 $C$ 交换，因此其逆 $S_F$ 与 $C$ 交换。
\end{proof}

\begin{example}[最小 Laplace 算子]
在有界开集 $G$ 上，$S=-\Delta$、$D(S)=C_c^\infty(G)$ 为正对称算子。平移形式 $q_1(u,v)=\int\nabla u\cdot\overline{\nabla v}+\int u\bar v$ 的闭包定义域为 $H_0^1(G)$。其 Friedrichs 扩张满足
\[
 D(S_F)=\{u\in H_0^1(G):-\Delta u\in L^2(G)\},\qquad S_Fu=-\Delta u,
\]
其中导数按分布理解。只有另有边界正则性时，才能再把定义域写成 $H^2(G)\cap H_0^1(G)$。
\end{example}

\begin{exercise}[型范数与图范数]\label{ex:y11-form-graph}
对 $Ae_n=n^2e_n$，比较型域 $D(A^{1/2})$ 与算子域 $D(A)$。
\end{exercise}
\begin{solution}
型域要求 $\sum n^2|x_n|^2<\infty$，算子域要求 $\sum n^4|x_n|^2<\infty$。取 $x_n=n^{-2}$，前者成立、后者不成立。闭型空间不是算子图范数完备化。
\end{solution}

\begin{exercise}[扩张不唯一]\label{ex:y11-friedrichs-nonunique}
在 $L^2(0,1)$ 上解释 $-\dd^2/\dd x^2$ 的最小算子为何同时有 Dirichlet 与 Neumann 自伴扩张，但 Friedrichs 扩张选择前者。
\end{exercise}
\begin{solution}
两种定义域分别为 $H^2\cap H_0^1$ 与 $\{u\in H^2:u'(0)=u'(1)=0\}$；分部积分的边界型和非实边值问题给出自伴性，也可分别用正弦、余弦基对角化。测试函数的能量型闭包是 $H_0^1$，而非 $H^1$，故 Friedrichs 构造选 Dirichlet 扩张。
\end{solution}

\Needspace{10\baselineskip}
\section[自伴算子的谱、Rayleigh 原理和 Krylov--Weinstein 定理、谱的重数]{自伴算子的谱、Rayleigh 原理和\\Krylov--Weinstein 定理、谱的重数}\label{sec:yosida-11-08}
\noindent\textit{原书起始页：中译本 332；英文版 319。}

\begin{theorem}[谱、原子与残差]\label{thm:y11-spectral-tests}
设 $A=A^*$。实数 $\lambda\in\sigma(A)$ 当且仅当每个 $\varepsilon>0$ 都有 $E((\lambda-\varepsilon,\lambda+\varepsilon))\ne0$；并且
\[
 E(\{\lambda\})H=\ker(A-\lambda I),\qquad \sigma_r(A)=\varnothing.
\]
对每个单位向量 $x\in D(A)$ 及实数 $a$，
\[
 \operatorname{dist}(a,\sigma(A))\le\|(A-aI)x\|.
\]
\end{theorem}
\begin{proof}
若某邻域投影为零，$(t-\lambda)^{-1}$ 在测度意义下有界，谱积分给出有界逆；反之，任意小邻域有非零投影时，选其中单位向量 $x$，它属 $D(A)$ 且 $\|(A-\lambda)x\|\le\varepsilon$，与有界逆不相容。原子断言由平方范数积分直接得到。又 $\overline{R(A-\lambda)}^\perp=\ker(A-\lambda)$，故单射时值域必稠密，没有残余谱。最后对 $d=\operatorname{dist}(a,\sigma(A))$，有 $\|(A-a)x\|^2=\int|t-a|^2\,\dd\mu_x\ge d^2$。
\end{proof}

\begin{theorem}[Rayleigh 与 Krylov--Weinstein 估计]\label{thm:y11-rayleigh}
对自伴 $A$，以扩展实数理解，有
\[
 \inf\sigma(A)=\inf_{\|x\|=1,\,x\in D(A)}(Ax,x),\quad
 \sup\sigma(A)=\sup_{\|x\|=1,\,x\in D(A)}(Ax,x).
\]
若 $a=(Ax,x)$、$\|x\|=1$、$r=\|Ax-ax\|$，则 $\sigma(A)$ 中有一点距 $a$ 至多为 $r$；而且
\[
 r^2=\|Ax\|^2-a^2.
\]
\end{theorem}
\begin{proof}
Rayleigh 商是概率测度 $\mu_x$ 的均值，故受谱上下界约束。反向在谱边缘任意小的有界区间取非零投影中的单位向量，其均值就在该区间，逼近谱边缘；无界时取任意远的谱区间，得到扩展实数结论。残差恒等式由展开平方得到，距离界用上一项。实轴上的非空闭谱到固定点的距离可达到，故可把原书带任意正误差的区间估计加强为闭区间估计。
\end{proof}

残差只保证附近有谱，不能无条件保证附近有特征值。只有额外知道该区间的谱离散时，才能把它用于本征值误差估计。

\begin{definition}[简单谱与重数]
若存在 $x$ 使 $\overline{\operatorname{span}\{E(B)x:B\text{ 为 Borel 集}\}}=H$，称 $A$ 有简单谱，$x$ 为循环向量。对于纯点谱，特征值的重数为 $\dim E(\{\lambda\})H$。连续谱处的重数需用循环分解或直积分中的纤维维数描述，不能用单点投影的维数代替。
\end{definition}

例如 $M_t$ 作用于 $L^2(0,1)$ 有简单谱，但所有单点投影均为零；$M_t\oplus M_t$ 在每个频率有两个分量，因而不是简单谱。一般可分情形的测度类与重数函数分类不在本节完整证明，第 \ref{sec:yosida-11-14} 节证明循环情形。

\begin{exercise}[残差不保证特征值]\label{ex:y11-residual}
对 $A=M_t$ 在 $L^2(0,1)$ 上，取归一化指示函数 $x_\varepsilon$ 支于 $(a-\varepsilon,a+\varepsilon)\subset(0,1)$。计算 Rayleigh 商及残差。
\end{exercise}
\begin{solution}
均值为 $a$，方差为 $\varepsilon^2/3$，故残差为 $\varepsilon/\sqrt3$。谱为 $[0,1]$，却没有任何特征值，因为单点测度为零。
\end{solution}

\begin{exercise}[循环与重数二]\label{ex:y11-multiplicity}
证明 $M_t\oplus M_t$ 作用于 $L^2(0,1;\mathbb C^2)$ 不存在循环向量。
\end{exercise}
\begin{solution}
给定 $x=(f,g)$，若 $x=0$ 则显然不循环；否则 $y=(\bar g,-\bar f)$ 非零且属于同一空间。对每个标量 Borel 乘子 $h$，逐点有 $hf\,\overline{\bar g}+hg\,\overline{-\bar f}=0$，故 $y$ 与所有 $h x$ 正交，尤其与全部 $E(B)x$ 正交。
\end{solution}

\Needspace{10\baselineskip}
\section{一般展开定理、连续谱不出现的一个条件}\label{sec:yosida-11-09}
\noindent\textit{原书起始页：中译本 336；英文版 323。}

\begin{theorem}[一般展开与预解式反演]\label{thm:y11-stone-formula}
若 $A=A^*$，则对每个 $x\in H$，$x=\lim_{R\to\infty}E([-R,R])x$。更具体地，对有限 $a<b$，令 $Q(z)=(A-zI)^{-1}$，则
\[
 \lim_{\varepsilon\downarrow0}\frac1{2\pi i}
 \int_a^b\bigl(Q(t+i\varepsilon)-Q(t-i\varepsilon)\bigr)x\,\dd t
 =\bigl(E((a,b))+\tfrac12E(\{a\})+\tfrac12E(\{b\})\bigr)x.
\]
极限为强极限；积分端点如有原子，各计一半。
\end{theorem}
\begin{proof}
对固定 $\varepsilon>0$，有界预解式的 Bochner 积分可与谱积分交换，所得标量乘子为
\[
 k_\varepsilon(s)=\frac1\pi\int_a^b
 \frac{\varepsilon}{(s-t)^2+\varepsilon^2}\,\dd t
 =\frac1\pi\left(\arctan\frac{b-s}{\varepsilon}-\arctan\frac{a-s}{\varepsilon}\right).
\]
它在 $0$ 与 $1$ 之间，逐点趋于开区间指示函数加上两个端点的半质量。用谱积分平方范数公式和控制收敛得到强极限。一般展开来自 $\mu_x([-R,R])\to\mu_x(\R)=\|x\|^2$。
\end{proof}

注意本节 $Q(z)=(A-zI)^{-1}$ 与第八章 $R(z,A)=(zI-A)^{-1}$ 差一个负号；反演公式中的次序也随之固定。

\begin{theorem}[紧自伴展开与紧预解式]\label{thm:y11-compact-expansion}
紧自伴算子 $K$ 的空间可正交分解为其核与所有非零特征空间，且
\[
 x=P_0x+\sum_{\lambda\ne0}P_\lambda x,\qquad
 Kx=\sum_{\lambda\ne0}\lambda P_\lambda x.
\]
后一式按 $|\lambda|>\varepsilon$ 截断时算子范数误差至多 $\varepsilon$。若自伴 $A$ 的一个非实预解式紧，则 $A$ 有完备特征向量系，各特征空间有限维，有界实区间内仅有有限多个特征值；无限特征值序列只能逃向无穷远。
\end{theorem}
\begin{proof}
紧性和第十章说明非零谱离散、各特征空间有限维。其正交和的补空间约化 $K$，且限制算子的谱只能含零；否则非零谱投影会再给出特征向量。自伴性给出 $\|K|_{H_0}\|=r(K|_{H_0})=0$，所以补空间正是核。范数余项由谱乘子在 $|\lambda|\le\varepsilon$ 上至多为 $\varepsilon$ 得到。

令 $Q=(A-iI)^{-1}$，它是紧正规算子且单射。紧正规算子的同一结论由 $\|N\|=r(N)$ 和正交谱投影证明，其非零特征空间正交张成全空间，因为剩余部分范数为零且核为零。$Q$ 的特征值 $\mu\ne0$ 对应 $A$ 的特征值 $i+\mu^{-1}$，且 $Q$ 的特征向量在 $D(A)$ 中。紧性使非零 $\mu$ 只能聚于零，故 $A$ 的特征值只会逃向无穷远。若最初紧的是另一非实预解式，预解恒等式将紧性传到 $i$。
\end{proof}

\begin{example}[有界域上的 Dirichlet Laplace 算子]
第 \ref{sec:yosida-11-07} 节的 $A=-\Delta_D$ 满足 $(A+I)^{-1}:L^2(G)\to H_0^1(G)$ 有界；由 Rellich 紧嵌入，作为 $L^2$ 上算子它紧。因此 $A$ 有正交本征展开。无须在此预先证明 $H^2$ 边界正则性；有界域与零边界型空间已经给出紧性。
\end{example}

\begin{exercise}[端点半质量]\label{ex:y11-stone-endpoints}
若 $A=aI$，求反演公式在区间 $[a,b]$、$b>a$ 上的极限，解释为何不是 $I$。
\end{exercise}
\begin{solution}
谱为集中于 $a$ 的单原子，核乘子在该点趋于 $1/2$，故极限为 $I/2$。要恢复含端点全质量的投影，应另加端点的一半投影，或选避开原子的端点。
\end{solution}

\begin{exercise}[一维本征展开]\label{ex:y11-dirichlet-basis}
对 $(0,\pi)$ 上的 Dirichlet Laplace 算子，写出归一化本征函数、本征值和算子定义域的序列条件。
\end{exercise}
\begin{solution}
$e_n(x)=\sqrt{2/\pi}\sin(nx)$，$\lambda_n=n^2$。写 $u=\sum u_ne_n$，则 $D(A)=\{u:\sum n^4|u_n|^2<\infty\}$，$Au=\sum n^2u_ne_n$。型域只要求 $\sum n^2|u_n|^2<\infty$。
\end{solution}

\Needspace{10\baselineskip}
\section{Peter--Weyl--Neumann 理论}\label{sec:yosida-11-10}
\noindent\textit{原书起始页：中译本 339；英文版 326。}

本节完整证明紧 Hausdorff 群 $G$ 的 Peter--Weyl 定理。采用外部前提：$G$ 有唯一归一化 Haar 概率测度 $m$，它左右不变，且每个非空开集有正测度。原书还对全有界群和抽象群上的殆周期函数使用不变均值；那一推广在此仅作衔接。

\begin{definition}[有限维表示与矩阵系数]
连续酉表示是连续同态 $\pi:G\to U(V_\pi)$，$V_\pi$ 有限维。函数 $g\mapsto(\pi(g)v,w)$ 称矩阵系数，其有限线性组合构成代表函数代数 $\mathcal R(G)$。常数、乘积和共轭分别来自平凡、张量积和共轭表示。
\end{definition}

\begin{lemma}[Schur 正交关系]\label{lem:y11-schur}
每个有限维酉表示是不可约表示的正交直和。取不可约代表 $\pi,\rho$ 的正交基，有
\[
 \int_G\pi_{ij}(g)\overline{\rho_{kl}(g)}\,\dd m(g)
 =\begin{cases}d_\pi^{-1}\delta_{ik}\delta_{jl},&\pi=\rho,\\0,&\pi\not\cong\rho.\end{cases}
\]
\end{lemma}
\begin{proof}
不变子空间的正交补仍不变，按维数归纳得到完全可约性。若 $T$ 是两个不可约表示间的交织算子，其核与像不变，所以非零时为同构；同一不可约表示的交织算子取一个复特征值，$T-\lambda I$ 的核非零，故 $T=\lambda I$。这就是 Schur 引理。

给定线性映射 $B:V_\rho\to V_\pi$，平均 $\int\pi(g)B\rho(g)^{-1}\,\dd m(g)$ 为交织算子。不同不可约时它为零；相同时它为 $(\operatorname{tr}B/d_\pi)I$，因取迹不改变。对矩阵单位 $B$ 比较各矩阵元，便得公式。
\end{proof}

\begin{theorem}[Peter--Weyl]\label{thm:y11-peter-weyl}
$\mathcal R(G)$ 在 $C(G)$ 中一致稠密；各不可约表示的归一化系数 $\sqrt{d_\pi}\,\pi_{ij}$ 构成 $L^2(G,m)$ 的完备正交系。尤其有限维连续酉表示分离 $G$ 中的点。
\end{theorem}
\begin{proof}
先构造紧的平滑算子。取连续实函数 $k\ge0$，$k(g^{-1})=k(g)$，支于单位元的对称小邻域，且 $\int k=1$。定义
\[
 T_kf(x)=\int_G k(y^{-1}x)f(y)\,\dd m(y).
\]
核连续，故 $T_k:L^2(G)\to C(G)$ 有界：Cauchy--Schwarz 给出一致界，核的平移在 $L^2$ 中连续给出连续性。它在 $L^2$ 上为紧自伴算子，并与左平移 $L_af(x)=f(a^{-1}x)$ 交换。所有非零特征空间有限维、左平移不变，其中的函数连续。因此左平移在这些空间上给出有限维连续酉表示，而其中每个函数通过在单位元处取值可写成该表示或其对偶表示的矩阵系数。

若 $v\in L^2$ 与所有有限维系数正交，紧自伴展开使 $T_kv=0$。当单位元邻域缩小时，$T_kv\to v$ 于 $L^2$：先对连续函数用一致连续性，再由它们在 $L^2$ 中稠密及 $\|T_k\|\le1$ 推广。故 $v=0$，证明矩阵系数 $L^2$ 稠密。结合 Schur 关系即得完备正交系。

最后取 $f\in C(G)$，先选 $k$ 使 $\|T_kf-f\|_\infty<\varepsilon$，再选 $p\in\mathcal R(G)$ 使 $\|f-p\|_2<\varepsilon/\|T_k\|_{2\to\infty}$。卷积一个矩阵系数仍为同一有限维表示的系数线性组合，所以 $T_kp\in\mathcal R(G)$，且 $\|T_kp-f\|_\infty<2\varepsilon$。这证明一致稠密。若表示不分离两个点，所有代表函数均在两点取相同值，与连续函数分离点及一致稠密矛盾。
\end{proof}

不可分 $L^2(G)$ 的正交基可不可数，展开按有限子集的网理解。圆群结论只保证三角多项式一致稠密，不保证普通 Fourier 部分和一致收敛。

\begin{exercise}[紧交换群的不可约表示]\label{ex:y11-abelian-representations}
证明紧交换群的有限维不可约酉表示都是一维；对圆群写出这些表示。
\end{exercise}
\begin{solution}
交换性使每个 $\pi(g)$ 都是交织算子，Schur 引理使它为标量。若空间维数大于一，任意一维子空间都不变，违反不可约性。圆群的一维连续特征为 $e^{i\theta}\mapsto e^{in\theta}$，$n\in\mathbb Z$，由实轴特征及 $2\pi$ 周期性得到。
\end{solution}

\begin{exercise}[有限群的维数公式]\label{ex:y11-finite-group}
若 $G$ 有限，证明 $\sum_{\pi\in\widehat G}d_\pi^2=|G|$。
\end{exercise}
\begin{solution}
归一化计数测度下 $L^2(G)$ 的维数为 $|G|$；Peter--Weyl 的正交基每个不可约表示贡献 $d_\pi^2$ 个系数，数基向量便得等式。
\end{solution}

\Needspace{10\baselineskip}
\section{关于非交换紧群的 Tannaka 对偶性定理}\label{sec:yosida-11-11}
\noindent\textit{原书起始页：中译本 345；英文版 332。}

交换紧群可用一维特征描述，非交换紧群则需要所有有限维表示及其乘法关系。下面采用与原书同型的代表函数环表述，避免把“任意一组酉矩阵”误认为群的一个点。

\begin{theorem}[Tannaka 的代表函数表述]\label{thm:y11-tannaka}
设 $G$ 为紧 Hausdorff 群，$\mathcal R(G)$ 为代表函数代数。每个含单位、保持复共轭的复代数同态
\[
 \chi:\mathcal R(G)\longrightarrow\mathbb C
\]
都是唯一的点值泛函 $\chi(f)=f(g)$。在所有此类 $\chi$ 上赋逐点收敛拓扑及下述群运算，点值映射给出 $G$ 到该空间的拓扑群同构。这里没有预先假定 $\chi$ 对一致范数连续。
\end{theorem}
\begin{proof}
对每个不可约酉表示 $\pi$，矩阵 $U_\pi(\chi)=(\chi(\pi_{ij}))$ 是酉矩阵：直接把 $\pi(g)^*\pi(g)=I$ 的各系数恒等式送入保持共轭的 $\chi$。代表函数由这些系数线性张成，所以 $\chi$ 由所有 $U_\pi$ 决定。线性关系、乘法关系、共轭及单位关系均为闭条件，故所有这样的 $\chi$ 组成乘积 $\prod_\pi U(d_\pi)$ 的闭子集 $K$，从而紧。

代表函数有代数余乘法 $\Delta f(g,h)=f(gh)$，其值在 $\mathcal R(G)\otimes\mathcal R(G)$ 中，因为 $\Delta\pi_{ij}=\sum_k\pi_{ik}\otimes\pi_{kj}$。定义
\[
 (\chi\psi)(f)=(\chi\otimes\psi)(\Delta f),\qquad
 \chi^{-1}(f)=\chi(f\circ\mathrm{inv}).
\]
这些公式在系数上就是矩阵乘法与取逆，代数恒等式保证它们仍为含单位、保持共轭的同态。于是 $K$ 为紧拓扑群。点值映射 $G\to K$ 连续，Peter--Weyl 分离点使它单射；紧性使其像为闭子群。

将 $f\in\mathcal R(G)$ 变成 $\widetilde f(\chi)=\chi(f)$。这些函数在 $K$ 上含常数、分离点且对共轭封闭，所以一致稠密于 $C(K)$。每个 $\pi$ 的矩阵 $U_\pi$ 是 $K$ 上的连续酉表示；它限制到 $G$ 后是不可约 $\pi$，所以在 $K$ 上也不可约，且非平凡表示仍非平凡。对 $K$ 和 $G$ 分别使用 Haar 平均：平凡表示系数的平均为 $1$，其他不可约系数的平均均为零。因此两种平均对每个 $\widetilde f$ 相同，再由一致稠密性对所有 $C(K)$ 函数相同。

若 $G$ 的像是真闭子群，Urysohn 引理给出非负连续函数在其上为零、在像外某点为正。它对 $G$ 的积分为零，对 $K$ 的积分为正，因为非空开集 Haar 测度为正，矛盾。故 $G=K$，得到满射；唯一性及拓扑同构随分离点和紧 Hausdorff 性得到。
\end{proof}

等价地，群的一个点作用在每个有限维表示上，且与所有交织算子、直和、张量积、共轭及平凡表示相容。代表函数中的代数关系正编码这些条件。仅要求每个表示对应一个酉矩阵，没有这些相容条件，会得到远大于 $G$ 的乘积空间。

\begin{exercise}[圆群的重构]\label{ex:y11-tannaka-circle}
对 $G=\mathbb T$，用 Laurent 多项式代数证明上述重构。
\end{exercise}
\begin{solution}
$\mathcal R(\mathbb T)=\mathbb C[z,z^{-1}]$，共轭满足 $z^*=z^{-1}$。同态由 $w=\chi(z)$ 决定，保持共轭及 $zz^{-1}=1$ 给出 $|w|=1$，于是 $\chi(p)=p(w)$，恰是点值。
\end{solution}

\begin{exercise}[张量相容性]\label{ex:y11-tannaka-tensor}
圆群的第 $n$ 个特征记为 $\pi_n$。若为每个 $n$ 任取 $u_n\in\mathbb T$，写出来自同一群元素所需的条件。
\end{exercise}
\begin{solution}
需 $u_0=1$、$u_{m+n}=u_mu_n$、$u_{-n}=\overline{u_n}$。这些条件使 $u_n=u_1^n$，于是由唯一的 $u_1\in\mathbb T$ 决定。独立选择各个 $u_n$ 通常不满足乘法关系。
\end{solution}

\Needspace{20\baselineskip}
\section{自伴算子的函数}\label{sec:yosida-11-12}
\noindent\textit{原书起始页：中译本 351；英文版 338。}

令 $A=\int t\,\dd E(t)$。对有限值 Borel 函数 $f$，定义 $f(A)$ 为定理 \ref{thm:y11-spectral-integral} 中的闭算子。允许在 $E$-零集上修改函数。其定义域取决于向量的谱测度，而不仅取决于形式表达式。

\begin{theorem}[无界函数演算的运算规则]\label{thm:y11-borel-calculus}
设 $f,g$ 为有限值 Borel 函数。则 $f(A)^*=\bar f(A)$，并且
\begin{align*}
 D(g(A)f(A))&=D(f(A))\cap D((gf)(A)),\\
 g(A)f(A)x&=(gf)(A)x\quad(x\text{ 属于左侧定义域}),\\
 (f+g)(A)x&=f(A)x+g(A)x\quad(x\in D(f(A))\cap D(g(A))).
\end{align*}
$f(A)$ 正规；$f$ 实值时自伴。复合满足 $h(f(A))=(h\circ f)(A)$，其中复值 $f$ 时左侧用正规算子的 Borel 演算，并取对应的最大定义域。
\end{theorem}
\begin{proof}
对谱投影 $E(B)$ 和 $x\in D(f(A))$，截断后取极限得 $E(B)f(A)x=f(A)E(B)x$，且
\[
 \mu_{f(A)x}(B)=\int_B|f|^2\,\dd\mu_x.
\]
故 $f(A)x\in D(g(A))$ 等价于 $\int|gf|^2\,\dd\mu_x<\infty$，连同先验的 $x\in D(f(A))$ 得出定义域公式。等式对共同截断成立，再由平方积分极限推广。和的公式同理，但一般只能在定义域交上断言。

伴随等式已证。$f(A)^*f(A)$ 与 $f(A)f(A)^*$ 的定义域均为 $D(|f|^2(A))$，因为四次可积结合 $x\in H$ 蕴含二次可积；两者都等于 $|f|^2(A)$，故正规。$f(A)$ 的谱测度是 $E$ 经 $f$ 的推前，复合的定义域条件就是 $\int|h(f(t))|^2\,\dd\mu_x<\infty$，由截断给出复合公式。
\end{proof}

例如若 $f(t)=t$、$g(t)=1/t$（在零处另行有限定义），即使乘积在谱上等于 $1$，$g(A)f(A)$ 仍只在 $D(A)$ 上定义，未必是处处定义的恒等算子。它的闭包可以是 $I$，而算子本身不等于 $I$。

\begin{proposition}[与谱投影交换]\label{prop:y11-commutation}
有界算子 $B$ 若与自伴 $A$ 在定义域上交换，即 $BD(A)\subset D(A)$ 且 $ABx=BAx$，则与每个 $E(C)$ 交换；因此保持 $D(f(A))$ 并满足 $Bf(A)x=f(A)Bx$。
\end{proposition}
\begin{proof}
定义域上的关系先给出 $B(A\pm iI)^{-1}=(A\pm iI)^{-1}B$，从而与 Cayley 变换及其伴随交换。连续演算、再以有界 Borel 函数的强逼近得到 $BE(C)=E(C)B$。为处理无界 $f$，令 $P_n=E(\{|f|\le n\})$，则 $f(A)P_nBx=Bf(A)P_nx$，右侧收敛，左侧由闭性得到 $Bx\in D(f(A))$ 及交换等式。
\end{proof}

原书进一步证明 Neumann--Riesz--Mimura 的闭算子判别：在可分空间上，适当的交换子包含关系刻画 $T=f(A)$。这里证明正向的谱交换关系及循环模型中的乘法解释；一般闭算子判别的反向拼接证明尚未展开，不能把“与 $A$ 交换”本身当作“是 $A$ 的标量函数”的充分条件，重数大于一时尤其如此。

\begin{exercise}[定义域使恒等式失效]\label{ex:y11-product-domain}
在 $L^2(1,\infty)$ 上取 $A=M_t$。比较 $A^{-1}A$ 与 $AA^{-1}$ 的定义域。
\end{exercise}
\begin{solution}
$A^{-1}=M_{1/t}$ 有界且值域为 $D(A)$。所以 $AA^{-1}=I$ 处处定义；$A^{-1}A=I|_{D(A)}$，只在 $D(A)$ 定义。其闭包才是 $I$。
\end{solution}

\begin{exercise}[实函数与复函数]\label{ex:y11-functional-examples}
对 $A\ge0$ 定义 $A^{1/2}$ 与 $e^{-tA}$，证明前者自伴，后者在 $t\ge0$ 时为压缩算子且强连续。
\end{exercise}
\begin{solution}
使用乘子 $\sqrt\lambda$ 与 $e^{-t\lambda}$。前者实值，定义域要求 $\int\lambda\,\dd\mu_x<\infty$，故自伴。后者绝对值至多 $1$，给出压缩；$t\to s\ge0$ 时乘子逐点收敛且有界，控制收敛给出强连续。
\end{solution}

\Needspace{10\baselineskip}
\section{Stone 定理和 Bochner 定理}\label{sec:yosida-11-13}
\noindent\textit{原书起始页：中译本 358；英文版 345。}

\begin{theorem}[Stone 定理的谱积分形式]\label{thm:y11-stone}
若 $A=A^*$，则
\[
 U(t)=e^{itA}=\int_\R e^{it\lambda}\,\dd E(\lambda)
\]
为强连续酉群，生成元为 $iA$，包括定义域相等。每个强连续酉群都唯一以此形式表示。
\end{theorem}
\begin{proof}
乘子乘法和 $|e^{it\lambda}|=1$ 给出群律与酉性。控制收敛给出强连续。对 $x\in D(A)$，$|(e^{it\lambda}-1)/t|\le|\lambda|$，差商在平方积分中趋于 $i\lambda$，故生成元包含 $iA$。

反之，若差商在 $H$ 中收敛，则其范数有界。沿 $t_n\to0$ 用 Fatou 引理，得到 $\int\lambda^2\,\dd\mu_x<\infty$，所以 $x\in D(A)$，生成元正好为 $iA$。对任意强连续酉群，第九章定理 \ref{thm:y9-stone} 已在不使用谱定理的前提下证明生成元为某个 $iA$。本节积分群与它生成元相同，由生成元唯一性两群相同。
\end{proof}

\begin{definition}[正定函数]
函数 $\varphi:\R\to\mathbb C$ 称为正定，若对任意有限组 $t_j\in\R$、$c_j\in\mathbb C$，
\[
 \sum_{j,k}c_j\overline{c_k}\varphi(t_j-t_k)\ge0.
\]
该条件蕴含 $\varphi(0)\ge0$、$\varphi(-t)=\overline{\varphi(t)}$ 及 $|\varphi(t)|\le\varphi(0)$，可由一阶与二阶矩阵检验得到。
\end{definition}

\begin{theorem}[Bochner 表示]\label{thm:y11-bochner}
$\varphi$ 是连续正定函数，当且仅当存在唯一有限正 Borel 测度 $\mu$ 使
\[
 \varphi(t)=\int_\R e^{it\lambda}\,\dd\mu(\lambda),\qquad
 \mu(\R)=\varphi(0).
\]
\end{theorem}
\begin{proof}
若存在测度，正定性由积分 $|\sum_jc_je^{it_j\lambda}|^2$ 得到，连续性由控制收敛得到。

反向在有限形式和 $\sum c_j\delta_{t_j}$ 上定义半内积 $(\delta_s,\delta_t)=\varphi(s-t)$。正定性保证非负；Cauchy--Schwarz 使零范数向量构成可商掉的子空间。商后完备化为 Hilbert 空间 $K$，平移 $V(a)\delta_s=\delta_{s+a}$ 保内积，逆为 $V(-a)$。又
\[
 \|\delta_{s+a}-\delta_s\|^2=2\varphi(0)-2\operatorname{Re}\varphi(a)\longrightarrow0,
\]
故先在稠密有限和上、再由等距性在全空间上得到强连续酉群。记 $v=\delta_0$ 的商像，则 $\varphi(t)=(V(t)v,v)$。对 $V(t)$ 使用 Stone 谱积分，取 $\mu(B)=(E(B)v,v)$ 即得表示。

若两有限测度有同一 Fourier 变换，对任意 Schwartz 测试函数应用 Fourier 反演与 Fubini，得到它们对所有 $C_c^\infty$ 函数的积分相同；有限测度的正则性使它们相等。因此表示唯一。
\end{proof}

这里正定与逐点非负不同；概率测度对应额外条件 $\varphi(0)=1$。积分指数取正号，与本书函数 Fourier 变换的负号约定配套。

\begin{exercise}[含负值的正定函数]\label{ex:y11-positive-definite}
证明 $\cos t$ 正定，求其表示测度，并说明正定不意味着逐点非负。
\end{exercise}
\begin{solution}
$\cos t=(e^{it}+e^{-it})/2$，对应正测度 $(\delta_1+\delta_{-1})/2$，所以正定；但 $\cos\pi=-1$。
\end{solution}

\begin{exercise}[有界生成元与范数连续]\label{ex:y11-unitary-norm-continuity}
证明 $e^{itA}$ 在 $t=0$ 范数连续当且仅当自伴 $A$ 有界。
\end{exercise}
\begin{solution}
有界时指数级数给出范数连续。若 $A$ 无界，谱中存在 $|\lambda_n|\to\infty$，取 $t_n=\pi/|\lambda_n|$，则 $|e^{it_n\lambda_n}-1|=2$。连续演算的等距性或谱支集判别给出 $\|e^{it_nA}-I\|=2$，故不范数连续。
\end{solution}

\Needspace{10\baselineskip}
\section{具有简单谱的自伴算子的标准型}\label{sec:yosida-11-14}
\noindent\textit{原书起始页：中译本 360；英文版 347。}

\begin{theorem}[循环谱的乘法模型]\label{thm:y11-cyclic-model}
若自伴 $A$ 有循环向量 $v$，令 $\mu(B)=(E(B)v,v)$。则
\[
 W:L^2(\R,\mu)\longrightarrow H,\qquad Wf=f(A)v
\]
先对有界简单函数定义，再延拓为酉映射，并满足
\[
 W^{-1}AW=M_t,\qquad
 D(M_t)=\{f\in L^2(\mu):tf\in L^2(\mu)\}.
\]
谱投影变成乘 $1_B$。此处 $v$ 不必属于 $D(A)$。
\end{theorem}
\begin{proof}
对简单函数 $f,g$，正交投影关系给出 $(Wf,Wg)=\int f\bar g\,\dd\mu$。简单函数在 $L^2(\mu)$ 中稠密，故可等距延拓；其像闭，且包含所有 $E(B)v$，循环性使其满射。又 $E(B)Wf=W(1_Bf)$，从而 $Wf$ 的谱测度为 $|f|^2\mu$。因此 $Wf\in D(A)$ 恰等价于 $tf\in L^2(\mu)$；对谱截断先验证 $AWf=W(tf)$，再由闭性得全定义域上的等式。
\end{proof}

\begin{proposition}[简单谱的酉不变量]\label{prop:y11-measure-class}
两个有限正测度的坐标乘法算子 $M_t$ 具有循环向量 $1$，且它们酉等价当且仅当这两个测度相互绝对连续。因而谱集相同不足以确定自伴算子的酉等价类。
\end{proposition}
\begin{proof}
常数 $1$ 循环，因为简单函数稠密。若 $\mu\sim\nu$，以 $f\mapsto\sqrt{\dd\mu/\dd\nu}\,f$ 定义酉映射；它与乘 $t$ 交换，并保持对应定义域。反之，若 $U:L^2(\mu)\to L^2(\nu)$ 交织两个自伴乘法算子，它也交织谱投影。置 $h=U1$，则 $\mu(B)=\|U1_B\|^2=\int_B|h|^2\,\dd\nu$。因为 $h$ 循环，它不能在一个正 $\nu$ 测度集合上为零，故 $\mu\sim\nu$。
\end{proof}

例如 $[0,1]$ 上的 Lebesgue 测度与 $\nu=\sum_n2^{-n}\delta_{q_n}$（$q_n$ 枚举互异有理数）均有全支集 $[0,1]$。两个乘法算子的谱都为 $[0,1]$，但前者无点谱、后者为纯点谱，不酉等价。

\begin{exercise}[换循环向量]\label{ex:y11-change-cyclic}
在 $L^2(\mu)$ 的坐标乘法模型中，证明 $h$ 为循环向量当且仅当 $h\ne0$ 几乎处处。
\end{exercise}
\begin{solution}
若 $h$ 在正测度集上为零，所有 $1_Bh$ 也为零，无法稠密。反之，对任意 $f\in L^2(\mu)$，令 $f_n=f1_{\{|f|\le n,\,1/n\le|h|\le n\}}$。则 $f_n/h$ 有界，故 $f_n$ 属于有界 Borel 乘子作用于 $h$ 的闭包；$f_n\to f$，故循环。
\end{solution}

\begin{exercise}[离散简单谱]\label{ex:y11-discrete-cyclic}
对 $Ae_n=ne_n$ 作用于 $\ell^2$，给出一个循环向量及对应有限谱测度。
\end{exercise}
\begin{solution}
取 $v=\sum_{n\ge1}2^{-n}e_n$。每个 $E(\{n\})v=2^{-n}e_n$，故循环；对应 $\mu=\sum_{n\ge1}4^{-n}\delta_n$。此例还说明循环向量不要求范数为 $1$；测度总质量为 $\|v\|^2$。
\end{solution}

\Needspace{10\baselineskip}
\section{对称算子的亏指数、广义单位分解}\label{sec:yosida-11-15}
\noindent\textit{原书起始页：中译本 362；英文版 349。}

设 $S$ 闭、稠密定义且对称，沿用第七章
\[
 N_+=\ker(S^*-iI),\quad N_-=\ker(S^*+iI),\quad
 n_\pm=\dim N_\pm.
\]
其 Cayley 变换从 $M_+=R(S+iI)$ 等距映到 $M_-=R(S-iI)$，且 $M_\pm^\perp=N_\pm$。亏指数是 Hilbert 维数，允许无穷；取相反 Cayley 约定会交换两者，必须同时改动符号。

\begin{theorem}[自伴扩张与最大对称性]\label{thm:y11-extensions}
$S$ 在同一 $H$ 中有自伴扩张当且仅当 $n_+=n_-$。其自伴扩张与所有酉映射 $V:N_+\to N_-$ 一一对应，定义为
\begin{align*}
 D(S_V)&=D(S)\mathbin{\dot+}\{h+Vh:h\in N_+\},\\
 S_V(u+h+Vh)&=Su+ih-iVh.
\end{align*}
$S$ 为最大对称算子当且仅当 $n_+=0$ 或 $n_-=0$；最大对称未必自伴。$S$ 自伴当且仅当两亏指数皆为零。
\end{theorem}
\begin{proof}
任何对称扩张对应 Cayley 等距映射的扩张。它为自伴恰对应定义域与值域均为 $H$ 的酉扩张，所以须以酉映射配对缺少的两个正交补；这恰等价于亏指数相等。任一补空间的酉配对都给出酉扩张 $\widetilde U$。因为 $R(I-\widetilde U)$ 包含稠密的 $D(S)$，$\widetilde U$ 不含不动向量，逆 Cayley 公式有效。将补空间配对取为 $-V$，代入 $(I-\widetilde U)h=h+Vh$、$i(I+\widetilde U)h=ih-iVh$，得到所写定义域及作用式。反向每个自伴扩张都如此取得。

若两个缺陷空间均非零，可在各取一维等距配对，扩张原 Cayley 映射，得到严格的对称扩张。若其中一个为零，则原等距映射的定义域或值域已经为全空间，无法再扩大定义域而保持等距性，故最大对称。最后一项就是第七章自伴判据。
\end{proof}

\begin{theorem}[扩空间与广义谱测度]\label{thm:y11-generalized-resolution}
每个闭稠密对称算子 $S$ 都有某个更大 Hilbert 空间 $\widetilde H\supset H$ 中的自伴扩张 $\widetilde A$。若 $P_H$ 为正交投影，则
\[
 F(B)=P_HE_{\widetilde A}(B)|_H
\]
是正算子值测度，$F(\R)=I_H$，并且对 $x\in D(S)$，
\[
 Sx=\int t\,\dd F(t)x,\qquad
 \int t^2\,\dd(F(t)x,x)=\|Sx\|^2.
\]
一般 $F(B)$ 不是投影；二阶积分有限可能只保证 $x\in D(\widetilde A)\cap H$，未必 $x\in D(S)$。若 $S$ 最大对称，则该条件等价于 $x\in D(S)$。
\end{theorem}
\begin{proof}
取 $\widetilde H=H\oplus H$，在其上考虑 $S\oplus(-S)$。亏指数均为 $n_++n_-$，故有自伴扩张 $\widetilde A$。把原 $H$ 视为第一分量便得到所需扩张。压缩谱测度保持正性、单位质量和强可加性。对 $x\in D(S)$，$\widetilde Ax=Sx\in H$，所以压缩谱积分给出作用式，二阶矩等于 $\|\widetilde Ax\|^2$。

对于任意 $x\in H$，二阶矩有限等价于 $x\in D(\widetilde A)\cap H$。算子 $T=P_H\widetilde A|_{D(\widetilde A)\cap H}$ 对称且扩张 $S$；若 $S$ 最大对称，则 $T=S$，从而定义域相等。一般情形没有这个最大性推论。
\end{proof}

\begin{example}[区间与半直线上的动量]
对 $S=-i\,\dd/\dd x$，$D(S)=H_0^1(0,1)$，伴随定义域为 $H^1(0,1)$。解 $S^*u=\pm iu$ 得 $e^{-x}$ 与 $e^x$，均平方可积，所以亏指数为 $(1,1)$。自伴扩张可写成边界条件 $u(1)=e^{i\theta}u(0)$，本征值为 $\theta+2\pi n$。

在半直线 $(0,\infty)$ 上，$e^{-x}$ 可积而 $e^x$ 不可积，故同一符号约定下亏指数为 $(1,0)$，算子最大对称但无同空间自伴扩张。可把它扩到整直线的自伴动量算子，再压缩其谱测度。
\end{example}

\begin{exercise}[压缩未必为投影]\label{ex:y11-compressed-pvm}
在 $\mathbb C^2$ 上取 $\widetilde A=\operatorname{diag}(-1,1)$，把 $H$ 取为 $(1,1)$ 张成的直线。求 $F(\{1\})$ 并检验幂等性。
\end{exercise}
\begin{solution}
以 $(1,1)/\sqrt2$ 为 $H$ 的单位基，$F(\{1\})=\tfrac12I_H$，其平方为 $\tfrac14I_H$，不是投影。这说明广义单位分解应理解为正算子值测度。
\end{solution}

\begin{exercise}[区间边界相位]\label{ex:y11-boundary-phase}
对 $u(1)=e^{i\theta}u(0)$ 的动量算子写出归一化本征函数，并解释不同 $\theta$ 给出不同扩张。
\end{exercise}
\begin{solution}
$u_n(x)=e^{i(\theta+2\pi n)x}$，范数为 $1$。乘 $e^{i\theta x}$ 将周期 Fourier 基变成该基，故完备。不同模 $2\pi$ 的相位给出不同边界定义域及不同本征值序列；它们都包含 $H_0^1$ 上的最小算子。
\end{solution}

\Needspace{10\baselineskip}
\section{群环 \texorpdfstring{$L^1$}{L1} 及 Wiener 的 Tauber 定理}\label{sec:yosida-11-16}
\noindent\textit{原书起始页：中译本 367；英文版 354。}

在 $L^1(\R)$ 上取卷积乘法和 $f^*(x)=\overline{f(-x)}$。Fubini 给出 $\|f*g\|_1\le\|f\|_1\|g\|_1$，所以这是交换 Banach 代数，但没有卷积单位：形式上的单位 $\delta_0$ 不属于 $L^1$。以下用 $\widetilde{L^1}=\mathbb C1\oplus L^1$ 表示单位化。

\begin{theorem}[群环的特征与半单性]\label{thm:y11-group-characters}
$L^1(\R)$ 的所有非零特征为
\[
 \chi_\xi(f)=\widehat f(\xi)=\int_\R e^{-it\xi}f(t)\,\dd t,\qquad \xi\in\R.
\]
单位化的特征另外包含 $\chi_\infty(\lambda1+f)=\lambda$，其他特征为 $\lambda+\widehat f(\xi)$。该代数半单。
\end{theorem}
\begin{proof}
非零特征延为单位化的含单位特征，故自动连续。选 $g$ 使 $\chi(g)\ne0$，令 $\tau_tg(x)=g(x-t)$，$c(t)=\chi(\tau_tg)/\chi(g)$。恒等式 $(\tau_tg)*h=g*(\tau_th)$ 说明此比值与所选非零特征值函数无关，且 $\chi(\tau_th)=c(t)\chi(h)$ 对所有 $h$ 成立。于是 $c(t+s)=c(t)c(s)$，$c(0)=1$。平移的 $L^1$ 连续性给出 $c$ 连续，等距性给出全轴有界。群律使 $|c(t)|=1$，故 $c(t)=e^{-it\xi}$，其中实数 $\xi$ 唯一。

Bochner 积分恒等式 $f*g=\int f(t)\tau_tg\,\dd t$ 使
\[
 \chi(f)\chi(g)=\chi(g)\int f(t)c(t)\,\dd t.
\]
消去 $\chi(g)$ 得 Fourier 公式。反向由 Fubini 验证这些都是特征；单位化中限制到 $L^1$ 为零时只剩 $\chi_\infty$。若所有特征在 $\lambda1+f$ 上为零，则先由 $\chi_\infty$ 得 $\lambda=0$，再由 Fourier 变换的单射性得 $f=0$，故根为零。
\end{proof}

\begin{theorem}[Wiener 平移稠密定理]\label{thm:y11-wiener}
对 $f\in L^1(\R)$，其所有平移的线性张成在 $L^1(\R)$ 中稠密，当且仅当 $\widehat f(\xi)\ne0$ 对所有实数 $\xi$ 成立。
\end{theorem}
\begin{proof}
必要性：若 $\widehat f(\xi_0)=0$，连续非零泛函 $\chi_{\xi_0}$ 消灭所有平移及其闭包，故不稠密。

充分性：先证明闭卷积理想 $I=\overline{f*L^1}$ 为全空间。取 $\varphi$ 满足 $\widehat\varphi\in C_c^\infty(\R)$。选 $h$ 使 $\widehat h\in C_c^\infty$、$0\le\widehat h\le1$，且在 $\operatorname{supp}\widehat\varphi$ 上恒为 $1$；由 Fourier 反演 $h\in\Sch\subset L^1$。
\par\Needspace{4\baselineskip}
在单位化代数中令
\[
 d=1-h+f*f^*.
\]
每个有限特征给出 $1-\widehat h(\xi)+|\widehat f(\xi)|^2>0$，而无穷特征值为 $1$。极大理想判据使 $d$ 可逆。置 $g=d^{-1}*(f^* *\varphi)\in L^1$。在 $\widehat\varphi$ 的支集上 $\widehat h=1$，所以 $\widehat{f*g}=\widehat\varphi$；支集外两者均为零。Fourier 单射性给出 $f*g=\varphi$，故所有这类 $\varphi$ 均在 $I$ 中。

这些 $\varphi$ 在 $L^1$ 中稠密。具体地，先把任意 $L^1$ 函数逼近为 Schwartz 函数 $v$，再取 $\eta\in\Sch$ 满足 $\widehat\eta\in C_c^\infty$、$\widehat\eta(0)=1$。$\eta_\varepsilon(x)=\varepsilon^{-1}\eta(x/\varepsilon)$ 为有界 $L^1$ 近似单位，故 $v*\eta_\varepsilon\to v$ 于 $L^1$，而其 Fourier 变换光滑且紧支。因此 $I=L^1$。

最后 $f*g$ 是平移 $f$ 的 Bochner 加权积分，所以属于平移张成的闭包；可由简单函数逼近积分验证。反向每个 $\tau_af$ 都由 $f*(\tau_a\eta_\varepsilon)$ 逼近。因此闭理想恰为平移张成的闭包，结论成立。
\end{proof}

这里使用的是 $f$ 的平移稠密性，不是 $f$ 在 $L^1$ 卷积代数中可逆。事实上 $\widehat f(\xi)\to0$ 于无穷远，不能把它的倒数当作某个 $L^1$ 函数的 Fourier 变换。

\begin{corollary}[有界函数的 Tauber 传递]\label{cor:y11-tauber-transfer}
若 $h\in L^\infty(\R)$、$\widehat f$ 处处非零，并且
\[
 (f*h)(x)\longrightarrow C\int f\quad(x\to+\infty),
\]
则对每个 $g\in L^1$，$(g*h)(x)\to C\int g$。
\end{corollary}
\begin{proof}
结论对 $f$ 的平移成立，因为 $(\tau_af*h)(x)=(f*h)(x-a)$，故对它们的有限线性组合成立。用 Wiener 稠密定理逼近 $g$，并估计 $\|(g-p)*h\|_\infty\le\|g-p\|_1\|h\|_\infty$、$|\int(g-p)|\le\|g-p\|_1$，即可先控制逼近误差，再令 $x\to\infty$。
\end{proof}

原书末尾还讨论 Littlewood 的 Abel--Tauber 推论。本节未展开其证明；通常级数收敛还需 $na_n$ 有界等条件，不能仅凭 Abel 可和性推出。

\begin{exercise}[高斯平移]\label{ex:y11-gaussian-translates}
证明 $f(x)=e^{-x^2}$ 的有限平移线性组合在 $L^1(\R)$ 中稠密。
\end{exercise}
\begin{solution}
由本书 Fourier 约定，$\widehat f(\xi)=\sqrt\pi e^{-\xi^2/4}$，处处非零，因此应用 Wiener 定理。结论允许复系数的有限线性组合，不要求正系数。
\end{solution}

\begin{exercise}[一个零点就足够阻碍稠密]\label{ex:y11-zero-obstruction}
取 $f=1_{[0,1]}$。证明其平移张成不在 $L^1(\R)$ 中稠密，并写出一个消灭该闭包的非零连续泛函。
\end{exercise}
\begin{solution}
$\widehat f(2\pi)=\int_0^1e^{-2\pi ix}\,\dd x=0$。泛函 $g\mapsto\int e^{-2\pi ix}g(x)\,\dd x$ 非零且有界，并在每个平移 $f$ 上取零，故其核包含整个闭包。
\end{solution}


\Needspace{6\baselineskip}
\section*{本章参考说明与阅读衔接}
\phantomsection\addcontentsline{toc}{section}{本章参考说明与阅读衔接}
直接来源为吉田第十一章 \cite{yosida-zh,yosida-en}。本章采用循环子空间构造谱测度，以紧群卷积和 Fourier 局部逆组织后半章证明；与原书的范畴集、殆周期均值及赋范环路线作适用范围对照。原书所引 Gelfand、Neumann、Friedrichs、Tannaka、Wiener 等原始文献未计为直接阅读。闭算子的可测函数判别一般版本、非可分谱重数分类及 Littlewood 的最终 Tauber 推论等未展开范围见 COVERAGE。下一章转向凸集端点、向量格及其表示。

% !TeX encoding = UTF-8
% 吉田第六版：中译本第 12 章；此文件当前仅为目录与正文占位。
\SetChapterInfo{subtitle={吉田耕作第六版主线 · 正文待撰写}}
\chapter{线性空间中其他表示定理}\label{ch:yosida-12}
\YosidaChapterPlaceholder{12}{Other Representation Theorems in Linear Spaces}{377--396}{362--378}

\section{端点、Krein--Milman 定理}\label{sec:yosida-12-01}
\YosidaSectionPlaceholder{377}{362}
% TODO: 定义与记号；定理与完整证明；例子；习题；中英文核对。

\section{向量格}\label{sec:yosida-12-02}
\YosidaSectionPlaceholder{379}{364}
% TODO: 定义与记号；定理与完整证明；例子；习题；中英文核对。

\section{B-格和 F-格}\label{sec:yosida-12-03}
\YosidaSectionPlaceholder{385}{369}
% TODO: 定义与记号；定理与完整证明；例子；习题；中英文核对。

\section{Banach 收敛定理}\label{sec:yosida-12-04}
\YosidaSectionPlaceholder{386}{370}
% TODO: 定义与记号；定理与完整证明；例子；习题；中英文核对。

\section{向量格的点函数表示}\label{sec:yosida-12-05}
\YosidaSectionPlaceholder{388}{372}
% TODO: 定义与记号；定理与完整证明；例子；习题；中英文核对。

\section{向量格的集合函数表示}\label{sec:yosida-12-06}
\YosidaSectionPlaceholder{391}{375}
% TODO: 定义与记号；定理与完整证明；例子；习题；中英文核对。

% !TeX encoding = UTF-8
\SetChapterInfo{subtitle={从时间平均到扩散生成元与抽象位势},keywords={Markov; 遍历; 熵; Brownian; Feller; 跳跃; 位势},prereq={平均遍历; 正算子; 强连续半群; 向量格; 测度表示},goals={区分平均与逐点收敛; 解释局部与跳跃生成元; 控制位势定义域}}
\chapter[遍历理论和扩散理论]{遍历理论\\和扩散理论}\label{ch:yosida-13}
\ChapterSource{按 \cite{yosida-zh,yosida-en} 第十三章九节组织，2026 年 10 月 27 日补写。直接查看英文印刷 379--418 页全部扫描图，第 418 页上部仍为本章结尾、下部开始第十四章，故修正本章英文范围。中译本范围为 397--436 页，选页与证明边界详见本章 COVERAGE。}

遍历理论研究长时间平均，扩散理论研究短时间变化；半群把两端联系起来。不变测度产生各 $L^p$ 空间上的收缩，极大不等式使平均收敛提升为几乎处处收敛；局部矩与跳跃测度决定生成式，无折扣预解式极限则产生位势与 Poisson 方程。

遍历部分先用概率空间与固定一步算子 $T=T_1$；连续时间结论另列假设。记 $A_n=n^{-1}(I+T+\cdots+T^{n-1})$，它与生成元 $A$ 的记号按上下文区分。扩散部分采用 $b\cdot\nabla+\tfrac12Q:D^2$，热核模型的生成元为 $c\Delta$。抽象位势统一取 $V=-A^{-1}$，以 $\lambda\downarrow0$ 的 Abel 极限定义。

\Needspace{10\baselineskip}
\section{具有不变测度的 Markov 过程}\label{sec:yosida-13-01}
\noindent\textit{原书起始页：中译本 397；英文版 379。}

\begin{definition}[转移核]
可测空间 $(S,\Sigma)$ 上的概率转移核 $P$ 满足：$P(x,\cdot)$ 为概率测度，$P(\cdot,B)$ 可测。时间齐次转移族满足 $P_0(x,\cdot)=\delta_x$ 以及 Chapman--Kolmogorov 方程
\[
 P_{t+s}(x,B)=\int_S P_s(y,B)P_t(x,\dd y).
\]
\end{definition}

测度 $m$ 不变，指 $\int P_t(x,B)\,m(\dd x)=m(B)$。有限非零不变测度可归一化，本章遍历部分约定 $m(S)=1$。记
\[
 (T_tf)(x)=\int_S f(y)P_t(x,\dd y).
\]

本章主要研究转移核与其算子。转移族的存在不自动等于已构造出具有连续样本路径的随机过程；样本路径构造所需的概率空间与正则性应另外说明。

\begin{proposition}[Markov 收缩]\label{prop:y13-markov-contraction}
若 $m$ 不变，则 $T_t$ 在每个 $L^p(m)$（$1\le p\le\infty$）上诱导正收缩，满足 $T_t1=1$、$\int T_tf\,\dd m=\int f\,\dd m$ 以及半群性质。
\end{proposition}
\begin{proof}
对 $m$ 零集 $N$，$\int P_t(x,N)\,\dd m=0$，故 $P_t(x,N)=0$ 几乎处处；算子因而与代表元无关。对有界函数由核积分得到正性、单位与保积分；单调收敛延伸到非负可积函数。Jensen 不等式给出 $|T_tf|^p\le T_t|f|^p$，积分并用不变性即得有限 $p$ 的估计。$p=\infty$ 用概率加权平均估计。截断及完备性将算子延拓到各 $L^p$；Fubini 与 Chapman--Kolmogorov 方程给出半群性质。
\end{proof}

\begin{theorem}[有限不变测度下的平均遍历]\label{thm:y13-markov-mean}
固定一步转移 $T=T_1$，置 $A_n=n^{-1}\sum_{k=0}^{n-1}T^k$。对每个 $f\in L^p(m)$、$1\le p<\infty$，$A_nf$ 在 $L^p$ 中收敛到 $Qf$，满足
\[
 TQf=Qf,\quad Q^2=Q,\quad Q1=1,\quad
 \int Qf\,\dd m=\int f\,\dd m.
\]
$Q$ 为正收缩投影。若 $T_t$ 在 $L^p$ 上强连续，相同结论适用于 $t^{-1}\int_0^tT_sf\,\dd s$，其极限投到全半群的共同不动空间。
\end{theorem}
\begin{proof}
在 $L^2$ 上应用第八章的 Hilbert 空间平均遍历定理。对 $L^1$，有界函数在其中稠密，且 $\|g\|_1\le\|g\|_2$。先对有界函数得到 $L^1$ 收敛，再用 $\sup_n\|A_n\|_1\le1$ 延拓到全部 $L^1$。对有界 $g$，$A_ng$ 一致有界，故其 $L^1$ 收敛通过 $|h|^p\le C^{p-1}|h|$ 提升到任意有限 $p$；再次使用有界函数的稠密性完成证明。

$TA_n-A_n=(T^n-I)/n$，所以极限不动；不动元素被每个平均保持，故 $Q^2=Q$。其余性质由极限得到。连续时间在 $L^2$ 上使用连续平均遍历，再作同样稠密延拓；端段积分估计给出对所有 $T_s$ 的不变性。
\end{proof}

\begin{example}[确定性过程与有限状态链]
若可测变换 $\tau$ 保持 $m$，则 $P(x,B)=\mathbf1_B(\tau x)$，$Tf=f\circ\tau$，是 Markov 核的特殊情形。有限状态时 $P=(p_{ij})$ 为行随机矩阵，不变概率是行向量 $\pi P=\pi$；算子在函数列向量上左乘 $P$，不要把两个方向混用。
\end{example}

\begin{exercise}[周期链]\label{ex:y13-periodic-chain}
对 $P=\left(\begin{smallmatrix}0&1\\1&0\end{smallmatrix}\right)$ 求不变概率和 $A_nf$ 的极限，并说明 $P^nf$ 未必收敛。
\end{exercise}
\begin{solution}
不变概率为 $(1/2,1/2)$。常数方向特征值为 $1$，$(1,-1)$ 方向特征值为 $-1$，故 $A_n(f_1,f_2)\to((f_1+f_2)/2,(f_1+f_2)/2)$。在第二个方向，$P^n$ 交替变号。平均收敛与各时刻的分布收敛是不同命题。
\end{solution}

\begin{exercise}[不变性不可省略]\label{ex:y13-noninvariant}
在两点空间取均匀概率，令两行转移概率均为 $(1,0)$。证明该算子保持常数，却不是 $L^1$ 收缩。
\end{exercise}
\begin{solution}
$T(a,b)=(a,a)$，故 $T1=1$。取 $f=(1,0)$，则 $\|f\|_1=1/2$ 而 $\|Tf\|_1=1$。均匀概率不满足 $\pi P=\pi$，不能从 Markov 性单独推得相对于给定测度的 $L^1$ 收缩。
\end{solution}

\Needspace{10\baselineskip}
\section{个别遍历定理及其应用}\label{sec:yosida-13-02}
\noindent\textit{原书起始页：中译本 401；英文版 383。}

本节 $T$ 为上一节的 Markov 收缩，$m(S)=1$，$A_n$ 从 $k=0$ 起求和。将原书从 $k=1$ 起的记法改到此处约定，不改变下述极限。证明采用有限极大不等式；它把稠密子空间上的逐点收敛延拓到 $L^1$，避免误把稠密集视为非第一纲集。

\begin{lemma}[Hopf 有限极大不等式]\label{lem:y13-hopf-maximal}
对实 $f\in L^1$，令 $S_kf=\sum_{j=0}^{k-1}T^jf$，$E_N=\{\max_{1\le k\le N}S_kf>0\}$。则 $\int_{E_N}f\,\dd m\ge0$。因而对任意 $f\in L^1$ 和 $\alpha>0$，
\begin{equation}\label{eq:y13-maximal-weak}
 m\{\sup_{n\ge1}|A_nf|>\alpha\}\le\|f\|_1/\alpha.
\end{equation}
\end{lemma}
\Needspace{4\baselineskip}
\begin{proof}
设 $g_N=\max(0,S_1f,\ldots,S_Nf)$、$g_0=0$。由于 $T$ 正，$S_kf=f+TS_{k-1}f\le f+Tg_{N-1}$。在 $E_N$ 上取最大值得 $g_N\le f+Tg_{N-1}$，故
\[
 \int_{E_N}f\ge\int g_N-\int_{E_N}Tg_{N-1}
 \ge\int g_N-\int g_{N-1}\ge0.
\]
对 $|f|-\alpha$ 应用此式；$T1=1$ 使正部分和事件恰为 $\max_{k\le N}A_k|f|>\alpha$。于是 $\alpha m(E_N)\le\|f\|_1$。用 $|A_nf|\le A_n|f|$ 并令 $N\to\infty$，得到所需估计。
\end{proof}

\begin{theorem}[Markov 算子的个别遍历定理]\label{thm:y13-individual}
对每个 $f\in L^1(m)$，$A_nf$ 几乎处处收敛到有限函数 $Qf$，且在 $L^1$ 中收敛。对 $f\in L^p$、$1\le p<\infty$，极限也与上一节的 $L^p$ 极限相同。
\end{theorem}
\begin{proof}
在 Hilbert 空间 $L^2$ 中，收缩性给出 $\ker(I-T)=\ker(I-T^*)$，因此
\[
 L^2=\ker(I-T)\oplus\overline{(I-T)L^2}.
\]
为说明第一个等式，若 $Th=h$，则 $\|T^*h\|\le\|h\|$ 且 $\langle T^*h,h\rangle=\|h\|^2$，Cauchy--Schwarz 的等号情形给出 $T^*h=h$；反向同理。由于有界函数在 $L^2$ 中稠密，
\[
 \mathcal D=\ker(I-T)|_{L^2}+(I-T)L^\infty
\]
在 $L^2$ 及 $L^1$ 中稠密。对不动函数，平均本来不变；对 $g=(I-T)h$、$h\in L^\infty$，有 $A_ng=(h-T^nh)/n\to0$ 几乎处处，甚至本质一致。因此 $\mathcal D$ 中每个元素都有有限逐点极限。

先取实 $f\in L^1$。极大估计在 $\alpha\to\infty$ 时说明 $\sup_n|A_nf|<\infty$ 几乎处处。其上、下极限之差记为 $\omega(f)$。任取 $g\in\mathcal D$，由于 $A_ng$ 收敛，
\[
 \omega(f)\le2\sup_n A_n|f-g|,
 \qquad m\{\omega(f)>2\alpha\}\le\|f-g\|_1/\alpha.
\]
令 $g\to f$，得到 $\omega(f)=0$ 几乎处处。复函数分别处理实虚部。上一节已有 $L^1$ 极限；抽取几乎处处收敛的子序列后，极限的唯一性识别该函数为 $Qf$。有限 $p$ 情形亦然。
\end{proof}

这与第十二章的 Banach 收敛思想一致：关键是先证明控制收敛集闭性的估计。仅在一个稠密子空间验证收敛并不足够。对连续时间，本节结论先用于一步平均 $g=\int_0^1T_sf\,\dd s$ 得整数时刻的积分平均；一般几乎处处连续时间版本还需共同可测代表与余项控制，此处不据此省略这些条件。

\Needspace{7\baselineskip}
\paragraph{原书的比值推广。}
Chacon--Ornstein 比值遍历定理对一般正 $L^1$ 收缩 $T$ 断言：若 $f,p\in L^1$、$p\ge0$，则
\[
 \frac{\sum_{k=0}^{n-1}T^kf}{\sum_{k=0}^{n-1}T^kp}
\]
在 $\{\sum_{k\ge0}T^kp>0\}$ 上几乎处处有有限极限；分母最终为正后才定义比值。本节完整证明的是上述保概率 Markov 情形；一般比值定理及原书的消去引理作为明确标注的外部推广保留，不将极大不等式的证明说成已证明整个比值理论。

\begin{exercise}[上同调差的平均]\label{ex:y13-coboundary}
若 $g=(I-T)h$ 且 $h\in L^\infty$，求 $A_ng$ 的显式误差界。说明同一公式在 $h\in L^1$ 时只立即给出哪种收敛。
\end{exercise}
\begin{solution}
$A_ng=(h-T^nh)/n$，故 $\|A_ng\|_\infty\le2\|h\|_\infty/n$。若只有 $h\in L^1$，同样得到 $\|A_ng\|_1\le2\|h\|_1/n$，但该范数估计本身不推出全序列几乎处处收敛；本节定理另给出后者。
\end{solution}

\begin{exercise}[保持变换的时间平均]\label{ex:y13-birkhoff}
若 $\tau$ 保持概率测度且 $f\in L^1$，把本节定理写成轨道函数公式，并说明例外集对一个给定 $f$ 的含义。
\end{exercise}
\begin{solution}
取 $Tf=f\circ\tau$，则 $n^{-1}\sum_{k=0}^{n-1}f(\tau^kx)\to Qf(x)$ 几乎处处，并且积分保持。可为这个 $f$ 选一个零测例外集；可数个给定函数可以统一例外集。不能在任意不可数函数族上直接取零测集的并，声称得到所有可积函数同时逐点收敛。
\end{solution}

\Needspace{10\baselineskip}
\section{遍历假设和 \texorpdfstring{$H$}{H}-定理}\label{sec:yosida-13-03}
\noindent\textit{原书起始页：中译本 407；英文版 389。}

\begin{definition}[遍历与混合]
保持概率 $m$ 的 Markov 算子称为遍历的，若每个 $f\in L^1$ 的时间平均均趋于常数 $\int f\,\dd m$。集合 $B$ 称为不变集，若 $T\mathbf1_B=\mathbf1_B$ 几乎处处。不变集构成的完备化 $\sigma$ 代数记为 $\mathcal I$。若对每对 $f,g\in L^2$ 都有
\[
 \langle T^nf,g\rangle\longrightarrow
 \left(\int f\,\dd m\right)\overline{\left(\int g\,\dd m\right)},
\]
则称一步算子混合；连续半群用 $t\to\infty$ 作相应定义。
\end{definition}

\begin{theorem}[遍历性的等价表述]\label{thm:y13-ergodic-equivalence}
下列条件等价：遍历性；$\ker(I-T)|_{L^2}$ 只有常数；每个不变集的概率为 $0$ 或 $1$；对任意可测集 $B,C$，
\[
 \frac1n\sum_{k=0}^{n-1}\int_B P^k(x,C)\,m(\dd x)
 \longrightarrow m(B)m(C).
\]
一般情形下，平均投影是条件期望 $Qf=\mathbb E_m(f\mid\mathcal I)$。
\end{theorem}
\begin{proof}
前两条件由平均投影的值域及 $L^2$ 在 $L^1$ 中稠密得到。若所有不动函数为常数，不变集的示性函数只能为零或一。反向先考虑实不动函数 $f\in L^2$。Jensen 及不变性给出
\[
 T(f^2)-(Tf)^2\ge0,\qquad \int[T(f^2)-f^2],\dd m=0.
\]
故对几乎每个 $x$，$P(x,\cdot)$ 集中于 $f(y)=f(x)$ 的集合。于是每个水平集 $\{f>a\}$ 都不变。若不变集概率只能为零或一，$f$ 必为常数。水平集判别也说明一般不动空间恰为 $L^2(\mathcal I)$；不变集在核下几乎必然保持其内外，故确实对补及可数并封闭。

第一节证明的 $\ker(I-T)=\ker(I-T^*)$ 表明不变示性函数也被 $T^*$ 保持。因此 $\int_B A_nf=\int_B f$ 对 $B\in\mathcal I$ 成立，取极限得到条件期望的刻画，先在 $L^2$ 后在 $L^1$ 成立。最后取 $f=\mathbf1_C$、$g=\mathbf1_B$ 得集合相关式；反向由简单函数稠密性及收缩性，将该式延伸到任意 $L^2$ 函数，识别平均投影只有常数。
\end{proof}

混合蕴含遍历，因为普通极限必有相同的 Ces\`aro 极限；反向不成立。还须区分固定一步 $T_1$ 的遍历与全连续半群的遍历：二者的不动空间可能不同，离散采样可产生共振。

\begin{example}[无理旋转遍历但不混合]
圆周上 $\tau(x)=x+\alpha\pmod1$，$\alpha$ 无理。Fourier 模 $e_k(x)=e^{2\pi ikx}$ 满足 $Te_k=e^{2\pi ik\alpha}e_k$，不动函数只有常数，故遍历。但 $|\langle T^ne_1,e_1\rangle|=1$，所以不混合。更一般地，遍历 Markov 算子的单位模本征函数满足 $|f|$ 为常数：$T|f|\ge|Tf|=|f|$，积分相等迫使等号，再用不动空间的刻画。这正是原书角变量的来源。
\end{example}

\Needspace{14\baselineskip}
\begin{theorem}[$H$-定理：熵的单调性]\label{thm:y13-entropy}
若 $f\ge0$、$\int f\,\dd m=1$ 且 $\int f\log^+f\,\dd m<\infty$，定义
\[
 \mathcal H(f)=-\int f\log f\,\dd m,\qquad 0\log0=0.
\]
则 $\mathcal H(T_tf)\ge\mathcal H(f)$，且对 $t_2\ge t_1$ 有 $\mathcal H(T_{t_2}f)\ge\mathcal H(T_{t_1}f)$。相反号的相对熵 $\int f\log f$ 单调不增。
\end{theorem}
\begin{proof}
$\Phi(r)=r\log r$ 凸且下界为 $-e^{-1}$，核的 Jensen 不等式给出 $\Phi(T_tf)\le T_t\Phi(f)$。通过截断凸函数或积分形式的 Jensen 可处理无界 $f$；右侧可积，且不变性给出 $\int\Phi(T_tf)\le\int\Phi(f)$。再用半群性质得到随时间的单调性。该结论只利用保积分与 Markov 性，不依赖遍历性，也不自动给出熵趋于最大值。
\end{proof}

这里证明的是 $T_t$ 对非负函数的 Jensen 不等式。若把 $f$ 解释为给定过程的分布密度，其前向演化通常由核的对偶决定，须在相应对偶空间中另行核对上述性质。

\begin{exercise}[采样共振]\label{ex:y13-sampling}
圆周上的连续平移 $T_tf(x)=f(x+t)$ 的共同不动函数有哪些？比较 $T_1$ 的不动空间。
\end{exercise}
\begin{solution}
所有 $t\ge0$ 的平移共同不变时，Fourier 系数中只有零频可以非零，故为常数。但 $T_1=I$，其不动空间为整个 $L^2$。连续时间平均投到常数，整数时间采样平均却保持原函数。
\end{solution}

\begin{exercise}[熵不保证混合]\label{ex:y13-entropy-rotation}
若 $T$ 来自保测变换，证明每个具有有限熵的概率密度满足 $\mathcal H(Tf)=\mathcal H(f)$。结合无理旋转说明熵单调性的限度。
\end{exercise}
\begin{solution}
因为 $Tf=f\circ\tau$，$\Phi(Tf)=\Phi(f)\circ\tau$，保测性使两者积分相等。无理旋转遍历却不混合，熵在每一步都保持；熵单调并不意味着严格增长、分布收敛或混合。
\end{solution}

\Needspace{10\baselineskip}
\section{带局部紧相空间的 Markov 过程的遍历分解}\label{sec:yosida-13-04}
\noindent\textit{原书起始页：中译本 411；英文版 393。}

以下先固定一步核 $P$。设 $S$ 为局部紧、第二可数 Hausdorff 空间，$C_0(S)$ 表示无穷远消失的连续函数。假定 $T C_0(S)\subset C_0(S)$；这是一项 Feller 性质，不由核的可测性自动推出。$C_c(S)$ 则表示紧支集连续函数，勿与旧文献中相近的记号混用。

\begin{theorem}[不变测度与不逃逸的平均质量]\label{thm:y13-invariant-existence}
存在非零有限不变正测度，当且仅当某个点 $x$ 与非负函数 $g\in C_0(S)$ 满足
\[
 \limsup_n A_ng(x)>0.
\]
若对每个 $x$ 和每个 $f\in C_0(S)$ 都有 $A_nf(x)\to0$，则不存在这样的不变测度。
\end{theorem}
\begin{proof}
令 $\eta_n=n^{-1}\sum_{k=0}^{n-1}P^k(x,\cdot)$。这些概率测度作为 $C_0(S)$ 上的泛函，其范数为一。$C_0(S)$ 可分，用可数稠密集及对角抽取得到某子序列在所有 $C_0$ 函数上收敛到正泛函 $L$；相应正则测度 $\eta$ 满足 $\eta(S)\le1$。选择子序列时保留 $\eta_n(g)$ 的正下界，故 $\eta\ne0$。由于 $Tf\in C_0$ 且
\[
 \eta_n(Tf)-\eta_n(f)=\frac{T^nf(x)-f(x)}n\longrightarrow0,
\]
$\eta$ 不变。这里使用 Riesz--Markov 表示及测度由 $C_0$ 积分确定的性质。

反之，若 $\eta$ 为非零有限不变测度，选非负 $g\in C_0$ 使 $\eta(g)>0$。第二节应用于归一化的 $\eta$，给出 $A_ng\to Qg$ 几乎处处且 $\int Qg\,\dd\eta=\eta(g)>0$，所以至少有一点的极限为正。最后的断言也可直接由控制收敛与不变性推出。
\end{proof}

非紧空间中上述子序列只有弱星即模糊收敛，极限可能丢失总质量；故必须先证明它非零，再归一化。紧空间时常数 $1\in C(S)$，极限质量仍为一，这就是 Krylov--Bogoliubov 平均构造的离散版本。

\begin{theorem}[紧可度量 Feller 情形的遍历分解]\label{thm:y13-ergodic-decomposition}
设 $S$ 紧可度量、$TC(S)\subset C(S)$，$m$ 为不变概率。存在可测概率核 $x\mapsto m_x$，使对 $m$ 几乎每个 $x$，$m_x$ 为遍历不变概率，并有
\[
 Qf(x)=\int f(y)m_x(\dd y)\quad(f\in C(S)),\qquad
 m(B)=\int m_x(B)m(\dd x).
\]
第一式可选共同的满测集，对所有连续函数同时成立。
\end{theorem}
\begin{proof}
取包含常数的可数实函数族，其线性包在 $C(S,\mathbb R)$ 中稠密。第二节给出一个共同满测集 $N$，在其中所有这些函数的平均收敛。由于 $\|A_n\|_\infty\le1$，对每个 $x\in N$，极限延拓为 $C(S)$ 上的正单位泛函 $L_x$。Riesz--Markov 给出概率 $m_x$。同样的端项估计说明 $L_x(Tf)=L_x(f)$，故 $m_x$ 不变。对连续函数，$x\mapsto m_x(f)$ 可测；单调类论证将其延伸为 Borel 概率核。积分平均保持性给出重心公式，也说明 $m_x(N)=1$ 对几乎每个 $x$ 成立。零测补集上的核可任取一点质量。

还须证明分量遍历，不能只从重心公式推断。记可数稠密族为 $f_j$，$q_j=Qf_j$，取上述极限代表。由第三节 $Q=\mathbb E(\cdot\mid\mathcal I)$，故
\[
 Qq_j=q_j,\qquad Q(q_j^2)=q_j^2.
\]
重心核在所有有界 Borel 函数上代表条件期望，因此
\[
 \int(q_j(y)-q_j(x))^2m_x(\dd y)=0
\]
对几乎每个 $x$、所有 $j$ 成立。于是 $m_x$ 集中在 $N$ 中使全部 $q_j(y)=q_j(x)$ 的共同水平集上。对其中的 $y$，$A_nf_j(y)\to q_j(x)=m_x(f_j)$。由于一致有界，控制收敛给出相对于 $m_x$ 的 $L^2$ 收敛到常数。连续函数在 $L^2(m_x)$ 中稠密，平均收缩使全部 $L^2(m_x)$ 函数的平均都趋于其积分，故 $m_x$ 遍历。
\end{proof}

一般局部紧空间中的完整“耗散部分加遍历部分”的逐点划分还涉及模糊极限的质量及共同例外集。上面完整证明了不变测度的存在判别与紧可度量情形的概率分解；原书更一般的全局划分及连续时间共同分解不在此处无条件推广。

\begin{exercise}[向无穷远逃逸]\label{ex:y13-escape}
在 $S=[0,\infty)$ 上取 $\tau(x)=x+1$。证明它保持 $C_0(S)$，且对每个 $f\in C_0(S)$、每个 $x$ 有 $A_nf(x)\to0$，故无不变概率。
\end{exercise}
\begin{solution}
$Tf(x)=f(x+1)$ 连续且无穷远消失。固定 $x$，$f(x+k)\to0$，其 Ces\`aro 平均也趋零。定理\ref{thm:y13-invariant-existence}排除非零有限不变测度。这说明非紧空间中的平均概率质量可能全部逃逸。
\end{solution}

\begin{exercise}[两个封闭类]\label{ex:y13-closed-classes}
四点链的转移矩阵为两个 $\left(\begin{smallmatrix}0&1\\1&0\end{smallmatrix}\right)$ 的直和。求全部不变概率、遍历不变概率及分解。
\end{exercise}
\begin{solution}
不变概率为 $(a/2,a/2,(1-a)/2,(1-a)/2)$，$0\le a\le1$。两个封闭类上的均匀概率 $m_1,m_2$ 分别遍历，全部不变概率为 $am_1+(1-a)m_2$。条件分量 $m_x$ 在第一个类等于 $m_1$，第二个类等于 $m_2$；周期性不妨碍各类的平均遍历性。
\end{solution}

\Needspace{10\baselineskip}
\section{齐次 Riemann 空间上的 Brown 运动}\label{sec:yosida-13-05}
\noindent\textit{原书起始页：中译本 416；英文版 398。}

Riemann 空间在本章理解为 Riemann 流形。转移核与二阶微分算子的联系，来自短时间内位移的一阶、二阶矩，而不是单靠“随机运动”一词。

\begin{definition}[局部性条件与空间齐次]
若对每个 $\varepsilon>0$、固定 $x$ 有
\[
 t^{-1}P_t(x,\{y:d(x,y)>\varepsilon\})\longrightarrow0,
\]
称满足 Lindeberg 型小时间条件；一致版本要求在所讨论的 $x$ 集合上一致。若传递等距群 $G$ 满足 $P_t(gx,gB)=P_t(x,B)$，称转移在空间上齐次。该条件使算子与群作用交换。
\end{definition}

\begin{proposition}[从局部矩得到生成元]\label{prop:y13-local-moments}
在点 $x$ 附近取坐标 $z^1,\ldots,z^d$，令 $\Delta z^i=z^i(y)-z^i(x)$。设上述小时间条件成立，截断一阶矩及二阶矩满足
\[
 t^{-1}\int_{d(x,y)<\varepsilon}\Delta z^i\,P_t(x,\dd y)\to b^i(x),\qquad
 t^{-1}\int_{d(x,y)<\varepsilon}\Delta z^i\Delta z^j\,P_t(x,\dd y)\to q^{ij}(x),
\]
极限不依赖充分小的 $\varepsilon$，且局部二阶矩除以 $t$ 一致有界。若 $f$ 有界、在 $x$ 附近为 $C^2$，则点态生成元为
\[
 Af(x)=\sum_i b^i(x)\partial_i f(x)+\tfrac12\sum_{i,j}q^{ij}(x)\partial_i\partial_j f(x).
\]
矩阵 $q$ 半正定。要把点态极限识别为某 Banach 空间生成元，还需相应范数收敛与定义域条件。
\end{proposition}
\begin{proof}
将积分分成小球内外。外部的差商绝对值不超过 $2\|f\|_\infty P_t(x,B_\varepsilon^c)/t$，趋零。内部作二阶 Taylor 展开，余项为 $o(|\Delta z|^2)$。先取 $t\downarrow0$，再令 $\varepsilon\downarrow0$，二阶矩的有界性使余项消失，主项即上式。对实向量 $\xi$，截断积分 $t^{-1}\int(\sum_i\xi_i\Delta z^i)^2\,\dd P_t\ge0$，取极限得到半正定性。原书把二阶系数直接记作 $b^{ij}$；此处以 $q/2$ 保留协方差的通常约定。
\end{proof}

\begin{example}[欧氏 Brownian 半群的完整模型]\label{exm:y13-gaussian-semigroup}
在 $\mathbb R^d$ 上取
\[
 p_t(z)=(4\pi ct)^{-d/2}e^{-|z|^2/(4ct)},\qquad T_tf=p_t*f,\quad c>0.
\]
Gaussian 卷积给出 $p_t*p_s=p_{t+s}$，总质量为一，故为正收缩半群。在 $C_0(\mathbb R^d)$ 上，近似单位估计给出强连续性；在 $C_c^\infty$ 上 Taylor 展开与 Gaussian 尾部估计给出 $Af=c\Delta f$，且差商在一致范数中收敛。位移均值为零、协方差为 $2ctI$；尾概率以指数速度衰减，满足 Lindeberg 条件。此为生成元 $c\Delta$ 的约定；常见概率论约定 $\tfrac12\Delta$ 对应 $c=1/2$。
\end{example}

\begin{proposition}[各向同性下的唯一微分形式]\label{prop:y13-isotropic}
若球面 $S^d$ 上的保守 Markov 生成元具有局部二阶形式，且与所有等距变换交换，则它为 $c\Delta_{S^d}$，$c\ge0$。在欧氏空间中相同结论适用于与全部欧氏等距变换交换的局部扩散生成元。
\end{proposition}
\Needspace{4\baselineskip}
\begin{proof}
在一点的法坐标中，稳定该点的等距群在切空间上包含全正交群。其不变向量只能为零，故漂移为零；其不变对称二阶型只能是单位型的数倍。正最大值性质使该倍数非负。传递性使系数在所有点相同，保常数性排除零阶项，故得到 Laplace--Beltrami 算子。
\end{proof}

一般齐次流形上的不变二阶型未必只有一个，因此“齐次”本身不能推出唯一的 $c\Delta$。一般流形热核的存在、保总质量及连续路径构造属于额外几何/PDE/概率理论，本节不把局部矩计算当成这些存在性定理的证明。紧连通流形上标准热半群保持归一化体积；其长期行为可再结合 Laplace 谱论分析。

\begin{exercise}[生成元的归一化]\label{ex:y13-brownian-scale}
一维核 $(2\pi t)^{-1/2}e^{-z^2/(2t)}$ 的生成元是什么？比较其二阶矩与 $p_t$ 在 $c=1$ 时的二阶矩。
\end{exercise}
\begin{solution}
该核方差为 $t$，故生成元为 $\tfrac12\,\dd^2/\dd x^2$。$c=1$ 的核方差为 $2t$，生成元为 $\dd^2/\dd x^2$。Taylor 展开的二分之一正是两个约定相差的来源。
\end{solution}

\begin{exercise}[平移不等于各向同性]\label{ex:y13-anisotropic}
对常向量 $b$ 和对称正定矩阵 $Q$，说明 $Af=b\cdot\nabla f+\tfrac12 Q:D^2f$ 的 Gaussian 半群为何平移齐次；何时还与全部正交变换交换？
\end{exercise}
\begin{solution}
转移为从 $x$ 出发、均值 $x+tb$、协方差 $tQ$ 的 Gaussian 分布，只依赖位移，故平移齐次。与所有正交变换交换要求 $Rb=b$、$RQR^{\mathsf T}=Q$ 对所有 $R$ 成立，所以 $b=0$、$Q=\sigma^2I$。平移齐次允许漂移与各向异性。
\end{solution}

\Needspace{10\baselineskip}
\section{W. Feller 的广义 Laplace 算子}\label{sec:yosida-13-06}
\noindent\textit{原书起始页：中译本 422；英文版 403。}

本节在连续函数空间上理解生成元，差商的一致范数收敛保证可逐点计算；常数不在空间中时，保守性由转移核表述。

\begin{proposition}[最大值性质与局部性]\label{prop:y13-local-maximum}
保常数的 Markov 生成元在其定义域上满足正最大值原理：若 $f$ 在 $x_0$ 达到非负全局最大值，则 $Af(x_0)\le0$。若另有上一节的 Lindeberg 条件，且 $f$ 在 $x_0$ 的邻域恒为零，则 $Af(x_0)=0$；因此达到局部最大值时也有相应不等式。
\end{proposition}
\begin{proof}
全局情形由 $T_tf(x_0)\le f(x_0)$ 后除以 $t$ 得到。局部为零时，$|T_tf(x_0)|\le\|f\|_\infty P_t(x_0,B_\varepsilon^c)=o(t)$。局部最大值时，同样将小球内的差分与小球外的 $o(t)$ 分开。若空间包含常数，保守性还给出 $1\in D(A)$ 时 $A1=0$；在非紧空间 $C_0$ 中，常数通常不属于该空间，须在核或适当扩张上理解此式。
\end{proof}

\begin{theorem}[光滑系数的尺度与速度表示]\label{thm:y13-scale-speed}
在开区间 $I=(l,r)$ 上考虑 $Lf=a(x)f''+b(x)f'$，其中 $a>0$ 连续、$b$ 连续。固定 $x_0\in I$，令
\[
 s'(x)=\exp\left(-\int_{x_0}^x\frac{b(u)}{a(u)}\,\dd u\right),\qquad
 m'(x)=\frac1{a(x)s'(x)}.
\]
则 $s$ 严格递增，且
\[
 L=D_mD_s,\qquad D_sf=\frac{f'}{s'},\qquad D_mg=\frac{g'}{m'}.
\]
调和方程 $Lf=0$ 的解为 $\alpha+\beta s$。尺度 $s$ 允许正仿射变换；相应速度测度需作逆比例调整。
\end{theorem}
\begin{proof}
由 $(\log s')'=-b/a$，直接计算
\[
 \frac1{m'}\left(\frac{f'}{s'}\right)'
 =a f''-a\frac{s''}{s'}f'=a f''+bf'.
\]
若该式为零，则 $D_sf$ 为常数，所以解恰为上述形式。将 $s$ 换为 $\gamma s+\delta$、$\gamma>0$ 时，$\dd m$ 换为 $\gamma^{-1}\dd m$，乘积算子不变。
\end{proof}

\paragraph{广义速度测度。}
令 $s$ 连续严格递增，$m$ 为局部有限且在每个开子区间有正质量的测度。广义表达式 $Lf=g$ 可用积分关系定义：$D_sf$ 有局部有界变差的适当代表，并满足
\[
 D_sf(v)-D_sf(u)=\int_{(u,v]}g\,\dd m.
\]
当 $m$ 有原子时，$D_sf$ 可以跳跃，故 $D_mD_s$ 不必是普通光滑系数的微分式。这是 Feller 广义 Laplace 表示容纳黏滞等现象的原因。原书从局部性、最大值原理及非退化条件反向构造尺度和速度；其完整反向表示定理不由本节的光滑计算替代，列入进一步阅读范围。

\paragraph{边界类型与定义域。}
选 $c\in(l,r)$，定义左端测试量
\[
 \Sigma_l=\int_{(l,c)}[s(x)-s(l)]\,m(\dd x),\qquad
 N_l=\int_{(l,c)}[s(c)-s(x)]\,m(\dd x),
\]
端点尺度取单侧极限，允许无穷值。右端相应使用 $s(r)-s(x)$ 与 $s(x)-s(c)$。Feller 边界分类定理在其一维正则扩散假设下给出
\[
 \begin{array}{c|cc}
 &N_l<\infty&N_l=\infty\\\hline
 \Sigma_l<\infty&\text{正则}&\text{流出}\\
 \Sigma_l=\infty&\text{流入}&\text{自然}
 \end{array}
\]
本节采用该分类作为外部结构定理，并计算典型例子；不声称已证明一般边界可达性的概率刻画。正则端点还须选择吸收、反射或其他允许的边界条件，才能确定闭生成元。内部的 $D_mD_s$ 表达式本身不决定半群。

\begin{example}[三种左端行为]
在 $(0,\infty)$ 上，$L=f''$ 有 $s=x$、$m'=1$，故零端两个测试量均有限，为正则端点。$L=xf''$ 有 $s=x$、$m'=1/x$，于是 $\Sigma_0<\infty$、$N_0=\infty$，零端为流出端点。对三维 Bessel 表达式 $L=f''+2f'/x$，有 $s=-1/x$、$m'=x^2$，所以 $\Sigma_0=\infty$、$N_0<\infty$，零端为流入端点。
\end{example}

\begin{exercise}[自然端点]\label{ex:y13-natural-boundary}
对 $Lf=x^2f''+xf'$、$x>0$，求尺度、速度密度与零端类型，并解释对数坐标的作用。
\end{exercise}
\begin{solution}
可取 $s=\log x$、$m'=1/x$。由于 $s(0)=-\infty$，有 $\Sigma_0=\infty$，且
\[
 N_0=\int_0^c\log(c/x)\,\frac{\dd x}{x}=\infty.
\]
所以零端为自然端点。若 $f(x)=g(\log x)$，则 $Lf=g''(\log x)$，零端对应实轴 Brownian 运动的负无穷远。
\end{solution}

\begin{exercise}[边界条件改变转移]\label{ex:y13-boundary-kernels}
令 $p_t(z)=(4\pi t)^{-1/2}e^{-z^2/(4t)}$。比较半直线上 $p_t(y-x)-p_t(y+x)$ 与 $p_t(y-x)+p_t(y+x)$ 的总质量和零端边界条件。
\end{exercise}
\begin{solution}
对 $x>0$，差核在 $y\ge0$ 上质量小于一，给出到达零端后被杀死的次 Markov 半群，边界为 Dirichlet。和核总质量为一，给出反射半群；其关于 $x$ 的零端导数为零，对应 Neumann 条件。二者内部生成式均为二阶导数，保守性与状态空间的边界处理不同。
\end{solution}

\Needspace{10\baselineskip}
\section{扩散算子的一种推广}\label{sec:yosida-13-07}
\noindent\textit{原书起始页：中译本 426；英文版 408。}

去掉 Lindeberg 条件后，短时间内发生远距离跳跃的概率可以是一阶小量。生成元因而出现积分项，通常不再具有局部性。

\begin{proposition}[有限跳跃测度的生成元]\label{prop:y13-compound-poisson}
设 $\nu$ 为 $\mathbb R^d\setminus\{0\}$ 上有限正测度，$r=\nu(\mathbb R^d)>0$、$\pi=\nu/r$。定义 $Kf(x)=\int f(x+z)\pi(\dd z)$。则 $B=r(K-I)$ 是 $C_0(\mathbb R^d)$ 上的有界生成元，
\[
 e^{tB}f=e^{-rt}\sum_{n=0}^\infty\frac{(rt)^n}{n!}K^nf,
 \qquad Bf(x)=\int[f(x+z)-f(x)]\nu(\dd z).
\]
这是保守的正收缩半群，且为强连续。
\end{proposition}
\begin{proof}
$K$ 正且范数至多一；连续性由控制收敛，无穷远消失性由先把 $\pi$ 的质量限制在大球再估计尾部得到。级数在算子范数中收敛，系数为总和等于一的非负 Poisson 权重，所以正且收缩。半群性质由指数级数得到，差商在算子范数中趋于 $B$。核可理解为以 Poisson 次数抽取独立位移；这里公式本身已构成转移核，无需先援引样本路径构造。
\end{proof}

与常系数漂移 Gaussian 半群复合，因各部分都是平移卷积而交换，得到跳跃扩散，其生成式为
\[
 Af=b\cdot\nabla f+\tfrac12Q:D^2f+
 \int[f(x+z)-f(x)]\nu(\dd z),\qquad Q\ge0.
\]
在测试函数上可直接对复合半群求导验证此式。

\begin{definition}[补偿后的跳跃项]
允许 $\nu$ 在原点附近有无限质量，但要求
\[
 \int_{\mathbb R^d\setminus\{0\}}(1\wedge|z|^2)\nu(\dd z)<\infty.
\]
对 $f\in C_c^\infty$，定义
\[
 \mathcal Jf(x)=\int\bigl[f(x+z)-f(x)-\mathbf1_{\{|z|\le1\}}z\cdot\nabla f(x)\bigr]\nu(\dd z).
\]
一阶补偿使小位移余项为 $O(|z|^2)$，所以积分绝对收敛。截断函数的改变须由漂移项的相应调整补偿，不能单独更改积分项。
\end{definition}

\paragraph{一般结构定理的边界。}
对 $\mathbb R^d$ 上具有 $C_c^\infty\subset D(A)$ 的 Feller 生成元，正最大值原理的结构表示通常写为
\begin{align*}
 Af(x)={}&-\kappa(x)f(x)+b(x)\cdot\nabla f(x)+\tfrac12Q(x):D^2f(x)\\
 &+\int\bigl[f(x+z)-f(x)-\mathbf1_{\{|z|\le1\}}z\cdot\nabla f(x)\bigr]\nu(x,\dd z),
\end{align*}
其中 $\kappa\ge0$、$Q\ge0$，每个 $x$ 的跳跃测度满足上述二阶可积条件；保守且单位函数可合法纳入定义域时杀死项为零。这是广义扩散的标准形式。原书在固定时空点的截断测度有界、紧性及差商存在条件下给出相应的子序列表示。本节将一般结构表示列为外部定理，完整构造并证明的是有限跳跃与常系数复合模型；任意给定的变系数组合是否生成 Feller 半群还需要额外存在性条件。

\begin{proposition}[最大值原理并不等于局部性]
上述跳跃积分在 $f$ 的全局最大点非正，而 $f$ 在该点邻域恒为零时积分仍可非零。对有限非零 $\nu$，相应过程不满足 Lindeberg 条件。
\end{proposition}
\begin{proof}
最大点处梯度为零，每个差分非正。若 $f$ 远处为正，且跳跃测度能跳到该处，则 $\mathcal Jf(x)>0$，显示非局部性。选 $\varepsilon>0$ 使 $\nu\{|z|>\varepsilon\}>0$，Poisson 级数的一次跳跃项给出
\[
 \liminf_{t\downarrow0}t^{-1}P_t(x,\{|y-x|>\varepsilon\})
 \ge\nu\{|z|>\varepsilon\}>0.
\]
故局部扩散所需的小时间条件失败。
\end{proof}

\begin{exercise}[对称跳跃的符号]\label{ex:y13-jump-symbol}
在实线上取 $\nu=\tfrac r2(\delta_a+\delta_{-a})$、$a\ne0$。写出 $A$，并求它在有界 Fourier 模 $e^{i\xi x}$ 上的形式作用。
\end{exercise}
\begin{solution}
$Af(x)=\tfrac r2[f(x+a)+f(x-a)-2f(x)]$。形式代入得到 $Ae^{i\xi x}=r(\cos(a\xi)-1)e^{i\xi x}$。这些模不在 $C_0$ 中，等式在有界函数扩张或符号计算意义下理解；乘以测试截断后仍可用于辨认跳跃项，不能当作 $C_0$ 的本征向量结论。
\end{solution}

\begin{exercise}[小跳跃的补偿]\label{ex:y13-small-jumps}
实线上 $\nu(\dd z)=\mathbf1_{0<|z|<1}|z|^{-1-\alpha}\,\dd z$，$0<\alpha<2$。验证二阶可积条件，并说明何时未补偿的一阶绝对估计足够。
\end{exercise}
\begin{solution}
$\int z^2\,\dd\nu=2\int_0^1r^{1-\alpha}\,\dd r=2/(2-\alpha)<\infty$，所以补偿后的 Taylor 余项可积。一阶绝对估计要求 $\int|z|\,\dd\nu=2\int_0^1r^{-\alpha}\,\dd r<\infty$，恰在 $\alpha<1$ 时成立。对称性可允许主值抵消，但主值与绝对收敛不能混称。
\end{solution}

\Needspace{10\baselineskip}
\section{Markov 过程和位势}\label{sec:yosida-13-08}
\noindent\textit{原书起始页：中译本 429；英文版 410。}

设 $T_t$ 为 Banach 空间上的有界强连续半群，生成元 $A$，此节先假定为正收缩。记
\[
 R_\lambda=(\lambda I-A)^{-1}
 =\int_0^\infty e^{-\lambda t}T_t\,\dd t,\qquad\lambda>0.
\]
右侧对每个向量是 Bochner 积分。核表示下它也是折扣占用算子。$h\in D(A)$ 满足 $Ah=0$ 称为调和元素；等价地 $T_th=h$，或 $\lambda R_\lambda h=h$ 对所有 $\lambda>0$ 成立。

\Needspace{12\baselineskip}
\begin{proposition}[位势与 Poisson 方程]\label{prop:y13-poisson}
若 $R_\lambda g$ 在 $\lambda\downarrow0$ 时有范数极限 $v$，则 $v\in D(A)$ 且 $Av=-g$。特别地，若 $\int_0^\infty\|T_tg\|\,\dd t<\infty$，则
\[
 v=\int_0^\infty T_tg\,\dd t=\lim_{\lambda\downarrow0}R_\lambda g.
\]
反向，$Av=-g$ 的解可以相差调和元素，不能仅凭方程断言唯一。
\end{proposition}
\begin{proof}
$AR_\lambda g=\lambda R_\lambda g-g\to-g$，而 $R_\lambda g\to v$，生成元闭性给出结论。绝对可积时，Bochner 控制收敛给出积分极限。两个解之差属于 $\ker A$，故最后的限制必要。
\end{proof}

\begin{theorem}[正格中的 Riesz 型分解]\label{thm:y13-riesz-potential}
设空间为 Dedekind $\sigma$ 完备 Banach 格，范数序列序连续。若 $u\in D(A)$、$u\ge0$、$Au\le0$，则
\[
 h=\lim_{\lambda\downarrow0}\lambda R_\lambda u,\qquad
 p=u-h=\lim_{\lambda\downarrow0}R_\lambda(-Au)
\]
均以范数收敛，$h\ge0$ 调和，$p\ge0$ 为位势。$h$ 是所有不超过 $u$ 的调和元素中的最大者。相应地，若次调和 $x\in D(A)$、$Ax\ge0$ 有调和上界，则它有最小调和上界，并分解为该上界加一个非正位势。
\end{theorem}
\begin{proof}
恒等式 $\lambda R_\lambda u-u=R_\lambda Au\le0$ 给出 $0\le\lambda R_\lambda u\le u$。对 $\lambda>\mu>0$，预解式恒等式给出
\[
 (\lambda R_\lambda-\mu R_\mu)u
 =-(\lambda-\mu)R_\lambda R_\mu Au\ge0.
\]
故 $\lambda\downarrow0$ 时该族递减。先沿 $\lambda_n\downarrow0$ 取格下确界，序连续范数使序列趋于 $h$；任意中间参数由单调夹逼得到同一范数极限。由于
\[
 A(\lambda R_\lambda u)=\lambda(\lambda R_\lambda u-u)\longrightarrow0,
\]
闭性给出 $Ah=0$。再用 $R_\lambda(-Au)=u-\lambda R_\lambda u$ 得 $p$ 及其位势表示。若 $k\le u$ 调和，则 $k=\lambda R_\lambda k\le\lambda R_\lambda u$，取极限得 $k\le h$。次调和情形对 $u=H-x\ge0$ 应用以上结论，其中 $H$ 是任一调和上界，即得最小调和上界。
\end{proof}

此定理是原书次调和分解的一个明确 Banach 格版本；一般局部凸序空间的弱单调完备假设不在此处另作推广。特别地，$C_0(S)$ 的一致范数通常不满足序连续假设，不能把此定理未经检查就用于所有连续函数位势。

\begin{example}[欧氏热核的暂留与再返界限]\label{exm:y13-green}
取生成元 $\Delta$。当 $d\ge3$ 且 $g\in C_c(\mathbb R^d)$，有 $\int_0^\infty\|T_tg\|_\infty\,\dd t<\infty$：小时间用收缩性，大时间用 $(4\pi t)^{-d/2}\|g\|_1$。因而
\[
 Vg(x)=\int_{\mathbb R^d}G_d(x-y)g(y)\,\dd y,\qquad
 G_d(z)=\frac{|z|^{2-d}}{(d-2)|\mathbb S^{d-1}|}.
\]
该核由 $\int_0^\infty(4\pi t)^{-d/2}e^{-|z|^2/(4t)}\,\dd t$ 换元计算，$-\Delta Vg=g$ 按生成元及分布意义成立；不据此声称一般连续 $g$ 的位势有连续二阶导数。对 $d=1,2$ 和非零非负 $g\in C_c$，固定 $x$ 时大时间的 $T_tg(x)$ 与 $t^{-d/2}\int g$ 同阶，故占用积分发散。该例说明无折扣位势可能不存在，虽每个 $R_\lambda$ 都存在。
\end{example}

\begin{exercise}[平衡态的障碍]\label{ex:y13-potential-obstruction}
在保概率 Markov 半群上，证明 $g=1$ 没有有限范数 Abel 位势；并证明 $Av=-g$ 有可积定义域中的解时必有 $\int g\,\dd m=0$。
\end{exercise}
\begin{solution}
$R_\lambda1=1/\lambda$，范数随 $\lambda\downarrow0$ 发散。若 $v\in D(A)\subset L^1$，保积分性对差商取极限给出 $\int Av=0$，所以 $Av=-g$ 要求 $\int g=0$。非零均值是 Poisson 方程的平衡态障碍。
\end{solution}

\begin{exercise}[有限链上的调和与位势]\label{ex:y13-two-state-potential}
令 $A=\left(\begin{smallmatrix}-a&a\\b&-b\end{smallmatrix}\right)$、$a,b>0$。求不变概率、平均投影，并在零均值子空间求位势。
\end{exercise}
\begin{solution}
不变概率为 $\pi=(b,a)/(a+b)$，$Qf=(\pi f)1$。零均值子空间由 $(a,-b)$ 张成，$A$ 在其上为 $-(a+b)I$，故位势 $Vg=g/(a+b)$。全空间上常数调和分量阻止无折扣逆，但投到零均值空间后可逆。
\end{solution}

\Needspace{10\baselineskip}
\section{抽象位势算子和半群}\label{sec:yosida-13-09}
\noindent\textit{原书起始页：中译本 430；英文版 411。}

本节不要求有次序或正性。$X$ 为实或复 Banach 空间，$T_t$ 为收缩强连续半群，$A$ 为其生成元。把连续对偶伴随写为 $V^*$，配对按对偶线性泛函理解。

\noindent\textit{原书校注：}末尾与 Hunt 理论比较的段落，把预解式位势的极限下标排成趋于无穷，与本节式 (7) 的零端极限不符。本稿统一采用 $\lambda\downarrow0$；实际上收缩半群满足 $\|R_\lambda\|\le1/\lambda$，若令 $\lambda\to\infty$，未乘 $\lambda$ 的预解式只能趋于零。

\begin{theorem}[零端 Abel 极限与抽象位势]\label{thm:y13-abstract-potential}
下列条件等价：$A$ 单射且 $\overline{R(A)}=X$；对每个 $x\in X$ 有 $\lambda R_\lambda x\to0$（$\lambda\downarrow0$）。此时
\[
 V=-A^{-1},\quad D(V)=R(A),\quad R(V)=D(A)
\]
为闭、定义域及值域均稠密的算子，称为抽象位势，满足
\[
 D(V)=\{x:\lim_{\lambda\downarrow0}R_\lambda x\text{ 在范数中存在}\},\qquad
 Vx=\lim_{\lambda\downarrow0}R_\lambda x.
\]
\end{theorem}
\begin{proof}
若 $x=Ay$，则 $\lambda R_\lambda x=\lambda(\lambda R_\lambda y-y)\to0$。因 $\|\lambda R_\lambda\|\le1$，值域稠密时此极限扩展到全部 $X$。反向若全部极限为零，则 $Ah=0$ 蕴含 $\lambda R_\lambda h=h=0$；又 $x-\lambda R_\lambda x=-AR_\lambda x\in R(A)$ 趋于 $x$，故值域稠密。

取 $x=-Ay$，则 $R_\lambda x=y-\lambda R_\lambda y\to y=Vx$。反之，若 $R_\lambda x\to y$，闭性与 $AR_\lambda x=\lambda R_\lambda x-x\to-x$ 给出 $Ay=-x$，所以 $x\in R(A)$。逆算子的图由 $A$ 的图交换坐标得到，故 $V$ 闭；两个稠密性由假设及生成元定义域稠密得到。
\end{proof}

\begin{theorem}[位势的双侧范数刻画]\label{thm:y13-potential-characterization}
设 $V$ 为闭线性算子，$D(V)$ 与 $R(V)$ 均在 $X$ 中稠密。则 $V$ 是某收缩强连续半群的抽象位势，当且仅当对每个 $\lambda>0$ 有
\begin{align}
 \|\lambda Vx\|&\le\|x+\lambda Vx\|&& (x\in D(V)),\label{eq:y13-potential-primal}\\
 \|\lambda V^*\phi\|&\le\|\phi+\lambda V^*\phi\|&& (\phi\in D(V^*)).\label{eq:y13-potential-dual}
\end{align}
此时 $A=-V^{-1}$，且 $R_\lambda=V(I+\lambda V)^{-1}$。
\end{theorem}
\begin{proof}
先证明充分性。令 $B_\lambda=I+\lambda V$。第一不等式还给出 $\|x\|\le2\|B_\lambda x\|$，所以 $B_\lambda$ 单射。若 $B_\lambda x_n$ 收敛，同一估计使 $x_n,Vx_n$ 均 Cauchy，闭性使极限仍在图中。因此 $R(B_\lambda)$ 闭。其零化子为 $\ker(I+\lambda V^*)$；第二不等式使该核为零。Hahn--Banach 于是给出值域稠密，结合闭性得满射。

令 $J_\lambda=VB_\lambda^{-1}$，则 $\|\lambda J_\lambda\|\le1$。直接在 $D(V)$ 上代入并利用 $V$ 与其有界预解算子的交换关系可得
\[
 J_\lambda-J_\mu=(\mu-\lambda)J_\lambda J_\mu.
\]
对 $y=Vx$，$\lambda J_\lambda y=y-J_\lambda x\to y$（$\lambda\to\infty$），因 $R(V)$ 稠密，这个极限对全部 $y$ 成立。伪预解式的核与参数无关，故其核为零；又 $R(J_\lambda)=R(V)$ 稠密。因此第八章的伪预解式构造给出稠密定义的闭算子 $A$，其预解式为 $J_\lambda$。$\|\lambda J_\lambda\|\le1$ 蕴含各幂的 Hille--Yosida 估计，故 $A$ 生成收缩强连续半群。

由 $J_\lambda(B_\lambda x)=Vx$ 得 $AVx=-x$。$V$ 单射也可从 $\ker J_\lambda=0$ 得到，故 $A=-V^{-1}$、$R(A)=D(V)$ 稠密。上一条定理将 $V$ 识别为该半群的抽象位势。

必要性方面，若 $V=-A^{-1}$，预解式给出
\[
 B_\lambda^{-1}=I-\lambda R_\lambda,\qquad
 VB_\lambda^{-1}=R_\lambda.
\]
所以 $\lambda VB_\lambda^{-1}$ 为收缩，得到第一不等式。闭稠密算子的伴随逆关系给出 $B_\lambda^{*-1}=(B_\lambda^{-1})^*$，且 $\lambda V^*B_\lambda^{*-1}=(\lambda VB_\lambda^{-1})^*$ 仍为收缩，得到第二不等式。这里未假定伴随半群在整个 $X^*$ 上强连续。
\end{proof}

\begin{corollary}[Hilbert 空间的增生表述]
在 Hilbert 空间中，上述两条范数不等式分别等价于
\[
 \operatorname{Re}\langle Vx,x\rangle\ge0,\qquad
 \operatorname{Re}\langle V^*y,y\rangle\ge0.
\]
因此闭、定义域和值域稠密的 $V$ 是抽象位势，当且仅当 $V,V^*$ 均增生。
\end{corollary}
\begin{proof}
展开 $\|x+\lambda Vx\|^2-\|\lambda Vx\|^2=\|x\|^2+2\lambda\operatorname{Re}\langle Vx,x\rangle$。它对所有 $\lambda>0$ 非负恰好等价于增生性；伴随同理。
\end{proof}

\begin{example}[乘法位势的定义域]
在 $L^2(0,1)$ 上取 $(T_tf)(s)=e^{-ts}f(s)$。则 $Af=-sf$，它单射但逆不有界，
\[
 (Vg)(s)=g(s)/s,\qquad D(V)=\{g:g(s)/s\in L^2(0,1)\}.
\]
$R_\lambda g=g/(\lambda+s)$；$\lambda R_\lambda g\to0$ 由控制收敛成立。只有满足上述定义域条件的 $g$ 才有位势范数极限，不能把所有有界预解式的零端极限称为处处定义的有界算子。
\end{example}

\begin{exercise}[Abel 位势不要求普通积分收敛]\label{ex:y13-abel-oscillation}
在 $X=\mathbb C$ 上令 $T_tz=e^{it}z$。求 $A,V$，并比较 $\int_0^R T_tz\,\dd t$ 与 $\int_0^\infty e^{-\lambda t}T_tz\,\dd t$ 的极限。
\end{exercise}
\begin{solution}
$A=iI$、$V=-1/i=iI$。普通积分为 $(e^{iR}-1)z/i$，对非零 $z$ 随 $R\to\infty$ 不收敛；折扣积分为 $z/(\lambda-i)$，随 $\lambda\downarrow0$ 趋于 $iz$。故抽象位势容纳 Abel 可和但普通占用积分不收敛的情形，也不要求格意义下的正性。
\end{solution}

\begin{exercise}[乘法位势的临界指数]\label{ex:y13-potential-domain}
在上例取 $g(s)=s^\beta$、$\beta>-1/2$。判断何时 $g\in D(V)$，并说明不在定义域时预解式范数的行为。
\end{exercise}
\begin{solution}
$g\in D(V)$ 等价于 $\int_0^1s^{2\beta-2}\,\dd s<\infty$，即 $\beta>1/2$。若 $\beta\le1/2$，则 $\|R_\lambda g\|_2^2=\int_0^1s^{2\beta}/(\lambda+s)^2\,\dd s\uparrow\infty$，不可能有范数极限，虽对每个正 $\lambda$ 都是良定义的预解式值。
\end{solution}


\Needspace{6\baselineskip}
\section*{本章参考说明与阅读衔接}
直接来源为吉田中英文第六版。个别遍历证明以有限极大不等式重组；不变测度分解在紧可度量 Feller 情形展开；扩散部分以局部矩、尺度速度与显式跳跃模型解释生成元。一般比值遍历、非紧空间的完整逐点分解、广义尺度速度的反向刻画、一般 Feller 边界分类及变系数跳跃表示的证明边界均明确记录，不把模型计算等同于全部结构定理。下一章将转向发展方程的积分与初值问题。

% !TeX encoding = UTF-8
% 吉田第六版：中译本第 14 章；此文件当前仅为目录与正文占位。
\SetChapterInfo{subtitle={吉田耕作第六版主线 · 正文待撰写}}
\chapter{发展方程的积分}\label{ch:yosida-14}
\YosidaChapterPlaceholder{14}{The Integration of the Equation of Evolution}{437--486}{418--465}

\section{在 \texorpdfstring{$L^2(\R^m)$}{L2(Rm)} 中扩散方程的积分}\label{sec:yosida-14-01}
\YosidaSectionPlaceholder{438}{419}
% TODO: 定义与记号；定理与完整证明；例子；习题；中英文核对。
相关模型：第 \ref{ch:pde} 章的热方程。\par

\section{紧 Riemann 空间中扩散方程的积分}\label{sec:yosida-14-02}
\YosidaSectionPlaceholder{444}{425}
% TODO: 定义与记号；定理与完整证明；例子；习题；中英文核对。

\section{Euclid 空间 \texorpdfstring{$\R^m$}{Rm} 中波动方程的积分}\label{sec:yosida-14-03}
\YosidaSectionPlaceholder{446}{427}
% TODO: 定义与记号；定理与完整证明；例子；习题；中英文核对。

\section{B-空间中非时齐发展方程的积分}\label{sec:yosida-14-04}
\YosidaSectionPlaceholder{449}{430}
% TODO: 定义与记号；定理与完整证明；例子；习题；中英文核对。

\section{Tanabe 和 Sobolevski 的方法}\label{sec:yosida-14-05}
\YosidaSectionPlaceholder{457}{438}
% TODO: 定义与记号；定理与完整证明；例子；习题；中英文核对。

\section{非线性发展方程 1（Komura--Kato 方法）}\label{sec:yosida-14-06}
\YosidaSectionPlaceholder{464}{445}
% TODO: 定义与记号；定理与完整证明；例子；习题；中英文核对。

\section{非线性发展方程 2（立足于 Crandall--Liggett 收敛定理的方法）}\label{sec:yosida-14-07}
\YosidaSectionPlaceholder{473}{454}
% TODO: 定义与记号；定理与完整证明；例子；习题；中英文核对。


% 第二部分为已写成的专题扩展，章号顺接为 15--22。
% 吉田原书中的广义函数相关节保留在第一部分，正文可交叉引用本部分。
\part{广义函数专题扩展}
\SetChapterInfo{subtitle={从积分泛函走向连续对偶},keywords={Dirac 质量; 弱极限; 分部积分},prereq={Lebesgue 积分; 多元微积分},goals={理解测试函数的作用与拓扑的必要性}}
\chapter{为什么需要广义函数}\label{ch:motivation}
\ChapterSource{主要对应施瓦尔茨 \cite{schwartz} 第 I 章“函数概念的推广”与第 II 章求导的基本思想。下面的例子按本书的配对约定重新推导。}

\section{函数也可以看成泛函}
通常把 $f$ 看作映射 $x\mapsto f(x)$。另一种观察方式是把它看作加权测量：给定探针 $\varphi$，测得 $\int f\varphi$。探针的紧支集使观察局限在有限区域；探针的光滑性使任意阶分部积分成为可能。
\begin{definition}[正则分布的原型]
对 $f\in L^1_\Loc(\Omega)$，定义 $T_f:C_c^\infty(\Omega)\to\mathbb C$ 为
\[
 \Pair{T_f}{\varphi}=\int_\Omega f\varphi\,\dd x.
\]
若 $\supp\varphi\subset K\Subset\Omega$，则
\begin{equation}\label{eq:regular-bound}
 |\Pair{T_f}{\varphi}|\le \left(\int_K|f|\,\dd x\right)\sup_K|\varphi|.
\end{equation}
\end{definition}
这一估计已包含未来连续性的全部信息：固定观察区域后，只须控制测试函数的零阶大小。几乎处处相等的函数给出同一泛函；反过来，若 $T_f=T_g$，则 $f=g$ 几乎处处。后一事实将在正则化建立后证明。

\section{集中质量与 Dirac 分布}
\begin{definition}[Dirac 分布]
对 $a\in\Omega$，令 $\Pair{\delta_a}{\varphi}=\varphi(a)$。这是点值探针的对偶对象；它也是单位点质量测度所给的分布。
\end{definition}
Dirac 分布不是局部可积函数。若 $T_f=\delta_a$，对所有支集避开 $a$ 的测试函数积分都为零，便推出 $f=0$ 于 $\Omega\setminus\{a\}$ 几乎处处。单点的 Lebesgue 测度为零，故 $f=0$ 几乎处处，矛盾于取 $\varphi(a)=1$。

取非负 $\rho\in C_c^\infty(\R^n)$，满足 $\supp\rho\subset\overline{B(0,1)}$ 与 $\int\rho=1$。例如将
\[
 \rho_0(x)=\begin{cases}\exp[-1/(1-|x|^2)],&|x|<1,\\0,&|x|\ge1\end{cases}
\]
乘上归一化常数。令 $\rho_\varepsilon(x)=\varepsilon^{-n}\rho(x/\varepsilon)$。
\begin{proposition}[集中质量的极限]\label{prop:delta-approx}
当 $\varepsilon\downarrow0$ 时，$T_{\rho_\varepsilon(\cdot-a)}\to\delta_a$，即对每个 $\varphi\in\Test(\R^n)$，其配对趋于 $\varphi(a)$。
\end{proposition}
\begin{proof}
变量替换 $x=a+\varepsilon y$ 给出
\[
 \int\rho_\varepsilon(x-a)\varphi(x)\,\dd x
 =\int\rho(y)\varphi(a+\varepsilon y)\,\dd y.
\]
在 $\supp\rho$ 上 $\varphi(a+\varepsilon y)$ 一致趋于 $\varphi(a)$，而 $\rho$ 可积，故结论成立。
\end{proof}
这里函数的峰值通常趋于无穷，其 $L^1$ 范数恒为 1，且对 $x\ne a$ 最终为零。逐点极限不能记录保留下来的总质量；逐探针极限却能记录它。

\section{把导数移到测试函数上}
对光滑 $f$ 和紧支集 $\varphi$，分部积分给出
\[
 \int (\partial_j f)\varphi=-\int f\partial_j\varphi.
\]
右边只要求 $f$ 局部可积。由此可以把求导定义在更大范围：
\begin{equation}\label{eq:derivative-preview}
 \Pair{\partial_jT}{\varphi}=-\Pair{T}{\partial_j\varphi}.
\end{equation}
紧支集消除边界项；负号来自分部积分，不是任意约定。
\begin{example}[跳跃的导数]
令 $H(x)=0$ 于 $x<0$，$H(x)=1$ 于 $x>0$；零点值任取。则
\[
 \Pair{H'}{\varphi}=-\int_0^\infty\varphi'(x)\,\dd x=\varphi(0),
 \qquad H'=\delta_0.
\]
同理 $\operatorname{sgn}'=2\delta_0$，$|x|'=\operatorname{sgn}x$，$|x|''=2\delta_0$。
\end{example}

\section{弱收敛记录什么信息}
在有界区间上，$f_k(x)=\sin(kx)$ 的振荡幅度并未变小。但对每个测试函数，
\[
 \int\sin(kx)\varphi(x)\,\dd x
 =k^{-1}\int\cos(kx)\varphi'(x)\,\dd x\longrightarrow0.
\]
因此 $T_{f_k}\to0$。同时 $f_k^2=(1-\cos(2kx))/2$ 在分布意义趋于 $1/2$。这个例子说明线性弱极限可以控制线性方程，却不能自动穿过乘法这样的非线性运算。
\begin{pitfall}[符号计算必须有定义域]
“$\delta$ 在零点无穷大、其余点为零”只是一种不精确的图像；定义是 $\Pair{\delta}{\varphi}=\varphi(0)$。同样，$\delta^2$ 并没有由上述定义得到合法乘法。后续一切运算须先证明测试函数侧的映射连续。
\end{pitfall}

\section{习题}
\begin{exercise}[质量与缩放]
设 $g\in L^1(\R^n)$，$\int g=c$。证明 $\varepsilon^{-n}g(\cdot/\varepsilon)\to c\delta_0$ 于分布意义；$g$ 不需要非负。
\end{exercise}
\begin{solution}
对固定 $\varphi$，配对等于 $\int g(y)\varphi(\varepsilon y)\,\dd y$。被积函数趋于 $g(y)\varphi(0)$，且其绝对值不超过 $\|\varphi\|_\infty|g(y)|$。由控制收敛定理得 $c\varphi(0)$。
\end{solution}
\begin{exercise}[边界质量]
在 $\R$ 上证明 $\mathbf1_{(a,b)}'=\delta_a-\delta_b$，再求它的二阶分布导数。
\end{exercise}
\begin{solution}
$-\int_a^b\varphi'=\varphi(a)-\varphi(b)$，故一阶导数为 $\delta_a-\delta_b$。再求导得 $\delta_a'-\delta_b'$，其作用为 $-\varphi'(a)+\varphi'(b)$。
\end{solution}

\SetChapterInfo{subtitle={半范数、Fréchet 空间与 LF 拓扑},keywords={局部凸; 紧支集; 有界集},prereq={拓扑向量空间的基本概念},goals={区分局部光滑收敛与测试函数收敛}}
\chapter{测试函数空间及其拓扑}\label{ch:test}
\ChapterSource{主要对应施瓦尔茨 \cite{schwartz} 第 I 章 §2、第 III 章 §§1--3。归纳极限的有界集定理作为标准 LF 空间结论引用，其后给出本书需要的具体推论。}

\section{局部凸空间与有限半范数估计}
\begin{definition}[半范数]
向量空间 $E$ 上的函数 $p:E\to[0,\infty)$ 若满足 $p(ax)=|a|p(x)$ 和 $p(x+y)\le p(x)+p(y)$，称为半范数。给定分离点的半范数族 $\{p_\lambda\}$，以有限交
\[
 \{x:p_{\lambda_1}(x)<\varepsilon_1,\ldots,p_{\lambda_r}(x)<\varepsilon_r\}
\]
作为零邻域基，可定义 Hausdorff 局部凸拓扑。
\end{definition}
\begin{proposition}[连续线性泛函的估计]\label{prop:seminorm-est}
线性泛函 $L:E\to\mathbb C$ 连续，当且仅当存在有限多个半范数和 $C>0$，使
\[
 |L(x)|\le C\max_{1\le j\le r}p_{\lambda_j}(x).
\]
\end{proposition}
\begin{proof}
连续性给出某个基本零邻域上 $|L|<1$。取 $q(x)=\max_jp_{\lambda_j}(x)/\varepsilon_j$。若 $q(x)>0$，把 $x$ 缩放到 $q<1$ 的邻域内，得到 $|L(x)|\le q(x)$。若 $q(x)=0$，则任意倍数 $tx$ 都在该邻域内，故 $L(x)=0$。逆向由估计直接得到连续性。
\end{proof}
因此，无限多个半范数描述空间的拓扑，一个连续泛函却只需要有限多个半范数控制。

\section{固定紧集上的测试空间}
对 $K\Subset\Omega$ 令
\[
 \Test_K=\{\varphi\in C_c^\infty(\Omega):\supp\varphi\subset K\},\qquad
 p_{K,m}(\varphi)=\max_{|\alpha|\le m}\sup_{x\in K}|\partial^\alpha\varphi(x)|.
\]
若 $K$ 没有足够的内部，$\Test_K$ 可能为零；实际构造穷竭时取具有内部的紧集。由于支集固定，导数在 $K$ 外为零，上式也可在 $\Omega$ 上取上确界。
\begin{proposition}[Fréchet 性]
$\Test_K$ 是完备可度量的局部凸空间。其拓扑可由平移不变距离
\[
 d(\varphi,\psi)=\sum_{m=0}^\infty2^{-m-1}
 \frac{p_{K,m}(\varphi-\psi)}{1+p_{K,m}(\varphi-\psi)}
\]
描述。
\end{proposition}
\begin{proof}
该距离的收敛等价于各个 $p_{K,m}$ 收敛。若序列为 Cauchy，则每个导数序列在 $\Omega$ 上一致 Cauchy，因而有连续极限 $g_\alpha$。在任意包含于 $\Omega$ 的小矩形上，对原函数应用微积分基本定理并取极限，得到 $\partial_jg_\alpha=g_{\alpha+e_j}$。故 $g_0$ 光滑且导数为 $g_\alpha$。在 $\Omega\setminus K$ 上各函数均为零，极限也为零。于是 $g_0\in\Test_K$，且在所有半范数下收敛。
\end{proof}

\section{从固定紧集到全部紧支集}
选取紧穷竭 $K_1\subset\operatorname{int}K_2\subset\cdots\Subset\Omega$，其并为 $\Omega$，且任意紧集最终包含于某个 $K_j$。则
\[
 \Test(\Omega)=\bigcup_{j=1}^\infty\Test_{K_j}.
\]
\begin{definition}[标准测试函数拓扑]
在上述并上取使所有包含映射 $\Test_{K_j}\to\Test(\Omega)$ 连续的最强局部凸拓扑，称为局部凸归纳极限拓扑。它不依赖所选穷竭。由于各层为 Fréchet 空间，且相邻层的包含映射给出闭子空间及其原拓扑，这里是严格 LF 空间。
\end{definition}
\begin{theorem}[测试空间的有界集与序列收敛]\label{thm:D-convergence}
集合 $B\subset\Test(\Omega)$ 有界，当且仅当它包含于某个 $\Test_K$，且对每个 $m$，$\sup_{\varphi\in B}p_{K,m}(\varphi)<\infty$。特别地，$\varphi_j\to\varphi$ 当且仅当存在共同紧集 $K\Subset\Omega$ 包含全部支集，并且对每个多重指标 $\alpha$，
\[
 \sup_K|\partial^\alpha(\varphi_j-\varphi)|\longrightarrow0.
\]
\end{theorem}
有界集结论引用标准严格 LF 空间定理；这是本章唯一需要较一般归纳极限理论的地方。序列推论的理由是：收敛序列连同极限在任意拓扑向量空间中有界，因而落在同一层；该层上的诱导拓扑恰为各阶一致收敛拓扑。反向由包含映射的连续性立即得到。有限多个初始项可以吸收到更大紧集中。

\begin{example}[支集逃向无穷远]
取非零 $\eta\in\Test(\R)$，$\varphi_j(x)=\eta(x-j)$。对每个固定紧区间，各阶导数最终为零，因此 $\varphi_j\to0$ 于 $C^\infty(\R)$ 的局部拓扑。但是它在 $\Test(\R)$ 中不趋于零，因为不存在共同紧支集。即使改成 $e^{-j^2}\eta(x-j)$，仍然不能满足测试空间的收敛条件。
\end{example}
\begin{pitfall}[序列判别的使用边界]
$\Test(\Omega)$ 的标准拓扑一般不可度量。上述序列收敛判别是其特殊结构的结论，不能由此推断任意映射的连续性都等价于序列连续。定义分布时应使用逐紧集连续性或有限阶估计。
\end{pitfall}

\section{光滑函数空间与 Schwartz 空间}
$\Smooth(\Omega)=C^\infty(\Omega)$ 由 $p_{K_j,m}$ 构成 Fréchet 拓扑，收敛只要求每个紧集上的各阶一致收敛，不要求共同紧支集。取切断函数 $\chi_j\in\Test(\Omega)$，在 $K_j$ 附近为 1，则 $\chi_jf\to f$ 于 $\Smooth$，故 $\Test$ 在 $\Smooth$ 中稠密。

\begin{definition}[Schwartz 空间]
$\Sch(\R^n)$ 是所有满足
\[
 q_{N,m}(\varphi)=\max_{|\alpha|\le m}\sup_{x\in\R^n}
 (1+|x|)^N|\partial^\alpha\varphi(x)|<\infty
 \quad(N,m\in\N_0)
\]
的光滑函数组成的空间，拓扑由这些半范数定义。它也可用 $\sup|x^\beta\partial^\alpha\varphi|$ 这族等价半范数描述。
\end{definition}
高斯函数 $e^{-|x|^2}$ 属于 $\Sch$，而 $1$ 不属于。各导数比任意负幂衰减；这比仅仅趋于零强得多。用与 $\Test_K$ 相似的导数极限论证可证 $\Sch$ 完备：各带权导数在加权一致范数下有极限，这些极限相容并满足所有权重估计。因此它是 Fréchet 空间。

包含映射 $\Test\to\Sch\to\Smooth$ 都连续。在固定紧集上，权重有界，给出第一种连续性；在固定紧集上去掉权重得到第二种。又取 $\chi\in\Test$ 在单位球附近为 1，令 $\chi_R(x)=\chi(x/R)$，则 $\chi_R\varphi\to\varphi$ 于 $\Sch$。证明时对每个导数用 Leibniz 公式，利用 $|x|\gtrsim R$ 的尾部衰减，并以更高的权重吸收切断函数的导数。因此 $\Test$ 在 $\Sch$ 中稠密。

\section{习题}
\begin{exercise}[零阶范数不够]
证明导数泛函 $L(\varphi)=\varphi'(0)$ 在 $\Test(\R)$ 连续，但无法在含零点内部的固定紧区间上由 $C\|\varphi\|_\infty$ 控制。
\end{exercise}
\begin{solution}
它由一阶半范数控制。取 $\eta\in\Test$ 在零点附近为 1，令 $\varphi_k(x)=\eta(x)\sin(kx)$，则 $\|\varphi_k\|_\infty\le\|\eta\|_\infty$，但 $\varphi_k'(0)=k$，故零阶控制不可能。
\end{solution}
\begin{exercise}[固定支集与快速衰减]
证明 $j^{-1}\eta(x-j)\to0$ 于 $\Smooth(\R)$，但不趋于零于 $\Sch(\R)$；设 $\eta$ 非零且紧支集。
\end{exercise}
\begin{solution}
局部收敛来自支集最终离开每个固定紧集。取 $y_0$ 使 $\eta(y_0)\ne0$，在 $x=j+y_0$ 处，二阶权重给出 $q_{2,0}\ge j^{-1}(1+|j+y_0|)^2|\eta(y_0)|\to\infty$，故不在 Schwartz 拓扑中收敛。
\end{solution}

\SetChapterInfo{subtitle={有限阶估计、支集与对偶收敛},keywords={连续对偶; 分布阶数; 局部结构},prereq={第 2 章; Hahn--Banach 定理},goals={用估计判定分布并理解奇异支集}}
\chapter{分布空间的定义与结构}\label{ch:distributions}
\ChapterSource{对应施瓦尔茨 \cite{schwartz} 第 I 章 §§2--3、第 III 章 §§3、6--7、10。尤其参见原书印刷第 10、49、61、70 页。}

\section{连续对偶与有限阶判别}
\begin{definition}[分布]
开集 $\Omega$ 上的分布是 $\Test(\Omega)$ 上的连续复线性泛函。全体分布记为 $\Dist(\Omega)$。泛函 $T$ 在各层 $\Test_K$ 上连续，等价于它在归纳极限上连续。
\end{definition}
\begin{theorem}[分布的局部有限阶判别]\label{thm:finite-order}
线性泛函 $T:\Test(\Omega)\to\mathbb C$ 为分布，当且仅当对每个 $K\Subset\Omega$，存在整数 $m_K\ge0$ 和常数 $C_K$，使
\begin{equation}\label{eq:finite-order}
 |\Pair{T}{\varphi}|\le C_K\max_{|\alpha|\le m_K}\sup_K|\partial^\alpha\varphi|,
 \qquad\varphi\in\Test_K.
\end{equation}
\end{theorem}
\begin{proof}
在 $\Test_K$ 上用命题 \ref{prop:seminorm-est}。因半范数随 $m$ 递增，有限多个半范数可合并为最大的一个。反过来，估计保证每个限制映射连续，归纳极限的定义保证 $T$ 连续。
\end{proof}
正则分布由式 \eqref{eq:regular-bound} 满足零阶估计，Radon 测度 $\mu$ 也满足 $|\int\varphi\,\dd\mu|\le|\mu|(K)\|\varphi\|_\infty$。而 $\partial^\alpha\delta_a$ 的作用为 $(-1)^{|\alpha|}\partial^\alpha\varphi(a)$，需要导数控制。

称 $T$ 在 $K$ 上的阶不超过 $m$，是指式 \eqref{eq:finite-order} 可取该 $m$。若同一整数 $m$ 对所有紧集均可使用，则称其全局阶不超过 $m$，常数仍可依赖紧集。每个分布局部有限阶，但未必全局有限阶。例如
\[
 T=\sum_{j=1}^\infty\delta_j^{(j)}\quad\text{于 }\R
\]
是合法分布，因为每个紧集只遇到有限项；在整数 $j$ 附近的阶至少为 $j$。后者可用缩放测试函数 $\varepsilon^m\eta((x-j)/\varepsilon)$ 检验，其不超过 $m$ 阶的导数有界，而第 $j>m$ 阶点值可以发散。

\section{弱对偶拓扑与强对偶拓扑}
\begin{definition}[两种对偶拓扑]
弱对偶拓扑 $\sigma(\Dist,\Test)$ 由半范数 $T\mapsto|\Pair{T}{\varphi}|$ 定义。强对偶拓扑 $\beta(\Dist,\Test)$ 则由
\[
 P_B(T)=\sup_{\varphi\in B}|\Pair{T}{\varphi}|,
 \qquad B\subset\Test\text{ 为有界集}
\]
定义。强收敛要求在每个有界测试函数集上配对一致收敛。
\end{definition}
本书写 $T_j\to T$ 时默认弱对偶收敛。这给出容易使用的判别：每个固定探针的读数都收敛。强对偶拓扑通常更精细，定义也不同；不能因某些序列有相同极限就把两种拓扑混为一谈。
\begin{theorem}[逐点极限仍是分布]\label{thm:distribution-limit}
若 $T_j\in\Dist(\Omega)$ 且对每个 $\varphi\in\Test(\Omega)$，$\Pair{T_j}{\varphi}$ 有有限极限，则该极限定义一个分布 $T$。此外，对每个 $K$，存在共同的 $C_K,m_K$ 使式 \eqref{eq:finite-order} 对所有 $T_j$ 成立。
\end{theorem}
\begin{proof}
极限显然线性。固定 $K$，各泛函 $T_j|_{\Test_K}$ 逐点有界。$\Test_K$ 为 Fréchet 空间，故可用一致有界原理得到等度连续性。具体地，闭集
\[
 A_r=\{\varphi:\sup_j|\Pair{T_j}{\varphi}|\le r\}\quad(r\in\N)
\]
覆盖 $\Test_K$。Baire 定理保证某个 $A_r$ 有内点；相减给出零邻域上的共同控制，再缩放得到共同的有限阶估计。对 $j$ 取极限，该估计也控制 $T$，故 $T$ 是分布。
\end{proof}
如果同时 $\varphi_j\to\varphi$ 于 $\Test$，则 $\Pair{T_j}{\varphi_j}\to\Pair{T}{\varphi}$。因为
\[
 \Pair{T_j}{\varphi_j}-\Pair{T}{\varphi}
 =\Pair{T_j}{\varphi_j-\varphi}+\Pair{T_j-T}{\varphi},
\]
第一项由共同有限阶估计趋于零，第二项由弱收敛趋于零。

\begin{remark}[关于序列的进一步结论]
在这里弱收敛的分布序列也强收敛。理由是：每个有界 $B\subset\Test$ 落在固定层，并由各阶导数界、Arzelà--Ascoli 定理和对角抽取得到相对紧性。上面的等度连续性配合有限网格覆盖，将逐点收敛提升为在 $B$ 上一致收敛。这是测试空间的 Montel 性所带来的特殊现象；关于一般网或一般对偶空间，不能照搬该结论。
\end{remark}

\section{限制、粘合与支集}
对开集 $U\subset\Omega$，测试函数零延拓给出 $\Test(U)\to\Test(\Omega)$ 的连续映射，故可定义 $T|_U$。若一族局部分布在重叠处一致，可借助从属于开覆盖的光滑单位分解进行粘合：$\Pair{T}{\varphi}=\sum_i\Pair{T_i}{\chi_i\varphi}$。因为 $\varphi$ 紧支集，局部有限的和实际只含有限项；重叠一致性保证定义不依赖单位分解。
\begin{definition}[支集与奇异支集]
若 $T|_U=0$，称 $T$ 在 $U$ 上为零。$\supp T$ 是所有这样的开集之并在 $\Omega$ 中的补集。奇异支集 $\operatorname{sing\,supp}T$ 是 $T$ 不能局部由光滑函数表示的点集。
\end{definition}
两者均为相对闭集。$\supp\delta_a=\operatorname{sing\,supp}\delta_a=\{a\}$。对 Heaviside 分布，$\supp H=[0,\infty)$，而 $\operatorname{sing\,supp}H=\{0\}$。任何非零光滑函数都可以有非空支集而奇异支集为空。
若 $\supp\varphi$ 与 $\supp T$ 不交，单位分解将 $\varphi$ 拆成有限个支于 $T$ 为零的开集中的函数，故 $\Pair{T}{\varphi}=0$。

\section{紧支集分布与光滑函数的对偶}
\begin{theorem}[$\Smooth'$ 与紧支集分布]\label{thm:Eprime}
$\Smooth(\Omega)$ 的连续对偶可自然识别为 $\Omega$ 上的紧支集分布。任意紧支集分布都有有限全局阶。
\end{theorem}
\begin{proof}
若 $L$ 在 $\Smooth$ 上连续，则存在紧集 $K$ 和有限阶 $m$ 使 $|L(f)|\le Cp_{K,m}(f)$。限制到 $\Test$ 得分布，且所有支于 $\Omega\setminus K$ 的探针均被消去，因此支集包含于 $K$。

反向，若 $T$ 支于紧集 $A$，选 $\chi\in\Test(\Omega)$ 在 $A$ 的邻域为 1，定义 $\Pair{T}{f}=\Pair{T}{\chi f}$。两种切断函数的差在 $A$ 的邻域为零，故定义与切断无关。对 $\chi f$ 使用固定紧集有限阶估计和 Leibniz 公式，可得 $|\Pair{T}{f}|\le Cp_{K,m}(f)$。该估计也给出 $T$ 的有限全局阶。唯一性来自 $\Test$ 在 $\Smooth$ 中稠密。
\end{proof}

\section{点支集分布的结构}
\begin{theorem}[有限点射流]\label{thm:point-support}
若 $\supp T\subset\{a\}$，则存在某个整数 $m$ 和唯一一族常数 $c_\alpha$，使
\begin{equation}\label{eq:point-support}
 T=\sum_{|\alpha|\le m}c_\alpha\partial^\alpha\delta_a.
\end{equation}
\end{theorem}
\begin{proof}
平移令 $a=0$。由紧支集性取共同阶 $m$。若 $\partial^\alpha\varphi(0)=0$ 对所有 $|\alpha|\le m$ 成立，取 $\chi_\varepsilon(x)=\chi(x/\varepsilon)$ 在原点附近为 1，则 $T(\varphi)=T(\chi_\varepsilon\varphi)$。Taylor 余项给出在 $|x|\le C\varepsilon$ 上
\[
 \partial^\gamma\varphi(x)=o(\varepsilon^{m-|\gamma|})\quad(|\gamma|\le m).
\]
Leibniz 公式中切断函数的 $\beta$ 阶导数产生 $\varepsilon^{-|\beta|}$，故 $p_{K,m}(\chi_\varepsilon\varphi)=o(1)$。有限阶估计推出 $T(\varphi)=0$。

因此 $T$ 只依赖有限个数 $\partial^\alpha\varphi(0)$，在这个有限维射流空间上是线性泛函，故可写成 $\sum b_\alpha\partial^\alpha\varphi(0)$。令 $c_\alpha=(-1)^{|\alpha|}b_\alpha$ 得结论。用 $\chi(x)x^\beta/\beta!$ 分别提取射流坐标，得到系数唯一性。
\end{proof}

\section{局部结构：测度导数的有限和}
\begin{proposition}[测度表示]\label{prop:measure-structure}
设 $U\Subset\Omega$，则存在有限阶 $m$ 和有限复 Radon 测度 $\mu_\alpha$，使
\[
 T|_U=\sum_{|\alpha|\le m}\partial^\alpha\mu_\alpha.
\]
这里各测度可以选为支于同一紧集，并在 $U$ 上解释等式。
\end{proposition}
\begin{proof}
取 $\chi\in\Test(\Omega)$ 在 $\overline U$ 附近为 1，令 $S=\chi T$。取包含 $\supp\chi$ 的紧集 $K$，由前述结论可得
$|S(\varphi)|\le C\max_{|\alpha|\le m}\|\partial^\alpha\varphi\|_{C(K)}$。
把 $\varphi$ 映入有限乘积 $\prod_{|\alpha|\le m}C(K)$ 的射流向量。估计保证 $S$ 在像上是有界泛函。Hahn--Banach 定理将它延拓到整个乘积，Riesz 表示定理将各坐标泛函表示为测度 $\nu_\alpha$。于是
\[
 S(\varphi)=\sum_{|\alpha|\le m}\int_K\partial^\alpha\varphi\,\dd\nu_\alpha.
\]
令 $\mu_\alpha=(-1)^{|\alpha|}\nu_\alpha$ 得到分布导数表示；在 $U$ 上 $S=T$。
\end{proof}
这说明分布的奇异性仍具有有限阶的局部控制。施瓦尔茨还证明了局部分布可表示为连续函数的有限阶导数；这里证明的测度版本已足以支持本稿的结构讨论，不把更强版本作为已证明结论使用。

\section{习题}
\begin{exercise}[零阶与正性]
说明每个零阶分布是局部 Radon 测度。证明若分布 $T$ 对每个非负实测试函数有 $T(\varphi)\ge0$，则 $T$ 为正 Radon 测度。
\end{exercise}
\begin{solution}
零阶估计使泛函在每个固定紧集上按一致范数延拓，Riesz 表示得到局部测度；不同局部表示由测试函数稠密性唯一而相容。正性情形，固定 $K$，选 $\chi\ge0$ 且在 $K$ 上为 1。实测试函数满足 $-\|\varphi\|_\infty\chi\le\varphi\le\|\varphi\|_\infty\chi$，因此 $|T(\varphi)|\le T(\chi)\|\varphi\|_\infty$。对复函数分解实虚部得到零阶估计；表示测度的正性再由非负探针检验。
\end{solution}
\begin{exercise}[无穷远的点质量]
在 $\R$ 上证明 $\delta_j\to0$ 于 $\Dist$，并说明 $e^{j^2}\delta_j$ 也有相同弱极限。
\end{exercise}
\begin{solution}
对固定紧支集 $\varphi$，当 $j$ 足够大时 $\varphi(j)=0$，因此两个序列的配对都最终为零。这不与有限阶估计矛盾：每个固定紧集只遇到有限多个 $j$。
\end{solution}

\SetChapterInfo{subtitle={转置原理与合法运算},keywords={分布导数; 光滑乘子; 主值},prereq={分布有限阶估计},goals={证明运算连续并处理跳跃与奇性}}
\chapter{分布的导数与基本运算}\label{ch:operations}
\ChapterSource{对应施瓦尔茨 \cite{schwartz} 第 II 章、第 III 章 §§4--5 及第 V 章。主值和跳跃公式以下从定义直接验证。}

\section{转置原理}
若 $A:\Test(\Omega_1)\to\Test(\Omega_2)$ 连续线性，则
\[
 \Pair{{}^tA T}{\varphi}=\Pair{T}{A\varphi}
\]
定义连续对偶间的线性映射。若 $T_j\to T$，则 $\Pair{{}^tAT_j}{\varphi}\to\Pair{{}^tAT}{\varphi}$，故转置对弱收敛连续。这是分布运算的统一来源。要定义对象侧的运算，先确定探针侧的运算，并检查它保持测试空间和连续性。

\section{任意阶求导及其连续性}
\begin{definition}[分布导数]
对多重指标 $\alpha$，定义
\begin{equation}\label{eq:distribution-derivative}
 \Pair{\partial^\alpha T}{\varphi}=(-1)^{|\alpha|}\Pair{T}{\partial^\alpha\varphi}.
\end{equation}
\end{definition}
求导不扩大测试函数的支集，且 $p_{K,m}(\partial^\alpha\varphi)\le p_{K,m+|\alpha|}(\varphi)$。因此导数仍是分布，并满足
\[
 \partial^\alpha\partial^\beta T=\partial^{\alpha+\beta}T,
 \quad\supp(\partial^\alpha T)\subset\supp T,
 \quad T_j\to T\Rightarrow\partial^\alpha T_j\to\partial^\alpha T.
\]
若 $f\in C^{|\alpha|}$，反复分部积分得到 $\partial^\alpha T_f=T_{\partial^\alpha f}$。如果 $f$ 只是局部可积，则导数存在但未必仍由函数表示。

\begin{proposition}[导数为零的分布]\label{prop:zero-derivative}
在连通开区间 $I\subset\R$ 上，若 $T'=0$，则 $T$ 是常数分布。
\end{proposition}
\begin{proof}
若 $\varphi\in\Test(I)$ 且 $\int\varphi=0$，则其原函数 $\psi(x)=\int_{-\infty}^x\widetilde\varphi(t)\,\dd t$ 支于 $I$ 内某个紧区间，属于 $\Test(I)$。因此 $T(\varphi)=T(\psi')=-T'(\psi)=0$。选 $\eta\in\Test(I)$ 满足 $\int\eta=1$，把任意探针写成
\[
 \varphi=\left(\int\varphi\right)\eta+
 \left[\varphi-\left(\int\varphi\right)\eta\right].
\]
方括号积分为零，故 $T(\varphi)=T(\eta)\int\varphi$。
\end{proof}
由此归纳得，在区间上 $T^{(m)}=0$ 当且仅当 $T$ 为次数不超过 $m-1$ 的多项式分布。高维连通开集上，若所有一阶偏导都为零，类似的局部论证与粘合也得到常数分布。

\section{光滑乘子与 Leibniz 公式}
\begin{definition}[与光滑函数相乘]
对 $a\in C^\infty(\Omega)$，定义 $\Pair{aT}{\varphi}=\Pair{T}{a\varphi}$。
\end{definition}
在固定紧集上，Leibniz 公式给出 $p_{K,m}(a\varphi)\le C_{a,K,m}p_{K,m}(\varphi)$，故乘法合法；并且 $\supp(aT)\subset\supp a\cap\supp T$。
\begin{proposition}[乘积求导]
对光滑 $a$ 和分布 $T$，
\[
 \partial_j(aT)=(\partial_ja)T+a\partial_jT,
 \qquad
 \partial^\alpha(aT)=\sum_{\beta\le\alpha}\binom\alpha\beta
 (\partial^\beta a)\partial^{\alpha-\beta}T.
\]
\end{proposition}
\begin{proof}
$-T(a\partial_j\varphi)=-T(\partial_j(a\varphi))+T((\partial_ja)\varphi)$，正是右端对 $\varphi$ 的作用。高阶公式由归纳得到。
\end{proof}
特别地，在一维
\[
 a\delta_0=a(0)\delta_0,\qquad
 a\delta_0'=a(0)\delta_0'-a'(0)\delta_0.
\]
例如 $x\delta_0=0$，$x\delta_0'=-\delta_0$。不能把 $a\delta_0'$ 简化为 $a(0)\delta_0'$，因为导数会读取乘子的导数。

\section{跳跃函数与界面项}
\begin{proposition}[一维跳跃公式]
设 $f$ 在 $a$ 两侧分片 $C^1$，且有有限单侧极限。令 $f'_{\mathrm{reg}}$ 为两侧的普通导数组成的局部可积函数，则
\begin{equation}\label{eq:jump}
 (T_f)'=T_{f'_{\mathrm{reg}}}+[f]_a\delta_a,
 \qquad [f]_a=f(a+)-f(a-).
\end{equation}
\end{proposition}
\begin{proof}
把 $-\int f\varphi'$ 分成 $(-\infty,a)$ 与 $(a,\infty)$ 两部分，在探针支集范围内分别分部积分。左右边界项之和为 $(f(a+)-f(a-))\varphi(a)$，其余积分为 $\int f'_{\mathrm{reg}}\varphi$。
\end{proof}
若两侧分片 $C^2$，还有 $f''=f''_{\mathrm{reg}}+[f']_a\delta_a+[f]_a\delta_a'$。多界面时把各跳跃项求和。

在高维，若 $D$ 有光滑边界，外法向为 $\nu$，则散度定理给出
\[
 \Pair{\partial_j\mathbf1_D}{\varphi}
 =-\int_D\partial_j\varphi\,\dd x
 =-\int_{\partial D}\varphi\nu_j\,\dd S.
\]
故 $\nabla\mathbf1_D=-\nu\,\delta_{\partial D}$。这解释了守恒律和界面方程中表面质量的出现，也固定了法向符号。

\section{平移、缩放与坐标变换}
定义平移 $(\tau_hf)(x)=f(x-h)$，对分布采用
\[
 \Pair{\tau_hT}{\varphi}=\Pair{T}{\varphi(\cdot+h)}.
\]
于是 $\tau_h\delta_a=\delta_{a+h}$。对可逆线性变换 $A$，$f(Ax)$ 的分布推广为
\[
 \Pair{T(A\cdot)}{\varphi}=|\det A|^{-1}\Pair{T}{\varphi(A^{-1}\cdot)}.
\]
由此 $\delta_0(A\cdot)=|\det A|^{-1}\delta_0$。这里的 Jacobian 不能遗漏。
更一般地，对光滑微分同胚 $F:U\to V$，
\[
 \Pair{F^*T}{\varphi}=\Pair{T}{
 (\varphi\circ F^{-1})|\det DF^{-1}|}.
\]
微分同胚把紧集送到紧集，链式法则给出各阶半范数估计，因此拉回合法。对任意非可逆光滑映射，拉回分布则需要额外条件，本稿不作一般定义。

\section{主值分布与奇性消去}
\begin{definition}[Cauchy 主值]
在 $\R$ 上定义
\[
 \Pair{\pv\frac1x}{\varphi}
 =\lim_{\varepsilon\downarrow0}\int_{|x|>\varepsilon}\frac{\varphi(x)}x\,\dd x.
\]
\end{definition}
取 $R$ 包含探针支集，利用对称性将该式写成
\[
 \int_0^R\frac{\varphi(x)-\varphi(-x)}x\,\dd x.
\]
均值定理使绝对值不超过 $2R\|\varphi'\|_\infty$，所以极限存在且定义分布。又因 $x\varphi$ 消去分母，$x\,\pv(1/x)=1$。
\begin{proposition}[对数的分布导数]
局部可积函数 $\log|x|$ 满足 $(\log|x|)'=\pv(1/x)$。
\end{proposition}
\begin{proof}
在 $(-\infty,-\varepsilon)$ 与 $(\varepsilon,\infty)$ 分部积分，有
\[
 -\int_{|x|>\varepsilon}\log|x|\,\varphi'(x)\,\dd x
 =\int_{|x|>\varepsilon}\frac{\varphi(x)}x\,\dd x
 +\log\varepsilon[\varphi(\varepsilon)-\varphi(-\varepsilon)].
\]
最后一项为 $O(\varepsilon|\log\varepsilon|)$，而被删去区间上的左端积分趋于零，故得结论。
\end{proof}
\begin{pitfall}[为什么不能任意相乘]
用 $\rho_\varepsilon$ 逼近 $\delta_0$ 时，
\[
 \int\rho_\varepsilon(x)^2\varphi(x)\,\dd x
 =\varepsilon^{-1}\int\rho(y)^2\varphi(\varepsilon y)\,\dd y
\]
在 $\varphi(0)>0$ 时发散。因此该正则化不能定义 $\delta_0^2$。这个例子展示了直接乘法的障碍；它本身不代替关于各种广义乘法体系的完整不可能性定理。
\end{pitfall}

\section{习题}
\begin{exercise}[除法的不唯一性]
求 $\R$ 上方程 $xT=1$ 的全部分布解。
\end{exercise}
\begin{solution}
一个特解为 $\pv(1/x)$。若 $xS=0$，则在避开零点的开集上可乘以光滑函数 $1/x$，得到 $S=0$，故 $S$ 为点支集分布。由定理 \ref{thm:point-support} 写成有限个 $\delta_0^{(k)}$ 的和，再用 $x\delta_0^{(k)}=-k\delta_0^{(k-1)}$，可知除 $k=0$ 外系数皆零。故 $T=\pv(1/x)+c\delta_0$。
\end{solution}
\begin{exercise}[阶数与支集]
证明 $\supp(\partial_jT)\subset\supp T$，并给出严格包含的例子。
\end{exercise}
\begin{solution}
若 $T$ 在开集 $U$ 上为零，$\partial_j\varphi$ 仍支于 $U$，故导数也在 $U$ 上为零。严格包含的例子是常数分布 $T=1$：其支集为整个空间，而导数为零，支集为空。
\end{solution}

\SetChapterInfo{subtitle={平滑逼近与运算的支集条件},keywords={卷积; 正则化; 张量积},prereq={分布导数; 紧支集分布},goals={说明卷积何时有定义并构造光滑逼近}}
\chapter{卷积、张量积与正则化}\label{ch:convolution}
\ChapterSource{对应施瓦尔茨 \cite{schwartz} 第 IV 章及第 VI 章，特别是两个分布卷积的定义与紧支集条件（原书印刷第 111 页）。}

\section{分布与测试函数的卷积}
在本章先取全空间 $\R^n$。对 $T\in\Dist(\R^n)$ 与 $\rho\in\Test(\R^n)$，定义
\begin{equation}\label{eq:convolution-function}
 (T*\rho)(x)=\Pair{T_y}{\rho(x-y)}.
\end{equation}
固定 $x$ 时右侧探针关于 $y$ 紧支撑，因此有意义。若 $x$ 在某个紧集变化，全部探针的支集仍落在一个共同紧集中；对 $x$ 求差商时，差商在该层的所有半范数下收敛。这允许把任意阶导数移入配对。
\begin{proposition}[卷积平滑性]\label{prop:smoothing}
$T*\rho\in C^\infty(\R^n)$，且
\[
 \partial^\alpha(T*\rho)=T*(\partial^\alpha\rho)
 =(\partial^\alpha T)*\rho.
\]
此外 $\supp(T*\rho)\subset\supp T+\supp\rho$。
\end{proposition}
\begin{proof}
第一等式由上述测试空间中的差商极限得到。第二等式中 $\partial_y^\alpha\rho(x-y)=(-1)^{|\alpha|}\partial_x^\alpha\rho(x-y)$，与分布导数的负号抵消。如果 $x$ 不在支集之和中，$\rho(x-\cdot)$ 的支集与 $\supp T$ 不交，配对为零。由于 $\supp\rho$ 紧，闭集与它的和仍闭，故得到支集包含关系。
\end{proof}
例如 $\delta_a*\rho=\rho(\cdot-a)$，$\delta_0'*\rho=\rho'$。平滑过程保留分布信息，但将奇性转成窄峰和高速变化。

\section{近似恒等与局部正则化}
\begin{theorem}[平滑逼近]\label{thm:mollification}
取第 \ref{ch:motivation} 章的 $\rho_\varepsilon$，则
\[
 T*\rho_\varepsilon\longrightarrow T\quad\text{于 }\Dist(\R^n).
\]
这里将左侧光滑函数识别为其正则分布。
\end{theorem}
\begin{proof}
令 $\Reflected\rho(x)=\rho(-x)$。对测试函数 $\varphi$，将紧支集参数积分移入连续配对，得
\[
 \Pair{T*\rho_\varepsilon}{\varphi}
 =\Pair{T}{\Reflected\rho_\varepsilon*\varphi}.
\]
可以先用 Riemann 和验证这一步：参数范围紧，探针族在固定 $\Test_K$ 层中连续，故积分极限在每个半范数下成立。
而
\[
 (\Reflected\rho_\varepsilon*\varphi)(y)
 =\int\rho(z)\varphi(y+\varepsilon z)\,\dd z.
\]
对每个 $\alpha$，$\partial^\alpha\varphi$ 一致连续，所以上式的各阶导数一致趋于 $\partial^\alpha\varphi$。当 $0<\varepsilon\le1$ 时支集落在固定紧集 $\supp\varphi+\overline B(0,1)$。故探针收敛发生在 $\Test$，最后利用 $T$ 的连续性。
\end{proof}
一般开集上，在 $\Omega_\varepsilon=\{x:\operatorname{dist}(x,\R^n\setminus\Omega)>\varepsilon\}$ 可用同样配对定义 $T*\rho_\varepsilon$，并在每个固定内紧集上得到逼近。另一种做法是先用切断函数将局部分布化为全空间紧支集分布，再卷积；不同切断在较小区域上给出相同局部极限。

若 $f\in L^p_\Loc$，$1\le p<\infty$，则 $f*\rho_\varepsilon\to f$ 于 $L^p$ 的局部范数。证明先局部切断，再用
\[
 \|f*\rho_\varepsilon-f\|_p
 \le\int|\rho(z)|\|f(\cdot-\varepsilon z)-f\|_p\,\dd z
\]
和 $L^p$ 平移的连续性。对于 $p=\infty$，任意函数的范数收敛不成立，例如 $H$ 的跳跃无法由连续函数一致逼近。

\begin{corollary}[正则分布的唯一性]
若 $f\in L^1_\Loc(\Omega)$ 且 $T_f=0$，则 $f=0$ 几乎处处。
\end{corollary}
\begin{proof}
在内紧集上，所有足够小 $\varepsilon$ 的卷积都为零；它们又在 $L^1$ 局部收敛到 $f$，故 $f$ 在每个内紧集上几乎处处为零。取可数穷竭即得全域结论。
\end{proof}

\section{张量积：分离不同变量}
对 $S\in\Dist(\R^m)$、$T\in\Dist(\R^n)$，定义
\[
 \Pair{S\otimes T}{\Phi}
 =\Pair{S_x}{\Pair{T_y}{\Phi(x,y)}},
 \qquad\Phi\in\Test(\R^{m+n}).
\]
内层配对关于 $x$ 为光滑紧支集函数，其支集包含 $\supp\Phi$ 在 $x$ 方向的投影。固定矩形紧集上，先用 $T$ 的有限阶估计，再用 $S$ 的有限阶估计，可得
\[
 |(S\otimes T)(\Phi)|\le C
 \max_{|\alpha|\le r,|\beta|\le s}
 \sup|\partial_x^\alpha\partial_y^\beta\Phi|.
\]
因此张量积是分布。两个配对次序相同：在固定矩形内，有限和形式 $\sum_jf_j(x)g_j(y)$ 在光滑测试函数拓扑中稠密，可用切断后的 Fourier 级数逼近证明；在这些有限和上两种次序显然一致，连续性给出一般结论。这一步是分布版本的 Fubini 原理。
特别地
\[
 (S\otimes T)(f\otimes g)=S(f)T(g),\qquad
 \supp(S\otimes T)\subset\supp S\times\supp T.
\]
张量积在不同变量中相乘，不是同一点上任意两个分布的乘积。比如 $\delta_0\otimes\delta_0$ 有定义，但这并不定义一维的 $\delta_0^2$。

\section{两个分布的卷积}
想写 $(S*T)(\varphi)=(S\otimes T)(\varphi(x+y))$，但 $\varphi(x+y)$ 通常不紧支撑：沿 $x+y$ 固定的方向可延伸到无穷远。因此必须补上支集条件。
\begin{theorem}[一个因子紧支撑的卷积]\label{thm:compact-convolution}
若 $S,T\in\Dist(\R^n)$ 且至少一个具有紧支集，则 $S*T\in\Dist(\R^n)$ 有定义，满足
\begin{align*}
 S*T&=T*S,\\
 \partial^\alpha(S*T)&=(\partial^\alpha S)*T=S*(\partial^\alpha T),\\
 \supp(S*T)&\subset\supp S+\supp T.
\end{align*}
\end{theorem}
\begin{proof}
设 $A=\supp S$ 紧，$K=\supp\varphi$。当 $x\in A$ 且 $x+y\in K$ 时，有 $y\in K-A$，也是紧集。于是相关支集交落在紧集 $A\times(K-A)$ 内，可以选紧支集光滑切断函数 $\theta(x,y)$ 在该紧集的邻域为 1，令
\[
 (S*T)(\varphi)=(S\otimes T)(\theta(x,y)\varphi(x+y)).
\]
两种切断的差与 $\supp(S\otimes T)$ 及 $\supp\varphi(x+y)$ 的相关部分不相交，故作用相同。固定 $K$ 时可统一选切断，链式法则与张量积的有限阶估计保证连续性。交换性来自张量积配对次序可交换；导数公式来自分部积分；支集包含由相关交集为空时配对为零得到。
\end{proof}
若三重加法 $(x,y,z)\mapsto x+y+z$ 在三个支集的乘积上适当，则可用同一个三重张量积配对解释两种括号，因而结合律成立。例如三因子中至少两个紧支撑，或一维三因子均支于 $[0,\infty)$，都满足该条件。不能只凭几个形式卷积符号分别出现，就省去共同三重配对的合法性检查。

更一般地，只要加法映射 $(x,y)\mapsto x+y$ 在 $\supp S\times\supp T$ 上是适当映射，即每个紧集的原像与该支集的交是紧的，同样定义成立。一维中两个分布都支于 $[0,\infty)$ 就满足此条件。这对因果基本解尤为重要。
\begin{example}[半轴上的因果卷积]
$H*H=xH(x)$，因为在 $x\ge0$ 时积分范围是 $0\le y\le x$，积分值为 $x$；在 $x<0$ 时为零。对其求导得到 $(H*H)'=H$，与 $H'*H=\delta_0*H=H$ 一致。
\end{example}

\section{习题}
\begin{exercise}[卷积单位]
证明 $\delta_0*T=T$ 与 $\delta_a*T=\tau_aT$。解释这些等式为何对任意分布都有意义。
\end{exercise}
\begin{solution}
点质量紧支撑，故卷积存在。对 $\varphi$，先对点质量配对，得到 $T(\varphi(a+\cdot))$，正是平移的定义。取 $a=0$ 得单位元。
\end{solution}
\begin{exercise}[常数卷积的障碍]
说明两个常数分布 $1$ 的形式卷积在 $\R$ 上为何不能用普通积分定义；与 $1*\rho$ 比较。
\end{exercise}
\begin{solution}
形式值为 $\int_\R1\,\dd y=\infty$，加法在全支集上也不适当。若 $\rho\in\Test$，则第二因子紧支撑，$1*\rho=\int\rho$ 是合法的常数分布。这体现支集或衰减条件的作用。
\end{solution}

\SetChapterInfo{subtitle={Schwartz 对偶与微分的乘子化},keywords={温和分布; Fourier 自同构; Sobolev 空间},prereq={测试函数拓扑; 分布运算与卷积; 经典 Fourier 反演},goals={统一常数并使用对偶定义 Fourier 变换}}
\chapter{温和分布与傅里叶变换}\label{ch:fourier}
\ChapterSource{对应施瓦尔茨 \cite{schwartz} 第 VII 章 §§2--6，特别是原书印刷第 169、172 页。全章使用式 \eqref{eq:fourier-convention}，与原书频率归一化不同。}

\section{温和分布的增长控制}
\begin{definition}[温和分布]
Schwartz 空间 $\Sch(\R^n)$ 的连续线性对偶记为 $\Temp(\R^n)$。其元素称为温和分布，也称缓增分布。
\end{definition}
由有限半范数估计，$T\in\Temp$ 当且仅当存在 $N,m,C$ 使
\begin{equation}\label{eq:tempered-bound}
 |T(\varphi)|\le Cq_{N,m}(\varphi)\qquad(\varphi\in\Sch).
\end{equation}
$\Test\to\Sch$ 连续，故限制得到 $\Temp\to\Dist$；由于 $\Test$ 在 $\Sch$ 中稠密，这一映射单射。也就是说，温和分布是普通分布中的特殊一类，它们能连续地作用于非紧支撑但快速衰减的探针。

\begin{proposition}[基本例子]
以下对象都是温和分布：$L^p(\R^n)$ 函数（$1\le p\le\infty$）；具有至多多项式增长的局部可积函数；紧支集分布；$\delta_a$ 的任意阶导数。
\end{proposition}
\begin{proof}
对 $L^p$ 函数用 Hölder 不等式，再以足够高权重控制测试函数的 $L^{p'}$ 范数。若 $|f(x)|\le C(1+|x|)^M$ 几乎处处，取 $N>M+n$ 得
\[
 |\int f\varphi|\le Cq_{N,0}(\varphi)
 \int(1+|x|)^{M-N}\,\dd x<\infty.
\]
紧支集分布的有限阶估计由某个 $q_{0,m}$ 控制，故也温和。
\end{proof}
更一般，若 $\int|f(x)|(1+|x|)^{-N}\,\dd x<\infty$，则 $T_f$ 温和。对带符号或复值的振荡函数，这只是充分条件，不能误当成必要条件。

\begin{example}[普通分布并非都温和]
函数 $e^x$ 在 $\R$ 上局部可积，给出分布但不是温和分布。若它满足式 \eqref{eq:tempered-bound}，取非负非零 $\eta\in\Test$ 并令 $\eta_j(x)=\eta(x-j)$，则配对为 $e^j\int e^y\eta(y)\,\dd y$，而 $q_{N,m}(\eta_j)$ 至多按 $j^N$ 增长，矛盾。
\end{example}
求导和多项式乘法保持 $\Temp$。更一般，若光滑 $a$ 的每个导数均有某种多项式增长，则 $\varphi\mapsto a\varphi$ 连续作用于 $\Sch$，其转置给出温和分布乘法。

\section{Schwartz 空间上的傅里叶变换}
\begin{theorem}[Fourier 自同构]\label{thm:schwartz-fourier}
按式 \eqref{eq:fourier-convention} 定义的 $\FT$ 是 $\Sch(\R^n)$ 的连续线性自同构，其逆为
\[
 \FT^{-1}\psi(x)=(2\pi)^{-n}\int e^{ix\cdot\xi}\psi(\xi)\,\dd\xi.
\]
\end{theorem}
\begin{proof}
微分可移入积分，分部积分无边界项，故
\begin{equation}\label{eq:fourier-identities}
 \widehat{\partial_j\varphi}=i\xi_j\widehat\varphi,
 \qquad\widehat{x_j\varphi}=i\partial_{\xi_j}\widehat\varphi.
\end{equation}
任意 $\xi^\beta\partial_\xi^\alpha\widehat\varphi$ 因此可表示成 $x^\alpha\varphi$ 的有限阶导数的 Fourier 变换，差一个常数因子。其上确界由对应函数的 $L^1$ 范数控制，而该范数由有限个 $q_{N,m}(\varphi)$ 控制。故 $\FT$ 连续映入 $\Sch$。逆号积分同理连续映入 $\Sch$。使用经典 Schwartz 函数 Fourier 反演公式，两个映射互为逆。
\end{proof}
经典反演也可以用高斯截断验证：在逆积分中加入 $e^{-\varepsilon|\xi|^2}$，得到与高斯核的卷积；高斯核构成近似恒等，令 $\varepsilon\downarrow0$ 得到原函数。这里所需 Fubini 与控制收敛均由 Schwartz 衰减保证。

\section{温和分布的傅里叶变换}
\begin{definition}[对偶定义]
对 $T\in\Temp$，定义
\begin{equation}\label{eq:fourier-distribution}
 \Pair{\widehat T}{\varphi}=\Pair{T}{\widehat\varphi},\qquad\varphi\in\Sch.
\end{equation}
\end{definition}
注意右侧是探针的正向变换。这是因为核 $e^{-ix\cdot\xi}$ 在两个变量中对称，配对采用双线性约定。若 $T=T_f$ 且 $f\in L^1$，Fubini 定理给出该定义与普通 $\widehat f$ 一致。连续自同构的转置也是连续自同构，因此 $\widehat T\in\Temp$ 且 Fourier 变换可逆。
\begin{keyformula}[常数与基本对象]
\begin{align*}
 \widehat{\delta_0}&=1,&\widehat1&=(2\pi)^n\delta_0,\\
 \widehat{\delta_a}&=e^{-ia\cdot\xi},&
 \widehat{\partial^\alpha T}&=(i\xi)^\alpha\widehat T,\\
 \widehat{x_jT}&=i\partial_{\xi_j}\widehat T,&
 \widehat{\tau_hT}&=e^{-ih\cdot\xi}\widehat T.
\end{align*}
\end{keyformula}
例如 $\widehat{\delta_0}(\varphi)=\widehat\varphi(0)=\int\varphi$。又由反演 $\int\widehat\varphi=(2\pi)^n\varphi(0)$，得到常数的变换。对于导数，
\[
 \widehat{\partial_jT}(\varphi)
 =-T(\partial_{x_j}\widehat\varphi)
 =T(\widehat{i\xi_j\varphi})
 =\widehat T(i\xi_j\varphi).
\]
这直接验证 $i$ 与负号，无须靠记忆。

\begin{example}[高斯]
对 $t>0$，令
\[
 G_t(x)=(4\pi t)^{-n/2}e^{-|x|^2/(4t)}.
\]
一维高斯积分与张量积计算给出 $\widehat G_t(\xi)=e^{-t|\xi|^2}$。由于 $\int G_t=1$，$G_t\to\delta_0$ 于分布意义；变换侧则 $e^{-t|\xi|^2}\to1$ 于温和分布意义。
\end{example}

\Needspace{10\baselineskip}
\section{卷积定理的合法范围}
对 Schwartz 函数 $f,g$，Fubini 直接给出 $\widehat{f*g}=\widehat f\widehat g$。若 $S\in\CompactDist$、$T\in\Temp$，则 $S*T\in\Temp$，并且
\begin{equation}\label{eq:convolution-fourier}
 \widehat{S*T}=\widehat S\,\widehat T.
\end{equation}
紧支集分布的 Fourier 变换是光滑函数
\[
 \widehat S(\xi)=\Pair{S_x}{e^{-ix\cdot\xi}},
\]
这里以紧支集切断解释非紧支撑探针。对 $\xi$ 求导产生 $(-ix)^\alpha$；有限阶估计表明每个导数至多多项式增长。因此它是合法 Schwartz 乘子，右端有意义。卷积定理可由对探针的配对与第 \ref{ch:convolution} 章的切断定义验证。

任意两个温和分布的卷积未必存在，它们的 Fourier 变换之积也未必有定义。例如 $1*1$ 无定义，相应右侧形式上含 $\delta_0^2$。在允许的函数类中，乘积定理为 $\widehat{fg}=(2\pi)^{-n}\widehat f*\widehat g$；其中常数因子来自反演。

\section{Sobolev 空间：分布的 Hilbert 尺度}
\begin{definition}[Bessel--Sobolev 空间]
对 $s\in\R$，$H^s(\R^n)$ 由那些 Fourier 变换可用可测函数表示、且满足下列有限范数条件的 $u\in\Temp$ 组成：
\[
 \|u\|_{H^s}^2=(2\pi)^{-n}\int_{\R^n}
 (1+|\xi|^2)^s|\widehat u(\xi)|^2\,\dd\xi<\infty.
\]
\end{definition}
Plancherel 定理给出 $H^0=L^2$。对非负整数 $m$，权重 $(1+|\xi|^2)^m$ 与 $\sum_{|\alpha|\le m}|\xi^\alpha|^2$ 等价，故 $H^m$ 等价于所有不超过 $m$ 阶弱导数属于 $L^2$ 的空间。并有
\[
 \partial_j:H^s\to H^{s-1}\quad\text{连续},\qquad
 s_1\ge s_2\Rightarrow H^{s_1}\subset H^{s_2}.
\]
在本稿的双线性分布配对下，$H^{-s}$ 与 $H^s$ 的连续线性对偶的对应明确为
\[
 \langle u,v\rangle=(2\pi)^{-n}\int_{\R^n}
 \widehat u(\xi)\widehat v(-\xi)\,\dd\xi,
 \qquad |\langle u,v\rangle|\le\|u\|_{H^{-s}}\|v\|_{H^s}.
\]
先在 Schwartz 函数上由反演得到公式，再用加权 Cauchy--Schwarz 延拓。反向将 $v\mapsto\langle\xi\rangle^s\widehat v(\xi)$ 作为到 $L^2$ 的同构，对任意连续泛函用 Riesz 表示，再反射频率并共轭表示函数，即得到唯一 $u\in H^{-s}$，且泛函范数等于 $\|u\|_{H^{-s}}$。因此这里的识别是线性的；若改写成 Hilbert 内积形式，必须同时调整共轭位置。
\begin{example}[点质量的正则性门槛]
因 $\widehat{\delta_0}=1$，有
\[
 \delta_0\in H^s(\R^n)\quad\Longleftrightarrow\quad s<-\frac n2.
\]
积分在原点附近有限，在无穷远需 $\int_1^\infty r^{2s+n-1}\,\dd r<\infty$。类似地，$\partial^\alpha\delta_0\in H^s$ 当且仅当 $s<-n/2-|\alpha|$，可在球面上一块多项式不为零的区域给出必要性。
\end{example}

\section{从 Sobolev 范数到普通函数}\label{sec:topic-sobolev-embedding}
\begin{proposition}[全空间的连续嵌入]\label{prop:topic-sobolev-embedding}
若 $k\in\N_0$ 且 $s>k+n/2$，则 $H^s(\R^n)$ 中每个元素有唯一的 $C^k$ 代表元，各阶 $|\alpha|\le k$ 的导数有界且在无穷远趋于零，并满足
\[
 \max_{|\alpha|\le k}\|\partial^\alpha u\|_\infty
 \le C_{s,k,n}\|u\|_{H^s}.
\]
\end{proposition}
\begin{proof}
加权 Cauchy--Schwarz 不等式给出
\[
 \int |\xi|^{|\alpha|}|\widehat u(\xi)|\,\dd\xi
 \le \left(\int(1+|\xi|^2)^s|\widehat u|^2\,\dd\xi\right)^{1/2}
 \left(\int |\xi|^{2|\alpha|}(1+|\xi|^2)^{-s}\,\dd\xi\right)^{1/2}.
\]
最后的积分由 $s>|\alpha|+n/2$ 保证有限。因此 $u$ 及这些导数均由绝对收敛的逆 Fourier 积分表示；控制收敛保证连续性与逐阶微分，Riemann--Lebesgue 引理保证无穷远消失。该函数代表原分布，连续代表元的唯一性由几乎处处相等推出。
\end{proof}
这把分布意义的正则性变成点值和导数估计，与主线第六章 §7 的 Sobolev 引理衔接。严格不等号是本证明中积分收敛的条件，不能直接改为临界等号。区域上的嵌入还需要延拓或局部化；边界迹是另一类结论，见书后补充说明第 \pageref{sec:supp-trace} 页。

\Needspace{12\baselineskip}
\section{习题}
\begin{exercise}[高阶导数的有界性]
设 $u\in H^s(\R^n)$ 且 $s>1+n/2$。证明 $u$ 的连续代表元为全局 Lipschitz 函数，并估计其 Lipschitz 常数。
\end{exercise}
\begin{solution}
上面的嵌入给出 $u\in C^1$ 且 $\|\nabla u\|_\infty\le C\|u\|_{H^s}$。沿 $x$ 到 $y$ 的线段积分，有 $|u(x)-u(y)|\le\|\nabla u\|_\infty|x-y|$。
\end{solution}
\begin{exercise}[二次变换]
证明 $\FT^2T=(2\pi)^n\Reflected T$，其中 $\Reflected T(\varphi)=T(\varphi(-\cdot))$。
\end{exercise}
\begin{solution}
对探针用经典反演得到 $\FT^2\varphi(x)=(2\pi)^n\varphi(-x)$，故 $\FT^2T(\varphi)=T(\FT^2\varphi)=(2\pi)^n\Reflected T(\varphi)$。
\end{solution}
\begin{exercise}[一阶基本解的变换]
令 $E(x)=H(x)e^{-x}$。验证 $(\partial_x+1)E=\delta_0$，并求 $\widehat E$。
\end{exercise}
\begin{solution}
由光滑乘积公式 $E'=\delta_0-E$。又 $E\in L^1$，故 $\widehat E(\xi)=\int_0^\infty e^{-(1+i\xi)x}\,\dd x=(1+i\xi)^{-1}$。变换侧 $(1+i\xi)\widehat E=1$ 与原方程一致。
\end{solution}

\SetChapterInfo{subtitle={基本解、弱解与柯西问题},keywords={Poisson 方程; 热半群; 适定性},prereq={卷积与 Fourier 变换; Hilbert 空间},goals={明确解空间并区分存在、唯一和稳定}}
\chapter{广义函数与微分方程}\label{ch:pde}
\ChapterSource{基本解部分对应施瓦尔茨 \cite{schwartz} 第 VI 章 §10；柯西问题的视角对应盖尔芳特、希烙夫 \cite{gelfand} 第 II--III 章，特别是初值在广义函数意义下取极限的论述（印刷第 108 页）。本章选择可完整核验的模型方程展开。}

\section{分布解与算子的形式转置}
对光滑系数微分算子 $P=\sum_{|\alpha|\le m}a_\alpha(x)\partial^\alpha$，定义
\[
 Pu=\sum_{|\alpha|\le m}a_\alpha\partial^\alpha u,
 \qquad {}^tP\varphi=\sum_{|\alpha|\le m}(-1)^{|\alpha|}
 \partial^\alpha(a_\alpha\varphi).
\]
则 $\Pair{Pu}{\varphi}=\Pair{u}{{}^tP\varphi}$。这里是双线性配对的形式转置，若使用 Hilbert 内积讨论伴随算子，还要对系数取复共轭。
\begin{definition}[分布解]
$u\in\Dist(\Omega)$ 是 $Pu=f$ 的分布解，是指对每个 $\varphi\in\Test(\Omega)$ 有 $u({}^tP\varphi)=f(\varphi)$。
\end{definition}
它不要求 $u$ 具有普通导数。然而，分布解也不自动具有边界值；边界条件须在具有迹算子的函数空间中解释，或另行使用适合的分布框架。

\Needspace{10\baselineskip}
\section{基本解与卷积求解}
\begin{definition}[基本解]
常系数算子 $P(\partial)$ 的基本解是满足 $P(\partial)E=\delta_0$ 的分布 $E$。
\end{definition}
若 $f$ 紧支撑，则 $u=E*f$ 有定义，导数与卷积交换给出 $Pu=(PE)*f=\delta_0*f=f$。这构造了一个解；全部解还需加上任意齐次解。基本解也可相差齐次解，唯一性须附加辐射、因果或增长条件。
\begin{example}[一阶因果基本解]
对 $P=\partial_x+a$，$E_a=H(x)e^{-ax}$ 在分布意义满足 $PE_a=\delta_0$，$a$ 可为复数。若源 $f$ 与 $E_a$ 都支于右半轴，即使都不紧支撑，加法适当性也保证因果卷积合法。当 $\operatorname{Re}a<0$ 时 $E_a$ 指数增长，通常不是温和分布，但仍是普通分布。
\end{example}

\section{Laplace 算子的基本解}
设 $\omega_{n-1}$ 为单位球面面积。对于算子 $-\Delta$，取
\begin{equation}\label{eq:laplace-fundamental}
 E_n(x)=\begin{cases}
 \displaystyle\frac{|x|^{2-n}}{(n-2)\omega_{n-1}},&n\ge3,\\[4pt]
 \displaystyle-\frac1{2\pi}\log|x|,&n=2,\\[4pt]
 \displaystyle-\frac12|x|,&n=1.
 \end{cases}
\end{equation}
这些函数在原点附近局部可积，且至多多项式或对数增长，故定义温和分布。
\begin{theorem}[Newton 势基本解]
式 \eqref{eq:laplace-fundamental} 满足 $-\Delta E_n=\delta_0$。
\end{theorem}
\begin{proof}
$n=1$ 时来自 $|x|''=2\delta_0$。当 $n\ge2$ 时，$E_n$ 在原点外调和，且其径向导数满足
\[
 \partial_rE_n(r)=-\frac1{\omega_{n-1}r^{n-1}}.
\]
对 $\varphi\in\Test$，在 $B_R\setminus\overline B_\varepsilon$ 应用第二 Green 公式，外边界选在探针支集之外。内边界相对于穿孔区域的外法向为 $-e_r$，故
\[
 \int_{|x|>\varepsilon}E_n(-\Delta\varphi)\,\dd x
 =\int_{|x|=\varepsilon}
 [E_n\partial_r\varphi-\varphi\partial_rE_n]\,\dd S.
\]
第一项的绝对值为 $O(\varepsilon)$（$n\ge3$）或 $O(\varepsilon|\log\varepsilon|)$（$n=2$），趋于零。第二项等于球面上 $\varphi$ 的平均值，趋于 $\varphi(0)$。局部可积性允许补回小球中的积分，得到结论。
\end{proof}
例如三维 $E_3=1/(4\pi|x|)$，故紧支撑源的 Newton 势 $u=E_3*f$ 满足 $-\Delta u=f$。无限空间的基本解并不直接给出有界区域的 Dirichlet Green 函数，后者还需要修正边界值。

\section{弱解与能量空间}
在有界 Lipschitz 区域 $\Omega$，令 $H_0^1(\Omega)$ 为 $C_c^\infty(\Omega)$ 在 $H^1$ 范数下的闭包。它提供零边界迹的解释。定义 $H^{-1}(\Omega)$ 为 $H_0^1(\Omega)$ 的连续对偶。
\begin{definition}[Poisson 方程的能量弱解]
对 $f\in H^{-1}(\Omega)$，寻找 $u\in H_0^1(\Omega)$ 使
\begin{equation}\label{eq:weak-poisson}
 \int_\Omega\nabla u\cdot\overline{\nabla v}\,\dd x
 =\Pair{f}{\overline v},\qquad v\in H_0^1(\Omega).
\end{equation}
右侧配对按双线性分布约定写成 $f(\overline v)$，故对 $v$ 共轭线性。
\end{definition}
Poincaré 不等式使 $\|u\|_{H_0^1}$ 可等价地取为 $\|\nabla u\|_2$。双线性型在复空间中是 sesquilinear 型，连续且强制。Lax--Milgram 定理给出唯一解，并有
\[
 \|u\|_{H_0^1}\le C\|f\|_{H^{-1}}.
\]
这同时说明解对数据稳定。取 $v=\overline\varphi$ 可得 $\int\nabla u\cdot\nabla\varphi=f(\varphi)$，即 $-\Delta u=f$ 于分布意义。反过来，若 $u\in H_0^1$ 满足分布方程，则先对光滑紧支集探针分部积分，再由稠密性延拓，得到能量弱解。

\begin{methodnote}[能量估计中的两个前提]
有界性已足以给出 $H_0^1$ 上的 Poincaré 估计。把 $\Omega$ 放进边长 $L$ 的立方体，对 $\varphi\in\Test(\Omega)$ 零延拓，沿第一坐标从立方体外侧积分，Cauchy--Schwarz 与 Fubini 给出 $\|\varphi\|_2\le L\|\partial_1\varphi\|_2$；再由闭包定义推广到 $H_0^1$。这里无需先使用边界迹定理。

Lipschitz 边界使迹算子有定义且 $H_0^1=\ker\gamma$；这是标准迹定理的零迹刻画，此处作为外部结论使用。书后\hyperref[sec:supp-trace]{“迹与边界值”节}证明其半空间基本估计，并说明光滑区域的局部化。任意分布没有自动定义的边界值。
\end{methodnote}

\begin{pitfall}[分布解与能量解的条件差别]
Dirac 源并非在任何维数都属于 $H^{-1}$。在内点附近，$\delta_a\in H^{-1}$ 需要 $1>n/2$，所以仅在一维成立；二维临界情形也不成立。高维点源可有分布基本解，但不能据此直接套用上述 $H_0^1$ 弱解定理。
\end{pitfall}

\Needspace{10\baselineskip}
\section{热方程与温和分布初值}
考虑 $\partial_tu=\Delta u$，$u(0)=u_0\in\Temp(\R^n)$。Fourier 变换把方程化为 $\partial_t\widehat u=-|\xi|^2\widehat u$，自然令
\begin{equation}\label{eq:heat-solution}
 \widehat u(t)=e^{-t|\xi|^2}\widehat u_0,\qquad t\ge0.
\end{equation}
乘子 $e^{-t|\xi|^2}$ 的各阶导数都有多项式界，因此对 Schwartz 测试空间连续，公式合法。对于 $t>0$，也可写
\[
 u(t,x)=\Pair{(u_0)_y}{G_t(x-y)}.
\]
这里 $G_t(x-\cdot)\in\Sch$，故与温和分布卷积有定义；这使用 Gaussian 衰减，而不是第 \ref{ch:convolution} 章的紧支集假设。对 $x$ 求导可移入配对，得到 $u(t,\cdot)$ 光滑，且每个固定 $t>0$ 下各导数有至多多项式增长。

\begin{theorem}[热半群的基本性质]
由式 \eqref{eq:heat-solution} 定义的 $u$ 属于 $C([0,\infty);\Temp)$，在 $t>0$ 时可微并满足热方程，且 $u(t)\to u_0$ 于温和分布意义。此外 $G_t*G_s=G_{t+s}$。
\end{theorem}
\begin{proof}
对固定 $\varphi\in\Sch$，$e^{-t|\xi|^2}\varphi\to\varphi$ 于 $\Sch$。可先在大球内使用光滑收敛，再以更高权重统一控制尾部；对 $t$ 在紧区间内，乘子的每个导数具有统一多项式界。故对偶配对在 $t=0$ 连续。对 $t>0$ 在 Schwartz 半范数下求导，得到 $\partial_t\widehat u=-|\xi|^2\widehat u$。逆变换给出方程。半群公式由 Fourier 侧 $e^{-t|\xi|^2}e^{-s|\xi|^2}=e^{-(t+s)|\xi|^2}$ 与变换单射性得到。
\end{proof}
在 $C([0,T];\Temp)\cap C^1((0,T];\Temp)$ 解类中，这个初值解唯一。可在 Fourier 侧取紧支集探针，对时间依赖探针 $e^{t|\xi|^2}\psi(\xi)$ 求配对的导数；紧支集保证乘法合法，得到配对恒定。由初值得到式 \eqref{eq:heat-solution} 在所有紧支集探针上成立，再用稠密性确定温和分布。这种证明避免把指数增长乘子当成全 Schwartz 空间上的合法算子。

当 $u_0\in L^2$ 时，Plancherel 给出 $\|u(t)\|_2\le\|u_0\|_2$。当 $u_0=\delta_0$ 时，$u(t)=G_t$，点质量立即被扩散为光滑密度。

\section{唯一性与稳定性为什么依赖空间}
Hadamard 适定性要求在选定数据空间和解空间中，解存在、唯一、且连续依赖数据。仅写出形式解并不能同时得到三者。
反向热方程 $\partial_tu=-\Delta u$ 的频率因子是 $e^{t|\xi|^2}$。在圆周 $\mathbb T$ 上，取 $u_0^{(k)}(x)=e^{-kt_0}e^{ikx}$，$t_0>0$ 固定。对任意固定 $s$，
\[
 \|u_0^{(k)}\|_{H^s(\mathbb T)}\asymp
 e^{-kt_0}(1+k^2)^{s/2}\longrightarrow0,
\]
而时刻 $t_0$ 的解范数为 $e^{(k^2-k)t_0}(1+k^2)^{s/2}\to\infty$。每个单模解都存在，但解映射在该 Sobolev 空间中不连续。

在 $\R^n$ 上，反向热方程以 $\delta_0$ 为初值也不能产生任意正时刻的温和分布解：局部 Fourier 方程强制其变换为光滑函数 $e^{t|\xi|^2}$，它不能满足温和分布估计。若初值 Fourier 变换带有足够 Gaussian 衰减，则有限时间内可能仍有解。由此可见，改变测试空间或增长条件会改变存在性与稳定性。

\section{常系数系统的 Fourier 视角}
对系统 $\partial_tu=P(\partial_x)u$，形式解为
\[
 \widehat u(t,\xi)=e^{tP(i\xi)}\widehat u_0(\xi).
\]
若该矩阵乘子的各阶导数在每个有限时间区间具有统一多项式增长，就能在 $\Temp$ 上定义连续演化。若增长更快，则需要更强衰减的测试空间和相应对偶。盖尔芳特、希烙夫对这些增长、唯一性类和适定性类进行了系统研究；其一般分类定理需要更多加权空间工具，不能仅由这条形式公式推出。

\Needspace{10\baselineskip}
\section{一个完全可解的椭圆正则性模型}\label{sec:topic-bessel-regularity}
\begin{proposition}[带质量项的全空间方程]\label{prop:topic-bessel-inverse}
对任意 $s\in\R$ 和 $f\in H^s(\R^n)$，方程 $(1-\Delta)u=f$ 有唯一温和分布解，且
\[
 u\in H^{s+2}(\R^n),\qquad \|u\|_{H^{s+2}}=\|f\|_{H^s}.
\]
\end{proposition}
\begin{proof}
令 $\widehat u=(1+|\xi|^2)^{-1}\widehat f$。该光滑乘子及其各阶导数有多项式界，故在 $\Temp$ 上合法；代入加权范数即得等式。若齐次解 $v\in\Temp$ 满足 $(1+|\xi|^2)\widehat v=0$，再乘逆乘子得 $v=0$，故唯一。
\end{proof}
这说明椭圆算子在该模型中增加两阶正则性。唯一性依赖温和增长类：在一维，普通分布 $e^x$ 满足 $(1-\partial_x^2)e^x=0$，故不能把结论改成所有 $\Dist$ 中唯一。对 $-\Delta$，逆符号在零频率奇异；变系数和有边界的问题另需局部估计与边界条件，不能套用本命题。

\Needspace{29\baselineskip}
\section{习题}
\begin{exercise}[迭代的正则性增益]
设 $m\ge1$ 为整数，$f\in H^s(\R^n)$。求 $(1-\Delta)^mu=f$ 的唯一温和分布解，并写出它属于 $C^k$ 的一个充分条件。
\end{exercise}
\begin{solution}
$\widehat u=(1+|\xi|^2)^{-m}\widehat f$，故 $\|u\|_{H^{s+2m}}=\|f\|_{H^s}$。唯一性用逆乘子。若 $s+2m>k+n/2$，命题 \ref{prop:topic-sobolev-embedding} 给出 $C^k$ 代表元及导数的一致估计。
\end{solution}
\begin{exercise}[三维点源]
求 $-\Delta u=\delta_a$ 在 $\R^3$ 上的一个分布解，并说明为什么不能仅由这个方程得到唯一性。
\end{exercise}
\begin{solution}
取 $u(x)=1/(4\pi|x-a|)$。平移基本解即得方程。任意调和分布 $h$ 都可相加，$u+h$ 仍为解；例如任意常数。需附加无穷远行为等条件来消除齐次自由度。
\end{solution}
\begin{exercise}[能量唯一性]
若 $u,v\in H_0^1(\Omega)$ 是同一 Poisson 方程的能量弱解，证明 $u=v$。
\end{exercise}
\begin{solution}
相减后取测试函数 $u-v$，得 $\int|\nabla(u-v)|^2=0$。Poincaré 不等式推出 $\|u-v\|_2=0$，故两者代表同一 $H_0^1$ 元素。
\end{solution}
\begin{exercise}[输运的点质量]
对常向量 $b\in\R^n$，求解 $\partial_tu+b\cdot\nabla u=0$，$u(0)=\delta_a$。
\end{exercise}
\begin{solution}
取 $u(t)=\delta_{a+tb}$。对探针 $\varphi$，时间导数为 $b\cdot\nabla\varphi(a+tb)$，而 $b\cdot\nabla u$ 的作用为其负值，故方程成立。与热方程不同，质量集中移动而不扩散。
\end{solution}

\SetChapterInfo{subtitle={从增长控制到连续谱},keywords={W 型空间; Young 对偶; 广义固有函数},prereq={Schwartz 对偶; Fourier 变换},goals={理解更强测试空间的动机与谱展开的分布意义}}
\chapter{加权测试空间与广义固有函数}\label{ch:weighted}
\ChapterSource{主要对应盖尔芳特、希烙夫 \cite{gelfand} 第 I 章 §§1--3（印刷第 1--22 页）与第 IV 章 §§4--6（印刷第 177--198 页）。本章区分原书的 $W$ 型基本空间、这里给出的加权模型，以及可直接验证的 Fourier 谱展开。}

\section{为什么要改变测试空间}
$\Test$ 只观察局部，因此允许指数乃至更快增长的正则分布，但普通 Fourier 变换未必存在。$\Sch$ 可以使用 Fourier 自同构，却要求对偶对象能由多项式权重估计。若需要处理更快增长的对象，可以使测试函数衰减更快，让相应积分有意义。

但空间的集合关系与拓扑关系都要检查。若连续嵌入 $E_1\hookrightarrow E_2$ 稠密，则限制给出单射 $E_2'\hookrightarrow E_1'$。因此更小且更精细的测试空间通常对应更大的对偶对象类。若稠密性不成立，限制可能丢失信息，不能把两个对偶空间自然等同。

\section{凸权重与 \texorpdfstring{$W_M$}{W-M} 型空间}
按原书思想，取递增连续函数 $\mu:[0,\infty)\to[0,\infty)$，满足 $\mu(0)=0$、$\mu(r)\to\infty$，定义
\[
 M(r)=\int_0^r\mu(s)\,\dd s\quad(r\ge0),\qquad M(-r)=M(r).
\]
则 $M$ 为偶凸函数，并有 $M(r)/r\to\infty$。凸性和 $M(0)=0$ 给出
\[
 M(r)+M(s)\le M(r+s),\qquad 2M(r)\le M(2r)\quad(r,s\ge0).
\]
例如 $M(r)=r^p/p$，$p>1$，相应 $\mu(r)=r^{p-1}$。

\begin{definition}[本书使用的加权模型]
对 $a>0$，定义
\[
 E_{M,a}=\{\varphi\in C^\infty(\R):
 p_{a,q}(\varphi)=\sup_xe^{M(a|x|)}|\varphi^{(q)}(x)|<\infty
 \text{ 对每个 }q\ge0\}.
\]
在每一层取这些半范数的 Fréchet 拓扑，再令
\[
 E_M=\bigcup_{a>0}E_{M,a}
 =\bigcup_{j\ge1}E_{M,1/j}
\]
并赋予局部凸归纳极限拓扑。
\end{definition}
每层的完备性由带权各导数一致 Cauchy 与导数极限相容性证明。这一模型突出原书 $W_M$ 中的估计
\[
 |\varphi^{(q)}(x)|\le C_qe^{-M(a|x|)},
\]
其中 $a$ 对一个函数的所有阶导数共同使用。原书还引入 $W_{M,a}$ 中允许 $a-\delta$ 的估计及其他整函数基本空间，并仔细规定其收敛性；因此这里的 $E_M$ 作为明确拓扑的教学模型使用，不宣称逐项等同于原书全部 $W$ 空间。

由于 $M(ar)$ 比任何线性函数增长更快，$E_{M,a}\subset\Sch$ 连续。任意紧支集光滑函数属于每个 $E_{M,a}$。若 $M(r)=r^p/p$，则选择整数 $k$ 使 $2k\ge p$ 后，$e^{-x^{2k}}$ 也属于某个 $E_{M,a}$：每阶导数只是多项式乘以同一指数，而适当选 $a$ 能留出指数衰减余量。这说明模型含有非紧支撑测试函数。

\section{快速增长函数如何作用于加权探针}
设 $|f(x)|\le Ce^{b|x|}$。对 $\varphi\in E_{M,a}$，
\[
 |\int f\varphi|\le Cp_{a,0}(\varphi)
 \int_\R e^{b|x|-M(a|x|)}\,\dd x<\infty.
\]
每一层上都有这样的连续估计，故 $f$ 给出 $E_M$ 上的连续泛函。例如 $e^x$ 虽非温和分布，却能作用于这个加权测试空间。这并不自动给出 $\Sch'$ 中的 Fourier 变换；它只说明已选对偶框架扩大了可处理的增长范围。

对于更快增长 $e^{M(b|x|)}$，当 $a>b$ 时它能作用于足够强衰减的固定层。却不能据此声称它在所有 $a>0$ 的并上都能积分。量词“存在衰减参数”与“对任意衰减参数”会给出不同的对偶空间，正是原书细分空间的重要原因。

\section{Young 对偶与复频率增长}
\begin{definition}[凸共轭]
定义 $M^*(s)=\sup_{r\ge0}(rs-M(r))$，$s\ge0$。则
\begin{equation}\label{eq:young-conjugate}
 rs\le M(r)+M^*(s)
\end{equation}
称为 Young 不等式。
\end{definition}
若 $M(r)=r^p/p$，最大值由 $s=r^{p-1}$ 给出，故 $M^*(s)=s^{p'}/p'$，其中 $1/p+1/p'=1$。若 $\mu$ 严格递增，可用其反函数积分表示 $M^*$；原书用这类对偶权重描述 Fourier 变换前后的增长。

取 $\varphi\in E_{M,a}$，其 Fourier 积分可延拓到复数 $z=\xi+i\eta$：
\[
 \widehat\varphi(z)=\int_\R e^{-ixz}\varphi(x)\,\dd x.
\]
因为 $|e^{-ixz}|\le e^{|\eta||x|}$，超线性衰减保证在每个复紧集上一致可积。各阶复导数也由附加的 $x^k$ 权重控制，因此它是整函数。Young 不等式进一步给出
\[
 |\eta||x|\le\tfrac12M(a|x|)+\tfrac12M^*(2|\eta|/a),
\]
从而
\begin{equation}\label{eq:entire-growth}
 |\widehat\varphi(\xi+i\eta)|\le C_a p_{a,0}(\varphi)
 \exp\!\left[\tfrac12M^*(2|\eta|/a)\right],
 \quad C_a=\int_\R e^{-M(a|x|)/2}\,\dd x<\infty.
\end{equation}
这条可直接证明的估计说明：实轴方向的衰减权重决定复频率虚部方向的增长权重。要证明两种完整基本空间互为 Fourier 像，还必须控制实部方向的衰减、复导数以及逆变换；单独的估计 \eqref{eq:entire-growth} 不等于完整空间同构定理。

\section{Hilbert 空间中的连续谱}
有限维自伴矩阵有正交本征向量基。无限维自伴算子可能有连续谱，此时 Hilbert 空间内未必存在对应本征向量。例如 $A=-i\partial_x$ 在 $L^2(\R)$ 上，取定义域 $H^1(\R)$，Fourier 变换将其变为乘法算子 $\xi$。若 $Au=\lambda u$，则 $(\xi-\lambda)\widehat u=0$，$L^2$ 函数只能支于单点时必为零，故不存在非零 $L^2$ 本征向量。

然而 $e^{i\lambda x}$ 是温和分布，满足
\[
 -i\partial_xe^{i\lambda x}=\lambda e^{i\lambda x}.
\]
因此它可作为广义固有函数。它不是单位范数向量；“正交性”只能以分布核或谱变换的意义解释。

\section{装配空间与可核验的广义展开}
考虑连续稠密嵌入
\[
 \Sch(\R^n)\subset L^2(\R^n)\subset\Temp(\R^n).
\]
第二个映射为 $f\mapsto T_f$，由 Cauchy--Schwarz 连续；使用双线性配对时这是线性嵌入。若改为通常的共轭线性反对偶约定，需要相应调整 Hilbert 空间的识别。这类三重结构称为装配 Hilbert 空间或 Gelfand 三元组的典型模型。

对 $u\in\Sch$，Fourier 反演可写成
\begin{equation}\label{eq:generalized-expansion}
 u=(2\pi)^{-n}\int_{\R^n}\widehat u(\xi)E_\xi\,\dd\xi,
 \qquad E_\xi(x)=e^{ix\cdot\xi}.
\end{equation}
等式在 $\Temp$ 中的弱意义是：对每个 $\varphi\in\Sch$，
\[
 \Pair{u}{\varphi}=(2\pi)^{-n}
 \int\widehat u(\xi)\widehat\varphi(-\xi)\,\dd\xi.
\]
所有积分绝对收敛，故 Fubini 合法。每个 $E_\xi$ 同时是 $-i\partial_j$ 的广义本征分布，本征值为 $\xi_j$，也是 $-\Delta$ 的广义本征分布，本征值为 $|\xi|^2$。

对 $u\in L^2$，展开需要用 Fourier 截断在 $L^2$ 中取极限，不宜声称逐点绝对积分总成立。Plancherel 定理提供的“连续正交性”是
\[
 \|u\|_2^2=(2\pi)^{-n}\|\widehat u\|_2^2.
\]
形式核 $\int e^{ix\cdot(\xi-\eta)}\,\dd x=(2\pi)^n\delta(\xi-\eta)$ 则必须按分布恒等式解释，不能把左侧作为普通可积函数计算。

\section{一般谱展开需要什么假设}
盖尔芳特、希烙夫通过适当基本空间、其 Hilbert 完备化和谱测度的微分来构造固有泛函。一般装配 Hilbert 空间的谱定理通常要求基本空间具备核性等额外性质、算子保持该空间并连续，且其 Hilbert 实现满足自伴性等条件；还涉及相对于谱测度“几乎处处”定义的广义本征泛函。不同版本的假设并不相同。

本稿已证明的具体模型是式 \eqref{eq:generalized-expansion} 的 Fourier 展开。一般自伴算子的广义本征展开在此作为后续研究方向，不把“任意自伴算子在任意测试空间都有完整广义本征基”作为结论。对于有界区域的椭圆算子，紧预解式常带来离散谱；对于全空间算子，连续谱与广义固有函数则更自然。

\section{习题}
\begin{exercise}[Gaussian 权重的对偶]
令 $M(r)=r^2/2$，求 $M^*$，并由式 \eqref{eq:entire-growth} 写出虚部方向的增长估计。
\end{exercise}
\begin{solution}
$rs-r^2/2$ 在 $r=s$ 取最大值，因此 $M^*(s)=s^2/2$。于是 $\frac12M^*(2|\eta|/a)=|\eta|^2/a^2$，得到 $|\widehat\varphi(\xi+i\eta)|\le C_a p_{a,0}(\varphi)e^{|\eta|^2/a^2}$。
\end{solution}
\begin{exercise}[乘法算子的广义本征向量]
对 $M_xf=xf$，说明 $\delta_\lambda$ 是广义本征分布，但并非 $L^2(\R)$ 中的向量。
\end{exercise}
\begin{solution}
对探针 $\varphi$，$(x\delta_\lambda)(\varphi)=\lambda\varphi(\lambda)$，故 $x\delta_\lambda=\lambda\delta_\lambda$。Dirac 分布不是局部可积函数，当然也不是 $L^2$ 函数。该例与微分算子的平面波本征分布通过 Fourier 变换联系起来。
\end{solution}

\appendix
\chapter{吉田第六版目录与版本对照}\label{app:yosida-map}
\section{主线章节与参考版本}
本书第一部分保留中译本第零章至第十四章的原章次。以下页码为原书印刷页码；本地两个 PDF 的正文第 1 页均位于 PDF 第 16 页，因此正文部分的 PDF 页序等于印刷页码加 15。前置页不使用这一换算。中译本目录位于 PDF 第 10--15 页，英文版目录位于 PDF 第 11--15 页。
\begin{center}\small
\begin{tabular}{p{0.08\linewidth}p{0.53\linewidth}p{0.14\linewidth}p{0.14\linewidth}}
\toprule 章 & 主题 & 中译本页码 & 英文版页码\\\midrule
0 & 预备知识 & 1--22 & 1--22 \\
1 & 半范数 & 23--70 & 23--67 \\
2 & Baire--Hausdorff 定理的应用 & 71--84 & 68--80 \\
3 & 正交投影及 F. Riesz 表示定理 & 85--106 & 81--101 \\
4 & Hahn--Banach 定理 & 107--124 & 102--119 \\
5 & 强收敛和弱收敛 & 125--150 & 119--145 \\
6 & Fourier 变换和微分方程 & 151--200 & 145--192 \\
7 & 对偶算子 & 201--218 & 193--208 \\
8 & 预解式和谱 & 219--242 & 209--231 \\
9 & 半群的分析理论 & 243--286 & 231--274 \\
10 & 紧算子 & 287--306 & 274--293 \\
11 & 赋范环和谱表示 & 307--376 & 294--362 \\
12 & 线性空间中其他表示定理 & 377--396 & 362--378 \\
13 & 遍历理论和扩散理论 & 397--436 & 379--418 \\
14 & 发展方程的积分 & 437--486 & 418--465 \\
\bottomrule\end{tabular}\end{center}
\section{目录核对说明}
章节名与各节主题以中译本为准，统一连字符、标点与少数人名拼写的排版。不使用英文版页码替代中译本页码。第五章附录的四个子目在中译本总目录中未展开，依据英文版目录补足，已在第五章附录中保留。空间记号表依据英文版书后目录建立辅助占位。
英文章名及页码用于辅助核对目录，不替代中译本页码。正文复查确认几处共页：英文第 119 页上部为第四章结尾、下部开始第五章；第 145 页由第五章接到第六章；第 231 页由第八章接到第九章；第 274 页由第九章接到第十章。本次又确认第 362 页上部仍为第十一章 Littlewood 推论证明结尾，下部才开始第十二章，故第十一章英文范围修正为 294--362 页，第十二章仍从 362 页起。本轮确认英文第 418 页上部仍为第十三章抽象位势说明、下部开始第十四章，故第十三章改为 379--418 页，第十四章起始 418 不变。这些范围允许重叠，编号节起始页未改。

目前第零至第十三章共 125/132 个编号节已有首轮正文，配 218 道附解答习题，第五、十章的章内附录亦已写入。第十一章涵盖交换赋范代数、谱测度、Friedrichs 扩张、紧群表示及 Wiener 稠密定理，配三十二道习题。第十二章六节及十四道习题连接凸集端点、序收敛与格的表示。第十三章九节与十八道习题连接遍历、扩散与位势。各章的实际查阅页、数学自查、外部前提与尚未展开的证明均分别记录于对应 COVERAGE，不把已有正文或编译通过等同于独立数学审校。第十四章的七节仍保留占位。

\section{广义函数专题的接入位置}
第二部分为第十五章至第二十二章，沿用现有八个专题文件。第一部分中 I.8、I.12--14、II.3、III.11、VI.1--3 等主题仍保留其原书位置；相关节用标签与对应专题章衔接。这使主线目录保持完整，也便于后续逐节写入吉田原书的特定证明与论述。
\section{占位完成后的更新规则}
正文写成后删除占位提示，把页码转入出处说明；保留章、节标题及标签。新增定理与习题使用独立标签，章号随全书顺序自动更新。

\chapter{参考书对照与公式核查}\label{app:references}
\section{本稿与两本主要参考书的对应}
下列对照依据本地扫描 PDF 的目录与所检查的正文。印刷页码指原书页面中的页码；PDF 页序从 1 计数，包括前置页。两者切勿混用。
本附录专用于第二部分“广义函数专题扩展”；吉田主线的版本对照见附录 \ref{app:yosida-map}。

\begin{center}\small
\begin{tabular}{p{0.25\linewidth}p{0.43\linewidth}p{0.22\linewidth}}
\toprule
本稿主题 & 原书位置 & 页码线索\\\midrule
动机、分布定义 & 施瓦尔茨，第 I 章 §§1--3 & 印刷 1--16；章首 PDF 24\\
测试空间拓扑 & 施瓦尔茨，第 III 章 §§1--3 & 印刷 43--53；章首 PDF 66\\
求导与局部结构 & 施瓦尔茨，第 II 章；第 III 章 §§4--10 & 第 II 章 PDF 40；点支集印刷 70 / PDF 93\\
张量积与光滑乘法 & 施瓦尔茨，第 IV--V 章 & 章首 PDF 96、106\\
卷积与基本解 & 施瓦尔茨，第 VI 章 §§2、10 & 卷积印刷 111 / PDF 134\\
温和分布与 Fourier & 施瓦尔茨，第 VII 章 §§2--6 & 印刷 169、172 / PDF 192、195\\
凸权重与 $W$ 型空间 & 盖尔芳特、希烙夫，第 I 章 §§1--3 & 印刷 1--22 / PDF 5--26\\
柯西问题与适定性 & 盖尔芳特、希烙夫，第 II--III 章 & 章首 PDF 30、107\\
广义固有函数 & 盖尔芳特、希烙夫，第 IV 章 §§4--6 & 印刷 177--198 / PDF 181--202\\
\bottomrule
\end{tabular}
\end{center}
盖尔芳特、希烙夫扫描目录的第一章标题显示为 $M$ 型空间，正文标题使用 $W$ 型空间，故本稿优先遵从正文符号。该卷本身不是广义函数初步教程，它多处调用前两卷结果。本稿将其作为应用和拓展的主要参考，基础内容以施瓦尔茨的系统叙述为主。

\section{运算条件速查}
\begin{center}\small
\begin{tabular}{p{0.23\linewidth}p{0.43\linewidth}p{0.24\linewidth}}
\toprule
运算 & 本稿使用的充分条件 & 所得对象\\\midrule
$\partial^\alpha T$ & $T\in\Dist(\Omega)$ & 分布\\
$aT$ & $a\in C^\infty(\Omega)$ & 分布\\
$T*\rho$ & $T\in\Dist(\R^n),\rho\in\Test$ & 光滑函数\\
$S*T$ & 至少一个因子紧支集 & 分布\\
$S*T$ & 加法在支集乘积上适当 & 分布\\
$\widehat T$ & $T\in\Temp(\R^n)$ & 温和分布\\
$S*T$ 的 Fourier 公式 & $S\in\CompactDist,T\in\Temp$ & 温和分布\\
$F^*T$ & $F$ 是光滑微分同胚 & 分布\\
$G_t*T$ & $t>0,T\in\Temp$ & 光滑函数与温和分布\\
\bottomrule
\end{tabular}
\end{center}
这些都是明确可用的充分条件，不宣称列尽所有可能的合法情形。

\section{检查证明时的五个问题}
\begin{enumerate}
\item 测试空间是什么？要求共同紧支集，还是全空间带权衰减？
\item 配对是否线性？是否混入了 Hilbert 内积的复共轭？
\item 对偶运算由哪个测试函数映射的转置得到？该映射是否连续？
\item 移动导数、极限或积分时，支集和半范数是否有统一控制？
\item 唯一性和稳定性在哪个解空间与数据拓扑中成立？
\end{enumerate}

\Needspace{14\baselineskip}
\section{与已完成主线的衔接}
测试空间与紧支集分布衔接第一章 §§8、12--14；分布极限衔接第二章 §3；局部结构衔接第三章 §11。Fourier、卷积及 Sobolev 估计与第六章 §§1--9 互补；本稿第六章 §4 已有 Paley--Wiener 内容，不再把它整体列为未写专题。紧嵌入、Dirichlet 弱解和核空间分别见第十章 §§3、6 及其附录。一般自伴谱测度在第十一章，半群及发展方程在第九、十四章。

本轮新增的全空间嵌入、带质量项椭圆模型和加权空间稠密性，分别见命题 \ref{prop:topic-sobolev-embedding}、\ref{prop:topic-bessel-inverse}、\ref{prop:weighted-density}；边界迹与反向密度演化见书后补充说明。本稿尚未证明一般 Schwartz 核定理、完整 $W$ 型空间 Fourier 同构、一般装配空间谱定理以及非光滑区域的完整迹与正则性定理。查阅这些主题时应保留各节注明的外部前提，不能把有限模型的证明提升为一般结论。

\backmatter
\chapter{补充说明}\label{back:yosida-notes}
本章依据吉田第六版中译本印刷第 487--490 页（PDF 502--505）、英文版第 466--468 页（PDF 481--483）整理。按原书主题顺序记录与本稿的联系，并补写迹估计、Markov 密度演化及非时齐唯一性的推导。原书所说的“近期发展”指该版出版时的文献状况；书目线索不表示已直接阅读所列原始论文。

\section*{广义函数、超函数与 Sobolev 嵌入}
\phantomsection
原书第一、六章补注指向 Schwartz 的增订本及 Gelfand 系列英文版。本稿第二部分以分布、温和分布和加权探针构成具体阅读路线。Sato 的超函数及 Schapira 的相关工作属于另一理论框架，不能把任意超函数直接当成 $\Dist$ 元素；本稿未建立其定义或运算。

第六章补注指出原书 Sobolev 引理只是一般嵌入定理的特例，并列 Sobolev、Adams、Stein 的线索。本稿第六章 §7 与命题 \ref{prop:topic-sobolev-embedding} 给出全空间 $H^s\hookrightarrow C^k$ 的 Fourier 证明，条件是 $s>k+n/2$。一般 $W^{m,p}$ 嵌入、临界指数及非光滑区域的延拓不从这个特例自动得到。

\section*{迹与边界值}
\phantomsection\addcontentsline{toc}{section}{迹与边界值}\label{sec:supp-trace}
原书第十章补注在有界 $C^2$ 区域上陈述 $W^{1,2}$ 的 $L^2$ 边界迹，并提及沿法线的几乎处处极限。这里先证明用于能量方法的基本范数估计。
\begin{proposition}[半空间的迹估计]\label{prop:supp-trace}
令 $\R^n_+=\{(x',t):t>0\}$。对 $u\in C_c^\infty(\overline{\R^n_+})$，其中该记号指邻域内光滑函数的限制，有
\[
 \|u(\cdot,0)\|_2^2
 \le 2\|u\|_{L^2(\R^n_+)}\|\partial_tu\|_{L^2(\R^n_+)}
 \le \|u\|_{H^1(\R^n_+)}^2.
\]
因此边界限制唯一连续延拓为 $\gamma:H^1(\R^n_+)\to L^2(\R^{n-1})$。
\end{proposition}
\begin{proof}
微积分基本定理给出
\[
 |u(x',0)|^2=-2\operatorname{Re}\int_0^\infty
 \partial_tu(x',t)\overline{u(x',t)}\,\dd t.
\]
对 $x'$ 积分，用 Cauchy--Schwarz 和 $2ab\le a^2+b^2$ 即得估计。光滑至边界且紧支撑的函数在半空间 $H^1$ 中稠密：先切断，再沿内法向平移后用半径小于平移距离的光滑核正则化；平移及卷积在 $L^2$ 中逼近函数和弱导数。故估计延伸到闭包，延拓与逼近序列无关。
\end{proof}
在有界 $C^2$ 区域上，用有限边界图展平边界，再用光滑单位分解；坐标 Jacobian 及其逆有界，链式法则控制 $H^1$ 范数，于是同一估计给出
\[
 \|\gamma u\|_{L^2(\partial\Omega)}\le C_\Omega\|u\|_{H^1(\Omega)}.
\]
内部片无边界贡献。标准边界图延拓与正则化提供这里的光滑稠密性；Lipschitz 区域的完整迹定理及 $H_0^1=\ker\gamma$ 的逆向包含仍作为外部结论使用。此处直接可见 $H_0^1\subset\ker\gamma$，因为紧支集光滑逼近的迹全为零。

迹是 $L^2(\partial\Omega)$ 等价类，不是任意代表元的边界点值。原书更强的沿法线极限表述须选择合适代表元并使用切片定理；上面的范数延拓证明没有单独证明该几乎处处极限。第 \ref{ch:pde} 章的能量弱解使用连续迹与零迹空间。

\section*{核空间、收缩算子与凸表示的线索}
\phantomsection
第十章附录补注把 Minlos 的柱分布延拓指向 Schwartz 的统一处理。第 \pageref{app:yosida-10} 页已定义核空间并证明到 Banach 空间的核映射判别；这与概率测度延拓之间仍隔着独立的存在性定理，不能只用有限维相容性代替连续性及可数可加性。

第十一章补注推荐 Sz.-Nagy--Foias 的 Hilbert 空间收缩算子理论。第九章的收缩半群和第十一章的酉谱表示提供本稿入口；一般收缩的酉扩张与函数模型未在这里证明。第十二章补注以 Phelps 为 Choquet 理论线索：Krein--Milman 的闭凸包结论不等于已经构造了集中于端点集的表示测度，可测性、紧性及唯一性须另立条件。

\section*{不变测度下的反向密度演化}
\phantomsection\addcontentsline{toc}{section}{不变测度下的反向密度演化}\label{sec:supp-markov}
沿用第十三章 §1 的核 $P_t$、不变概率 $m$ 和 $T_tg(x)=\int g(y)P_t(x,\dd y)$。
\begin{proposition}[密度的对偶 Markov 半群]\label{prop:supp-markov-density}
对每个 $f\in L^1(m)$，存在唯一 $U_tf\in L^1(m)$，使对所有可测 $B$，
\[
 \int_B U_tf\,\dd m=\int_S f(x)P_t(x,B)\,m(\dd x).
\]
$U_t$ 是正收缩，保积分且 $U_t1=1$，满足半群律，以及
\[
 \int (U_tf)g\,\dd m=\int f(T_tg)\,\dd m
 \qquad(f\in L^1,\ g\in L^\infty).
\]
\end{proposition}
\begin{proof}
右端定义有限复测度 $\nu_f$，可数可加性由控制收敛得到。对有限可测分割 $\{B_i\}$，
\[
 \sum_i|\nu_f(B_i)|\le\int |f(x)|\sum_iP_t(x,B_i)\,m(\dd x)
 \le\|f\|_1,
\]
故全变差至多为 $\|f\|_1$。若 $m(B)=0$，不变性给出 $P_t(x,B)=0$ 几乎处处，故 $\nu_f\ll m$。Radon--Nikodym 定理得到唯一密度及收缩估计。正性、保积分和 $U_t1=1$ 分别由非负源、$B=S$ 和不变性得到。对偶公式先对示性函数，再对简单函数和有界极限成立。由 $T_{t+s}=T_tT_s$ 与对偶分离性得 $U_{t+s}=U_sU_t$；时间参数可交换，所以也是半群律。
\end{proof}
原书将 $U_tf$ 记作 $f_{-t}$。若 $P_t$ 来自可逆保测流 $\tau_t$，则 $T_tf=f\circ\tau_t$，$U_tf=f\circ\tau_{-t}$。一般随机转移中，$U_t$ 不是 $T_t$ 的逆；也不能据此推出细致平衡 $U_t=T_t$。有限状态、不变概率 $\pi_i>0$ 时，反向转移为
\[
 q_{ij}(t)=\frac{\pi_jp_{ji}(t)}{\pi_i},\qquad
 (U_tf)_i=\sum_jq_{ij}(t)f_j.
\]
例如三点单向循环链有不变均匀概率，但反向循环核与原核不同。

对一步算子 $U=U_1$，$n^{-1}\sum_{k=1}^nU^kf$ 在 $L^1$ 中收敛，极限保积分。证明与定理 \ref{thm:y13-markov-mean} 相同：对有界 $f$ 先在 $L^2$ 中平均，再用概率空间的 $L^2\subset L^1$；最后以 $L^1$ 收缩及有界函数稠密性推广。$U$ 的 $L^2$ 收缩性由其正性、$U1=1$ 和保积分的 Jensen 估计得到。求和从 $k=1$ 开始仅与从 $k=0$ 开始相差 $(U^nf-f)/n$。

\Needspace{16\baselineskip}
\section*{非时齐经典解的唯一性}
\phantomsection\addcontentsline{toc}{section}{非时齐经典解的唯一性}\label{sec:supp-uniqueness}
原书第十四章 §4 补注用每一时刻的收缩预解式证明唯一性。以下将该论证单独写出。
\begin{proposition}[冻结时刻的预解式估计]
设每个 $A(t)$ 为稠定线性算子，且所有充分小 $h>0$ 的 $J_h(t)=(I-hA(t))^{-1}$ 在全空间上存在、$\|J_h(t)\|\le1$。若 $x\in C^1([0,T];X)$、$x(t)\in D(A(t))$ 且 $x'=A(t)x$，则 $\|x(t)\|\le\|x(0)\|$。同一初值的经典解至多一个。
\end{proposition}
\begin{proof}
固定 $t$。对 $z\in D(A(t))$，$J_h(t)z-z=hJ_h(t)A(t)z\to0$；稠密性及一致有界性推广到所有 $z\in X$。可微性给出
\[
 x(t+h)=x(t)+hA(t)x(t)+o(h)=J_h(t)x(t)+o(h),
\]
因为 $J_h(t)x(t)=x(t)+hJ_h(t)A(t)x(t)$。于是范数的右上 Dini 导数至多为零，连续函数的 Dini 单调性引理推出范数不增。该引理可对函数减去 $\varepsilon t$，用紧区间极值点的右侧差商反证，再令 $\varepsilon\downarrow0$。两个解之差满足同一齐次方程，故唯一。
\end{proof}
第十四章 §4 的冻结乘积还承担存在性和演化族构造，需要共同定义域与稳定性等额外假设。本命题仅说明经典解已经存在时，唯一性可由上述较少条件得到。

\section*{发展方程与凸分析的后续阅读}
\phantomsection
第九、十四章补注列出 Krein、Lax--Phillips、Lions、Tr\`eves 的线性发展方程著作，以及 Barbu、Browder、Martin、Masuda、Miyadera、Tanabe 的线性或非线性理论；凸分析指向 Rockafellar 和 Ekeland--Temam。本稿第十四章已给出线性演化、极大单调算子及隐式 Euler 指数公式的具体入口。书目定位见参考文献后的对照表；原书有关译本“正在准备”的措辞不作为当前出版状态使用。

\PrintSolutions
\begin{thebibliography}{99}
\addcontentsline{toc}{chapter}{参考文献}
\bibitem{yosida-zh} 吉田耕作（K\={o}saku Yosida）.
\textbf{泛函分析}，第六版，中译本.
主线目录依据本地文件《泛函分析 第6版 (吉田耕作) 中译本.pdf》。
目录位于 PDF 第 10--15 页，正文印刷第 1 页位于 PDF 第 16 页。
\bibitem{yosida-en} K\={o}saku Yosida.
\textbf{Functional Analysis}, Sixth Edition. Springer.
英文辅助依据本地文件《泛函分析 第6版 英文原版.pdf》。
目录位于 PDF 第 11--15 页，正文印刷第 1 页位于 PDF 第 16 页。
\bibitem{schwartz} L. 施瓦尔茨（Laurent Schwartz）著，姚家燕译.
\textbf{广义函数论}，法兰西数学精品译丛.
本稿依据工作目录中同名中文扫描版；具体查阅位置见附录 \ref{app:references}。
\bibitem{gelfand} I. M. 盖尔芳特，G. E. 希烙夫（Gelfand--Shilov）.
\textbf{广义函数 III：微分方程理论的若干问题}.
本稿依据工作目录中同名中文扫描版。
\bibitem{alexandroff-hopf} P. Alexandroff (Alexandrov), H. Hopf.
\textbf{Topologie}, Vol.~I. Springer, 1935.
吉田第零章拓扑参考；原书总书目英文版印刷第 469 页。
\bibitem{bourbaki-topology} N. Bourbaki.
\textbf{Topologie G\'en\'erale}.
Actualit\'es Scientifiques et Industrielles, nos.~858, 916, 1029, 1045, 1084.
Hermann, 1940--1942（按吉田英文版印刷第 470 页的分册书目记录）。
现代合订重印本另见 \href{https://link.springer.com/book/10.1007/978-3-540-33982-3}{\emph{Topologie g\'en\'erale, Chapitres 1 \`a 4}, Springer, 2007}；重印本与原书引文的版次须区别。
\bibitem{kelley-topology} J. L. Kelley.
\textbf{General Topology}. D.~Van Nostrand, 1955.
吉田第零章拓扑参考；原书总书目英文版印刷第 475 页。
后由 Springer 再版为 Graduate Texts in Mathematics, Vol.~27, 1975；
\href{https://link.springer.com/book/9780387901251}{出版社书目页}注明原版出版于 1955 年。
\bibitem{halmos-measure} P. R. Halmos.
\textbf{Measure Theory}. D.~Van Nostrand, 1950.
吉田第零章测度参考；原书总书目英文版印刷第 473 页。
后由 Springer 重印，见 \href{https://link.springer.com/book/10.1007/978-1-4684-9440-2}{出版社书目页}。
\bibitem{saks-integral} S. Saks.
\textbf{Theory of the Integral}. Warszawa, 1937.
按吉田英文版印刷第 481 页的书目记录；原书章末列为 Saks [1]。
\end{thebibliography}
\section*{原书参考书目整理（占位）}
\phantomsection
\addcontentsline{toc}{section}{原书参考书目整理（占位）}
中译本印刷第 491--514 页（PDF 第 506--529 页）；
英文版印刷第 469--486 页（PDF 第 484--501 页）。
第零章涉及的 Alexandroff--Hopf、Bourbaki、Kelley、Halmos、Saks 五条已核对并收入上表，其余原书参考书目后续按实际引用逐条整理。

\chapter{索引}\label{back:index}
全书索引占位，后续随主线正文逐步建立术语、人名、定理与记号的索引条目。
中译本原书索引从印刷第 515 页起，英文版索引从印刷第 487 页起。
此处不手工填入尚未完成正文的索引页码，也不把原书索引页码当作本书页码。
% TODO: 正文写作开始后选用 makeidx 或 imakeidx，并同步更新编译脚本。

\chapter{空间记号表（英文版辅助）}\label{back:notation}
英文第六版书后目录列有空间记号表（Notation of Spaces，印刷第 501 页起），
本书在此建立辅助占位；它不是中译本总目录中另列的主线章节。
后续按本书实际采用的记号整理空间、对偶、算子和拓扑，并注明同一符号在不同参考书中的含义。
当前广义函数的核心记号与 Fourier 约定见前言。

\end{document}
